% !TeX root = ../Book2.tex
% 纲要标题：# 第三编　代数、有限性、半单性与根
% 纲要位置：工作记录/计划纲要.md，第 1404 行。

\BookPart{代数、有限性、半单性与根}{b2:part:03}
\BookPartGuide
环的乘法可以通过它在模上的作用来研究。有限维代数尤其适合这样处理：它的正则模既记录全部乘法，又满足有限长度条件。本编从代数、双模和幂等元出发，逐步区分可以拆成简单部分的半单结构，以及使这些简单部分连接起来的根与扩张。

\ProseParagraphBreak
\chapref{b2:ch:11}重整域上代数的语言。结构常数记录选基后的乘法，反环保证右作用与复合次序一致，包络代数把双模写成一个左模。幂等元的分块则说明：模的直和不自动成为环的直积，非中心幂等元之间仍有非零的交叉乘积。正则作用和两个张量因子的中心化子提供了可以逐式核验的例子。

\ProseParagraphBreak
\chapref{b2:ch:12}讨论链条件、长度与分解。Jordan--Hölder 定理\footnote{赫尔德（Otto Ludwig Hölder，1859-1937），德国数学家，研究群的合成列及分析学中的不等式；Jordan 已在本卷序言介绍。}比较合成因子，Krull--Schmidt 定理\footnote{克鲁尔（Wolfgang Krull，1899-1971），德国数学家，研究交换环的维数与理想理论。施密特（\OriginalName{Отто Юльевич Шмидт}，1891-1956），苏联数学家及地球物理学家，研究群的直积分解。}比较不可分解直和因子。前者逐层取商，后者要求实际拆开，两种操作所保留的信息不同。有限长度使自同态的核链与像链同时稳定，Fitting 引理\footnote{菲廷（Hans Fitting，1906-1938），德国数学家，研究群与环中的分解问题。}由此分出幂零作用与可逆作用的两个部分。在不可分解模上，这个直和只能留下其中一部分，进而得到局部自同态环；这一局部性支持 Krull--Schmidt 的唯一性证明。具体反例将说明仅有有限生成或 Noether 条件还不足以代替有限长度。

\ProseParagraphBreak
\chapref{b2:ch:13}转而考察所有子模都能分裂的半单情形。半单模是简单模的直和，子模都有补模；对环自身的正则模，单位元又迫使简单直和只有有限个分量。Artin--Wedderburn 定理把半单环写成有限个除环矩阵环的乘积，其中的反环来自左正则模自同态的复合次序。这给出了完全分裂时的结构，下一章则用简单模共同检测不到的元素研究一般环。有限除环的交换性还需用类方程和分圆多项式证明，不能只因有限就把一般自同态除环预先视为域。

\ProseParagraphBreak
\chapref{b2:ch:14}用 Jacobson 根刻画所有简单模共同零化的元素。根为零本身并不保证半单；在左正则模 \({}_RR\) 有限长度时，第十二章的链条件把 \(R/J(R)\) 中极大子模的零交缩为有限交，进而将它嵌入有限个简单模的直和并证明它半单。可逆性判据导出非交换 Nakayama 引理，\footnote{中山正（假名：なかやま ただし，罗马音：Tadashi Nakayama，1912-1964），日本数学家，研究环论与表示论。}使剩余模中的生成元能够提升。随后把简单模的自同态看成右除环标量，证明 Jacobson 稠密性；代数闭域上的有限维情形再给出 Burnside 矩阵代数定理\footnote{伯恩赛德（William Burnside，1852-1927），英国数学家，研究有限群及其表示，著有《有限阶群论》。}。半本原、本原与半单各自提出不同要求，正文会给出区分它们的例子。

\ProseParagraphBreak
\chapref{b2:ch:15}回到有限维代数 \(A\)，记 \(J=J(A)\)。根幂的降链先由有限维性稳定，再由 Nakayama 引理得到稳定项为零。半单商 \(A/J\) 和被 \(J\) 零化的各个 Loewy 层\footnote{洛伊（Alfred Loewy，1873-1935），德国数学家，研究线性代数与群表示；模的根滤过以他命名。}由此描述模的分层结构，各层之间的扩张仍可能不分裂。构造投射覆盖时，先将 \(A/J\) 中的本原幂等元提升到 \(A\)，由左理想 \(Ae\) 得到简单模的不可分解投射覆盖。对一般的有限维左 \(A\) 模 \(M\)，再将 \(M/JM\) 分成简单模，按其重数取这些覆盖的直和作为来源 \(P\)。投射性使 \(P\to M/JM\) 沿商映射提升成 \(p:P\to M\)；Nakayama 引理保证 \(p\) 满射，商上同构又使 \(\ker p\subseteq JP\)，因而核多余。这里提升的是投射模出发的映射，商中的简单分解并不把 \(M\) 直接分成简单模。

\begin{center}
\begin{tikzcd}[column sep=large,row sep=large]
\begin{gathered}\text{第11章}\\\text{代数、双模与幂等元}\end{gathered}\arrow[r]\arrow[dr]
&\begin{gathered}\text{第12章}\\\text{链条件与有限长度}\end{gathered}\arrow[r]\arrow[d]\arrow[dr]
&\begin{gathered}\text{第13章}\\\text{半单结构}\end{gathered}\arrow[d]\arrow[dl,crossing over]\\
&\begin{gathered}\text{第15章}\\\text{根滤过与投射覆盖}\end{gathered}
&\begin{gathered}\text{第14章}\\\text{根、生成元与稠密性}\end{gathered}\arrow[l]
\end{tikzcd}
\end{center}

图中列出主要的直接依赖。投射覆盖的构造使用\cref{b2:thm:noncommutative-nakayama}检验满射性，使用\cref{b2:thm:artin-wedderburn}分解半单商，并由\cref{b2:prop:corner-endomorphism}识别 \(\operatorname{End}_A(Ae)\)；幂等元的提升在第十五章就地证明。

\ProseParagraphBreak
掌握第一编的自由模、展示、正合列、张量积与投射模后，可以先学习本编的基本代数与模结构。包络代数 \(A^e=A\otimes_kA^{\mathrm{op}}\) 使用第六章的张量积构造，代数张量积的乘法及其良定义在\secref{b2:sec:1102}就地证明。\secref{b2:sec:1104}的张量作用例子还使用第八章已构造的交换作用与 Schur--Weyl 对偶背景；\secref{b2:sec:1402}关于交换环 Nakayama 引理的证明说明则调用第九章的伴随矩阵恒等式。有限维代数的模不一定有限维，正文出现“有限生成”“有限长度”时都应检查它修饰的是环还是所讨论的模。一般环不预设交换，左理想、右理想和双边理想也须分别辨认。具体计算可以从上三角矩阵代数、截断多项式代数和除环矩阵环开始，观察同一种抽象结论如何给出不同的分解与层结构。

% !TeX root = ../../Book2.tex
% 纲要标题：## 第11章　代数、反环与双中心化子
% 纲要位置：计划纲要.md，第 1406 行。

\chapter{代数、反环与双中心化子}
\label{b2:ch:11}
\BookChapterNumber{b2:ch:11}

\begin{chapterabstract}
代数同时具有线性结构与乘法，表示则把乘法实现为线性算子的复合。本章通过结构常数、群代数与矩阵代数说明这两种结构如何相容，再用反环把右作用写成左作用，用包络代数把双模写成左模。幂等元给出正则模的分解和角环中的乘法；中心化子由与给定作用交换的算子组成。通过正则作用和张量因子的显式计算，建立双中心化子的基本例子，并辨明其何时不能恢复原作用代数。
\end{chapterabstract}

\begin{learninggoals}
\begin{itemize}
\item 根据结构常数检查结合律和单位元，计算换基后的乘法，区分向量空间同构与代数同构。
\item 在群表示、群代数模、双模及包络代数模之间转换，并核对反环与复合的次序。
\item 从正交幂等元构造模分解和矩阵分块，解释中心性何时把模分解提升为代数乘积分解。
\item 计算给定作用的中心化子及双中心化子，区分交换作用、互为中心化子与不可约性。
\end{itemize}
\end{learninggoals}

前两编研究了模的生成、关系与分解，现在要进一步问：一个代数可以怎样作用在这些模上，哪些模自同态与作用交换，又能否从这些自同态恢复原代数？把矩阵代数 \(A=M_2(k)\) 看作自身的向量空间，左乘 \(L_a(x)=ax\) 与右乘 \(R_b(x)=xb\) 总是交换：\(a(xb)=(ax)b\)。不过右乘的复合顺序反过来，\(R_b\circ R_c=R_{cb}\)：先作用 \(R_c\)，再作用 \(R_b\)。若要把两侧作用同时写成代数作用，就须用反环记录右侧的乘法方向。

\ProseParagraphBreak

同一组矩阵还可以从另一面观察。幂等矩阵把正则模分成像与核，和所有左乘都交换的算子则组成右作用的中心化子。先在矩阵单位上逐项计算这些事实，再抽去坐标，才知道哪些由代数结构保证、哪些只是在这个特殊例子里成立。

\ProseParagraphBreak

本章使用\chapref{b2:ch:02}的模与自同态环、\chapref{b2:ch:06}的双模和张量积，以及\chapref{b2:ch:08}的结合代数与张量作用。底域始终写作 \(k\)，代数同态保持单位元；一般不假设 \(k\) 代数闭，也不默认代数交换或有限维。矩阵单位提供了可直接核验的计算方法；左乘、右乘与模自同态产生的投影需要分别辨认。

\ProseParagraphBreak
代数和表示的基本语言可参阅 Etingof 等人的《Introduction to Representation Theory》\cite[\S 2.2, pp. 8--9; \S 2.3, pp. 9--10]{EtingofEtAl2011RepresentationTheory}。该书在这一部分默认底域代数闭并排除零代数；这里允许任意底域和零代数。后面的半单性与稠密性理论将进一步解释双中心化子计算。

% !TeX root = ../../Book2.tex
% 纲要标题：### 11.1 域上的代数
% 纲要位置：计划纲要.md，第 1408 行。

\section{域上的代数}
\label{b2:sec:1101}

矩阵可以相加、数乘，也可以相乘；群元素的形式线性组合也具有这三种运算。两者都是代数：乘法与向量空间的加法、数乘相容。本章以域 \(k\) 为底域，除特别说明外，代数结合且含单位元，同态保持单位元；允许零代数，需要把底域嵌入代数或讨论非零简单对象时，再明确加上非零条件。

% 纲要要点（供目录核对）：
%
% * 中心标量、结合含幺代数及表示与模同态
% * 乘法映射的结构常数、结合方程及基变换
% * 群代数泛性质、增广与矩阵单位作用
%

\subsection{结合含幺代数与代数同态}
\label{b2:subsec:1101:01}

\begin{definition}\label{b2:def:field-algebra}
设 \(k\) 为域。若 \(k\) 向量空间 \(A\) 配有双线性乘法 \((a,b)\mapsto ab\) 和元素 \(1_A\)，且对任意 \(a,b,c\in A\) 满足
\[
(ab)c=a(bc),\qquad 1_Aa=a1_A=a.
\]
则称 \(A\) 为域 \(k\) 上的结合含幺代数。若 \(A,B\) 均为这样的代数，\(k\) 线性映射 \(f:A\to B\) 保持乘法与单位元，则称 \(f\) 为 \(k\) 代数同态。\(A\) 中与所有元素交换的元素组成其中心 \(Z(A)\)。
\end{definition}

这是\cref{b2:def:associative-graded-algebra}中交换底环代数的定义在底环为域 \(k\) 时的具体形式。双线性给出 \((\lambda1_A)a=\lambda a=a(\lambda1_A)\)，所以 \(\lambda\mapsto\lambda1_A\) 是到中心的保单位元环同态；反过来，这样的结构同态给出标量乘法及上述双线性。环结构与线性结构正是通过中心的标量元素相容的。若 \(A\ne0\)，这个同态单射，因为其核是域的真理想。以下直接把 \(\lambda\in k\) 视作 \(A\) 中的元素时，默认 \(A\ne0\)；零代数仍由结构同态解释标量作用。

\ProseParagraphBreak
例如 \(k[t]/(t^2)\) 的元素唯一写成 \(a+b\varepsilon\)，其中 \(\varepsilon^2=0\)。它作为向量空间与 \(k\times k\) 都是二维，但前者有非零幂零元，后者没有。代数同构保持幂零性，所以它们不同构。区分代数应当寻找这样的乘法性质，不能只比较向量空间维数或某组基下的乘法表。

\begin{definition}\label{b2:def:algebra-representation}
设 \(k\) 为域、\(A\) 为含幺结合 \(k\) 代数、\(V\) 为 \(k\) 向量空间。保单位元的 \(k\) 代数同态 \(\rho:A\to\operatorname{End}_k(V)\) 称为 \(A\) 在 \(V\) 上的表示。若 \(V\ne0\)，且 \(V\) 没有非零真 \(A\) 稳定子空间，就称该表示为\term[bukeyue-biaoshi]{不可约表示}[irreducible representation]；此时相应左 \(A\) 模称为\term[jiandan-mo]{简单模}[simple module]。
\end{definition}

本章用不可约表示和简单模的定义辨认基本作用；它们与半单模的结构将在\chapref{b2:ch:13}讨论。

\begin{proposition}\label{b2:prop:algebra-modules-representations}
设 \(k\) 为域，\(A\) 为含幺结合 \(k\) 代数，\(V\) 为 \(k\) 向量空间。在 \(V\) 上指定幺性左 \(A\) 模结构，并使底域作用与已有标量作用相同，等价于指定保单位元的 \(k\) 代数同态
\[
\rho:A\longrightarrow\operatorname{End}_k(V).
\]
对两个这样的左 \(A\) 模 \(V,W\)，将相应表示记为 \(\rho_V,\rho_W\)。模同态正是满足 \(f\rho_V(a)=\rho_W(a)f\)（\(a\in A\)）的 \(k\) 线性映射 \(f:V\to W\)。
\end{proposition}
\begin{proof}
由模结构令 \(\rho(a)(v)=av\)。分配律给出对 \(a\) 和 \(v\) 的可加性；底域在中心保证每个 \(\rho(a)\) 为 \(k\) 线性。模的结合律与幺性给出 \(\rho(ab)=\rho(a)\rho(b)\)、\(\rho(1)=\operatorname{id}\)，而 \(\rho(\lambda a)=\lambda\cdot\rho(a)\)。反向用同一公式定义作用，这些等式逐项给出模公理。最后的相容等式在 \(v\) 处就是 \(f(av)=a\cdot f(v)\)。
\end{proof}

一个子空间是 \(A\) 子模，正是说它在每个 \(\rho(a)\) 下稳定。令上面的 \(W=V\)，模自同态的条件就成为
\[
\operatorname{End}_A(V)
=\{\,T\in\operatorname{End}_k(V)\mid T\rho(a)=\rho(a)T\;\text{对所有}\;a\in A\,\}.
\]
它由与全部作用算子交换的算子组成，\secref{b2:sec:1104}将把这种集合称为中心化子。表示的核为 \(V\) 的零化子；核为零恰好表示作用忠实。一个表示可以使许多不同代数元素产生同一个线性算子，定义本身并不排除这种情况。

\subsection{结构常数及基变换}
\label{b2:subsec:1101:02}

设 \(A\) 为域 \(k\) 上的有限维结合代数，以 \(e_1,\ldots,e_n\) 为一组 \(k\) 基。乘法的双线性与\cref{b2:thm:tensor-universal}给出 \(k\) 线性映射 \(\mu:A\otimes_kA\to A\)、\(\mu(a\otimes b)=ab\)。逐对展开乘积，写成
\[
\mu(e_i\otimes e_j)=e_ie_j=\sum_{r=1}^n c_{ij}^{r}e_r.
\]
所得标量 \(c_{ij}^{r}\in k\) 称为 \(A\) 在这组基下的\term[jiegou-changshu]{结构常数}[structure constants]。两个下标记录乘法的两个输入，上标记录输出坐标。它们决定全部乘法，因为两个一般元素的乘积可由双线性展开。

\begin{proposition}\label{b2:prop:structure-constants}
设 \(V\) 是域 \(k\) 上以 \(e_1,\ldots,e_n\) 为基的向量空间。指定 \(k\) 中的常数 \(c_{ij}^{r}\)（\(1\le i,j,r\le n\)），并规定 \(e_ie_j=\sum_r c_{ij}^{r}e_r\)，就指定了 \(V\) 上唯一的双线性乘法。该乘法结合，当且仅当对所有 \(i,j,\ell,s\in\{1,\ldots,n\}\)，
\begin{equation}\label{b2:eq:structure-associativity}
\sum_r c_{ij}^{r}c_{r\ell}^{s}=\sum_r c_{j\ell}^{r}c_{ir}^{s}.
\end{equation}
元素 \(u=\sum_i u_i e_i\) 是单位元，当且仅当
\[
\sum_i u_i c_{ij}^{s}=\delta_{js},\qquad
\sum_i u_i c_{ji}^{s}=\delta_{js}
\quad\text{对所有}\;j,s.
\]
\end{proposition}
\begin{proof}
令两个线性组合的乘积为其所有基乘积按系数展开的有限和。坐标唯一使其良定义，并直接得到双线性。比较 \((e_ie_j)e_\ell\) 与 \(e_i(e_je_\ell)\) 的第 \(s\) 个坐标，便得到\eqref{b2:eq:structure-associativity}。两种三变量表达都三线性，故在所有基向量组上相等即在任意输入上相等。类似地，\(ue_j=e_j\) 与 \(e_ju=e_j\) 的坐标恰是两条单位条件；线性性再推广到任意元素。
\end{proof}

因此，固定有限维向量空间上的结合乘法由满足\eqref{b2:eq:structure-associativity}的结构常数描述；这些是关于 \(c_{ij}^{r}\) 的二次多项式方程。乘法确定以后，寻找单位元只需解关于 \(u_i\) 的线性方程；若同时把 \(u_i\) 与 \(c_{ij}^{r}\) 当作未知量，单位条件则是关于这两组变量的双线性关系。

\ProseParagraphBreak
若换基为 \(f_a=\sum_i p_{ia}e_i\)，记 \(P^{-1}=(q_{dr})\)，两个输入向量的坐标都先用 \(P\) 转到旧基，乘积的输出坐标再用 \(P^{-1}\) 转回新基。因此新结构常数为
\begin{equation}\label{b2:eq:structure-basis-change}
(c')_{ab}^{d}=\sum_{i,j,r}q_{dr}p_{ia}p_{jb}c_{ij}^{r}.
\end{equation}
具体地，将 \(f_a,f_b\) 分别展开后，得到 \(f_af_b=\sum_{i,j,r}p_{ia}p_{jb}c_{ij}^{r}e_r\)，再代入 \(e_r=\sum_dq_{dr}f_d\) 即得此式。两个 \(P\) 和一个 \(P^{-1}\) 分别对应两个输入与一个输出，不能按只有一个输入的算子来作相似变换。

\ProseParagraphBreak
例如，在双数代数中以 \(1,y\) 为新基，令 \(y=b+a\varepsilon\)、\(a\ne0\)。直接计算得到 \(y^2=2by-b^2\)。乘法表中虽然出现了常数项和一次项，元素 \(y-b=a\varepsilon\) 仍非零且平方为零。因此结构常数随基改变，而存在非零幂零元这一乘法性质不变。

\ProseParagraphBreak
结构常数回答固定有限维向量空间上可以怎样定义乘法。群代数和矩阵代数则已经给定乘法；研究它们在向量空间上的作用，要检验乘法是否对应算子复合，并找出作用下保持不变的子空间。

\subsection{群代数和矩阵代数}
\label{b2:subsec:1101:03}

群乘法只规定群元素之间怎样相乘。以这些元素为基，允许有限线性组合，再将乘法双线性延拓，就把群乘法变成代数乘法，同时保留了原来的每个群元素。

\begin{definition}\label{b2:def:group-algebra}
设 \(k\) 为域、\(G\) 为群。以 \(G\) 的元素为基构造 \(k\) 向量空间，其元素为有限和 \(\sum_g a_g g\)（\(a_g\in k\)）。把群乘法按 \(k\) 双线性延伸：
\[
\left(\sum_g a_g g\right)\left(\sum_h b_h h\right)
=\sum_{g,h}a_gb_h(gh).
\]
所得含幺 \(k\) 代数称为群 \(G\) 的\term[qun-daishu]{群代数}[group algebra]，记为 \(k[G]\)；其单位为群单位元对应的基向量。
\end{definition}
\symidx{\(k[G]\)：群 \(G\) 的群代数}

两个支集有限，乘积仍是有限和；三重乘积的系数由同一个有限三重和给出，群乘法的结合律保证结果相同。因此该式定义结合代数。它交换当且仅当群交换，因为不同群元素是不同的基向量。若 \(G\) 有限，则 \(\dim_k k[G]=|G|\)。

\begin{proposition}\label{b2:prop:group-algebra-representations}
设 \(k\) 为域、\(G\) 为群、\(B\) 为非零结合含幺 \(k\) 代数，\(B^\times\) 为其单位群。限制到基向量 \(g\in G\) 给出双射
\[
\operatorname{Hom}_{k\text{-alg}}(k[G],B)
\longrightarrow\operatorname{Hom}_{\mathrm{Grp}}(G,B^\times).
\]
其逆把群同态 \(\rho:G\to B^\times\) 线性延拓为
\[
\widetilde\rho\left(\sum_g a_g g\right)=\sum_g a_g\cdot\rho(g).
\]
此对应对 \(G\) 逆变、对 \(B\) 协变地自然：对群同态 \(\alpha:H\to G\) 及非零结合含幺 \(k\) 代数之间的同态 \(\psi:B\to C\)，预复合 \(k[\alpha]:k[H]\to k[G]\) 对应预复合 \(\alpha\)，后复合 \(\psi\) 对应后复合其单位群限制 \(\psi^\times:B^\times\to C^\times\)。这里 \(k[\alpha]\) 是 \(h\mapsto\alpha(h)\) 的线性延拓。

对任意 \(k\) 向量空间 \(V\)（允许 \(V=0\)），群表示 \(G\to\operatorname{GL}(V)\) 与 \(V\) 上底域作用为原标量作用的幺性左 \(k[G]\) 模结构一一对应；两个表示之间的等变 \(k\) 线性映射对应 \(k[G]\) 模同态。
\end{proposition}
\begin{proof}
对代数同态 \(f:k[G]\to B\)，群基向量在群代数中满足 \(gg^{-1}=g^{-1}g=1\)，故 \(f(g)f(g^{-1})=f(g^{-1})f(g)=1_B\)。所以 \(f(g)\) 确在 \(B^\times\) 中，且保持群乘法。反向给定 \(\rho:G\to B^\times\)，由于 \(G\) 是一组基，其线性延拓 \(\widetilde\rho\) 唯一。由有限和展开，
\[
\begin{aligned}
\widetilde\rho\left(\left(\sum_g a_g g\right)\left(\sum_h b_h h\right)\right)
&=\sum_{g,h}a_gb_h\cdot\rho(gh)\\
&=\left(\sum_g a_g\cdot\rho(g)\right)\left(\sum_h b_h\cdot\rho(h)\right).
\end{aligned}
\]
因此它保持乘法；又 \(\rho(1_G)=1_B\)，故它是保单位元的代数同态。限制与线性延拓互逆，因为线性映射在一组基上的取值决定整个映射。

群同态 \(\alpha\) 的线性延拓同样保持乘法与单位元，因而 \(k[\alpha]\) 是代数同态；\(\psi\) 把单位及其逆仍送到互逆元素，所以单位群限制 \(\psi^\times\) 定义良好。对每个 \(h\in H\)，复合的取值都为 \(\psi(\rho(\alpha(h)))\)。由基上的唯一性，
\[
\widetilde{\psi^\times\circ\rho\circ\alpha}
=\psi\circ\widetilde\rho\circ k[\alpha],
\]
这就证明了两个变量的自然性。

当 \(V\ne0\) 时，取 \(B=\operatorname{End}_k(V)\)，其单位群就是 \(\operatorname{GL}(V)\)，再由\cref{b2:prop:algebra-modules-representations}得到表示与模结构的对应。明确地，作用为 \((\sum_g a_g g)v=\sum_g a_g\cdot\rho(g)v\)；反向限制到 \(g\) 的作用算子，其逆为 \(g^{-1}\) 的作用，因为 \(g\cdot(g^{-1}\cdot v)=v=g^{-1}\cdot(g\cdot v)\)。若 \(V=0\)，\(\operatorname{GL}(0)\) 是单元素群，且零空间只有一个幺性左模结构，对应仍成立。对两个空间之间的 \(k\) 线性映射，与每个 \(g\) 的作用相容等价于与它们的任意有限线性组合的作用相容，故等变映射与模同态也一致。
\end{proof}

把每个群元素都送到 \(1\in k^\times\)，就忘掉了群元素之间的区别。对群 \(G\) 和域 \(k\)，由上述泛性质得到的代数同态 \(\epsilon:k[G]\to k\)、\(\sum_g a_g g\mapsto\sum_g a_g\) 只保留系数总和，称为\term[zengguang]{增广}[augmentation]。相应一维模上每个群元素都恒等地作用，正是平凡表示。其核由 \(g-1\)（\(g\in G\)）线性生成，因为系数和为零的元素可写成 \(\sum_g a_g(g-1)\)。

\ProseParagraphBreak
矩阵代数 \(M_n(k)\)、\(n\ge1\)，以矩阵单位 \(E_{ij}\) 为基，乘法满足
\[
E_{ij}E_{rs}=\delta_{jr}E_{is},\qquad 1=\sum_iE_{ii}.
\]
\symidx{\(M_n(k)\)：域上的 \(n\) 阶矩阵代数}
它在列向量空间 \(k^n\) 上自然作用。矩阵单位 \(E_{ij}\) 读取向量的第 \(j\) 个坐标，把这个值写入第 \(i\) 个坐标，其余坐标置零。任何非零稳定子空间若含非零向量 \(v\)，取其一个非零坐标 \(v_j\)，则 \(E_{ij}v=v_j e_i\)；再乘 \(v_j^{-1}\)，便知它含有每个标准基向量，因而为全空间。这是不可约作用的一个基本例子。

% !TeX root = ../../Book2.tex
% 纲要标题：### 11.2 反环与包络代数
% 纲要位置：计划纲要.md，第 1414 行。

\section{反环与包络代数}
\label{b2:sec:1102}

左乘与右乘都可以表示为线性算子，但右乘算子的复合会颠倒乘法次序。反环记录这个次序，包络代数则把两个相容的作用合并成一个左作用。以下仍固定底域 \(k\)；所说的代数双模，其左右两侧的 \(k\) 标量作用须一致。

% 纲要要点（供目录核对）：
%
% * 反环与右作用的复合次序
% * 双模作为包络代数上的模，双边理想作为子模
% * 正则模自同态、矩阵包络作用及多项式商模型
%

\subsection{反环与左右模转换}
\label{b2:subsec:1102:01}

设 \(k\) 为域、\(A\) 为结合含幺 \(k\) 代数。取与 \(A\) 相同的向量空间和单位元，并规定乘法 \(a\cdot_{\mathrm{op}}b=ba\)，所得 \(k\) 代数称为 \(A\) 的反代数 \(A^{\mathrm{op}}\)。它就是\subsecref{b2:subsec:0201:01}所定义的反环，带上原有的 \(k\) 结构。将右 \(A\) 模 \(M\) 改写为左 \(A^{\mathrm{op}}\) 模时，规定 \(a^{\mathrm{op}}m=ma\)；原右作用的结合律保证这个左作用结合。

\ProseParagraphBreak
若 \(D_a(m)=ma\)，先右乘 \(b\)、再右乘 \(a\) 给出
\[
m\longmapsto mb\longmapsto(mb)a=m(ba),
\qquad D_a\circ D_b=D_{ba}.
\]
原来的右作用把乘法次序反转，改用反代数的乘法后，\(a^{\mathrm{op}}\mapsto D_a\) 就是通常的代数同态。例如 \(A=M_2(k)\) 时，先右乘 \(E_{12}\)、再右乘 \(E_{21}\)，得到 \(D_{E_{21}}\circ D_{E_{12}}=D_{E_{11}}\)；反向复合则为 \(D_{E_{12}}\circ D_{E_{21}}=D_{E_{22}}\)。

\subsection{双模作为包络代数上的模}
\label{b2:subsec:1102:02}

对两个 \(k\) 代数 \(A,B\)，\(A\otimes_kB\) 的乘法由
\[
(a\otimes b)(c\otimes d)=ac\otimes bd
\]
双线性延伸。四变量公式分别线性，故对两次张量关系良定义；在纯张量上，结合律由 \(A,B\) 的结合律得到，单位为 \(1_A\otimes1_B\)。这里两个代数本身都不必交换，只用到底域标量在各自中心。

\ProseParagraphBreak
对一个双模，我们希望用一个标量同时记录左乘 \(a\) 与右乘 \(b\)，使它作用为 \(m\mapsto amb\)。若先施加左乘 \(c\)、右乘 \(d\)，再施加左乘 \(a\)、右乘 \(b\)，连续作用为
\[
m\longmapsto cmd\longmapsto acmdb.
\]
左侧合成为 \(ac\)，右侧却合成为 \(db\)。因此记录右作用的因子须使用反代数，因为 \(b^{\mathrm{op}}d^{\mathrm{op}}=(db)^{\mathrm{op}}\)；张量积再把对这两个参数的双线性作用化为一个线性作用。

\begin{definition}\label{b2:def:enveloping-algebra}
设 \(k\) 为域、\(A\) 为结合含幺 \(k\) 代数。\(A\) 的\term[baoluo-daishu]{包络代数}[enveloping algebra]定义为
\[
A^e=A\otimes_k A^{\mathrm{op}}.
\]
若 \(k\) 向量空间 \(M\) 上有幺性的左、右 \(A\) 作用，两者对 \(k\) 双线性、满足 \((am)b=a(mb)\)，且两侧的底域标量作用都等于原有的 \(k\) 作用，则称 \(M\) 为 \(A\) 代数双模 \({}_AM_A\)。
\end{definition}
\symidx{\(A^e=A\otimes_k A^{\mathrm{op}}\)：代数的包络代数}

两侧标量一致的要求对应张量平衡关系 \((\lambda1_A)\otimes1=1\otimes(\lambda1_A)^{\mathrm{op}}\)。这两个张量是同一个元素，作用后分别得到 \((\lambda1_A)m\) 与 \(m(\lambda1_A)\)，所以它们必须相同，并等于原来的 \(\lambda m\)；仅有左右作用交换还不够。

\ProseParagraphBreak
这里的包络代数把代数的左、右作用合在一起；Lie 代数的泛包络代数 \(U(\mathfrak g)\) 则由 Lie 括号的另一种泛性质定义。

\begin{theorem}\label{b2:thm:enveloping-bimodules}
设 \(k\) 为域、\(A\) 为结合含幺 \(k\) 代数、\(A^e=A\otimes_kA^{\mathrm{op}}\)，\(M\) 为 \(k\) 向量空间。在 \(M\) 上指定与 \(k\) 标量作用相容的幺性 \(A\) 代数双模结构，等价于指定与该标量作用相容的幺性左 \(A^e\) 模结构。相应作用为
\begin{equation}\label{b2:eq:enveloping-action}
(a\otimes b^{\mathrm{op}})m=amb.
\end{equation}
这两种结构的同态也完全相同。
\end{theorem}
\begin{proof}
从双模出发，\((a,b^{\mathrm{op}})\mapsto(m\mapsto amb)\) 是到 \(\operatorname{End}_k(M)\) 的双线性映射，因此诱导 \(A^e\) 上的线性作用。连续作用两个纯张量给出
\[
(a\otimes b^{\mathrm{op}})\bigl((c\otimes d^{\mathrm{op}})m\bigr)
=acmdb.
\]
反代数中的乘法次序给出
\[
(a\otimes b^{\mathrm{op}})(c\otimes d^{\mathrm{op}})
=ac\otimes(db)^{\mathrm{op}},
\]
其作用也得到 \(acmdb\)。纯张量生成性保证乘法相容，\(1\otimes1\) 则作用为恒等。

反过来，令 \(am=(a\otimes1)m\)、\(mb=(1\otimes b^{\mathrm{op}})m\)。第一项保持乘法，第二项因反环而给出右结合律；两种纯张量 \(a\otimes1\) 与 \(1\otimes b^{\mathrm{op}}\) 交换，所以两侧作用相容。张量平衡关系给出
\[
(\lambda1_A)\otimes1=\lambda(1\otimes1)
=1\otimes(\lambda1_A)^{\mathrm{op}},
\]
而 \(1\otimes1\) 作用为恒等，故两侧的标量作用都等于 \(\lambda m\)。两次构造恢复原作用。同态同时与两侧作用相容，当且仅当与它们的复合及线性组合相容，亦即 \(A^e\) 线性。
\end{proof}

正则双模 \(A\) 本身由左右乘法给出。它的 \(A^e\) 子模恰是 \(A\) 的双边理想：对 \(a\otimes1\) 与 \(1\otimes b^{\mathrm{op}}\) 的作用稳定，分别就是对左乘与右乘稳定；反过来，这两种稳定性又保证对所有张量的作用稳定。因此研究 \(A\) 的双边理想可以转为研究正则 \(A^e\) 模 \(A\) 的子模。当 \(A\ne0\) 时，没有非零真双边理想，也就等价于这个 \(A^e\) 模没有非零真子模。

\ProseParagraphBreak
例如 \(A=k[t]\)，则 \(A^e\cong k[x,y]\)：\(x\) 记录左乘 \(t\)，\(y\) 记录右乘 \(t\)，包络代数把这两个作用分别记录。在正则双模上，这两个作用相同，因此 \(x-y\) 零化整个模。更一般的双模上，它们只须交换，不必相同。

\subsection{自同态代数与复合作用}
\label{b2:subsec:1102:03}

\begin{proposition}\label{b2:prop:regular-bimodule-endomorphism}
设 \(k\) 为域、\(A\) 为结合含幺 \(k\) 代数。对左正则模 \({}_AA\)，有 \(k\) 代数同构
\[
A^{\mathrm{op}}\longrightarrow\operatorname{End}_A({}_AA),\qquad
b^{\mathrm{op}}\longmapsto R_b,
\quad R_b(x)=xb.
\]
对正则双模，还有 \(\operatorname{End}_{A^e}(A)\cong Z(A)\)。
\end{proposition}
\begin{proof}
\cref{b2:prop:regular-hom-endomorphism}的公式给出每个左模自同态恰为 \(x\mapsto xb\)，其中 \(b=f(1)\) 唯一；这里它自动 \(k\) 线性。又 \(R_b\circ R_c=R_{cb}\)，正好给出反环同构。这样的映射同时保持右乘，当且仅当 \(ab=ba\) 对每个 \(a\in A\) 成立：把右线性等式 \(f(xa)=f(x)a\) 先取 \(x=1\) 得必要性，任意 \(x\) 的等式由同一交换性得到。因此双模自同态恰对应中心元素。中心交换，复合仍对应其通常乘法。
\end{proof}

对任意左 \(A\) 模 \(M\)，记 \(E=\operatorname{End}_A(M)\)。其中的映射与所有左标量作用交换。若要把这些自同态再作为 \(M\) 的右标量，系数代数须取 \(E^{\mathrm{op}}\)：定义 \(m\cdot f=f(m)\)，在反代数中有 \(f\cdot_{\mathrm{op}}g=g\circ f\)，因而
\[
m\cdot(f\cdot_{\mathrm{op}}g)
=g\bigl(f(m)\bigr)=(m\cdot f)\cdot g.
\]
这给出右作用的结合律，恒等映射给出单位作用，而 \(f(am)=a\cdot f(m)\) 保证它与左 \(A\) 作用相容。\chapref{b2:ch:14}的稠密性定理将采用这个右除环约定。

\begin{proposition}\label{b2:prop:regular-enveloping-models}
设 \(k\) 为域，\(n\ge1\)，对以下各 \(k\) 代数 \(C\) 记 \(C^e=C\otimes_k C^{\mathrm{op}}\)。对 \(A=M_n(k)\)，正则双模的作用给出含幺 \(k\) 代数同构
\[
\Phi:A^e\longrightarrow\operatorname{End}_k(A),
\qquad a\otimes b^{\mathrm{op}}\longmapsto(x\mapsto axb).
\]
对 \(B=k[t]\)，通过 \(x\mapsto t\otimes1\)、\(y\mapsto1\otimes t^{\mathrm{op}}\) 识别 \(B^e\cong k[x,y]\) 后，正则双模 \(B\) 作为左 \(B^e\) 模有同构
\[
k[x,y]/(x-y)\longrightarrow B,
\qquad F+(x-y)\longmapsto F(t,t).
\]
\end{proposition}
\begin{proof}
由\cref{b2:thm:enveloping-bimodules}应用于正则双模 \(A\)，\(\Phi\) 是含幺代数同态。设 \(E_{ij}\) 为 \(M_n(k)\) 的矩阵单位，对 \(1\le i,j,\ell,m,r,s\le n\)，有
\[
\Phi(E_{ij}\otimes E_{\ell m}^{\mathrm{op}})(E_{rs})
=E_{ij}E_{rs}E_{\ell m}
=\delta_{jr}\delta_{s\ell}E_{im}.
\]
因此这个算子把一个指定的基向量 \(E_{j\ell}\) 送到 \(E_{im}\)，把其余基向量送到零。让四个指标遍历，所得算子正是 \(\operatorname{End}_k(A)\) 的全部矩阵单位，组成一组基。来源中的 \(E_{ij}\otimes E_{\ell m}^{\mathrm{op}}\) 也是一组张量基，所以 \(\Phi\) 将一组基双射地送到另一组基，既单射又满射。

由于 \(B\) 交换，\(B^{\mathrm{op}}=B\)。张量基 \(t^i\otimes t^j\) 与单项式基 \(x^iy^j\) 一一对应，且乘法相符，所以得到所述 \(B^e\cong k[x,y]\)。在正则双模上，\(x,y\) 都作用为乘 \(t\)。于是取值映射
\[
\theta:k[x,y]\longrightarrow B,\qquad F\longmapsto F(t,t)
\]
是满射，也是左 \(k[x,y]\) 模同态，其中目标上的作用为 \(F\cdot p(t)=F(t,t)p(t)\)。显然 \(x-y\in\ker\theta\)。反向把 \(F(x,y)\) 看作 \(k[y][x]\) 中的多项式，除以首一多项式 \(x-y\)，可唯一写成
\[
F(x,y)=(x-y)Q(x,y)+R(y).
\]
若 \(\theta(F)=0\)，就有 \(R(t)=0\) 于多项式环 \(k[t]\)，因此 \(R\) 是零多项式，\(F\in(x-y)\)。这证明 \(\ker\theta=(x-y)\)，由\cref{b2:thm:module-first-isomorphism}得到所列模同构。
\end{proof}

% !TeX root = ../../Book2.tex
% 纲要标题：### 11.3 幂等元与分解
% 纲要位置：计划纲要.md，第 1420 行。

\section{幂等元与分解}
\label{b2:sec:1103}

一个投影满足连续作用两次与作用一次相同，即 \(p^2=p\)。在含幺代数 \(A\) 内，若 \(e^2=e\)，右乘映射 \(R_e(a)=ae\) 就满足 \(R_e^2=R_e\)，并将正则左模投到左理想 \(Ae\)。环中的幂等元由此产生模的分解，但未必分离整个代数的乘法；区别来自这个幂等元是否位于中心。

% 纲要要点（供目录核对）：
%
% * 正交幂等元的模分解与中心幂等元的代数分解
% * 角环、Peirce块及其双模结构
% * 右乘投影、角块与模同态复合
%

\subsection{正交幂等元及单位分解}
\label{b2:subsec:1103:01}

\begin{definition}\label{b2:def:orthogonal-idempotents}
设 \(k\) 为域、\(A\) 为含幺结合 \(k\) 代数。若 \(e\in A\) 满足 \(e^2=e\)，则称 \(e\) 为\term[mideng-yuan]{幂等元}[idempotent]。若两个幂等元 \(e,f\in A\) 满足 \(ef=fe=0\)，则称它们正交。若有限族 \(e_1,\ldots,e_r\in A\) 两两正交且 \(\sum_i e_i=1_A\)，则称 \(1_A=e_1+\cdots+e_r\) 为单位元的正交幂等元分解。
\end{definition}

\begin{proposition}\label{b2:prop:regular-orthogonal-decomposition}
设 \(k\) 为域、\(A\) 为含幺结合 \(k\) 代数，\(1_A=e_1+\cdots+e_r\) 是正交幂等元分解。它给出正则模的分解
\[
{}_AA=\bigoplus_i Ae_i,\qquad A_A=\bigoplus_i e_iA.
\]
若每个 \(e_i\) 还位于中心，则它们给出代数同构 \(A\cong\prod_i Ae_i\)，其中因子 \(Ae_i\) 的单位为 \(e_i\)。
\end{proposition}
\begin{proof}
对任意 \(a\)，有 \(a=\sum_i ae_i\)。若 \(\sum_i a_i e_i=0\)，右乘 \(e_j\) 只留下 \(a_j e_j\)，所以各项为零，左模分解成立。左乘的同样计算给出右模分解。

当各 \(e_i\) 在中心时，\(Ae_i\) 是双边理想和有单位的子代数，不同因子的乘积为零。映射 \(a\mapsto(ae_i)_i\) 保持加法、标量与乘法，并把单位送到各因子的单位族。其逆为求和映射，由刚证的直和分解互逆，得到代数同构。
\end{proof}

非中心幂等元只保证模分解。例如 \(A=M_2(k)=AE_{11}\oplus AE_{22}\) 是按列的左模分解，但 \(E_{11}\in AE_{11}\)、\(E_{12}\in AE_{22}\)，且 \(E_{11}E_{12}=E_{12}\ne0\)。这两个左模因子之间仍有非零乘积，不能成为直积代数中互相乘得零的两个因子。

\subsection{角环与矩阵分块}
\label{b2:subsec:1103:02}

设 \(k\) 为域、\(A\) 为含幺 \(k\) 代数、\(e\in A\) 为幂等元。集合 \(eAe\) 对加法、标量及乘法封闭，且以 \(e\) 为单位元；它称为 \(A\) 对应于 \(e\) 的\term[jiao-huan]{角环}[corner ring]，在当前情形也可称角代数。若 \(e\ne1_A\)，包含 \(eAe\subseteq A\) 不保持单位元，因而不是本书约定的含幺代数同态。

\ProseParagraphBreak
在 \(A=M_n(k)\) 中，若 \(e=\operatorname{diag}(I_s,0)\)、\(1\le s\le n\)，两侧乘 \(e\) 就只保留左上 \(s\times s\) 块，其余各块为零，所以 \(eAe\cong M_s(k)\)。对单位元的正交分解 \(1=\sum_i e_i\)，\(e_i a e_j\) 同样选出第 \((i,j)\) 块；下面的分解说明这种矩阵分块计算在一般代数中也成立。

\begin{proposition}[幂等元分块]\label{b2:prop:peirce-decomposition}
设 \(k\) 为域，\(A\) 为含幺结合 \(k\) 代数，\(1_A=e_1+\cdots+e_r\) 为正交幂等元分解，则
\begin{equation}\label{b2:eq:peirce-decomposition}
A=\bigoplus_{i,j}e_iAe_j
\end{equation}
作为向量空间成立。乘法满足 \((e_iAe_j)(e_sAe_t)=0\) 当 \(j\ne s\)；当 \(j=s\) 时，乘积属于 \(e_iAe_t\)。
每个块 \(e_iAe_j\) 通过 \(A\) 内的乘法成为 \((e_iAe_i,e_jAe_j)\) 双模，左右单位分别为 \(e_i\) 与 \(e_j\)。
\end{proposition}
\begin{proof}
把 \(a=1a1\) 展开，得到 \(a=\sum_{i,j}e_i a e_j\)。若这类分量之和为零，左乘 \(e_s\)、右乘 \(e_t\) 就只留下第 \((s,t)\) 个分量，所以分解唯一。两分量相乘时中间出现 \(e_je_s\)，它在不同指标时为零，相同指标时为 \(e_j\)，给出所述乘法规则。角环 \(e_iAe_i\) 左乘及 \(e_jAe_j\) 右乘都保持 \(e_iAe_j\)，并由结合律满足左右作用的相容性；对其中的 \(x\)，有 \(e_ix=x=xe_j\)，所以这些作用幺性。
\end{proof}

这种分块称为\term[Peirce-fenjie]{Peirce 分解}[Peirce decomposition]。\footnote{本杰明·皮尔斯（Benjamin Peirce，1809-1880），美国数学家、天文学家，在《Linear Associative Algebra》中以幂等元研究结合代数的结构。}因此可以把 \(a\) 记作分块矩阵 \((e_i a e_j)\)。乘积第 \((i,j)\) 块为 \(\sum_s(e_i a e_s)(e_s b e_j)\)，正是矩阵乘法式。不同块可能是不同的向量空间，不能未经论证就把它们全当作同一个系数环；非对角块作为双模连接两端的角环。

\ProseParagraphBreak
在 \(M_n(k)\) 中取 \(e_i=E_{ii}\)，各角环为 \(kE_{ii}\)，各块为 \(kE_{ij}\)，于是恢复通常的矩阵条目。若改取上三角代数 \(T_2(k)\)，右上块 \(e_1Ae_2=kE_{12}\) 非零，左下块 \(e_2Ae_1=0\)。两个角环都同构于 \(k\)，但整个代数不是 \(k\times k\)：非零右上块仍参与乘法。

\subsection{模分解与环中幂等元}
\label{b2:subsec:1103:03}

前面模论中已经见过投影的像与核分解；在\cref{b2:prop:finite-projective-summand}中，有限生成投射模还由有限秩自由模上的幂等自同态刻画。这里重新写出任意模的投影分解，是为了与环中幂等元产生的右乘投影对应起来。

\begin{proposition}\label{b2:prop:idempotent-module-decomposition}
对任意含幺环 \(R\) 和幺性左 \(R\) 模 \(M\)，满足 \(p^2=p\) 的 \(p\in\operatorname{End}_R(M)\) 给出
\[
M=\operatorname{im}p\oplus\ker p.
\]
反过来，每个指定的直和分解都给出投影幂等元。一组有限两两正交且和为恒等的投影，对应有限个子模的内直和分解。
\end{proposition}
\begin{proof}
写 \(m=p(m)+\bigl(m-p(m)\bigr)\)，第二项在核中；若 \(x=p(y)\) 又满足 \(p(x)=0\)，则 \(x=p^2(y)=p(x)=0\)，所以和为直和。反向取沿其余分量投到指定分量的映射，便满足平方为自身。有限族时，和为恒等保证生成，两两正交保证从一个分量到另一个的投影为零，从而分解唯一。
\end{proof}

在正则左模上，\cref{b2:prop:regular-bimodule-endomorphism}说明模自同态是右乘，所以幂等元 \(e\) 给出的投影为 \(a\mapsto ae\)，像为 \(Ae\)，核为 \(A(1-e)\)。左乘 \(L_e(a)=ea\) 虽也满足 \(L_e^2=L_e\)，却未必是左模自同态：左线性要求 \(e(ra)=r(ea)\)，取 \(a=1\) 就得到 \(er=re\) 对所有 \(r\in A\) 成立。反过来，\(e\) 在中心时此条件确实成立。取 \(A=M_2(k)\)、\(e=E_{11}\)、\(r=E_{21}\)，则
\[
L_e(r)=E_{11}E_{21}=0,
\qquad rL_e(1)=E_{21}E_{11}=E_{21}\ne0.
\]
因此左正则模中的投影须用右乘；不能仅凭算子平方等于自身，就把它当作左模投影。

\begin{proposition}\label{b2:prop:corner-endomorphism}
设 \(A\) 为含幺环，\(e,f\in A\) 为幂等元。取值映射给出交换群同构
\[
\operatorname{Hom}_A(Ae,Af)\longrightarrow eAf,
\qquad h\longmapsto h(e).
\]
它的逆把 \(b\in eAf\) 送到 \(h_b:Ae\rightarrow Af\)、\(h_b(ae)=ab\)。若 \(g\in A\) 也是幂等元，\(b\in eAf\)、\(c\in fAg\)，则
\[
h_c\circ h_b=h_{bc}:Ae\longrightarrow Ag.
\]
特别地，有含幺环同构
\[
(eAe)^{\mathrm{op}}\longrightarrow\operatorname{End}_A(Ae),
\qquad b^{\mathrm{op}}\longmapsto h_b.
\]
若 \(k\) 为域、\(A\) 为含幺结合 \(k\) 代数，则上述取值同构为 \(k\) 线性同构，上述自同态环同构为含幺 \(k\) 代数同构。
\end{proposition}
\begin{proof}
若 \(h:Ae\to Af\) 左线性，因为 \(h(e)\in Af\)，可写成 \(h(e)=a_0f\)。又 \(h(e)=h(e^2)=e h(e)\)，所以
\[
h(e)=ea_0f\in eAf.
\]
对任意 \(a\in A\)，左线性又迫使 \(h(ae)=a\cdot h(e)\)，故 \(e\) 的像决定整个同态。

反过来，取 \(b\in eAf\)。若 \(ae=a'e\)，则 \((a-a')e=0\)，而 \(eb=b\)，所以 \((a-a')b=(a-a')eb=0\)，证明 \(h_b(ae)=ab\) 与代表元选择无关。由于 \(bf=b\)，有 \(ab\in Af\)；结合律给出左 \(A\) 线性。取值映射与 \(b\mapsto h_b\) 都保持加法，而 \(h_b(e)=eb=b\)，所以这两个加法群映射互为逆。

对 \(b\in eAf\)、\(c\in fAg\)，有 \(bc\in eAg\)，并且
\[
(h_c\circ h_b)(ae)=h_c(ab)=abc=h_{bc}(ae).
\]
这里 \(ab=(ab)f\) 作为 \(Af\) 中的元素，正可代入 \(h_c\) 的定义。取 \(f=g=e\)，便有 \(h_b\circ h_c=h_{cb}\)，所以参数乘法相对于通常的自同态复合恰好反转，得到 \((eAe)^{\mathrm{op}}\) 到自同态环的同构；\(h_e(ae)=ae\)，故角环的单位 \(e\) 对应恒等映射。当 \(A\) 还是 \(k\) 代数时，取值与其逆都保持 \(k\) 标量，故这些对应分别是所述的 \(k\) 线性同构与 \(k\) 代数同构。
\end{proof}

这个取值对应说明，Peirce 块 \(eAf\) 同时记录从左理想 \(Ae\) 到 \(Af\) 的所有左模同态，块乘法 \((eAf)(fAg)\subseteq eAg\) 对应先施加 \(h_b\)、再施加 \(h_c\) 的复合。取 \(f=e\)，便把正则模的分解细化到某个直和因子内部。后续判别 \(Ae\) 是否还能分解时，考察的是它的自同态环中有无非平凡幂等元，也就是角环 \(eAe\) 中的同一问题。

% !TeX root = ../../Book2.tex
% 纲要标题：### 11.4 中心化子
% 纲要位置：计划纲要.md，第 1426 行。

\section{中心化子}
\label{b2:sec:1104}

一个作用确定以后，可以再问哪些线性算子与它交换。这些算子组成中心化子；再取一次中心化子，则至少能找回原来的作用。但是否恰好找回原代数，是一个需要证明的问题。本节先给出完全可算的例子，再说明它与张量表示的联系。

% 纲要要点（供目录核对）：
%
% * 模自同态与中心化子，双中心化子的包含与幂等性
% * 正则左右作用及有限维张量因子的中心化子
% * 三角反例与 Schur–Weyl 对偶的动机及证明归属
%
% ---
%

\subsection{交换作用与中心化子}
\label{b2:subsec:1104:01}

设 \(k\) 为域、\(V\) 为 \(k\) 向量空间，\(\mathcal A\subseteq\operatorname{End}_k(V)\) 为含恒等算子的子代数。沿用\subsecref{b2:subsec:0804:03}的中心化子定义，记
\[
\mathcal A'=\bigl\{\,T\in\operatorname{End}_k(V):Ta=aT\text{ 对所有 }a\in\mathcal A\,\bigr\}.
\]
再记 \(\mathcal A''=(\mathcal A')'\)，称为\term[shuang-zhongxinhua-zi]{双中心化子}[double centralizer]。这些撇号在本节表示中心化子，不表示导数或对偶。
\symidx{\(\mathcal A'\)、\(\mathcal A''\)：给定作用的中心化子及双中心化子}

\begin{proposition}\label{b2:prop:commutant-endomorphism}
设 \(k\) 为域、\(A\) 为结合含幺 \(k\) 代数、\(V\) 为 \(k\) 向量空间，\(\rho:A\to\operatorname{End}_k(V)\) 为含幺表示。对 \(\mathcal C\subseteq\operatorname{End}_k(V)\)，用 \(\mathcal C'\) 表示 \(\operatorname{End}_k(V)\) 中与 \(\mathcal C\) 全部元素交换的算子所成的中心化子，则
\[
\rho(A)'=\operatorname{End}_A(V),\qquad \rho(A)\subseteq\rho(A)''.
\]
若 \(\mathcal A\subseteq\mathcal B\subseteq\operatorname{End}_k(V)\)，则
\[
\mathcal B'\subseteq\mathcal A',\qquad
\mathcal A''\subseteq\mathcal B''.
\]
对任意算子族 \(\mathcal C\subseteq\operatorname{End}_k(V)\)，还有
\[
\mathcal C\subseteq\mathcal C'',\qquad
\mathcal C'''=\mathcal C',\qquad
\mathcal C''''=\mathcal C''.
\]
\end{proposition}
\begin{proof}
\cref{b2:prop:algebra-modules-representations}中模同态的相容公式在定义域、陪域均为 \(V\) 时，就是 \(T\rho(a)=\rho(a)T\)。在向量 \(v\) 上，这个等式写成 \(T(av)=a\cdot T(v)\)，故第一个等式成立。任意 \(c\in\mathcal C\) 都与 \(\mathcal C'\) 的全部元素交换，所以 \(\mathcal C\subseteq\mathcal C''\)，这也给出 \(\rho(A)\subseteq\rho(A)''\)。

若一个算子与 \(\mathcal B\) 的全部元素交换，它当然与子集 \(\mathcal A\) 的全部元素交换，故 \(\mathcal B'\subseteq\mathcal A'\)。再取一次中心化子并反转包含，得到 \(\mathcal A''\subseteq\mathcal B''\)。由 \(\mathcal C\subseteq\mathcal C''\) 得 \(\mathcal C'''\subseteq\mathcal C'\)；把扩大性质用于 \(\mathcal C'\)，又得 \(\mathcal C'\subseteq\mathcal C'''\)。于是 \(\mathcal C'''=\mathcal C'\)，两侧再取中心化子即得 \(\mathcal C''''=\mathcal C''\)。
\end{proof}

与更多算子交换会施加更多限制，所以中心化子反转包含；连续取两次后又保持包含。上述等式还说明，双中心化子扩大原算子族，并且再取双中心化子不会继续增大。这是它作为一种闭包运算的意义。

\ProseParagraphBreak
若两个代数在同一空间上的作用彼此交换，其中一个便落在另一个的中心化子中。只有证明反向包含，才能称它们互为中心化子。维数有限时可以把交换条件写成线性方程组；这通常比直接猜测所有相容映射更可靠。

\subsection{双中心化子的基本例子}
\label{b2:subsec:1104:02}

\begin{theorem}[正则作用的双中心化子]\label{b2:thm:regular-double-centralizer}
设 \(k\) 为域、\(A\) 为结合含幺 \(k\) 代数。在 \(\operatorname{End}_k(A)\) 中，令
\[
L(A)=\{L_a:x\mapsto ax\mid a\in A\},\qquad
R(A)=\{R_a:x\mapsto xa\mid a\in A\}.
\]
映射 \(a\mapsto L_a\) 与 \(a^{\mathrm{op}}\mapsto R_a\) 分别给出 \(k\) 代数同构 \(A\cong L(A)\) 与 \(A^{\mathrm{op}}\cong R(A)\)。两个作用代数互为中心化子，因此 \(L(A)''=L(A)\)，不要求 \(A\) 有限维。
\end{theorem}
\begin{proof}
复合满足 \(L_aL_b=L_{ab}\)、\(R_aR_b=R_{ba}\)，所以右乘的复合顺序对应反代数乘法。两种映射都线性、保单位元，并且算子在 \(1\) 处的值恢复 \(a\)，故它们单射；由像的定义，它们满射到相应作用代数。

\(T\) 与全部左乘交换，当且仅当它左 \(A\) 线性。正则左模由 \(1\) 生成；\cref{b2:prop:regular-bimodule-endomorphism}给出这样的 \(T\) 由 \(b=T(1)\) 唯一确定，且 \(T(x)=xb\)，故 \(L(A)'=R(A)\)。若 \(T\) 与全部右乘交换，则对任意 \(a\)，
\[
T(a)=T\bigl(R_a(1)\bigr)=R_a\bigl(T(1)\bigr)=T(1)a.
\]
整个算子仍由 \(T(1)\) 决定，这次它是左乘 \(T(1)\)。反向左乘与右乘由结合律交换，得到 \(R(A)'=L(A)\)。
\end{proof}

另一个重要例子来自两个向量空间的张量。令 \(U,W\) 为非零有限维空间。代数 \(\operatorname{End}_k(U)\) 作用于第一因子，\(\operatorname{End}_k(W)\) 作用于第二因子，两种作用交换。选定 \(U\) 的一组 \(m\) 元基后，\(U\otimes_kW\) 是 \(W\) 的 \(m\) 个副本的直和；第一因子的矩阵作用调配这些副本。与所有这些调配交换的算子，必须在每个副本上作同一个变换。

\begin{proposition}\label{b2:prop:tensor-factor-centralizers}
设 \(k\) 为域，\(U,W\) 为非零有限维 \(k\) 向量空间。在 \(\operatorname{End}_k(U\otimes_kW)\) 中，
\[
\operatorname{End}_k(U)\otimes k\operatorname{id}_W,
\qquad k\operatorname{id}_U\otimes\operatorname{End}_k(W)
\]
互为中心化子。此外，映射
\[
\begin{aligned}
&\operatorname{End}_k(U)\otimes_k\operatorname{End}_k(W)
\longrightarrow\operatorname{End}_k(U\otimes_kW),\\
&F\otimes G\longmapsto
\bigl(u\otimes w\mapsto F(u)\otimes G(w)\bigr)
\end{aligned}
\]
是 \(k\) 代数同构。
\end{proposition}
\begin{proof}
选 \(U\) 的基 \(u_1,\ldots,u_m\)。任意线性算子 \(T\) 可唯一写成
\[
T(u_j\otimes w)=\sum_i u_i\otimes T_{ij}(w),
\qquad T_{ij}\in\operatorname{End}_k(W).
\]
\(E_{aa}\otimes\operatorname{id}_W\) 是到第 \(a\) 个副本 \(ku_a\otimes W\) 的投影。若 \(T\) 与这些投影交换，它就保留各个副本，故 \(T_{ij}=0\) 当 \(i\ne j\)。而 \(E_{ab}\otimes\operatorname{id}_W\) 把第 \(b\) 个副本移到第 \(a\) 个副本，在其中保留 \(w\) 不变。与这个算子交换，便要求两个副本上的变换相同：在 \(u_b\otimes w\) 上比较，得到 \(T_{aa}(w)=T_{bb}(w)\)。因此 \(T=\operatorname{id}_U\otimes S\)，其中 \(S\) 为公共对角块。反向这样的算子确实交换，得到第一个中心化子。交换两因子的角色，得到另一个。

再选 \(W\) 的基 \(w_1,\ldots,w_\ell\)。张量 \(E_{ij}\otimes E_{ab}\) 将基向量 \(u_j\otimes w_b\) 送到 \(u_i\otimes w_a\)，将其余基向量送到零，所以这些张量恰对应全部矩阵单位。它们在两边均为基，故作用映射是线性同构；逐因子复合保证它保持乘法与单位。
\end{proof}

例如 \(k^n\oplus k^n\cong k^n\otimes_k k^2\)（\(n\ge1\)）时，\(\operatorname{diag}(B,B)\)、\(B\in M_n(k)\)，作用于第一因子。其中心化子为 \(\operatorname{id}_{k^n}\otimes\operatorname{End}_k(k^2)\)，允许在两个相同的 \(k^n\) 副本之间任意混合，而在各副本内部保留原来的矩阵作用。\cref{b2:ex:11-repeated-matrix-blocks}将这个结论写成具体的分块矩阵。

\ProseParagraphBreak
这里自同态代数的张量同构不能无条件推广到无限维空间。取 \(U,W\) 分别以 \(u_1,u_2,\ldots\) 与 \(w_1,w_2,\ldots\) 为基，定义
\[
Q(u_i\otimes w_j)=
\begin{cases}
u_i\otimes w_i,&i=j,\\
0,&i\ne j.
\end{cases}
\]
这在张量积的一组基上规定了像，故给出线性算子。若 \(Q=\sum_{r=1}^s A_r\otimes B_r\)，记 \(\alpha_i:U\to k\) 为第 \(i\) 个坐标泛函，对输入 \(u_i\otimes w\) 的输出取第 \(i\) 个 \(U\) 坐标，就得到
\[
P_i(w)=\sum_{r=1}^s\alpha_i(A_ru_i)B_r(w),
\]
其中 \(P_i\in\operatorname{End}_k(W)\) 保留 \(w_i\) 坐标，并把其余坐标归零。因此所有 \(P_i\) 都在有限维空间 \(\operatorname{span}_k\{B_1,\ldots,B_s\}\) 中。但这些投影线性无关：任意有限线性关系分别作用于对应的 \(w_i\)，就迫使每个系数为零。矛盾。所以这个 \(Q\) 不在张量作用映射的像中。

\ProseParagraphBreak
双中心化子等式不能推广到每个子代数。取 \(\mathcal A=T_2(k)\subseteq M_2(k)\)。与 \(E_{11}\) 交换迫使矩阵为对角矩阵，再与 \(E_{12}\) 交换迫使两个对角元相等。因此 \(\mathcal A'=kI_2\)，而 \(\mathcal A''=M_2(k)\supsetneq\mathcal A\)。原代数保留非零真子空间 \(ke_1\)，故作用不是不可约的。这就给出了从中心化子仅含标量反推不可约性的反例。

\subsection{Schur–Weyl 对偶的环论背景}
\label{b2:subsec:1104:03}

\chapref{b2:ch:08}已在 \(V^{\otimes d}\) 上构造一般线性群的对角作用与 \(S_d\) 的置换作用，并在\cref{b2:prop:commuting-tensor-actions}中证明它们交换。因而两种作用的线性生成代数互相包含于对方的中心化子。对于有限维复空间，\cref{b2:claim:schur-weyl-preview}给出它们互为中心化子的等式；证明见第五卷第7.3节。

\ProseParagraphBreak
这里的情形与\cref{b2:prop:tensor-factor-centralizers}不同。后者允许独立选取 \(F\in\operatorname{End}_k(U)\) 与 \(G\in\operatorname{End}_k(W)\)，作用为 \(u\otimes w\mapsto F(u)\otimes G(w)\)；Schur--Weyl 情形则让同一个 \(g\) 同时作用在所有因子上，得到 \(v_1\otimes\cdots\otimes v_d\mapsto gv_1\otimes\cdots\otimes gv_d\)，另一个作用改变因子的排列。独立因子上的中心化子计算没有计算出同一个 \(g\) 的对角作用的全部中心化子，因而不能代替 Schur--Weyl 对偶中反向包含的证明。

\ProseParagraphBreak
若只考察交换两个因子的算子 \(\tau\)，在特征不为二的域上，第八章的两个投影 \((\operatorname{id}\pm\tau)/2\) 把张量积分成对称与交替两部分。与 \(\tau\) 交换的算子恰好分别保持这两部分，因而可以在每一部分内部独立作用。\cref{b2:ex:11-tensor-flip-centralizer}在 \(V=k^2\) 上计算这一中心化子及其双中心化子。

\ProseParagraphBreak
\chapref{b2:ch:14}将研究不可约作用的中心化子及稠密性；有限维且底域代数闭时，不可约性会迫使作用代数等于整个自同态代数。上三角例子则表明，只知中心化子由标量组成还不足以作这个推断。双中心化子问题的答案取决于整个作用，而不只取决于原代数的抽象乘法表。

\ProseParagraphBreak
若另具备Hilbert空间和有界算子的基础，可参阅\appref{b2:app:D}的\cref{b2:thm:operator-bicommutant}。该结果在含幺且对伴随封闭的条件下，把双交换子与强、弱算子闭包联系起来；这里的纯代数中心化子计算并不自动包含那些拓扑条件。


\begin{chaptersummary}
有限维代数在选定一组有限基后，其结构常数记录全部乘法，结合律与单位条件成为有限组方程；换基会改变这些常数，却保留代数结构。把代数作用写成到自同态代数的同态，可以同时处理群表示、矩阵作用和一般模。

\ProseParagraphBreak
右作用对应反环，两个相容的左右作用对应包络代数的一个左作用，双边理想因而成为正则包络代数模的子模。正则左模的自同态来自右乘，正则双模的自同态则由中心决定。环中的正交幂等元给出正则模与角块分解；角块 \(eAf\) 记录 \(Ae\to Af\) 的模同态，块乘法记录这些同态的复合。只有幂等元还满足中心性时，这种模分解才进一步给出代数乘积。

\ProseParagraphBreak
中心化子描述与一个作用交换的全部算子，再取中心化子总包含原作用代数，却未必等于它。正则作用及有限维张量因子的独立作用具有双中心化子等式，上三角作用则说明标量中心化子本身不足以保证不可约。后续结构理论要补充的，正是这些具体计算背后的有限性与模分解条件。
\end{chaptersummary}

\BookExercises

\begin{exercise}\label{b2:ex:11-two-dimensional-constants}
设 \(k\) 为域，非零结合含幺 \(k\) 代数 \(A\) 有基 \(1,x\)，其中 \(x^2=\alpha x+\beta\)、\(\alpha,\beta\in k\)。令 \(y=ux+v\)，其中 \(u\in k^\times\)、\(v\in k\)。求基 \(1,y\) 下的乘法关系，并证明 \(\Delta=\alpha^2+4\beta\) 在这种换基下变为 \(u^2\Delta\)。说明等式本身是否要求特征不为二。
\end{exercise}
\begin{solution}
展开并消去 \(x=u^{-1}(y-v)\)，得到
\[
y^2=(u\alpha+2v)y+(u^2\beta-u\alpha v-v^2).
\]
新系数为 \(\alpha'=u\alpha+2v\)、\(\beta'=u^2\beta-u\alpha v-v^2\)，所以
\[
(\alpha')^2+4\beta'=u^2\alpha^2+4u^2\beta=u^2\Delta.
\]
全部计算都是整系数恒等式，不要求特征不为二。这个变换律说明换基只会把 \(\Delta\) 乘以一个非零平方，但分类还须证明相应的标准形：特征不为二时，取 \(z=x-\alpha/2\) 得 \(z^2=\Delta/4\)，再分别辨认 \(\Delta=0\)、非零平方及非平方三种情形；特征二时，\(\Delta=\alpha^2\) 甚至不记录 \(\beta\)。另有 \(A\cong k[t]/(t^2-\alpha t-\beta)\)：送 \(t\) 到 \(x\) 给出满同态，商中 \(1,t\) 与目标的 \(1,x\) 分别为基，所以它单射。
\end{solution}

\begin{exercise}\label{b2:ex:11-associative-table}
设 \(k\) 为域、\(\lambda,\mu\in k\)。在有基 \(1,x,y\) 的 \(k\) 向量空间中，令 \(1\) 为单位，并规定交换的基乘法
\(x^2=y\)、\(xy=yx=\lambda y\)、\(y^2=\mu y\)。求这一双线性乘法满足结合律的充要条件，并在满足条件时给出单生成元商代数描述。
\end{exercise}
\begin{solution}
由 \((x^2)y=x(xy)\)，必须有 \(\mu y=\lambda^2y\)，所以 \(\mu=\lambda^2\)。若该条件成立，候选的单生成元关系是 \(t^3=\lambda t^2\)。考察结合代数 \(B=k[t]/(t^3-\lambda t^2)\)，它以 \(1,t,t^2\) 为基，且 \(t\cdot t^2=\lambda t^2\)、\((t^2)^2=\lambda^2t^2\)。基对应 \(1\mapsto1,t\mapsto x,t^2\mapsto y\) 因而保留所有给定基乘法，使 \(B\) 与题给双线性乘法结构同构。后者由此结合，且得到 \(A\cong B\)。这种构造也核对了充分性，而不只检查某一个必要等式。
\end{solution}

\begin{exercise}\label{b2:ex:11-nilpotent-representation}
设 \(k\) 为域、\(A=k[t]/(t^3)\)。在 \(V=k^3\) 上分别令 \(t\) 的剩余类作用为三阶幂零 Jordan 块 \(J_3\)，以及 \(J_2\oplus0\)。证明这两者都给出 \(A\) 的表示，并求表示的核。
\end{exercise}
\begin{solution}
两个算子都满足 \(T^3=0\)，所以多项式代入 \(f\mapsto f(T)\) 经过商代数 \(A\)，且保持乘法与单位。对 \(J_3\)，若 \(aI+bJ_3+cJ_3^2=0\)，依次比较主对角线、第一条及第二条上对角线，得到 \(a=b=c=0\)，所以表示忠实。对 \(J_2\oplus0\)，有 \(T\ne0,T^2=0\)，且 \(I,T\) 线性无关，所以代入核正好为 \(k\bar t^2=(\bar t^2)\)。同一代数作用在相同维数空间上，仍可有不同的零化子。
\end{solution}

\begin{exercise}\label{b2:ex:11-group-algebra-universal}
设 \(k\) 为域，\(G\) 为群，\(N\triangleleft G\)。令 \(I\) 为 \(k[G]\) 中由 \(n-1\)、\(n\in N\) 生成的双边理想。证明商群映射诱导代数同构
\[
k[G]/I\cong k[G/N].
\]
说明在群代数中加入这些关系，为什么正好把同一 \(N\) 陪集中的基元素识别，而不会产生额外的线性关系。
\end{exercise}
\begin{solution}
由\cref{b2:prop:group-algebra-representations}，\(g\mapsto gN\) 唯一延拓为满代数同态 \(q:k[G]\to k[G/N]\)。每个 \(n-1\) 都在核中，故 \(I\subseteq\ker q\)。若有限和 \(z=\sum_g a_gg\) 在核中，由于各陪集是 \(k[G/N]\) 的基，每个陪集内的系数和都为零。对与 \(z\) 的支集相交的每个陪集 \(C\)，取代表元 \(r_C\)，便有
\[
\sum_{g\in C}a_gg
=\sum_{g\in C}a_g(g-r_C).
\]
这里所有和都只取原有限支集中的元素。对 \(g\in C=r_CN\)，写 \(g=r_Cn\)，则 \(g-r_C=r_C(n-1)\in I\)。逐陪集相加得 \(z\in I\)，所以 \(\ker q=I\)，商代数同构随之成立。陪集内系数和为零恰好描述了这些识别关系的全部线性后果。
\end{solution}

\begin{exercise}\label{b2:ex:11-augmentation-generators}
设 \(k\) 为域，群 \(G\) 由子集 \(X\) 生成，\(\epsilon:k[G]\to k\) 为增广同态。证明增广理想 \(\ker\epsilon\) 作为左理想已由 \(x-1\)、\(x\in X\) 生成，不需要把所有群元素逐个列入生成元。说明同样结论也适用于右理想生成。
\end{exercise}
\begin{solution}
令 \(L\) 为这些元素生成的左理想。恒等式
\[
gh-1=(g-1)+g(h-1),\qquad x^{-1}-1=-x^{-1}(x-1)
\]
说明若已知两个群元素减一属于 \(L\)，则其乘积减一也属于 \(L\)，并且每个生成元的逆减一也在 \(L\)。对生成字长度归纳，得到所有 \(g-1\in L\)。若 \(\sum_g a_gg\in\ker\epsilon\)，其系数和为零，故它等于 \(\sum_g a_g(g-1)\)，也属于 \(L\)。这证明 \(\ker\epsilon\subseteq L\)；反向包含由 \(\epsilon(x-1)=0\) 得到。右理想版本使用 \(gh-1=(g-1)h+(h-1)\) 与 \(x^{-1}-1=-(x-1)x^{-1}\)，其余归纳相同。
\end{solution}

\begin{exercise}\label{b2:ex:11-cyclic-two-units}
设 \(k\) 为域、\(C_2=\{1,g\}\)、\(g^2=1\)。对 \(a,b\in k\)，求 \(k[C_2]\) 中 \(a+bg\) 可逆的充要条件和逆元公式，允许 \(\operatorname{char}k=2\)。
\end{exercise}
\begin{solution}
左乘 \(a+bg\) 在基 \(1,g\) 下的矩阵为 \(\begin{pmatrix}a&b\\b&a\end{pmatrix}\)。若元素可逆，该矩阵可逆，故 \(a^2-b^2\ne0\)。反过来，若此数非零，直接相乘得到
\[
(a+bg)^{-1}=\frac{a-bg}{a^2-b^2}.
\]
这是双侧逆，所以条件充分。在特征二中，条件变为 \((a+b)^2\ne0\)，即 \(a+b\ne0\)，公式仍成立。可逆元不限于非零标量乘群元素，例如在 \(\mathbb Q[C_2]\) 中 \(2+g\) 也可逆。
\end{solution}

\begin{exercise}\label{b2:ex:11-matrix-opposite}
设 \(k\) 为域、\(n\ge1\)。证明转置给出代数同构 \(M_n(k)^{\mathrm{op}}\to M_n(k)\)，但在 \(n\ge2\) 时不是从 \(M_n(k)\) 到自身的代数同态。用两个矩阵单位写出具体失败等式。
\end{exercise}
\begin{solution}
来源反环中 \(A\cdot_{\mathrm{op}}B=BA\)，所以
\((A\cdot_{\mathrm{op}}B)^{\mathsf t}=A^{\mathsf t}B^{\mathsf t}\)。转置线性、保持单位且自身为逆，故是所述同构。在原乘法下取 \(A=E_{12},B=E_{21}\)，则 \((AB)^{\mathsf t}=E_{11}\)，而 \(A^{\mathsf t}B^{\mathsf t}=E_{21}E_{12}=E_{22}\)，两者不同。
\end{solution}

\begin{exercise}\label{b2:ex:11-matrix-enveloping}
设 \(k\) 为域，\(A=k\times k\)，\(e_1=(1,0)\)、\(e_2=(0,1)\)。计算正则双模作用 \(\Phi:A^e\to\operatorname{End}_k(A)\) 的像与核，其中 \(\Phi(a\otimes b)(x)=axb\)。说明左正则作用与右正则作用各自忠实，为什么仍不足以使整个包络代数的作用忠实。
\end{exercise}
\begin{solution}
因为 \(A\) 交换，\(A^{\mathrm{op}}=A\)，而 \(e_i\otimes e_j\)、\(1\le i,j\le2\) 构成 \(A^e\) 的基。由 \(e_ie_j=\delta_{ij}e_i\)，有
\[
\Phi(e_i\otimes e_j)(e_\ell)
=\delta_{i\ell}\delta_{j\ell}e_\ell.
\]
所以 \(e_1\otimes e_1\)、\(e_2\otimes e_2\) 给出两个坐标投影，非对角的两个张量作用为零。由此
\[
\operatorname{im}\Phi
=\left\{\begin{pmatrix}a&0\\0&b\end{pmatrix}:a,b\in k\right\},
\qquad
\ker\Phi=k(e_1\otimes e_2)\oplus k(e_2\otimes e_1).
\]
左乘或右乘 \(a\) 若为零，在 \(1\) 上求值便得 \(a=0\)，故两种单侧作用各自忠实。但包络代数把两侧分别记录，还含有上述非零张量，单侧的忠实性没有排除它们。与\cref{b2:prop:regular-enveloping-models}的矩阵代数相比，这里的左右乘有限和只能产生对角算子。
\end{solution}

\begin{exercise}\label{b2:ex:11-one-dimensional-bimodules}
设 \(k\) 为域，\(A=k[t]\)，\(\alpha,\beta,\gamma,\delta\in k\)。在一维空间 \(M_{\alpha,\beta}=k\) 上规定左乘 \(t\) 为乘 \(\alpha\)，右乘 \(t\) 为乘 \(\beta\)，两侧 \(k\) 标量都按通常方式作用。证明它总是 \(A\) 双模，并求
\(\operatorname{Hom}_{A^e}(M_{\alpha,\beta},M_{\gamma,\delta})\)。
\end{exercise}
\begin{solution}
左右作用分别为 \(f(t)m=f(\alpha)m\)、\(m\cdot g(t)=g(\beta)m\)。它们都保持多项式乘法、单位与底域标量，且两个标量作用交换，所以确为双模。任一线性映射是乘 \(c\in k\)，保持左、右作用分别要求 \(c\alpha=\gamma c\)、\(c\beta=\delta c\)。若两对参数相同，任意 \(c\) 都可行，Hom 为 \(k\)；只要有一项不同，域中消去非零差便得 \(c=0\)，Hom 为零。两个作用可以不同，这不违反双模的相容性。
\end{solution}

\begin{exercise}\label{b2:ex:11-dual-number-enveloping-matrices}
设 \(k\) 为域，\(A=k[\varepsilon]/(\varepsilon^2)\)，\(M=k^3\)，并以 \(E_{ij}\) 表示 \(M_3(k)\) 的矩阵单位。规定 \(\varepsilon\) 在 \(M\) 上的左作用为 \(E_{12}\)、右作用为 \(E_{13}\)，两侧标量均按通常方式作用。证明这些规定使 \(M\) 成为 \(A\) 双模。将 \(A^e\) 识别为 \(k[x,y]/(x^2,y^2)\)，写出 \(1,x,y,xy\) 在 \(M\) 上的表示矩阵，并说明如何从这些矩阵恢复左右 \(A\) 作用。
\end{exercise}
\begin{solution}
矩阵 \(E_{12}\)、\(E_{13}\) 的平方均为零，且 \(E_{12}E_{13}=E_{13}E_{12}=0\)，故它们分别满足 \(\varepsilon^2=0\) 并彼此交换。这给出相容的左、右幺性作用。令 \(x=\varepsilon\otimes1\)、\(y=1\otimes\varepsilon^{\mathrm{op}}\)。因 \(A\) 交换，\(A^e\) 由交换的 \(x,y\) 生成，满足 \(x^2=y^2=0\)；两边都以 \(1,x,y,xy\) 为 \(k\) 基，所以 \(A^e\cong k[x,y]/(x^2,y^2)\)。这四个元素在 \(M\) 上依次表示为 \(I_3,E_{12},E_{13},0\)。把所得表示分别限制到 \(A\otimes1\) 与 \(1\otimes A^{\mathrm{op}}\)，便恢复左、右 \(A\) 作用。
\end{solution}

\begin{exercise}\label{b2:ex:11-polynomial-bimodule-annihilator}
设 \(k\) 为域，\(n\ge2\)，\(A=k[t]\)，\(M=k[t]/(t^n)\) 以商中的乘法作左右 \(A\) 作用。把 \(A^e\) 识别为 \(k[x,y]\)，其中 \(x,y\) 分别记录左右乘 \(t\)。证明
\[
\operatorname{Ann}_{A^e}(M)=(x-y,y^n),
\qquad M\cong k[x,y]/(x-y,y^n)
\]
作为 \(A^e\) 模成立。说明为何也可把 \(y^n\) 换成 \(x^n\)，并求 \(\dim_kM\)。
\end{exercise}
\begin{solution}
用 \(x-y\) 对 \(F\in k[y][x]\) 作首一带余除法，写成 \(F=(x-y)Q+r(y)\)。在 \(M\) 上左右乘 \(t\) 相同，故 \(F\) 作用为乘以 \(r(t)+(t^n)\)。它零化 \(M\) 当且仅当零化 \(\bar1\)，即 \(t^n\mid r(t)\)。因此零化子正好是 \((x-y,y^n)\)。求值映射 \(F\mapsto F(t,t)+(t^n)\) 满射到 \(M\)，并保持 \(A^e\) 作用，给出所求商模同构。又
\[
x^n-y^n=(x-y)(x^{n-1}+x^{n-2}y+\cdots+y^{n-1}),
\]
所以 \((x-y,y^n)=(x-y,x^n)\)。商中 \(\bar1,\bar t,\ldots,\bar t^{n-1}\) 为基，故维数为 \(n\)。与正文正则双模的关系 \(x=y\) 相比，截断还加入了 \(t^n=0\)。
\end{solution}

\begin{exercise}\label{b2:ex:11-triangular-idempotents}
设 \(k\) 为域。求上三角代数 \(T_2(k)\) 的全部幂等元，并从中辨认中心幂等元。由此判断它能否分解为两个非零代数的直积。
\end{exercise}
\begin{solution}
写 \(e=\begin{pmatrix}a&b\\0&c\end{pmatrix}\)。\(e^2=e\) 等价于 \(a^2=a,c^2=c,(a+c-1)b=0\)。域中 \(a,c\in\{0,1\}\)。若 \(a=c\)，则 \(b=0\)，得到零与单位；若 \((a,c)=(1,0)\) 或 \((0,1)\)，则 \(b\) 任意。因此其余幂等元为
\[
\begin{pmatrix}1&b\\0&0\end{pmatrix},\qquad
\begin{pmatrix}0&b\\0&1\end{pmatrix}\quad(b\in k).
\]
与 \(E_{11}\) 交换要求 \(b=0\)，再与 \(E_{12}\) 交换要求 \(a=c\)，所以中心幂等元只有零与单位。若存在非零乘积分解，第一因子的单位 \((1,0)\) 会给出非平凡中心幂等元，矛盾。因此不存在这种分解，虽然正则模仍可按 \(E_{11},E_{22}\) 分成两个非零直和因子。
\end{solution}

\begin{exercise}\label{b2:ex:11-large-corner}
设 \(k\) 为域，\(A=M_3(k)\)，\(E_{ij}\) 为矩阵单位，\(e=E_{11}+E_{22}\)、\(f=E_{33}\)。写出四个空间 \(eAe,eAf,fAe,fAf\) 的矩阵形状与维数，并求 \(\operatorname{End}_A(Ae)\)。
\end{exercise}
\begin{solution}
按照前两个坐标与第三个坐标分块，四个空间依次是左上 \(2\times2\) 块、右上 \(2\times1\) 块、左下 \(1\times2\) 块、右下 \(1\times1\) 块，其他位置均为零；维数依次为 \(4,2,2,1\)，总和为九。角代数 \(eAe\cong M_2(k)\)，其单位是 \(e\)。由\cref{b2:prop:corner-endomorphism}，
\(\operatorname{End}_A(Ae)\cong M_2(k)^{\mathrm{op}}\)，所以其维数为四。这个自同态作用为右乘左上块，与 \(Ae\) 的六维底向量空间的全部自同态明显不同。
\end{solution}

\begin{exercise}\label{b2:ex:11-corner-hom}
设 \(k\) 为域，\(A=M_3(k)\)，\(E_{ij}\) 为矩阵单位，\(e=E_{11}+E_{22}\)、\(f=E_{33}\)。取
\[
b=E_{13}+E_{23}\in eAf,
\qquad c=E_{31}-E_{32}\in fAe.
\]
令 \(h_b:Ae\to Af\)、\(h_c:Af\to Ae\) 为对应的右乘同态。计算两个复合，证明其中一个为零，另一个是非零平方零自同态，并求后者的像与核的维数。复合仍按“先右边、后左边”的次序书写。
\end{exercise}
\begin{solution}
矩阵单位相乘给出
\[
bc=E_{11}-E_{12}+E_{21}-E_{22}\ne0,
\qquad cb=E_{33}-E_{33}=0.
\]
由\cref{b2:prop:corner-endomorphism}，\(h_c\circ h_b\) 是 \(Ae\) 上的右乘 \(bc\)，\(h_b\circ h_c\) 是 \(Af\) 上的右乘 \(cb\)，故后者为零。前者在 \(e\) 上取值为 \(bc\ne0\)，而 \((bc)^2=b(cb)c=0\)，所以非零且平方为零。若 \(X\in Ae\) 的前两列为 \(u,v\in k^3\)，则
\[
Xbc=\begin{pmatrix}u+v&-(u+v)&0\end{pmatrix}.
\]
其像由任意 \(w=u+v\in k^3\) 决定，维数为三；其核由 \(v=-u\) 决定，维数也为三，且像正好等于核。所有计算在特征二中仍成立。这里先作用 \(h_b\) 再作用 \(h_c\)，对应的是角块乘积 \(bc\)，与复合符号 \(h_c\circ h_b\) 的书写次序相反。
\end{solution}

\begin{exercise}\label{b2:ex:11-commuting-projections}
设 \(R\) 为含幺环，\(M\) 为幺性左 \(R\) 模，\(p,q\in\operatorname{End}_R(M)\) 都幂等且相互交换。证明四个算子
\(pq,p(1-q),(1-p)q,(1-p)(1-q)\) 是两两正交的幂等元且和为恒等，从而给出 \(M\) 的四项直和分解。给出不交换时其中某项不再幂等的例子。
\end{exercise}
\begin{solution}
因 \(p,q\) 交换，\(p,1-p\) 与 \(q,1-q\) 全部互相交换，故各乘积平方等于自身。两个不同乘积至少在一处分别含 \(p,1-p\) 或 \(q,1-q\)，相乘因此为零；展开 \(\bigl(p+(1-p)\bigr)\bigl(q+(1-q)\bigr)\) 得四项之和为恒等。由\cref{b2:prop:idempotent-module-decomposition}，四个像组成内直和，允许某些像为零。

在 \(k^2\) 上取 \(p=\begin{pmatrix}1&0\\0&0\end{pmatrix}\)、\(q=\begin{pmatrix}1&1\\0&0\end{pmatrix}\)。它们都幂等，但 \(pq=q\)、\(qp=p\)，不交换。此时 \(p(1-q)=\begin{pmatrix}0&-1\\0&0\end{pmatrix}\) 非零而平方为零，所以不是幂等元。原分解不能不加交换条件照搬。
\end{solution}

\begin{exercise}\label{b2:ex:11-cyclic-operator-centralizer}
设 \(k\) 为域，\(V\) 为 \(n\ge1\) 维 \(k\) 向量空间，\(T\in\operatorname{End}_k(V)\) 有循环向量 \(v\)，即 \(v,Tv,\ldots,T^{n-1}v\) 是 \(V\) 的基。所有中心化子均在 \(\operatorname{End}_k(V)\) 中取。证明 \(k[T]'=k[T]\)，从而 \(k[T]''=k[T]\)。写出三阶幂零 Jordan 块在 \(M_3(k)\) 中的中心化子。
\end{exercise}
\begin{solution}
若 \(S\) 与 \(T\) 交换，将 \(S(v)\) 按循环基展开，唯一写成 \(p(T)v\)，其中 \(\deg p<n\)。于是
\(S(T^jv)=T^jS(v)=p(T)T^jv\)，所以 \(S=p(T)\)。反过来，\(T\) 的任意多项式都与 \(T\) 交换，得到第一个等式；再取一次中心化子便得第二式。对标准上三角的 \(J_3\)，中心化子中的元素恰为
\[
aI+bJ_3+cJ_3^2
=\begin{pmatrix}a&b&c\\0&a&b\\0&0&a\end{pmatrix},\qquad a,b,c\in k.
\]
它包含非零幂零元，说明双中心化子等式本身不意味着作用代数没有幂零元。
\end{solution}

\begin{exercise}\label{b2:ex:11-repeated-matrix-blocks}
设 \(k\) 为域、\(n\ge1\)，在 \(V=k^n\oplus k^n\) 上令 \(\mathcal A=\bigl\{\operatorname{diag}(B,B):B\in M_n(k)\bigr\}\)。在 \(\operatorname{End}_k(V)\) 中求 \(\mathcal A'\) 与 \(\mathcal A''\)，并给出两者维数。
\end{exercise}
\begin{solution}
把算子写成四个 \(n\times n\) 块 \(T_{ij}\)。与每个 \(\operatorname{diag}(B,B)\) 交换，等价于每个 \(T_{ij}\) 都与全部 \(B\in M_n(k)\) 交换。逐个与矩阵单位比较，得每块都是标量矩阵，因此中心化子中的元素恰为
\[
T=\begin{pmatrix}aI_n&bI_n\\cI_n&dI_n\end{pmatrix},\qquad a,b,c,d\in k.
\]
与其中的两个块投影交换使双中心化子为块对角，再与交换两个副本的块矩阵交换使两对角块相等，故 \(\mathcal A''=\mathcal A\)。维数分别为四与 \(n^2\)。这个结果也符合\cref{b2:prop:tensor-factor-centralizers}在 \(k^n\otimes k^2\) 上的两个独立因子作用。
\end{solution}

\begin{exercise}\label{b2:ex:11-two-eigenvalue-blocks}
设 \(k\) 为域，\(\lambda,\mu\in k\)、\(\lambda\ne\mu\)，\(r,s>0\)，\(T=\operatorname{diag}(\lambda I_r,\mu I_s)\) 作用在 \(k^{r+s}\) 上。在 \(M_{r+s}(k)\) 中计算 \(k[T]'\) 和 \(k[T]''\)，并说明在 \(r>1\) 或 \(s>1\) 时，这与循环算子的例子有何不同。
\end{exercise}
\begin{solution}
分块交换方程使两个非对角块分别满足 \((\lambda-\mu)X=0\)、\((\mu-\lambda)Y=0\)，故为零；对角块任意，所以
\[
k[T]'=M_r(k)\times M_s(k)
\]
以块对角方式作用。其中心化子是每块各一个标量，故 \(k[T]''=kI_r\times kI_s\)。线性多项式可以在不同的 \(\lambda,\mu\) 处取任意两个给定值，所以后一代数恰为 \(k[T]\)。当有重数大于一时，第一次中心化子已经严格大于 \(k[T]\)；不过第二次中心化仍能恢复 \(k[T]\)，并不要求第一次中心化子等于自身。
\end{solution}

\begin{exercise}\label{b2:ex:11-tensor-flip-centralizer}
设 \(k\) 为特征不为二的域、\(V=k^2\)，令 \(\tau(v\otimes w)=w\otimes v\) 作用在 \(V\otimes_kV\) 上。在 \(\operatorname{End}_k(V\otimes_kV)\) 中求 \(k[\tau]'\) 与 \(k[\tau]''\)，并计算维数。把两个特征子空间与对称幂、外幂联系起来。
\end{exercise}
\begin{solution}
\(\tau\) 的 \(1\) 特征子空间有基
\(e_1\otimes e_1,e_2\otimes e_2,e_1\otimes e_2+e_2\otimes e_1\)，维数三；\(-1\) 特征子空间由 \(e_1\otimes e_2-e_2\otimes e_1\) 生成，维数一。特征不为二保证这些向量组成四维空间的一组基。与 \(\tau\) 交换恰好要求分别保持这两个特征子空间，因此中心化子为 \(M_3(k)\times k\)，维数十。再取中心化子，得到两个特征空间上的独立标量作用，即
\[
k\frac{1+\tau}{2}\oplus k\frac{1-\tau}{2}=kI\oplus k\tau,
\]
维数二。由\cref{b2:prop:symmetrization-characteristic}，\(\operatorname{Sym}^2_kV\) 通过 \(vw\mapsto(v\otimes w+w\otimes v)/2\) 与正特征子空间自然同构。映射 \(v\wedge w\mapsto(v\otimes w-w\otimes v)/2\) 由交替双线性关系良定义，取商到 \(\bigwedge^2_kV\) 又把它送回 \(v\wedge w\)，故它将外幂同构到负特征子空间。因而这里的两个块分别记录对称张量和交替张量上的独立算子。
\end{solution}

% !TeX root = ../../Book2.tex
% 纲要标题：## 第12章　链条件、长度与直和分解
% 纲要位置：工作记录/计划纲要.md:L1434

\chapter{链条件、长度与直和分解}
\label{b2:ch:12}
\BookChapterNumber{b2:ch:12}

\begin{chapterabstract}
链条件把有限维线性代数中“不能无限增加或减少维数”的现象推广到一般模。本章证明 Noether 条件与子模有限生成的等价性，说明升降链条件在短正合列中的传递及非交换环上的左右区别。合成列与模的 Jordan–Hölder 定理给出长度及其可加性；Fitting 引理把有限长度模分解为自同态作用的幂零分量与可逆分量，并由此证明不可分解模的自同态环局部。最后证明 Krull–Schmidt 分解的存在唯一性，用具体反例说明有限生成或单独的 Noether 条件为何不足。
\end{chapterabstract}

\begin{learninggoals}
\begin{itemize}
\item 区分有限生成、升链条件、降链条件与有限长度，能通过子模、商模及短正合列判断这些性质。
\item 证明模的 Jordan–Hölder 定理，计算长度和简单因子的重数，并区别长度与基域维数。
\item 证明 Fitting 分解，使用局部环判据辨认不可分解模，说明幂等元和投影细分的关系。
\item 在 Krull–Schmidt 唯一性证明中从恒等映射的有限和中找到可逆项，比较分量并交换补模，理解结论的适用条件及失效例子。
\end{itemize}
\end{learninggoals}

整数环 \(\mathbb Z\) 的每个理想都由一个整数生成，因而子模的生成问题十分简单；但
\[
\mathbb Z\supsetneq2\mathbb Z\supsetneq4\mathbb Z\supsetneq8\mathbb Z\supsetneq\cdots
\]
永不停止。它说明升链条件与降链条件不能混为一谈。有限维向量空间同时满足两者，是因为严格变化每次都改变维数；一般模需要分别研究这两种有限性。

\ProseParagraphBreak

当升降过程都能停止时，一个模可以被逐层剖开，记录简单的商层；但仅知道这些层是什么，并不能自动把它们重新拼成直和。Jordan–Hölder 定理确定这些商层的同构类型和重数，Krull–Schmidt 定理确定实际的不可分解直和分量，两者回答不同的问题。为证明后一种唯一性，还需要用链条件控制自同态反复作用后的核与像，由 Fitting 分解得到不可分解模的局部自同态环。

\ProseParagraphBreak

本章需要\chapref{b2:ch:02}的子模与同构定理、\chapref{b2:ch:03}的正合列与直和，以及\chapref{b2:ch:11}关于幂等自同态与模分解的结果。个别算例还会用到\chapref{b2:ch:04}的主理想整环模分类和\chapref{b2:ch:07}的投射模判据。局部环在这里由非单位理想定义。

\ProseParagraphBreak
交换环上的 Noether 模判据可参阅 Roman 的《Advanced Linear Algebra》\cite[Chapter 5, Theorem 5.7, pp. 133--134]{Roman2008AdvancedLinearAlgebra}。有限维表示的合成列与不可分解分解，可参阅 Etingof 等人的《Introduction to Representation Theory》\cite[Theorem 3.7.1, pp. 50--51; Theorem 3.8.1, pp. 51--52]{EtingofEtAl2011RepresentationTheory}。本章讨论一般环上的有限长度模，不预设代数闭底域。

% !TeX root = ../../Book2.tex
% 纲要标题：### 12.1 链条件
% 纲要位置：工作记录/计划纲要.md:L1436

\section{链条件}
\label{b2:sec:1201}

有限维向量空间的严格子空间链必在有限步后终止，因为每次严格包含都改变维数。一般模没有这样的维数，但仍可直接要求子模链终止。本节区分升链与降链两种条件，并说明它们怎样通过子模、商模和短正合列传递。除另外说明外，\(R\) 为含单位元的环，模均为幺性左 \(R\) 模，不假设环交换。

% 纲要要点（供目录核对）：
%
% * Noether 模与 Artin 模
% * 子模、商模与短正合列中的传递
% * 环的左右链条件
%

\subsection{Noether 模与 Artin 模}
\label{b2:subsec:1201:01}

\begin{definition}\label{b2:def:noetherian-artinian-module}
设 \(R\) 为含幺环、\(M\) 为幺性左 \(R\) 模。若 \(M\) 的每条子模升链
\[
N_1\subseteq N_2\subseteq\cdots
\]
都最终稳定，即存在 \(r\) 使 \(N_i=N_r\) 对所有 \(i\ge r\) 成立，则称 \(M\) 满足\term[sheng-lian-tiaojian]{升链条件}[ascending chain condition]，也称 \(M\) 为 \term[Noether-mo]{Noether 模}[Noetherian module]。若每条子模降链 \(N_1\supseteq N_2\supseteq\cdots\) 都最终稳定，则称 \(M\) 满足\term[jiang-lian-tiaojian]{降链条件}[descending chain condition]，也称 \(M\) 为 \term[Artin-mo]{Artin 模}[Artinian module]。两种条件分别简称 ACC 与 DCC。
\end{definition}

“最终稳定”不是说所有链都有一个共同的长度上界。例如 \(\mathbb Z\) 的每个子模为主理想，升链最终稳定；但
\[
0\subsetneq 2^n\mathbb Z\subsetneq2^{n-1}\mathbb Z
\subsetneq\cdots\subsetneq\mathbb Z
\]
可以任意长。这个模也不满足降链条件，因为 \(\mathbb Z\supsetneq2\mathbb Z\supsetneq4\mathbb Z\supsetneq\cdots\) 永不稳定。

\begin{proposition}\label{b2:prop:chain-maximal-minimal}
设 \(R\) 为含幺环、\(M\) 为幺性左 \(R\) 模。\(M\) 满足升链条件，当且仅当每个非空子模集合在包含关系下有极大元。把升链、极大元分别换成降链、极小元，也有同样的等价。
\end{proposition}
\begin{proof}
若某个非空子模集合没有极大元，从其中任取 \(N_1\)，再逐次取严格包含前一项的 \(N_{i+1}\)，得到不稳定的升链，矛盾。反之，升链的全部项组成非空集合，若其中某项 \(N_r\) 极大，则以后各项包含它，只能都等于它。降链的论证将包含方向反转即可：无极小元时逐次严格缩小；有极小项时后续项与之相等。这里的逐次选取沿用本书的选择公理约定。
\end{proof}

这里的极大元只要求不能在给定子模集合中继续严格增大，不要求包含这个集合中的所有成员。例如，域 \(k\) 上 \(k^2\) 中两个不同的一维子空间，在由它们组成的集合中都为极大元，却都不是最大元。在一条链中，各项可比，极大项才自动成为最大项；极小与最小的区别同样如此。

\begin{proposition}\label{b2:prop:noetherian-submodules-fg}
设 \(R\) 为含幺环、\(M\) 为幺性左 \(R\) 模。\(M\) 是 Noether 模，当且仅当它的每个子模都作为左 \(R\) 模有限生成。
\end{proposition}
\begin{proof}
若某个子模 \(N\) 不有限生成，就可以不断加入旧生成元无法表示的新元素。先取 \(x_1\in N\)，再从 \(N\setminus(Rx_1+\cdots+Rx_i)\) 取 \(x_{i+1}\)，得到严格升链
\[
Rx_1\subsetneq Rx_1+Rx_2\subsetneq\cdots.
\]
这与升链条件矛盾，故每个子模有限生成。零子模由空集生成，不需另选元素。

反过来，给定升链 \((N_i)\)，其并 \(N=\bigcup_iN_i\) 为子模：两个元素同属某个较大的 \(N_i\)，和与标量倍都仍在其中。取 \(N\) 的有限生成元 \(x_1,\ldots,x_s\)。每个生成元在链的某一项中，有限多个这样的项有共同的后继项 \(N_r\)，所以全部生成元已经同时落在 \(N_r\) 中；零子模的情形可任取 \(r\)。于是 \(N\subseteq N_r\subseteq N\)，后续各项只能等于 \(N_r\)，所以升链稳定。
\end{proof}

特别地，Noether 模本身有限生成，反向却依赖于环。例如 \(R=k[X_1,X_2,\ldots]\) 是循环左 \(R\) 模，但理想链
\((X_1)\subsetneq(X_1,X_2)\subsetneq\cdots\) 严格上升。严格性可令前 \(i\) 个不定元取零来核验：\(X_{i+1}\) 不可能属于前面的理想。有限生成只控制整个模的生成元，并不自动控制全部子模。

\ProseParagraphBreak
Artin 条件与升链条件的差别可在交换群中直接看到。固定素数 \(p\)，令
\[
C_{p^\infty}=\mathbb Z[1/p]/\mathbb Z,\qquad
C_{p^n}=\langle p^{-n}+\mathbb Z\rangle.
\]
每个元素的阶为某个 \(p\) 的幂。阶为 \(p^m\) 的元素可写成 \(a/p^m+\mathbb Z\)，其中分子与 \(p\) 互素，因而它生成 \(C_{p^m}\)。被 \(p^n\) 消去的全部元素恰好组成 \(C_{p^n}\)：这正要求约分后的分母不超过 \(p^n\)。所以，若子群 \(H\) 中的元素阶无界，对每个 \(n\) 都能找到阶至少为 \(p^n\) 的元素，其生成的循环群包含 \(C_{p^n}\)，于是 \(H=C_{p^\infty}\)。若非零子群 \(H\) 的元素阶有界，取其中达到最大阶 \(p^m\) 的元素，便有 \(C_{p^m}\subseteq H\)；同时 \(H\) 的全部元素都被 \(p^m\) 消去，所以 \(H\subseteq C_{p^m}\)。这就证明每个真子群都是某个有限循环群 \(C_{p^m}\) 或零群。

\ProseParagraphBreak
因此，\(C_{p^\infty}\) 中的降链一旦离开全群，便落在一个有限群中，此后必稳定，所以它是 Artin \(\mathbb Z\) 模。另一方面，\(C_p\subsetneq C_{p^2}\subsetneq\cdots\) 永不稳定，它不是 Noether 模；任何有限组元素又同时落在某个 \(C_{p^n}\) 中，不能生成全群，所以它也不是有限生成模。与上面 \(\mathbb Z\) 为 Noether 模而非 Artin 模的例子合起来可见，这两种链条件互不蕴含，Artin 条件也不保证有限生成。

\subsection{链条件在短正合列中的传递}
\label{b2:subsec:1201:02}

在子模和商模上同时检测一条链，可以把中间模的增长或缩小分成两部分。对于可比的子模 \(L\subseteq L'\)，若增长既没有在固定子模 \(A\) 内留下新部分，也没有在商模中留下新像，那么实际没有增长。下面的引理把这个比较写成精确的交与像条件。

\begin{lemma}\label{b2:lem:submodule-intersection-image}
设 \(R\) 为含幺环、\(B\) 为幺性左 \(R\) 模、\(A\subseteq B\) 为子模，\(q:B\rightarrow B/A\) 为商映射。若子模 \(L\subseteq L'\subseteq B\) 满足
\[
L\cap A=L'\cap A,\qquad q(L)=q(L'),
\]
则 \(L=L'\)。
\end{lemma}
\begin{proof}
给定 \(x\in L'\)，由像相同可取 \(y\in L\) 使 \(q(x)=q(y)\)。于是 \(x-y\in L'\cap A=L\cap A\subseteq L\)，所以 \(x=y+(x-y)\in L\)。这证明反向包含，得到相等。
\end{proof}

\begin{theorem}\label{b2:thm:chain-conditions-exact}
设 \(R\) 为含幺环，\(0\rightarrow A\rightarrow B\rightarrow C\rightarrow0\) 为幺性左 \(R\) 模的短正合列，则 \(B\) 为 Noether 模，当且仅当 \(A,C\) 均为 Noether 模；对 Artin 模也有相同结论。
\end{theorem}
\begin{proof}
将 \(A\) 识别为 \(B\) 的子模，\(C\) 识别为 \(B/A\)。若 \(B\) 的子模升链稳定，\(A\) 中的升链也是 \(B\) 中的升链，故稳定；商模 \(C\) 中的升链拉回为包含 \(A\) 的升链，由\cref{b2:thm:module-correspondence}的子模对应，其稳定性也成立。

反过来，给定 \(B_1\subseteq B_2\subseteq\cdots\)，同时观察它在 \(A\) 内的交链 \(B_i\cap A\) 与在 \(C\) 中的像链 \(q(B_i)\)。若 \(A,C\) 均为 Noether 模，取两条链稳定起点中的较大者 \(r\)。对每个 \(i\ge r\)，\(B_r\subseteq B_i\) 有相同交与像，\cref{b2:lem:submodule-intersection-image}给出 \(B_r=B_i\)，所以两端停止变化后，中间链也停止变化。

对降链，子模继承和商模拉回仍然有效。若两端满足降链条件，取交链和像链的共同稳定起点 \(r\)，对 \(i\ge r\) 应用该引理于 \(B_i\subseteq B_r\)，又得相等。因此 Artin 情形的两方向都成立。
\end{proof}

由 \(0\rightarrow M\rightarrow M\oplus N\rightarrow N\rightarrow0\)，有限直和满足某一种链条件，当且仅当每个分量都满足。这一结论不能推广为无限直和：非零模 \(M_i\) 的可数直和有严格升链 \(\bigoplus_{i\le n}M_i\)，也有严格降链 \(\bigoplus_{i\ge n}M_i\)。两种严格性都可用某个非零坐标元素核验。

\ProseParagraphBreak
引理中的“有相同交与像”必须配合包含关系使用。两个不同的补模可以与 \(A\) 都交于零、都满射到 \(B/A\)，却仍然不同。例如在 \(k^2\) 中，以 \(A=ke_1\) 为固定子空间，\(ke_2\) 与 \(k(e_1+e_2)\) 就有相同的交与像。链条件证明中，各项彼此可比，才保证引理适用。

\subsection{环的左右链条件}
\label{b2:subsec:1201:03}

\begin{definition}\label{b2:def:left-right-chain-ring}
设 \(R\) 为含幺环。若左正则模 \({}_RR\) 是 Noether 模，则称 \(R\) 为\term[zuo-Noether-huan]{左 Noether 环}[left Noetherian ring]；若 \({}_RR\) 是 Artin 模，则称 \(R\) 为\term[zuo-Artin-huan]{左 Artin 环}[left Artinian ring]。相应地，对右正则模 \(R_R\) 定义右 Noether 环和右 Artin 环。它们分别要求左理想链或右理想链满足相应条件。
\end{definition}

在交换环中，左右理想相同，通常省略方向。非交换环中则必须保留它。\cref{b2:prop:noetherian-submodules-fg}说明，左 Noether 环等价于每个左理想都作为左模有限生成；这没有自动给出右理想的有限生成性。

\begin{proposition}\label{b2:prop:ring-chain-finitely-generated}
设 \(R\) 为含幺环。\(R\) 是左 Noether 环，当且仅当每个有限生成幺性左 \(R\) 模都为 Noether 模。把 Noether 换成 Artin，同样成立。
\end{proposition}
\begin{proof}
若左正则模满足所讨论的链条件，\cref{b2:thm:chain-conditions-exact}说明有限直和 \(R^n\) 也满足。任意有限生成模都有满同态 \(R^n\rightarrow M\)，商模继承链条件，所以 \(M\) 满足。反向取有限生成左模 \(M=R\) 即得环本身的条件。两个论证对升链和降链分别有效。
\end{proof}

要让左右链条件不同，可以让同一块的左侧使用一个大标量域，右侧却只使用它的子域。令 \(k=\mathbb Q\)、\(K=\mathbb Q(t)\)，取三角矩阵环
\[
R=\left\{\,
\begin{pmatrix}a&u\\0&b\end{pmatrix}
\ \middle|\ a,u\in K,\ b\in k
\,\right\}.
\]
记 \(e_{11}=\begin{psmallmatrix}1&0\\0&0\end{psmallmatrix}\)、\(e_{22}=\begin{psmallmatrix}0&0\\0&1\end{psmallmatrix}\)，并令
\[
J=\left\{\,
\begin{pmatrix}0&u\\0&0\end{pmatrix}\ \middle|\ u\in K
\,\right\},
\]
将其中的元素记为 \(j(u)\)。对 \(r=\begin{psmallmatrix}a&v\\0&b\end{psmallmatrix}\in R\)，直接相乘得到
\[
rj(u)=j(au),\qquad j(u)r=j(ub).
\]
因此 \(J\) 左侧可以使用全部 \(K\) 标量，右侧却只能使用 \(k\) 标量。左模 \(Re_{11}\) 由左上位置的 \(K\) 组成，\(J\subseteq Re_{22}\) 也以 \(K\) 中的标量作左作用；这两个左模的任意非零元素都生成全模，所以没有非零真子模。商模 \(Re_{22}/J\cong k\) 只留下右下坐标，其左作用为乘 \(b\in k\)，也没有非零真子模。这三个模都满足两种链条件，故由短正合列传递和按列分解 \(R=Re_{11}\oplus Re_{22}\)，\(R\) 同时为左 Noether 环和左 Artin 环。

\ProseParagraphBreak
在右侧，上面的公式还说明 \(J_R\) 的子模恰好对应 \(K\) 的 \(k\) 线性子空间：任意 \(b\in k\) 都能作为右下坐标出现，所以右作用下封闭正是对这些标量倍封闭。取
\[
U_n=\operatorname{span}_k\{1,t,\ldots,t^n\},\qquad
W_n=\operatorname{span}_k\{t^j:j\ge n\},
\]
得到严格上升的 \(U_n\) 与严格下降的 \(W_n\)，因为不定元的不同幂在 \(k\) 上线性无关。由 \(j\) 将它们送入 \(J\)，便得到两条真正的右理想链，而不只是向量子空间链。因此这个 \(R\) 既不是右 Noether 环，也不是右 Artin 环；即使两种左链条件同时成立，也不能略去左右方向。若将 \(K/k\) 换成有限扩张，两侧都能具有有限的简单商层，但层数仍可能不同，\cref{b2:prop:triangular-left-right-length}将给出具体计算。

% !TeX root = ../../Book2.tex
% 纲要标题：### 12.2 有限长度
% 纲要位置：工作记录/计划纲要.md:L1442

\section{有限长度}
\label{b2:sec:1202}

升链条件阻止不断增加新子模，降链条件阻止不断缩小旧子模；同时具有这两种条件时，模可以在有限步内逐层取商，直到各商层都没有非零真子模。本节首先证明这些层的同构类型和重数不依赖于合成列的选择，再据此定义长度。长度计算的是简单商层，原模中相邻商层之间的扩张可以不分裂。

% 纲要要点（供目录核对）：
%
% * 合成列与 Jordan–Hölder 定理
% * 长度的可加性
% * 有限维代数的有限生成模与左右长度
%

\subsection{合成列与 Jordan–Hölder 定理}
\label{b2:subsec:1202:01}

\begin{definition}\label{b2:def:composition-series}
设 \(R\) 为含幺环。若非零幺性左 \(R\) 模 \(S\) 只有 \(0,S\) 两个子模，则称 \(S\) 为\term[jiandan-mo]{简单模}[simple module]。对幺性左 \(R\) 模 \(M\)，取有限子模链
\[
0=M_0\subsetneq M_1\subsetneq\cdots\subsetneq M_n=M,
\]
若其中每个商 \(M_i/M_{i-1}\) 都简单，则称这条链为 \(M\) 的\term[hecheng-lie]{合成列}[composition series]，这些商称为该列的\term[hecheng-yinzi]{合成因子}[composition factor]。若 \(M\) 存在合成列，则称它为\term[youxian-changdu-mo]{有限长度模}[module of finite length]；零模采用没有非零商层的空合成列。
\end{definition}

固定素数 \(p\)。例如 \(\mathbb Z/p\mathbb Z\) 为简单 \(\mathbb Z\) 模；而 \(\mathbb Z/p^2\mathbb Z\) 不简单，因为其中有非零真子模 \(p\mathbb Z/p^2\mathbb Z\)。不过
\[
0\subsetneq p\mathbb Z/p^2\mathbb Z\subsetneq\mathbb Z/p^2\mathbb Z
\]
是一条合成列，两层都同构于 \(\mathbb Z/p\mathbb Z\)。\((\mathbb Z/p\mathbb Z)^2\) 有同样的两个合成因子，却不与前者同构：前者有阶为 \(p^2\) 的元素，后者每个元素都被 \(p\) 零化。用\cref{b2:thm:split-exact-sequence}的语言来看，第一条列对应的短正合列
\[
0\longrightarrow\mathbb Z/p\mathbb Z
\xrightarrow{\ \bar a\mapsto pa+p^2\mathbb Z\ }
\mathbb Z/p^2\mathbb Z
\longrightarrow\mathbb Z/p\mathbb Z\longrightarrow0
\]
不分裂：中间模里被 \(p\) 消去的元素都在 \(p\mathbb Z/p^2\mathbb Z\) 中，取商后为零，不可能给出截面。而以 \((\mathbb Z/p\mathbb Z)^2\) 为中间模、把第一坐标作为子模的相应短正合列分裂，第二坐标就是补模。合成因子记录简单商层，不记录相邻层之间的扩张是否分裂，也不表示这些层都已成为原模的直和分量。

\begin{lemma}\label{b2:lem:composition-series-subquotient}
设 \(R\) 为含幺环、\(M\) 为有限长度幺性左 \(R\) 模。\(M\) 的每个子模和商模都有合成列。若 \(M\) 有一条包含 \(n\) 个简单商层的合成列，则所构造的子模与商模合成列各至多有 \(n\) 层。
\end{lemma}
\begin{proof}
给定合成列 \(0=M_0\subsetneq\cdots\subsetneq M_n=M\) 及子模 \(N\subseteq M\)。把各项与 \(N\) 相交，得到
\[
0=N\cap M_0\subseteq N\cap M_1\subseteq\cdots\subseteq N\cap M_n=N.
\]
包含映射诱导
\[
(N\cap M_i)/(N\cap M_{i-1})\longrightarrow M_i/M_{i-1}.
\]
它单射，因为映到零正好要求代表元落在 \(M_{i-1}\)。像为简单模的子模，所以或为零、或为全体。删去重复项，得到 \(N\) 的合成列。

再把各项映到 \(M/N\)，得到
\[
0=(M_0+N)/N\subseteq(M_1+N)/N\subseteq\cdots\subseteq(M_n+N)/N=M/N.
\]
相邻商是 \(M_i/M_{i-1}\) 的商：映射把 \(m+M_{i-1}\) 送到 \(m+(M_{i-1}+N)\)，在相应商层上满射。因此它或为零、或与原简单模同构；为零正是相邻项重复的情形。删去重复项，即得 \(M/N\) 的合成列。两种过程都不会增加原来的 \(n\) 个商层。
\end{proof}

比较两条合成列时，若它们紧邻全模的子模不同，可以先取这两个子模的交，再把两种最后的简单商交叉配对。这样剩余部分可以归纳比较，末端只交换两个简单因子的次序。

\begin{theorem}[Jordan–Hölder]\label{b2:thm:module-jordan-holder}
设 \(R\) 为含幺环、\(M\) 为有限长度幺性左 \(R\) 模。\(M\) 的任意两条合成列具有相同的层数，并且其合成因子在计入重数后，除排列次序外逐个同构。
\end{theorem}
\begin{proof}
对第一条合成列的层数 \(n\) 作强归纳。\(n=0\) 时 \(M=0\)，只有空列。设 \(n>0\)，记两条合成列中紧邻 \(M\) 的真子模为 \(A,B\)，所以 \(M/A,M/B\) 简单。

若 \(A=B\)，对 \(A\) 的两条合成列应用归纳假设，再在两侧加上相同的商 \(M/A\)，即得结论。现设 \(A\ne B\)。它们都是极大真子模，因此 \(A+B=M\)；否则 \(A+B\) 仍真且包含 \(A\)，会等于 \(A\)，从而 \(B\subseteq A\)，再由 \(B\) 极大得到 \(B=A\)，矛盾。

令 \(C=A\cap B\)。映射 \(A\rightarrow M/B\)、\(a\mapsto a+B\) 的核为 \(C\)，像由 \(A+B=M\) 为全商，因此\cref{b2:thm:module-first-isomorphism}给出 \(A/C\cong M/B\)。同样有 \(B/C\cong M/A\)。这两个商都简单。

由\cref{b2:lem:composition-series-subquotient}，\(C\subseteq A\) 有某条有限合成列，设这条列有 \(r\) 层。将其后接简单商 \(A/C\)，得到 \(A\) 的一条 \(r+1\) 层合成列。第一条原列在 \(A\) 中已有 \(n-1\) 层，故以原列的这一部分应用归纳假设，得到 \(r+1=n-1\)，且两种因子列表相同。于是所取 \(C\) 的合成列恰有 \(n-2\) 层。

把同一条 \(C\) 的合成列后接 \(B/C\)，得到 \(B\) 的一条 \(n-1\) 层合成列。以这条新列作为强归纳中的第一条列，就可与第二条原列在 \(B\) 中的部分比较。第一条 \(M\) 合成列的因子，因而由 \(C\) 的这些因子再加 \(A/C\cong M/B\) 与 \(M/A\) 组成；第二条则由同一组 \(C\) 因子再加 \(B/C\cong M/A\) 与 \(M/B\) 组成。它们只是最后两种因子的次序互换，故层数和全部重数相同。
\end{proof}

对含幺环 \(R\) 上的有限长度幺性左模 \(M\)，模的\term[mo-Jordan-Holder-dingli]{Jordan–Hölder 定理}[Jordan–Hölder theorem for modules]使以下记号有意义：
\[
\ell_R(M)=\text{\(M\) 的任意合成列的层数}.
\]
这个数称为 \(M\) 的\term[mo-changdu]{长度}[length of a module]；环已明确时简写 \(\ell(M)\)。它为零当且仅当 \(M=0\)。对于简单左 \(R\) 模 \(S\)，还可把它在合成列中出现的次数记为 \([M:S]\)，称为相应合成重数；只比较同构类型，不把具体嵌入方式计入这个数。
\symidx{\(\ell_R(M)\)：左 \(R\) 模的长度}
\symidx{\([M:S]\)：简单模 \(S\) 在 \(M\) 中的合成重数}

\subsection{长度的可加性}
\label{b2:subsec:1202:02}

\begin{theorem}\label{b2:thm:finite-length-chain-conditions}
设 \(R\) 为含幺环、\(M\) 为幺性左 \(R\) 模。\(M\) 有限长度，当且仅当它同时为 Noether 模和 Artin 模。
\end{theorem}
\begin{proof}
简单模没有可继续严格变化的子模链，故同时满足两种链条件。若 \(M\) 有合成列，逐层应用\cref{b2:thm:chain-conditions-exact}，便知全部层的扩张 \(M\) 同时满足两种链条件。

反过来，设 \(M\) 同时满足两种条件。升链条件用来保证每一步能找到简单商，降链条件则用来保证这个过程在有限步后结束。若 \(M\ne0\)，由\cref{b2:prop:chain-maximal-minimal}，全部真子模中有极大元 \(M^{(1)}\)，其商简单；若 \(M^{(1)}\ne0\)，再在它内部选取极大真子模。子模继承升链条件，因此每一步都可进行。若永不到零，就得到无限严格降链，违反 Artin 条件。所以经过有限步到达零，反向排列即为合成列。
\end{proof}

\begin{proposition}[长度的可加性]\label{b2:prop:length-additive}
设 \(R\) 为含幺环，\(0\rightarrow A\rightarrow B\rightarrow C\rightarrow0\) 为幺性左 \(R\) 模的短正合列。\(B\) 有限长度当且仅当 \(A,C\) 都有限长度。此时
\[
\ell(B)=\ell(A)+\ell(C),\qquad
[B:S]=[A:S]+[C:S]
\]
对每个简单幺性左 \(R\) 模 \(S\) 都成立，其中 \(\ell\) 表示合成列长度，\([X:S]\) 表示 \(S\) 在 \(X\) 的合成列中出现的重数。
\end{proposition}
\begin{proof}
若 \(B\) 有限长度，子模 \(A\) 与商模 \(C\) 由\cref{b2:lem:composition-series-subquotient}也有合成列。反过来，给定 \(A,C\) 的合成列，把 \(C\) 中的各项通过商映射拉回到 \(B\)，得到从 \(A\) 到 \(B\) 的链，其相邻商与 \(C\) 的各商同构；再在前面接上 \(A\) 的合成列，就得到 \(B\) 的合成列。新列恰好把两端的全部因子合并，Jordan–Hölder 定理保证所计算的长度与重数不依赖这些选择，于是两个公式成立。
\end{proof}

因此有限长度模的真子模和真商模长度都严格变小。对严格链 \(0=L_0\subsetneq L_1\subsetneq\cdots\subsetneq L_r=M\)，各 \(L_i/L_{i-1}\) 非零，长度至少为一，反复使用可加性得到 \(r\le\ell(M)\)。合成列达到这个上界，所以长度也是严格子模链可达到的最大层数。\secref{b2:sec:1201}中的 \(\mathbb Z\) 虽然满足升链条件，却有任意长的有限严格链；有限长度在此提供了真正统一的上界。

\begin{proposition}\label{b2:prop:finite-length-endomorphism}
设 \(R\) 为含幺环、\(M\) 为有限长度幺性左 \(R\) 模。对自同态 \(f:M\rightarrow M\)，单射、满射、可逆这三个条件等价。更一般地，
\[
\ell(M)=\ell(\ker f)+\ell(\operatorname{im}f)
\]
对任何到幺性左 \(R\) 模 \(N\) 的同态 \(f:M\to N\) 成立。
\end{proposition}
\begin{proof}
短正合列 \(0\rightarrow\ker f\rightarrow M\rightarrow\operatorname{im}f\rightarrow0\) 与长度可加性给出公式。若 \(f\) 是自同态且满射，像长度等于 \(\ell(M)\)，故核长度为零，核为零；若它单射，像长度等于 \(\ell(M)\)，再对像的包含取商，商长度为零，故像为全模。模的双射同态有线性逆，因而可逆。可逆显然同时给出单射和满射。
\end{proof}

这个长度公式对应线性代数中维数等于核维数与像维数之和；有限长度自同态的单射与满射等价，也对应有限维线性空间中的结论。例如，对素数 \(p\) 和 \(n\ge1\)，\(M=\mathbb Z/p^n\mathbb Z\) 的链 \(0\subsetneq p^{n-1}M\subsetneq\cdots\subsetneq M\) 有 \(n\) 个简单商层，所以长度为 \(n\)。乘 \(p\) 的核长度为一、像长度为 \(n-1\)，可加性给出 \(n=1+(n-1)\)。在 \(n>1\) 时，这个算子既不单射，也不满射。

\subsection{有限维代数的有限生成模}
\label{b2:subsec:1202:03}

域 \(k\) 上的简单模恰好是一维向量空间，所以对有限维 \(k\) 向量空间 \(V\)，有 \(\ell_k(V)=\dim_kV\)。一般代数上的简单模却可能有更大的 \(k\) 维数；要比较长度和维数，需要同时计算每个简单商层的维数。

\begin{proposition}\label{b2:prop:finite-dimensional-algebra-modules}
设 \(k\) 为域、\(A\) 为有限维结合含幺 \(k\) 代数、\(M\) 为幺性左 \(A\) 模。以下三个条件等价：\(M\) 有限生成；\(M\) 作为 \(k\) 向量空间有限维；\(M\) 作为 \(A\) 模有限长度。若 \(S_1,\ldots,S_r\) 为一条合成列的因子，则
\[
\dim_kM=\sum_{i=1}^r\dim_kS_i,\qquad \ell_A(M)\le\dim_kM.
\]
\end{proposition}
\begin{proof}
有限生成时存在满同态 \(A^n\rightarrow M\)，因 \(A\) 有限维，目标的 \(k\) 维数有限。有限维时，严格子模链也为严格子空间链，维数只能有限次增加或减少，所以 \(M\) 同时满足两种链条件，由\cref{b2:thm:finite-length-chain-conditions}得有限长度。有限长度又蕴含 Noether 条件，故由\cref{b2:prop:noetherian-submodules-fg}得 \(M\) 有限生成，完成等价。

每个简单因子 \(S_i\) 由任一非零元素生成，故是 \(A\) 的商模，具有有限且正的 \(k\) 维数。合成列逐层使用向量空间短正合列的维数公式，给出所列和式；每项至少为一，所以得到长度的不等式。
\end{proof}

长度计算简单商层的个数，维数则是这些商层的 \(k\) 维数之和，所以二者一般不同。取 \(A=M_n(k)\)、\(n\ge1\)，其标准列模 \(k^n\) 简单：任取非零向量，利用某个非零坐标和矩阵单位可以得到每个标准基向量，故它生成全模。因此这个模的 \(A\) 长度为一，而 \(k\) 维数为 \(n\)。对于交换环 \(k[t]\) 上的模 \(k[t]/(p^m)\)，其中 \(m\ge1\)、\(p\in k[t]\) 首一不可约，逐次乘 \(p\) 给出 \(m\) 个同构于 \(k[t]/(p)\) 的简单层；其长度为 \(m\)，\(k\) 维数为 \(m\deg p\)。

\ProseParagraphBreak
代数有限维是上述等价的实质条件。\(k[t]\) 作为自身的模由一生成，却有不稳定降链 \((t)\supsetneq(t^2)\supsetneq\cdots\)，所以不有限长度；域 \(k(t)\) 作为自身的模长度为一，但作为 \(k\) 向量空间无限维。研究有限维表示时，基域维数可以保证有限长度；研究一般环上的模时，应直接核验所需的链条件。

\ProseParagraphBreak
长度不仅取决于底层加法群，还取决于允许哪些标量作用。在\secref{b2:sec:1201}的三角矩阵环中，上右块的左作用使用 \(K\) 标量，右作用只使用 \(k\) 标量；将 \(K/k\) 改为有限扩张后，这个差别可以直接反映在左右正则模的长度上。

\begin{proposition}\label{b2:prop:triangular-left-right-length}
设 \(k\subseteq K\) 为有限域扩张，\(d=[K:k]\)，并取结合含幺环
\[
R=\left\{\,
\begin{pmatrix}a&u\\0&b\end{pmatrix}
\ \middle|\ a,u\in K,\ b\in k
\,\right\}.
\]
幺性左正则模 \({}_RR\) 与幺性右正则模 \(R_R\) 都有限长度，且
\[
\ell({}_RR)=3,\qquad \ell(R_R)=d+2.
\]
这里右模长度按相应的幺性左 \(R^{\mathrm{op}}\) 模的合成列定义。特别地，\(d>1\) 时，同一个环的左右正则模长度不同。
\end{proposition}
\begin{proof}
记矩阵单位 \(e_{11}=\begin{psmallmatrix}1&0\\0&0\end{psmallmatrix}\)、\(e_{22}=\begin{psmallmatrix}0&0\\0&1\end{psmallmatrix}\)，令
\[
j(u)=\begin{pmatrix}0&u\\0&0\end{pmatrix},\qquad J=j(K).
\]
对 \(r=\begin{psmallmatrix}a&v\\0&b\end{psmallmatrix}\in R\)，有 \(rj(u)=j(au)\)、\(j(u)r=j(ub)\)，所以 \(J\) 是双边理想。

按列分解有 \({}_RR=Re_{11}\oplus Re_{22}\)。\(Re_{11}\) 由左上坐标组成，其左作用为 \(K\) 上的标量乘法；\(J\) 上的左作用同样如此。这两个左模的任意非零元素都能用 \(K\) 标量生成全模，所以都简单。\(Re_{22}/J\) 只保留右下坐标，左作用为 \(k\) 上的标量乘法，故它也简单。因此 \(0\subsetneq J\subsetneq Re_{22}\) 是两层合成列，再由\cref{b2:prop:length-additive}得
\[
\ell({}_RR)=\ell(Re_{11})+\ell(Re_{22})=1+2=3.
\]

按行分解则有 \(R_R=e_{11}R\oplus e_{22}R\)。\(e_{22}R\) 只保留右下坐标，是按 \(k\) 标量右乘的简单右模。又 \(J\subseteq e_{11}R\)，商 \(e_{11}R/J\) 只保留左上坐标，其右作用是乘 \(a\in K\)；任意非零坐标都能生成整个 \(K\)，所以这个商也简单。

\(J\) 上的右作用却只用 \(b\in k\)。选 \(K/k\) 的一组基 \(\omega_1,\ldots,\omega_d\)，令 \(J_0=0\)、\(J_i=j(\operatorname{span}_k\{\omega_1,\ldots,\omega_i\})\)。于是
\[
0=J_0\subsetneq J_1\subsetneq\cdots\subsetneq J_d=J
\]
是右子模链，每个商 \(J_i/J_{i-1}\) 都是一维 \(k\) 空间，且右作用为右下坐标的任意 \(k\) 标量，故都简单。因此 \(\ell(J_R)=d\)。将这条列后接简单商 \(e_{11}R/J\)，得到 \(\ell(e_{11}R)=d+1\)；将右模视为幺性左 \(R^{\mathrm{op}}\) 模，再应用\cref{b2:prop:length-additive}，便得
\[
\ell(R_R)=\ell(e_{11}R)+\ell(e_{22}R)=d+1+1=d+2.
\]
\end{proof}

% !TeX root = ../../Book2.tex
% 纲要标题：### 12.3 不可分解模
% 纲要位置：工作记录/计划纲要.md:L1448

\section{不可分解模}
\label{b2:sec:1203}

合成列逐层取商，直和分解则要求每个分量在原模中有一个补模。两者的差别可由投影看清：一个直和分量对应一个满足 \(e^2=e\) 的自同态。本节先用自同态环描述不可分解性，再证明有限长度条件如何使模分解为自同态作用的幂零分量与可逆分量。

% 纲要要点（供目录核对）：
%
% * 局部自同态环的判别
% * Fitting 引理
% * 有限长度模上的幂等元分解
%

\subsection{局部自同态环的判别}
\label{b2:subsec:1203:01}

\begin{definition}\label{b2:def:indecomposable-module}
设 \(R\) 为含幺环，\(M\) 为非零幺性左 \(R\) 模。若不存在两个非零子模 \(A,B\) 使 \(M=A\oplus B\)，则称 \(M\) 为\term[bukefenjie-mo]{不可分解模}[indecomposable module]。即每次内部直和分解中，至少一个分量必须为零。
\end{definition}

简单模一定不可分解，但逆命题不成立。例如对素数 \(p\)，\(M=\mathbb Z/p^2\mathbb Z\) 的每个非零子群都包含唯一的阶 \(p\) 子群，所以两个非零子模不可能交于零。因此它不可分解，却不是简单模。这里并不否认其有两层合成列，而是说这两层不能在原模中分裂成直和。

\begin{corollary}\label{b2:prop:idempotent-direct-sum}
设 \(R\) 为含幺环，\(M\) 为非零幺性左 \(R\) 模。\(M\) 不可分解，当且仅当 \(\operatorname{End}_R(M)\) 中的幂等元只有 \(0,1\)。
\end{corollary}
\begin{proof}
\cref{b2:prop:idempotent-module-decomposition}已说明，幂等元 \(e\) 给出 \(M=\operatorname{im}e\oplus\ker e\)。若 \(e\ne0\)，像非零；若核为零，则 \(e\bigl(m-e(m)\bigr)=0\) 迫使 \(m=e(m)\) 对所有 \(m\) 成立，即 \(e=1\)。因此 \(e\ne0,1\) 恰好给出两个非零分量。反过来，两个非零分量的坐标投影显然既不为零，也不为恒等。
\end{proof}

比较两个直和分解时，会在某个分量的自同态环中得到恒等映射的有限和表达。若能从这个和中找出一个可逆项，就能比较它与另一分解中的某个分量。为此，我们需要一种环，使有限个非单位相加仍为非单位；这样它们的和便不可能是一。

\begin{definition}\label{b2:def:local-ring}
设 \(E\) 为非零含幺环。若 \(E\) 的全部非单位组成一个双边理想，则称 \(E\) 为\term[jubu-huan]{局部环}[local ring]。这里“单位”指具有两侧逆元的元素；不预设 \(E\) 交换。
\end{definition}

这里的局部性由环的非单位决定，与第六章的局部化构造是不同概念；局部化是一种使选定元素可逆的构造，并非每次局部化所得的环都是局部环。文献也常用“具有唯一极大左理想”定义局部环；在非交换情形，它与上述定义以及唯一极大右理想的条件等价。

\begin{proposition}\label{b2:prop:local-ring-criterion}
设 \(E\) 为非零含幺环。以下条件等价：\(E\) 为局部环；对每个 \(a\in E\)，\(a\) 与 \(1-a\) 中至少有一个为单位；\(E\) 有唯一极大左理想；\(E\) 有唯一极大右理想。在这些条件下，非单位组成的双边理想 \(J\) 是唯一极大左理想，也是唯一极大右理想，\(E/J\) 为除环。有限个非单位之和仍为非单位，特别地，若 \(1=x_1+\cdots+x_r\)，其中 \(r\ge1\)、\(x_i\in E\)，则至少一个 \(x_i\) 为单位。\(E\) 的幂等元只有 \(0,1\)。因此，若 \(R\) 为含幺环、\(M\) 为非零幺性左 \(R\) 模且 \(\operatorname{End}_R(M)\) 为局部环，则 \(M\) 不可分解。
\end{proposition}
\begin{proof}
若非单位构成理想 \(J\)，\(a\) 和 \(1-a\) 不能同时在 \(J\)，否则 \(1\in J\)，与一为单位矛盾。因此所述条件成立。

反过来，假设对每个 \(a\in E\)，\(a\) 或 \(1-a\) 为单位。先证明幂等元 \(e\) 只能为零或一：若 \(e\) 可逆，由 \(e^2=e\) 得 \(e=1\)；若 \(1-e\) 可逆，由 \(e(1-e)=0\) 得 \(e=0\)。这个限制还把一侧逆变成两侧逆。若 \(ab=1\)，则
\[
(ba)^2=b(ab)a=ba.
\]
幂等元 \(ba\) 不能为零，否则 \(b=bab=0\)，与 \(ab=1\) 矛盾。因此 \(ba=1\)，\(a,b\) 互为两侧逆。这一计算对只有左逆或只有右逆的情形都适用。

令 \(J\) 为非单位集合。若 \(a\in J\)、\(r\in E\)，\(ra\) 若可逆，则 \((ra)^{-1}r\) 是 \(a\) 的左逆；\(ar\) 若可逆，则 \(r(ar)^{-1}\) 是 \(a\) 的右逆。上段均会迫使 \(a\) 可逆，矛盾。所以 \(ra,ar\in J\)，并且 \(-a\in J\)。若 \(a,b\in J\) 而 \(u=a+b\) 可逆，则 \(u^{-1}a\) 与 \(1-u^{-1}a=u^{-1}b\) 都非单位，与假设矛盾。因此 \(a+b\in J\)，\(J\) 为不含一的双边理想。

对局部环的这个 \(J\)，任意非零剩余类由 \(a\notin J\) 代表，\(a\) 可逆，故 \(E/J\) 为除环。因此 \(J\) 在左右两侧均极大。任何真左理想或真右理想不能包含单位，故必包含于 \(J\)，所以它们各自的极大理想都只能是 \(J\)。

为从唯一极大左理想推出局部性，要用到每个真左理想都包含于一个极大左理想。对给定真左理想 \(L\)，考虑包含 \(L\) 的全部真左理想，按包含关系排序。它们组成非空偏序集；非空链的并仍是左理想，因为有限个元素可放进同一个链成员中检查减法与左乘，而这个并不含一，否则某个成员已经含一。空链以 \(L\) 为上界。由第一卷附录A定理A.1.4的 Zorn 引理，这个偏序集有极大元，即包含 \(L\) 的极大左理想。

现设 \(J\) 是 \(E\) 的唯一极大左理想。若 \(e^2=e\) 且 \(e\ne0,1\)，则 \(Ee\) 与 \(E(1-e)\) 都是真左理想：例如 \(re=1\) 会给出 \(e=(re)e=r e^2=1\)。由刚证的存在性，两者都包含于 \(J\)，便有 \(1=e+(1-e)\in J\)，矛盾。因此幂等元只有 \(0,1\)，前面关于 \(ba\) 的计算再次说明一侧逆必为两侧逆。

于是非单位 \(a\) 使 \(Ea\) 成为真左理想，必有 \(a\in J\)；反之 \(J\) 不含单位，所以 \(J\) 恰为非单位集合。它已对左乘封闭；若 \(a\in J\) 而 \(ar\) 可逆，则 \(a\) 有右逆，从而为单位，矛盾。因此 \(J\) 也对右乘封闭，是双边理想，\(E\) 为局部环。唯一极大右理想的情形在反环中应用同一论证。有限个非单位之和仍在理想 \(J\) 中，故不可能为单位；若它们之和为一，就至少有一项不在 \(J\) 中。幂等元结论与\cref{b2:prop:idempotent-direct-sum}还给出最后的不可分解性。
\end{proof}

例如，若 \(k\) 为域、\(n\ge1\)，则 \(k[t]/(t^n)\) 中元素可逆，当且仅当其常数项非零。确实，常数项非零的元素可写为 \(a(1-u)\)，其中 \(a\in k^\times\)、\(u\in(t)\) 幂零，逆由有限几何和给出；常数项为零的元素都在真理想 \((t)\)，不可能可逆。故非单位恰为 \((t)\)，这个环为局部环。但 \(\mathbb Z\) 不是局部环：非单位 \(2,-3\) 的和为单位 \(-1\)。虽然后者作为自身的模不可分解，反向局部性还需要下面的有限长度条件。

\subsection{Fitting 引理}
\label{b2:subsec:1203:02}

\begin{lemma}[Fitting]\label{b2:lem:fitting}
设 \(R\) 为含幺环，\(M\) 为有限长度幺性左 \(R\) 模，\(f\in\operatorname{End}_R(M)\)。对所有充分大的正整数 \(n\)，有
\[
M=\ker f^n\oplus\operatorname{im}f^n.
\]
在第一分量上 \(f\) 幂零，在第二分量上 \(f\) 可逆；两个分量都在 \(f\) 下稳定。可取任意 \(n\ge\max\bigl(1,\ell(M)\bigr)\)。
\end{lemma}
\begin{proof}
反复施加 \(f\) 产生两条链：
\[
\begin{aligned}
0&=\ker f^0\subseteq\ker f\subseteq\ker f^2\subseteq\cdots,\\
M&=\operatorname{im}f^0\supseteq\operatorname{im}f\supseteq\operatorname{im}f^2\supseteq\cdots.
\end{aligned}
\]
核记录逐次作用后变成零的元素，像记录仍能由反复作用得到的元素。有限长度同时限制严格升链与严格降链，所以两者都最终稳定。若某步 \(\ker f^r=\ker f^{r+1}\)，则对 \(f^{r+2}(x)=0\)，有 \(f(x)\in\ker f^{r+1}=\ker f^r\)，继而 \(x\in\ker f^{r+1}=\ker f^r\)。逐次归纳说明核链从此稳定。令 \(d=\ell(M)\)，从 \(\ker f^0=0\) 到 \(\ker f^d\) 若有一步相等，核链便从该步稳定；若前 \(d\) 步全都严格增加，则 \(\ell(\ker f^d)\ge d\)，所以 \(\ker f^d=M\)，此后也不再增加。因此从 \(r=d\) 起核链必稳定。由\cref{b2:prop:finite-length-endomorphism}，
\[
\ell(\operatorname{im}f^r)=\ell(M)-\ell(\ker f^r),
\]
故像链从同一位置起长度不变。相邻像有包含关系，长度相等时商长度为零，所以像也相等。

取满足上述界的 \(n\)，于是 \(\ker f^{2n}=\ker f^n\)、\(\operatorname{im}f^{2n}=\operatorname{im}f^n\)。要证明核与像交于零，考虑既经过 \(n\) 次作用、又会在 \(n\) 次作用后变成零的元素 \(y\)。写 \(y=f^n(x)\) 且 \(f^n(y)=0\)，就有 \(f^{2n}(x)=0\)。这正是比较 \(2n\) 与 \(n\) 的原因；由核稳定，\(x\in\ker f^{2n}=\ker f^n\)，从而 \(y=0\)。

对任意 \(x\in M\)，还要从像中减去一项，使余项落入核。由像稳定，\(f^n(x)\in\operatorname{im}f^{2n}\)，所以可取 \(z\in M\)，使 \(f^n(x)=f^{2n}(z)\)。减去 \(f^n(z)\) 后有
\[
f^n\bigl(x-f^n(z)\bigr)=0.
\]
因此 \(x-f^n(z)\in\ker f^n\)，而 \(f^n(z)\in\operatorname{im}f^n\)，得到所需分解。若 \(x\in\ker f^n\)，则 \(f^n(f(x))=f(f^n(x))=0\)，所以核在 \(f\) 下稳定；像也稳定，因为 \(f(\operatorname{im}f^n)=\operatorname{im}f^{n+1}=\operatorname{im}f^n\)。在核上 \(f^n=0\)；在像上，这个等式说明 \(f\) 满射，而 \(\ker f\subseteq\ker f^n\) 与已证的零交说明限制映射也单射，故可逆。
\end{proof}

这称为\term[Fitting-yinli]{Fitting 引理}[Fitting lemma]。\(\max(1,\ell(M))\) 是对所有自同态通用的指数上界，不一定是某个 \(f\) 的最小分解指数；\cref{b2:prop:fitting-integer-projections}中模的长度为五，乘二却在第三次作用时就得到分解。

\ProseParagraphBreak
对域 \(k\) 上有限维空间的算子 \(T\)，稳定核由那些被某个 \(T\) 的幂消去的向量组成，正是广义零特征子空间；稳定像是 \(T\) 可逆的部分。若最小多项式在 \(k\) 上分裂，\cref{b2:thm:jordan-canonical-form}的 Jordan 形把所有 \(J_r(0)\) 块放在第一部分，把非零特征值的块放在第二部分。Fitting 引理没有要求底环是域，也没有要求多项式分裂；有限长度提供的是核和像的同时稳定。

\begin{theorem}\label{b2:thm:indecomposable-local-end}
设 \(R\) 为含幺环，\(M\) 为非零幺性左 \(R\) 模。若 \(\operatorname{End}_R(M)\) 为局部环，则 \(M\) 不可分解。若 \(M\) 还有有限长度，则 \(M\) 不可分解，当且仅当 \(\operatorname{End}_R(M)\) 为局部环；在有限长度且不可分解情形，每个 \(f\in\operatorname{End}_R(M)\) 或者可逆，或者幂零，其全部非可逆自同态构成双边理想。
\end{theorem}
\begin{proof}
局部自同态环的幂等元只有零和一，由\cref{b2:prop:idempotent-direct-sum}，\(M\) 不可分解。这一方向没有使用长度假设。现设 \(M\) 有限长度且不可分解。对任意 \(f\)，Fitting 引理给出核与像的直和分解，不可分解性迫使
\[
\ker f^n=0\qquad\text{或}\qquad\operatorname{im}f^n=0.
\]
前一种情形使 \(f\) 单射，由\cref{b2:prop:finite-length-endomorphism}可逆；后一种情形使 \(f^n=0\)，即 \(f\) 幂零。

当 \(f\) 幂零时，\(1-f\) 可逆，因为若 \(f^n=0\)，有
\[
(1-f)(1+f+\cdots+f^{n-1})
=(1+f+\cdots+f^{n-1})(1-f)=1.
\]
因此每个 \(f\) 都使 \(f\) 或 \(1-f\) 可逆。\cref{b2:prop:local-ring-criterion}给出自同态环为局部环，同时保证非可逆元素对加法及左右乘法封闭。
\end{proof}

局部环一般只保证非单位组成理想，并不保证每个非单位都幂零。这里的幂零结论同时使用了有限长度与不可分解性；\cref{b2:prop:local-ring-nonnilpotent}将给出局部环中非幂零非单位的例子。

\ProseParagraphBreak

单有有限长度也不足以保证非可逆自同态幂零。例如对域 \(k\)，\(k^2\) 上的投影 \(\operatorname{diag}(1,0)\) 不可逆，也不幂零。相反，对循环 \(k[t]\) 模 \(k[t]/(p^n)\)，其中 \(p\) 不可约、\(n\ge1\)，任意两个非零子模均含有最小非零子模 \((p^{n-1})/(p^n)\)，所以它不可分解；其非可逆自同态因而都幂零。这与其自同态环 \(k[t]/(p^n)\) 中的非单位为 \((p)/(p^n)\) 相吻合。

\subsection{有限长度模上的幂等元分解}
\label{b2:subsec:1203:03}

\cref{b2:prop:idempotent-module-decomposition}也给出了有限版本：若自同态满足
\[
e_i^2=e_i,\qquad e_ie_j=0\ (i\ne j),\qquad
\sum_{i=1}^re_i=1,
\]
则 \(M=\bigoplus_i e_iM\)，反过来有限内部直和的坐标投影满足这些等式。现在进一步判别一个投影的像是否还能分解。

\begin{proposition}\label{b2:prop:orthogonal-idempotent-decomposition}
设 \(R\) 为含幺环，\(M\) 为幺性左 \(R\) 模。对非零幂等元 \(e\in\operatorname{End}_R(M)\)，其像 \(eM\) 不可分解，当且仅当 \(e\) 不能写成两个非零幂等元 \(u,v\) 之和且 \(uv=vu=0\)。若 \(M\) 有限长度，\(I\) 为集合，且 \((e_i)_{i\in I}\) 是 \(\operatorname{End}_R(M)\) 中满足 \(e_i\ne0\)、\(e_i^2=e_i\) 及 \(e_ie_j=0\ (i\ne j)\) 的一族元素，则 \(I\) 有限且 \(|I|\le\ell(M)\)。
\end{proposition}
\begin{proof}
若 \(e=u+v\) 如题述，则 \(eu=ue=u\)、\(ev=ve=v\)，所以两像都在 \(eM\) 中，并有 \(eM=uM\oplus vM\)，两项非零。反过来，若 \(eM=A\oplus B\) 且两项非零，先用 \(e\) 把 \(M\) 投到 \(eM\)，再投到 \(A,B\)，得到 \(u,v\)。它们分别在自身像上为恒等，在另一像及 \(\ker e\) 上为零，因而幂等、正交、非零，且 \(u+v=e\)。

对正交幂等元族，只需先看任意有限子集 \(F\subseteq I\)。若 \(x_i\in e_iM\) 且 \(\sum_{i\in F}x_i=0\)，施加 \(e_j\) 得 \(x_j=0\)，因为 \(e_j\) 在 \(e_jM\) 上为恒等，而在其余各项上为零。因此 \(\sum_{i\in F}e_iM\) 是内部直和。每个 \(e_i\ne0\) 使 \(e_iM\ne0\)，长度至少一，故由长度可加性，
\[
|F|\le\sum_{i\in F}\ell(e_iM)
=\ell\!\left(\bigoplus_{i\in F}e_iM\right)
\le\ell(M).
\]
若 \(I\) 有多于 \(\ell(M)\) 个元素，就能取到一个含 \(\ell(M)+1\) 个元素的有限子集，与此矛盾。所以 \(I\) 有限且满足所述界。整个论证只使用有限子族，没有对无限族定义幂等元之和。
\end{proof}

“正交”在这里指乘积为零，不涉及内积。若 \(M\) 有限长度，每个非零投影的像至少占一层长度，因此不能无限细分投影。如果 \(eM=A\oplus B\) 且两项非零，则 \(\ell(eM)=\ell(A)+\ell(B)\)，两个新分量的长度都严格小于原分量；继续分解时可对长度归纳。第十五章将把不能再写成两个非零正交幂等元之和的幂等元称为本原幂等元。

\ProseParagraphBreak
对内部直和 \(M=A\oplus B\)，自同态可写成由四个 Hom 群组成的分块矩阵；投影到 \(A\) 的幂等元是 \(\begin{pmatrix}1&0\\0&0\end{pmatrix}\)。在一般情形，它不在整个自同态环的中心，因为与含有非零交叉映射的分块矩阵不一定交换。模的直和分解需要的是幂等元，并不要求这些投影是中心幂等元；中心性对应更强的环论分解，后面另行讨论。

\begin{proposition}\label{b2:prop:fitting-integer-projections}
把 \(M=\mathbb Z/72\mathbb Z\) 看作幺性左 \(\mathbb Z\) 模，令 \(f(\bar x)=2\bar x\)。使 \(M=\ker f^n\oplus\operatorname{im}f^n\) 成立的最小正整数为 \(n=3\)，此时
\[
\ker f^3=\langle\bar9\rangle,\qquad
\operatorname{im}f^3=\langle\bar8\rangle.
\]
两个分量的投影分别为乘九与乘六十四，\(f\) 在第二分量上的逆为乘五。\(M\) 的长度为五。
\end{proposition}
\begin{proof}
\(f^3\) 为乘八，而 \(72=8\cdot9\)，故 \(8x\equiv0\pmod{72}\) 等价于 \(9\mid x\)。这给出所列核；像由 \(\bar8\) 生成。两分量分别有八与九个元素，交集为零，因为同时被八和九整除的整数被七十二整除。又 \(\bar9-\bar8=\bar1\)，所以它们共同生成 \(M\)，得到直和分解。这也就是中国剩余分解 \(\mathbb Z/72\mathbb Z\cong\mathbb Z/8\mathbb Z\oplus\mathbb Z/9\mathbb Z\)：乘二在前者上幂零，在后者上可逆。

乘九在 \(\langle\bar9\rangle\) 上为恒等，因为 \(9^2\equiv9\pmod{72}\)，在 \(\langle\bar8\rangle\) 上为零，因为 \(9\cdot8\equiv0\pmod{72}\)。因此它是第一分量的投影；另一分量的投影为 \(1-9=-8\equiv64\pmod{72}\)，即乘六十四。在阶九的第二分量上，乘二与乘五的复合为乘十，而 \(10\equiv1\pmod9\)，故乘五为所需逆。

当 \(n=1\) 时，\(\ker f=\langle\overline{36}\rangle\)，其中非零的 \(\overline{36}\) 也在 \(\operatorname{im}f=2M\) 中；当 \(n=2\) 时，\(\ker f^2=\langle\overline{18}\rangle\)，\(\overline{36}\) 又在 \(\operatorname{im}f^2=4M\) 中。因此前两步的核与像交集非零，最小指数确为三。最后，
\[
0\subsetneq\langle\overline{36}\rangle
\subsetneq\langle\overline{18}\rangle
\subsetneq\langle\bar9\rangle
\subsetneq\langle\bar3\rangle
\subsetneq M
\]
的相邻商依次为阶二、二、二、三、三的循环群，都是简单 \(\mathbb Z\) 模，所以这是长度五的合成列。
\end{proof}

\begin{proposition}\label{b2:prop:local-ring-nonnilpotent}
固定素数 \(p\)，令
\[
R=\mathbb Z_{(p)}
=\bigl\{\,a/b\in\mathbb Q\bigm|a,b\in\mathbb Z,\ p\nmid b\,\bigr\}.
\]
则 \(R\) 为交换局部环，非单位理想为 \(pR\)。其正则幺性左模 \({}_RR\) 不可分解，却不满足降链条件，因而不具有有限长度。非单位 \(p\in R\) 以及正则模上的乘 \(p\) 自同态都不幂零。
\end{proposition}
\begin{proof}
分母不被 \(p\) 整除的分数对加法、减法与乘法封闭，并含零与一，所以 \(R\) 是 \(\mathbb Q\) 的含幺子环，特别地为整环。对 \(a/b\in R\)，若 \(p\nmid a\)，其逆 \(b/a\) 仍在 \(R\) 中；若 \(p\mid a\)，则它在 \(R\) 中不可能有逆。否则存在 \(c/d\in R\) 满足 \(ac=bd\)，但左边被 \(p\) 整除、右边不被 \(p\) 整除，矛盾。因此非单位恰为 \(pR\)，这个集合是双边真理想，\(R\) 为局部环。由\cref{b2:prop:regular-hom-endomorphism}，\(\operatorname{End}_R({}_RR)\cong R^{\mathrm{op}}=R\)，故\cref{b2:prop:local-ring-criterion}推出正则模不可分解。

正则模的子模降链
\[
R\supsetneq pR\supsetneq p^2R\supsetneq\cdots
\]
处处严格：若 \(p^n\in p^{n+1}R\)，在分式域中消去 \(p^{n+1}\) 就得到 \(1/p\in R\)，矛盾。所以正则模不是 Artin 模，也不具有有限长度。由于 \(R\subseteq\mathbb Q\)，\(p^n\ne0\) 对每个 \(n\ge1\) 都成立；乘 \(p\) 自同态的第 \(n\) 次幂在元素一上取值为 \(p^n\)，同样非零。这给出局部自同态环中不幂零的非单位，而有限长度假设在这里恰好失败。
\end{proof}

% !TeX root = ../../Book2.tex
% 纲要标题：### 12.4 Krull–Schmidt 定理
% 纲要位置：工作记录/计划纲要.md:L1454

\section{Krull–Schmidt 定理}
\label{b2:sec:1204}

有限长度使一次非平凡直和分解的每个分量都变短，因而分解过程会终止。唯一性则更细：两个分解中的分量未必是同一个子模，要先证明某一对分量同构，再交换相应的补模，才能继续比较。局部自同态环恰好保证这一交换能够进行。

% 纲要要点（供目录核对）：
%
% * 有限长度情形的不可分解分解存在性
% * 局部自同态环、补模交换与不可分解因子的唯一性
% * 无限乘积的存在性失效及 Noether 模的非唯一分解
%
% ---
%

\subsection{有限长度模的分解存在性}
\label{b2:subsec:1204:01}

\begin{proposition}\label{b2:prop:finite-length-decomposition}
设 \(R\) 为含幺环，\(M\) 为有限长度幺性左 \(R\) 模。若 \(M\ne0\)，则 \(M\) 是有限多个非零不可分解左 \(R\) 模的直和，分量个数不超过 \(\ell(M)\)。零模对应空直和。
\end{proposition}
\begin{proof}
零模取空直和。对非零模的长度 \(\ell(M)\) 归纳。长度一时 \(M\) 简单，因而不可分解。一般地，若 \(M\) 不可分解，取它本身为唯一分量。否则可写成 \(M=A\oplus B\)，其中两项非零。由\cref{b2:prop:length-additive}，\(\ell(M)=\ell(A)+\ell(B)\)，所以两项长度都是严格小于 \(\ell(M)\) 的正整数，可以分别使用归纳假设，再将两个有限分解合并。沿每次继续分解的分量，长度都严格下降，因而分解过程不能无限继续。每个非零分量长度至少一，长度可加性还给出分量个数的上界。
\end{proof}

例如，若 \(k\) 为域，则 \(t^2\) 与 \((t-1)^3\) 互素，\(k[t]\) 模 \(k[t]/\bigl(t^2(t-1)^3\bigr)\) 通过\secref{b2:sec:0403}证明的中国剩余同构分成
\[
k[t]/(t^2)\oplus k[t]/\bigl((t-1)^3\bigr).
\]
这正是\cref{b2:thm:pid-elementary-divisors}中的两个循环主分量，分别对应不可约多项式 \(t\) 与 \(t-1\)。每个循环商的非零子模都含有其最小非零子模，分别为 \((t)/(t^2)\) 与 \(\bigl((t-1)^2\bigr)/\bigl((t-1)^3\bigr)\)，所以两个分量均不可分解，长度分别为二和三。它们各自仍有非平凡合成列：不可分解性不等于没有真子模。

\ProseParagraphBreak
分解存在性也可以表述为投影的逐步细分。若一个分量还能分解，就把对应幂等元拆成两个非零正交幂等元。\cref{b2:prop:orthogonal-idempotent-decomposition}说明这与模的操作一致；长度每次被分到两个正整数中，所以不可能无限继续。这个证明保证存在有限分解，却尚未比较不同选择得到的分量。

\subsection{不可分解因子的唯一性}
\label{b2:subsec:1204:02}

\begin{theorem}[Krull–Schmidt]\label{b2:thm:krull-schmidt}
设 \(R\) 为含幺环，\(M\) 为有限长度幺性左 \(R\) 模。若 \(M\) 有两个分解
\[
M=A_1\oplus\cdots\oplus A_r
=B_1\oplus\cdots\oplus B_s,
\]
其中各项非零且不可分解，则 \(r=s\)，并可重排下标，使 \(A_i\cong B_i\) 对所有 \(i\) 成立。因此不可分解直和分量的同构类型及重数由 \(M\) 唯一确定。
\end{theorem}
\begin{proof}
对 \(\ell(M)\) 归纳，零模只有空分解。非零时，令
\[
a_i:A_i\rightarrow M,\quad p_i:M\rightarrow A_i,\qquad
b_j:B_j\rightarrow M,\quad q_j:M\rightarrow B_j
\]
为两组包含和投影。复合 \(p_1b_jq_ja_1\) 先把 \(A_1\) 中的元素包含到 \(M\)，取出它的第 \(j\) 个 \(B\) 分量，将这个分量送回 \(M\)，最后取出其中的 \(A_1\) 分量。所有这些往返映射相加，应当恢复原元素：由 \(\sum_jb_jq_j=\operatorname{id}_M\) 和 \(p_1a_1=\operatorname{id}_{A_1}\)，在 \(\operatorname{End}_R(A_1)\) 中有
\begin{equation}\label{b2:eq:krull-schmidt-unit-sum}
\operatorname{id}_{A_1}=\sum_{j=1}^s p_1b_jq_ja_1.
\end{equation}
\(A_1\) 是非零不可分解有限长度模，故这个自同态环由\cref{b2:thm:indecomposable-local-end}为局部环。右侧只有有限项，左侧恒等自同态是单位；\cref{b2:prop:local-ring-criterion}中的有限和性质保证至少一项可逆。重排后设 \(p_1b_1q_1a_1\) 可逆。

记 \(u=q_1a_1:A_1\rightarrow B_1\)、\(v=p_1b_1:B_1\rightarrow A_1\)，于是 \(vu\) 可逆。要把 \(u(A_1)\) 分裂成 \(B_1\) 的直和因子，需要给 \(u\) 构造一个左逆。原映射 \(v\) 与 \(u\) 的复合为可逆自同态 \(vu\)，因此先用 \((vu)^{-1}\) 修正它，令
\[
w=(vu)^{-1}v:B_1\longrightarrow A_1.
\]
便有 \(wu=\operatorname{id}_{A_1}\)，所以 \(u\) 单射。幂等元 \(uw\) 在 \(B_1\) 上的像为 \(u(A_1)\)、核为 \(\ker w\)，由\cref{b2:prop:idempotent-module-decomposition}，
\[
B_1=u(A_1)\oplus\ker w.
\]
由于 \(A_1\ne0\)，第一项非零；\(B_1\) 不可分解，故 \(\ker w=0\)，上式只剩 \(B_1=u(A_1)\)。因此 \(u\) 满且单，得到 \(A_1\cong B_1\)。这里起作用的是：非零模分裂嵌入不可分解模后，其像必须占据整个目标模。

虽然 \(A_1\cong B_1\)，它们在 \(M\) 中的位置仍可能不同，尚不能据此约去两边的同构项。令 \(B'=\bigoplus_{j>1}B_j\)，要把实际子模 \(A_1\) 换入第二个分解，同时保留同一个补模 \(B'\)。因为 \(q_1|_{A_1}=u\) 为同构，对任意 \(m\in M\)，它的 \(B_1\) 坐标 \(q_1(m)\) 都对应唯一的 \(a\in A_1\)。于是 \(q_1(a)=q_1(m)\)，从而 \(m-a\in\ker q_1=B'\)，所以 \(A_1+B'=M\)。若 \(a\in A_1\cap B'\)，则 \(u(a)=0\)，所以 \(a=0\)。这就得到实际的直和分解 \(M=A_1\oplus B'\)。

令 \(A'=\bigoplus_{i>1}A_i\)。两个实际分解 \(M=A_1\oplus A'=A_1\oplus B'\) 都将其补模同构到商 \(M/A_1\)：限制商映射的核为零，且每个陪集都有补模中的代表。因此
\[
A'\cong M/A_1\cong B'.
\]
将 \(B'\) 的分解沿此同构转移到 \(A'\)，就得到同一个模上的两个不可分解分解。由长度可加性，\(\ell(A')=\ell(M)-\ell(A_1)<\ell(M)\)，故可用归纳假设，得到剩余项一一同构。加上 \(A_1\cong B_1\)，完成证明。
\end{proof}

这一\term[Krull-Schmidt-dingli]{Krull–Schmidt 定理}[Krull–Schmidt theorem]的唯一性在同构和排列意义下成立。分量作为子模的具体位置可以不同，例如 \(k^2=ke_1\oplus ke_2=ke_1\oplus k(e_1+e_2)\)，第二个分量改变了；两种分解的不可分解因子都为两个一维空间。

\begin{corollary}[有限长度模的消去]\label{b2:cor:finite-length-cancellation}
设 \(R\) 为含幺环，\(A,B,C\) 为有限长度幺性左 \(R\) 模。若 \(A\oplus C\cong B\oplus C\)，则 \(A\cong B\)。
\end{corollary}
\begin{proof}
将 \(A,B,C\) 分别分成有限个不可分解模。由\cref{b2:thm:krull-schmidt}，左右两侧每一种不可分解同构类型的出现次数相等。两侧减去来自 \(C\) 的相同有限重数，得到 \(A,B\) 的分量类型及重数相同，逐项同构的直和便给出 \(A\cong B\)。
\end{proof}

Jordan–Hölder 定理比较的是子模链中的简单商层，Krull–Schmidt 定理比较的是原模内部的不可分解直和因子。回看\subsecref{b2:subsec:1202:01}中比较过的 \(\mathbb Z/p^2\mathbb Z\) 与 \((\mathbb Z/p\mathbb Z)^2\)，其中 \(p\) 为素数：二者的合成因子都是两个 \(\mathbb Z/p\mathbb Z\)，但前者的任意非零子模都包含 \(p\mathbb Z/p^2\mathbb Z\)，所以自身不可分解；后者则已经分成两个简单直和因子。相同的简单商层不能说明它们怎样拼合，而不可分解分解会区分这两种拼合。对有限维代数的表示，二者的有限维版本可参阅 Etingof 等人\cite[Theorem 3.7.1, pp. 50--51; Theorem 3.8.1, pp. 51--52; Remark 3.8.6, p. 53]{EtingofEtAl2011RepresentationTheory}；本章用链条件与 Fitting 引理处理一般有限长度模。

\subsection{分解定理的条件与反例}
\label{b2:subsec:1204:03}

有限长度在证明中有两个作用。分解存在性用到的是非平凡分量的正整数长度严格下降，迫使分解过程终止；唯一性则用到核、像的两种链同时稳定，由 Fitting 引理得到不可分解模的自同态环局部，再从恒等映射的有限和中找到可逆项。去掉有限长度以后，这两步都可能失败，甚至有限生成或 Noether 条件也不足以替代它。

\ProseParagraphBreak
先看自同态的失效。整数模 \(\mathbb Z\) 不可分解，因为两个非零子群总有非零交；其自同态环却是非局部环 \(\mathbb Z\)。乘二的映射单射而不满，任意 \(n\ge1\) 都有 \(\ker(2^n)=0\)、\(\operatorname{im}(2^n)=2^n\mathbb Z\ne\mathbb Z\)，所以 Fitting 分解也不成立。缺少降链条件使像永远不能稳定。这个模已有以自身为唯一分量的不可分解分解；此处失去的是从不可分解性推出自同态环局部的机制。

\begin{example}[没有不可分解直和分解的循环模]\label{b2:ex:no-indecomposable-decomposition}
设 \(k\) 为域，令含幺环 \(R=\prod_{n\ge1}k\)，按坐标运算，单位元为 \((1,1,\ldots)\)。幺性左正则模 \(R\) 虽由一生成，却不能写成任何一族非零不可分解子模的内部直和。
\end{example}
\begin{proof}
任意直和因子的投影 \(p:R\rightarrow R\) 由 \(e=p(1)\) 决定，\(p(r)=re\)。投影条件给出 \(e^2=e\)，每个坐标因 \(k\) 为域只能是零或一。因此直和因子为 \(eR\)，对应一个坐标子集 \(S\)。

若 \(S\) 至少有两个元素，可分成两个非空不交子集，相应指示幂等元把 \(eR\) 分成两个非零直和因子。若 \(S\) 只有一个元素，\(eR\cong k\) 没有非零真子模，故不可分解。因此每个非零不可分解直和因子只能占据一个坐标。

若 \(R\) 是这些因子的内部直和，即使因子族无限，\cref{b2:prop:internal-direct-sum}也要求每个元素只有有限多个非零分量。因此单位元 \((1,1,\ldots)\) 必须是有限多个分量元素之和。这种和只能有有限个非零坐标，不可能等于单位元，矛盾。
\end{proof}

唯一性也可能真正失败，而不只是上述证明失去效力。以下例子给出两个不同的有限不可分解分解。

\begin{example}[Noether 模的非唯一分解]\label{b2:ex:krull-schmidt-nonunique}
令 \(s=\sqrt{-5}\)、\(R=\mathbb Z[s]\subseteq\mathbb C\)，\(I=(2,1+s)\)。以通常乘法为作用，将 \(R\) 与 \(I\) 视为幺性左 \(R\) 模。它们是不同构的不可分解 Noether 模，但
\[
R\oplus R\cong I\oplus I.
\]
\end{example}
\begin{proof}
先说明 \(I\) 是非零真理想。将整数模二、将 \(s\) 送到一，给出满同态 \(R\rightarrow\mathbb F_2\)，因为关系 \(s^2+5=0\) 在目标中成立。它将 \(2,1+s\) 送到零，故 \(I\) 真；又 \(2\in I\)，所以 \(I\ne0\)。

对 \(a+bs\in R\)，令 \(N(a+bs)=a^2+5b^2\)，这是乘共轭所得的非负整数，满足 \(N(uv)=N(u)N(v)\)。若 \(I=(u)\)，则 \(u\) 同时整除 \(2\) 和 \(1+s\)，故正整数 \(N(u)\) 同时整除四和六，只能为一或二。若为一，共轭元素就是 \(u\) 的逆，\(I=R\)，矛盾；若为二，方程 \(a^2+5b^2=2\) 无整数解。因此 \(I\) 非主，不能与 \(R\) 同构：一个同构会把一送到 \(I\) 的单个生成元。

为构造 \(R^2\) 与 \(I\oplus I\) 之间的同构，需要控制 \(I\) 中元素的乘积。这里 \(I^2\) 表示理想乘积，即由所有 \(xy\)、\(x,y\in I\) 生成的理想；\(I\oplus I\) 则以有序对为元素。计算理想 \(I^2\) 时，生成元 \(4,2(1+s),(1+s)^2=-4+2s\) 都在 \(2R\) 中。另一方面，\(s-1=(1+s)-2\in I\)，并且
\[
2=2(2\cdot2)+(1+s)(s-1),
\]
所以 \(2\in I^2\)，得到 \(I^2=2R\)。现在选一个条目都在 \(I\) 中的 \(2\times2\) 矩阵 \(B\)，它就把 \(R^2\) 送入 \(I\oplus I\)。二维伴随矩阵的条目仍在 \(I\) 中，故 \(\operatorname{adj}(B)\) 乘上 \(I\oplus I\) 中的向量后，分子各坐标都是两个 \(I\) 中元素的乘积之和，属于理想 \(I^2=2R\)。为了让逆矩阵 \((\det B)^{-1}\operatorname{adj}(B)\) 把这些向量送回 \(R^2\)，可以选择 \(\det B=2\)：这样分子中的二恰好抵消分母，所得坐标仍在 \(R\) 中。

为找到这样的 \(B\)，把刚才的 \(2=8+(s-1)(s+1)\) 组织成一个行列式：
\[
B=\begin{pmatrix}2&-(s-1)\\1+s&4\end{pmatrix},\qquad
\det B=2\cdot4+(s-1)(s+1)=2.
\]
\cref{b2:prop:adjugate-identity}给出
\[
B^{-1}=\frac12\operatorname{adj}(B)
=\frac12\begin{pmatrix}4&s-1\\-(1+s)&2\end{pmatrix}.
\]
于是对 \(x,y\in I\)，得到以下公式：
\[
\Phi(x,y)=
\left(2x+\frac{s-1}{2}y,\ -\frac{1+s}{2}x+y\right).
\]
这里的分式系数位于分式域中；但 \((s-1)y\) 与 \((1+s)x\) 都在 \(I^2=2R\) 中，所以两个除以二的表达确实属于 \(R\)，\(\Phi\) 取值于 \(R^2\)。
反向公式为
\[
\Psi(r,z)=\bigl(2r-(s-1)z,\ (1+s)r+4z\bigr).
\]
它确实取值于 \(I\oplus I\)，因为四个系数 \(2,s-1,1+s,4\) 均在 \(I\)。在分式域中写成矩阵，
\[
\begin{pmatrix}2&(s-1)/2\\-(1+s)/2&1\end{pmatrix}
\begin{pmatrix}2&-(s-1)\\1+s&4\end{pmatrix}
=\begin{pmatrix}1&0\\0&1\end{pmatrix};
\]
反序相乘也为单位矩阵，使用的恒等式为 \((s-1)(s+1)=-6\)。两映射都是分式域线性映射在相应子模上的限制，故为 \(R\) 线性；正反公式又已核对实际值域，因而给出 \(I\oplus I\cong R^2\) 的 \(R\) 模同构。

任意两个非零 \(R\) 子模 \(A,B\) 若都在 \(R\) 或 \(I\) 内，取非零 \(a\in A,b\in B\)。由于 \(a,b\) 同时属于标量环 \(R\)，有 \(ab=b\cdot a\in A\)、\(ab=a\cdot b\in B\)，且 \(R\) 为整环使 \(ab\ne0\)。因此这两个子模总有非零交，\(R,I\) 均不可分解。最后，\(R\) 作为整数模自由，基为 \(1,s\)。任意 \(R\) 子模都是 \(\mathbb Z^2\) 的子模，由\cref{b2:thm:pid-free-submodule}有限生成；一组整数生成元同时也是 \(R\) 生成元。所以 \(R\) 为 Noether 模，\(I\) 作为其子模也为 Noether 模。得到的两个分解均有限，却没有逐项同构的可能。
\end{proof}

这里的 \(R\) 与 \(I\) 都不满足 Artin 条件：\(R\supsetneq2R\supsetneq4R\supsetneq\cdots\) 严格下降；又因 \(2\in I\) 而 \(2\notin2I\)（否则 \(1\in I\)），\(I\supsetneq2I\supsetneq4I\supsetneq\cdots\) 也严格下降。因此这个反例不与有限长度情形的定理矛盾。

\ProseParagraphBreak

无限乘积环的循环正则模说明有限生成不保证不可分解分解的存在；\(R,I\) 的例子说明，即使 Noether 模已经有有限不可分解分解，其因子的同构类型也可能不唯一。Krull–Schmidt 还有适用于某些更大范围模的版本，但必须另外给出保证直和分量局部性和分解存在性的假设。


\begin{chaptersummary}
Noether 条件等价于每个子模有限生成，Artin 条件则限制降链；两者都能通过短正合列的两端判定。对非交换环，左理想与右理想给出的条件可能不同。两种链条件同时成立，才保证存在有限合成列；Jordan–Hölder 定理使其层数与简单因子重数成为模的不变量，长度可加性则把它们用于子模、商模和同态计算。

\ProseParagraphBreak
合成列并不自动分裂。不可分解只排除非平凡直和，不排除真子模或非分裂扩张；Fitting 引理使有限长度不可分解模的每个自同态或者可逆、或者幂零，从而使自同态环局部。Krull–Schmidt 唯一性利用这个局部性，从一个恒等自同态的有限和中找出可逆项，再交换对应的直和分量。结论确定不可分解类型和重数，不确定分量的具体嵌入位置。

\ProseParagraphBreak
无限乘积环的正则模说明，循环模也未必存在不可分解直和分解；\(\mathbb Z[\sqrt{-5}]\) 上的理想例子则说明，Noether 模可以有两个不同的有限不可分解分解。这些限制促使我们继续考察哪些环能使所有模或某类模的分解更简单。\chapref{b2:ch:13}的半单模把每个分量进一步要求为简单模，随后再由环的结构解释这种现象。
\end{chaptersummary}

\BookExercises

\begin{exercise}\label{b2:ex:12-vector-chain-conditions}
设 \(k\) 为域，\(V\) 为 \(k\) 向量空间。证明 \(V\) 作为 \(k\) 模满足升链条件、满足降链条件、有限维三者等价。说明仅将“向量空间”替换成“一般环上的模”后，哪些蕴含不再成立。
\end{exercise}
\begin{solution}
有限维时，严格包含使维数至少改变一，所以升链和降链都稳定。若维数无限，沿用本书的选择公理约定，逐次选择不在此前向量生成空间内的向量，得到线性无关序列 \(e_1,e_2,\ldots\)，其前 \(n\) 项的线性生成空间严格上升，其尾部 \(\operatorname{span}_k\{e_i:i\ge n\}\) 严格下降，所以两种链条件都失败。一般模上，\(\mathbb Z\) 是 Noether 而非 Artin 模，\(C_{p^\infty}\) 是 Artin 而非 Noether 模，故两种链条件不能互推。“有限维”也没有对任意环都适用的同一含义，合适的共同条件是\cref{b2:thm:finite-length-chain-conditions}中的有限长度。
\end{solution}

\begin{exercise}\label{b2:ex:12-directed-cover}
设 \(R\) 为含幺环，\(M\) 为幺性左 \(R\) 模。非空子模族 \((N_i)_{i\in I}\) 中若任意两个成员都包含于某个共同的成员，就称该族向上有向。证明 \(M\) 有限生成，当且仅当每个满足 \(M=\bigcup_iN_i\) 的向上有向子模族都含有等于 \(M\) 的一个成员。这个条件为何不足以直接说明 \(M\) 是 Noether 模？
\end{exercise}
\begin{solution}
若 \(M\) 由有限个 \(x_1,\ldots,x_r\) 生成，各元素分别落在族中某项；反复使用有向性，可找到同时包含它们的 \(N_j\)，故 \(N_j=M\)。反过来，取 \(M\) 的全部有限生成子模，它们的并为 \(M\)，任意两个的和仍有限生成，所以是向上有向族。假设迫使某个有限生成子模就是 \(M\)。这个条件只针对 \(M\) 自身的生成性；Noether 条件还要求每个子模有限生成，二者的区别由\cref{b2:prop:noetherian-submodules-fg}明确给出。
\end{solution}

\begin{exercise}\label{b2:ex:12-generator-extension}
设 \(R\) 为含幺环，\(0\rightarrow A\xrightarrow{i}B\xrightarrow{q}C\rightarrow0\) 为幺性左 \(R\) 模的短正合列。若 \(A\) 可由 \(r\) 个元素生成，\(C\) 可由 \(s\) 个元素生成，证明 \(B\) 可由至多 \(r+s\) 个元素生成，并说明在哪里使用了正合性。
\end{exercise}
\begin{solution}
取 \(A\) 的生成元 \(a_1,\ldots,a_r\)，\(C\) 的生成元 \(c_1,\ldots,c_s\)。由 \(q\) 满射，取 \(b_j\in B\) 使 \(q(b_j)=c_j\)。给定 \(b\in B\)，写 \(q(b)=\sum_j t_jc_j\)，则 \(b-\sum_jt_jb_j\in\ker q=\operatorname{im}i\)，所以又可写成 \(\sum_\nu u_\nu i(a_\nu)\)。于是 \(i(a_\nu),b_j\) 共同生成 \(B\)。满射性用于选取商生成元的原像，\(\ker q=\operatorname{im}i\) 用于把差送回 \(A\)；只知道两个映射可复合，并不足以得到这个结论。
\end{solution}

\begin{exercise}\label{b2:ex:12-one-sided-endomorphism}
设 \(R\) 为含幺环，\(M\) 为幺性左 \(R\) 模。证明 \(M\) 为 Noether 模时每个满射自同态可逆，\(M\) 为 Artin 模时每个单射自同态可逆。分别给出不能互换“满射”与“单射”的例子。
\end{exercise}
\begin{solution}
设 \(M\) Noether、\(f\) 满。取 \(n\) 使 \(\ker f^n=\ker f^{n+1}\)。若 \(x\in\ker f\)，由 \(f^n\) 满可取 \(y\) 使 \(f^n(y)=x\)。于是 \(y\in\ker f^{n+1}=\ker f^n\)，故 \(x=0\)，\(f\) 单且满。

设 \(M\) Artin、\(f\) 单。取 \(n\) 使 \(\operatorname{im}f^n=\operatorname{im}f^{n+1}\)。对 \(x\in M\)，可取 \(y\) 使 \(f^n(x)=f^{n+1}(y)\)。因 \(f^n\) 单，\(x=f(y)\)，故 \(f\) 满。整数模上的乘二单而不满，给出第一个反向失败；Artin 模 \(C_{p^\infty}\) 上的乘 \(p\) 满而不单，因为每个分式可以再除以 \(p\)，其核为非零的 \(C_p\)。
\end{solution}

\begin{exercise}\label{b2:ex:12-left-right-length}
设域扩张 \(k\subseteq K\) 的次数为有限整数 \(d\)，\(r\ge1\) 为整数。取由三元组 \((a,\mathbf u,b)\in K\times K^r\times k\) 组成的环 \(R\)，加法逐项进行，乘法为
\[
(a,\mathbf u,b)(a',\mathbf u',b')
=(aa',a\mathbf u'+\mathbf u b',bb').
\]
单位元为 \((1,\mathbf0,1)\)。分别计算左右正则模的长度，并比较 \(k=\mathbb R,K=\mathbb C,r=2\) 的情形。
\end{exercise}
\begin{solution}
这个乘法来自 \(K^r\) 的左 \(K\)、右 \(k\) 标量作用；两侧作用相容，展开三元组乘法即可验证结合律。令 \(e_1=(1,\mathbf0,0)\)、\(e_2=(0,\mathbf0,1)\)，则它们正交且和为一。由此左右正则模分别分解为 \(Re_1\oplus Re_2\) 与 \(e_1R\oplus e_2R\)。

记 \(J=\{\,(0,\mathbf u,0):\mathbf u\in K^r\,\}\)。左乘在 \(J\) 上只作用为 \(\mathbf u\mapsto a\mathbf u\)，所以其子模恰为 \(K^r\) 的 \(K\) 子空间，\(\ell({}_RJ)=r\)。\(Re_1\) 是由左上角作用的简单模 \(K\)，\(Re_2/J\) 是由右下角作用的简单模 \(k\)。长度可加性给出左正则长度 \(1+r+1=r+2\)。

右乘在 \(J\) 上只作用为 \(\mathbf u\mapsto\mathbf u b\)，其子模恰为 \(K^r\) 的 \(k\) 子空间，故 \(\ell(J_R)=rd\)。\(e_2R\) 是简单右模 \(k\)，\(e_1R/J\) 是简单右模 \(K\)，于是右正则长度为 \(rd+2\)。取实复扩张及 \(r=2\)，得到左右长度分别为四和六。相比\cref{b2:prop:triangular-left-right-length}，扩大中间的双模后，这一部分贡献的左侧简单层数为 \(r\)，右侧为 \(rd\)。
\end{solution}

\begin{exercise}\label{b2:ex:12-cyclic-length}
求整数模 \(M=\mathbb Z/360\mathbb Z\) 的长度及各简单因子的重数，并计算乘十二自同态的核与像的长度。
\end{exercise}
\begin{solution}
\(360=2^3\cdot3^2\cdot5\)。互素循环分解给出三种主分量，分别有三、二、一层简单因子，所以 \(\ell(M)=6\)，\(\mathbb Z/2,\mathbb Z/3,\mathbb Z/5\) 的重数分别为三、二、一。乘十二的核由 \(\bar{30}\) 生成，阶为十二；像由 \(\bar{12}\) 生成，阶为三十。因此两者分别同构于 \(\mathbb Z/12\mathbb Z\)、\(\mathbb Z/30\mathbb Z\)，长度均为三。长度可加性得到 \(6=3+3\)，但这两个长度三的模具有不同的简单因子重数。
\end{solution}

\begin{exercise}\label{b2:ex:12-truncated-polynomial-length}
设 \(k\) 为域，\(A=k[x,y]/(x,y)^3\)，把 \(A\) 看成自身的幺性左模。写出一条合成列，求其长度与合成因子。计算乘 \(x\) 的核和像，并核验长度可加性。
\end{exercise}
\begin{solution}
以 \(1,x,y,x^2,xy,y^2\) 的像为基，记 \(\mathfrak m=(x,y)\)。一条合成列是
\[
0\subsetneq kx^2\subsetneq kx^2+kxy\subsetneq\mathfrak m^2
\subsetneq\mathfrak m^2+kx\subsetneq\mathfrak m\subsetneq A.
\]
各项都为理想：二次部分被 \(x,y\) 零化，而乘一次部分所得都在 \(\mathfrak m^2\)。每个相邻商一维，且 \(x,y\) 作用为零，故都是简单模 \(A/\mathfrak m\cong k\)。于是长度为六。

乘 \(x\) 将 \(1,x,y\) 分别送到 \(x,x^2,xy\)，将三个二次基元素送到零。因此其像以 \(x,x^2,xy\) 为基，核为 \(\mathfrak m^2\)，两者均为三维。它们的合成因子都只能为上述 \(k\)，从列中的相应细分可得长度均为三，因此 \(6=3+3\)。
\end{solution}

\begin{exercise}\label{b2:ex:12-matrix-module-length}
设 \(k\) 为域，\(n\ge1,r\ge0\) 为整数。令 \(A=M_n(k)\)，以左矩阵乘法作用在 \(M_{n\times r}(k)\) 上。求该模的长度与一条合成列，说明长度与 \(k\) 维数的关系。
\end{exercise}
\begin{solution}
按列分解为 \(r\) 个 \(k^n\) 的直和，每个列空间都是简单左 \(A\) 模：从非零列向量的一个非零坐标出发，矩阵单位及标量可得到全部标准列向量。逐次加入前一、前二、直到全部 \(r\) 列，得到一条有 \(r\) 个简单商层的合成列。因此长度为 \(r\)，而 \(k\) 维数为 \(nr\)。只在 \(n=1\) 或零模情形，两者数值相同。
\end{solution}

\begin{exercise}\label{b2:ex:12-fitting-integers}
在整数模 \(M=\mathbb Z/600\mathbb Z\) 上令 \(f\) 为乘六。求最小的正整数 \(n\)，使 \(M=\ker f^n\oplus\operatorname{im}f^n\)。写出两个分量的投影，求 \(f\) 在可逆分量上的逆，并计算两分量的长度。
\end{exercise}
\begin{solution}
分解 \(600=24\cdot25=2^3\cdot3\cdot5^2\)。乘六在二、三主分量上幂零，在五主分量上可逆。因 \(\gcd(6^3,600)=24\)，有
\[
\ker f^3=\langle\bar{25}\rangle\cong\mathbb Z/24\mathbb Z,
\qquad
\operatorname{im}f^3=\langle\bar{24}\rangle\cong\mathbb Z/25\mathbb Z.
\]
它们交为零，且 \(25-24=1\) 保证其和为 \(M\)。当 \(n=1,2\) 时，非零元 \(\bar{300}\) 都在 \(\ker f^n\) 中；它又是六与十二的倍数，而这两数分别是 \(\gcd(6,600)\)、\(\gcd(6^2,600)\)，所以也在相应的像中。因此最小正整数为三。

乘二十五在阶二十四的分量上为恒等，在阶二十五的分量上为零，故为核分量投影；其余投影为乘 \(1-25\)，即乘五百七十六。它们是和为恒等的正交幂等元。可逆分量上乘二十一是 \(f\) 的逆，因为 \(6\cdot21\equiv1\pmod{25}\)。两分量的长度分别为 \(3+1=4\) 与二，合为 \(\ell(M)=6\)。
\end{solution}

\begin{exercise}\label{b2:ex:12-fitting-polynomial}
设 \(k\) 为域，\(p,q\in k[t]\) 为互不相伴的首一不可约多项式，\(a,b\ge1\) 为整数。在
\[
M=k[t]/(p^a)\oplus k[t]/(q^b)
\]
上令 \(f\) 为乘 \(p\)。求其 Fitting 分解，并求使分解成立的最小正整数。
\end{exercise}
\begin{solution}
第一分量上 \(f^a=0\)，而 \(f^{a-1}\ne0\)；第二分量上 \(p\) 与 \(q^b\) 互素，Bézout 等式给出乘 \(p\) 的逆。因此 \(n=a\) 时，核为第一分量，像为第二分量，得到分解。

若 \(1\le n<a\)，第一分量的核为 \((p^{a-n})/(p^a)\)，像为 \((p^n)/(p^a)\)，两者交含非零的 \(p^{a-1}+(p^a)\)，因为 \(a-1\ge n,a-n\)。所以此时不能为直和，最小整数正是 \(a\)。第二分量的指数 \(b\) 不影响这个最小值，因为 \(f\) 在其上从一开始就可逆。
\end{solution}

\begin{exercise}\label{b2:ex:12-local-not-finite-length}
设 \(k\) 为域，\(R=k[\![t]\!]\) 为形式幂级数环。证明 \(a(t)=\sum_{n\ge0}a_nt^n\) 是单位当且仅当 \(a_0\ne0\)，并在单位情形递推求出逆级数的系数。证明 \(R\) 局部，其左正则模有限生成且不可分解，却不有限长度；给出不幂零的非单位。
\end{exercise}
\begin{solution}
若 \(a(t)b(t)=1\)，常数项满足 \(a_0b_0=1\)，所以 \(a_0\ne0\)。反之，令
\[
b_0=a_0^{-1},\qquad
b_n=-a_0^{-1}\sum_{i=1}^n a_i b_{n-i}\quad(n\ge1).
\]
每一项只使用已确定的有限个系数，使乘积的常数项为一、其他系数为零，故得到形式幂级数逆。这里不涉及级数的解析收敛。

非单位恰为常数项为零的级数，即真理想 \(tR\)，所以 \(R\) 局部。左正则模由一生成；其自同态环由右乘法识别为 \(R^{\mathrm{op}}=R\)，因而局部，由\cref{b2:prop:local-ring-criterion}中幂等元判据及\cref{b2:prop:idempotent-direct-sum}得不可分解。另一方面，\(R\supsetneq tR\supsetneq t^2R\supsetneq\cdots\) 严格：\(t^n\) 的 \(n\) 次系数非零，而 \(t^{n+1}R\) 中该系数全为零。因此它不 Artin，也不有限长度。非单位 \(t\) 的每次幂都非零，故不幂零。
\end{solution}

\begin{exercise}\label{b2:ex:12-jordan-centralizer}
设 \(k\) 为域，\(n\ge1\)，\(J=J_n(0)\in M_n(k)\) 为单个上次对角线为一的幂零 Jordan 块。证明与 \(J\) 交换的全部矩阵恰为
\[
a_0I+a_1J+\cdots+a_{n-1}J^{n-1}.
\]
证明这个环局部，并确定其中的幂等元。
\end{exercise}
\begin{solution}
向量 \(e_n\) 在 \(J\) 下循环，因为 \(e_n,Je_n,\ldots,J^{n-1}e_n\) 是逆序标准基。代入映射 \(P(t)\mapsto P(J)e_n\) 因而给出 \(k[t]/(t^n)\cong k^n\)：\(J^n=0\)，而上述基保证次数小于 \(n\) 的非零多项式不在核中。特别地，\(\mu_J(t)=t^n\)，\(k^n\) 在作用代数 \(k[J]\) 下同构于其正则模。由\cref{b2:thm:regular-double-centralizer}，正则左作用的交换子来自右乘；\(k[J]\) 交换，右乘就是左乘，所以 \(\operatorname{Cent}(J)=k[J]\)。上述基也保证 \(I,J,\ldots,J^{n-1}\) 的线性组合表示唯一。

常数项 \(a_0\ne0\) 时，提出 \(a_0\) 后用幂零矩阵的有限几何和求逆；常数项为零时，矩阵为 \(J\) 的倍数，因而幂零、非可逆。非单位构成理想 \((J)\)，所以环局部。由\cref{b2:prop:local-ring-criterion}，幂等元只有零和单位矩阵。
\end{solution}

\begin{exercise}\label{b2:ex:12-nilpotent-sum}
设 \(k\) 为域。在 \(M_2(k)\) 中给出两个幂零矩阵，其和可逆。说明为什么这不违背有限长度不可分解模的局部自同态环结论。
\end{exercise}
\begin{solution}
取 \(U=E_{12}\)、\(V=E_{21}\)，有 \(U^2=V^2=0\)，但 \((U+V)^2=I\)，故其和可逆，任意特征下都成立。它们是 \(k^2\) 的自同态，而这个模可分解为两个一维子空间。定理要求有限长度与不可分解同时成立；在满足这两个条件的同一个模上，非可逆自同态才组成理想，幂零元素之和才由该理想性继续非可逆并因 Fitting 判据而幂零。
\end{solution}

\begin{exercise}\label{b2:ex:12-graph-projection}
设 \(R\) 为含幺环，\(M=A\oplus B\) 为幺性左 \(R\) 模的直和，\(u:A\rightarrow B\) 为 \(R\) 模同态。证明 \(u\) 的图
\(\Gamma_u=\Bigl\{\,\bigl(a,u(a)\bigr)\mid a\in A\,\Bigr\}\)
是以 \(B\) 为补模的直和因子，求对应投影，并证明它与投影到 \(A\) 的幂等元在 \(\operatorname{End}_R(M)\) 中共轭。
\end{exercise}
\begin{solution}
每个 \((a,b)\) 唯一写成 \(\bigl(a,u(a)\bigr)+\bigl(0,b-u(a)\bigr)\)，故 \(M=\Gamma_u\oplus B\)。对应投影为
\[
e_u=\begin{pmatrix}1&0\\u&0\end{pmatrix}.
\]
令 \(g=\begin{pmatrix}1&0\\u&1\end{pmatrix}\)，则 \(g^{-1}=\begin{pmatrix}1&0\\-u&1\end{pmatrix}\)，并且
\[
g\begin{pmatrix}1&0\\0&0\end{pmatrix}g^{-1}=e_u.
\]
所以不同图子模虽然作为原模中的子模位置不同，其投影通过这个明确自同构共轭，图本身也都同构于 \(A\)。
\end{solution}

\begin{exercise}\label{b2:ex:12-power-cancellation}
设 \(R\) 为含幺环，\(M,N\) 为有限长度幺性左 \(R\) 模，\(r\ge1\) 为整数。证明 \(M^{\oplus r}\cong N^{\oplus r}\) 蕴含 \(M\cong N\)。这里不能未经说明就逐次消去，因为两侧并未给定相同的直和因子。
\end{exercise}
\begin{solution}
先把 \(M,N\) 分解为有限个不可分解模。对每一种出现的同构类型，设其重数分别为 \(a,b\)。两边的 \(r\) 次直和中相应重数为 \(ra,rb\)，由\cref{b2:thm:krull-schmidt}必相等。正整数 \(r\) 可在整数等式中约去，所以 \(a=b\)。全部类型的重数都相同，逐项同构的直和便给出 \(M\cong N\)。这里约去的是整数重数，不是预先假定的模消去律。
\end{solution}

\begin{exercise}\label{b2:ex:12-finite-extension-descent}
设 \(L/k\) 为有限域扩张，\(A\) 为结合含幺 \(k\) 代数，\(M,N\) 为有限维幺性左 \(A\) 模。若 \(L\otimes_kM\cong L\otimes_kN\) 作为 \(L\otimes_kA\) 模，证明 \(M\cong N\) 作为 \(A\) 模。可对照有限扩张情形的证明方法\cite[Problem 3.8.4(i), pp. 52--53]{EtingofEtAl2011RepresentationTheory}。
\end{exercise}
\begin{solution}
令 \(d=[L:k]\)，选择 \(L\) 的 \(k\) 基 \(\lambda_1,\ldots,\lambda_d\)。限制为 \(A\) 模后，映射
\[
M^{\oplus d}\longrightarrow L\otimes_kM,\qquad
(m_1,\ldots,m_d)\longmapsto\sum_i\lambda_i\otimes m_i
\]
是同构：张量积的基展开给出唯一坐标，\(A\) 只作用在第二因子，所以映射 \(A\) 线性。对 \(N\) 同样成立，题设因而给出 \(M^{\oplus d}\cong N^{\oplus d}\)。两模有限维，其子模链由维数控制，故有限长度。分解为不可分解模后，Krull–Schmidt 定理使每一种类型的重数满足 \(d a=d b\)，故 \(a=b\)，得到 \(M\cong N\)。有限扩张条件用于把限制标量后的模写成有限个相同分量。
\end{solution}

\begin{exercise}\label{b2:ex:12-infinite-cancellation}
设 \(k\) 为域，\(k\) 向量空间 \(V\) 有可数无限基 \(e_0,e_1,\ldots\)。构造 \(k\oplus V\cong V\)，据此说明去掉有限长度以后，\(A\oplus C\cong B\oplus C\) 不再保证 \(A\cong B\)。
\end{exercise}
\begin{solution}
把左侧 \(k\) 的标准基向量送到 \(e_0\)，把 \(V\) 中的 \(e_i\) 送到 \(e_{i+1}\)，得到两组基之间的双射，线性延拓即为同构，逆映射将 \(e_0\) 送回首个分量，将其余基向量依次前移。取 \(A=k,B=0,C=V\)，则 \(A\oplus C\cong B\oplus C\)，但 \(k\not\cong0\)。无限公共分量能够吸收新增的一维空间，整数重数的有限消去论证不再适用。
\end{solution}

\begin{exercise}\label{b2:ex:12-nonprincipal-projective}
对\cref{b2:ex:krull-schmidt-nonunique}中的 \(R=\mathbb Z[\sqrt{-5}]\)、\(I=(2,1+\sqrt{-5})\)，证明 \(I\) 是有限生成投射模，却不是自由模。
\end{exercise}
\begin{solution}
\(I\) 由题述两个元素生成。正文给出的显式同构 \(I\oplus I\cong R^2\) 说明 \(I\) 是自由模的直和因子，由\cref{b2:thm:projective-characterizations}，它投射。若 \(I\) 自由，因为非零，其基非空。任意两个非零元素 \(a,b\in I\) 都满足 \(R\) 线性关系 \(b\cdot a+(-a)\cdot b=0\)。这里 \(a,b\) 首先是模元素，又因 \(I\subseteq R\) 而能作标量；系数 \(b,-a\) 均非零，所以二者不能同时属于一组基。基因而至多有一个元素。单元素基却意味着 \(I\) 为主理想，与正文已经证明的非主性矛盾。因此它不自由。
\end{solution}

\begin{exercise}\label{b2:ex:12-dual-numbers-decomposition}
设 \(k\) 为域，\(A=k[\varepsilon]/(\varepsilon^2)\)，\(M\) 为五维幺性左 \(A\) 模，乘 \(\varepsilon\) 的线性算子秩为二。确定 \(M\) 的不可分解直和分解及长度。
\end{exercise}
\begin{solution}
记算子为 \(T\)，则 \(T^2=0\)。由\cref{b2:thm:jordan-canonical-form}，其 Jordan 块只能为大小一或二的零特征值块，因为大小至少三的块平方不为零。每个二阶块贡献秩一，每个一阶块贡献秩零，所以恰有两个二阶块和一个一阶块。二阶块作为 \(A\) 模同构于正则模 \(A\)，一阶块同构于 \(A/(\varepsilon)\cong k\)，故
\[
M\cong A\oplus A\oplus k.
\]
\(A\) 的非零子模都包含 \(k\varepsilon\)，所以不可分解，\(k\) 简单也不可分解。长度分别为二、二、一，总长度为五，三个不可分解因子的个数则为三。
\end{solution}

\begin{exercise}\label{b2:ex:12-orthogonal-family-bound}
设 \(R\) 为含幺环，\(M\) 为长度 \(d>0\) 的幺性左 \(R\) 模。若 \(\operatorname{End}_R(M)\) 中有 \(d\) 个两两正交的非零幂等元 \(e_1,\ldots,e_d\)，证明它们的和自动为恒等，每个 \(e_iM\) 都简单。反之，若 \(M\) 是简单模的有限直和，证明正交非零幂等元的个数可以达到 \(d\)。比较整数模 \(\mathbb Z/p^2\mathbb Z\) 与 \((\mathbb Z/p\mathbb Z)^2\) 能达到的最大个数，其中 \(p\) 为素数。
\end{exercise}
\begin{solution}
令 \(e=\sum_i e_i\)。由正交性，\(e^2=e\)，且 \(eM=\bigoplus_i e_iM\)。长度可加性给出
\[
d\le\sum_{i=1}^d\ell(e_iM)=\ell(eM)\le\ell(M)=d.
\]
所以每项长度为一，因而简单；\(\ell(M/eM)=0\) 则迫使 \(eM=M\)。幂等元在自身像上为恒等，故 \(e=1_M\)。

若 \(M\) 为 \(s\) 个简单模的直和，其长度为 \(s=d\)，各分量投影便给出 \(d\) 个正交非零幂等元。对 \(\mathbb Z/p^2\mathbb Z\)，任何非零子模都包含其唯一的阶 \(p\) 子模；两非零直和因子不可能交于零，故它不可分解，自同态环只有零和恒等两个幂等元，最大个数为一。\((\mathbb Z/p\mathbb Z)^2\) 的两个坐标投影达到二，且\cref{b2:prop:orthogonal-idempotent-decomposition}的长度界禁止更多。这两个模的简单商层相同，但实际可分离的投影数不同。
\end{solution}

% !TeX root = ../../Book2.tex
% 纲要标题：## 第13章　半单模与 Artin–Wedderburn 理论
% 纲要位置：计划纲要.md，第 1462 行。

\chapter{半单模与 Artin–Wedderburn 理论}
\label{b2:ch:13}
\BookChapterNumber{b2:ch:13}

\begin{chapterabstract}
半单模可以分解为简单模的直和，等价地，它的每个子模都有补模。本章先证明这一判据及子商稳定性，再由正则模的有限简单分解建立 Artin–Wedderburn 定理。证明通过简单模之间的同态得到矩阵环，并逐步核对左模约定中的反环以及分解的唯一性。最后讨论有限维代数和群代数的基本例子，以乘法群的类方程和分圆多项式证明有限除环的交换性。
\end{chapterabstract}

\begin{learninggoals}
\begin{itemize}
\item 证明简单模直和、简单子模之和及所有子模有补三个条件等价，并处理无限直和的选择步骤。
\item 用 Schur 引理计算简单类型之间的 Hom，区分除环结论与代数闭域上的标量结论。
\item 从正则模推导矩阵环乘积分解，证明参数唯一性，并正确处理左右模与反环。
\item 比较半单性与半本原性，计算有限维和有限半单环的简单模及维数，证明 Wedderburn 小定理。
\end{itemize}
\end{learninggoals}

向量空间的每个子空间都有补空间，但 \(\mathbb Z/4\mathbb Z\) 中的子群 \(2\mathbb Z/4\mathbb Z\) 没有补子群：任何阶为二的子群仍是它自身。它的合成列有两个 \(\mathbb Z/2\mathbb Z\) 商因子，却不能把这两个因子拆成直和。\cref{b2:thm:module-jordan-holder}确定的是逐层出现的简单商及其重数；本章的半单性则要求简单模已作为原模内部的直和分量存在。二者的区别在于连接这些商层的扩张是否分裂。

\ProseParagraphBreak

从这个补模条件出发，可以先在模上辨认半单性，再问一个环何时使自身作为模半单。对于除环 \(D\)，矩阵环 \(M_n(D)\) 的左正则模可按列分解为 \(n\) 个同构的简单列模。一般半单环的分类要从正则模的简单类型及重数恢复这些矩阵块，并弄清每块所用的除环；这需要计算自同态环，不能只比较维数。

\ProseParagraphBreak

本章使用模的商、直和、正合列及自同态环；有限长度的术语沿用第十二章。无限直和的补模构造使用第一卷定理 A.1.4 的 Zorn 引理。Wedderburn 小定理还用到第一卷的有限群类方程和分圆多项式的整数因式分解。

\ProseParagraphBreak
有限维表示中的 Schur 引理和半单分解可参阅 Etingof 等人的《Introduction to Representation Theory》\cite[Proposition 2.3.9 and Corollary 2.3.10, pp. 11--12; Proposition 3.1.4 and Remark 3.1.5, pp. 42--43]{EtingofEtAl2011RepresentationTheory}；该书\cite[Proposition 3.5.8, p. 48]{EtingofEtAl2011RepresentationTheory}给出代数闭域上有限维代数的对应结构定理。本章先处理一般环与任意大小的半单模，再在有限维代数的应用中加上基域假设。

% !TeX root = ../../Book2.tex
% 纲要标题：### 13.1 简单模与半单模
% 纲要位置：计划纲要.md，第 1464 行。

\section{简单模与半单模}
\label{b2:sec:1301}

一个模可以有很短的合成列，却仍不能把各个商因子直接拆开。半单性要求更强：组成模本身的简单部分已经作为直和存在，所有子模都能找到补模。本节允许无限直和，不预设模有限生成。环仍含单位元，模仍为幺性左模；简单模总要求非零。

% 纲要要点（供目录核对）：
%
% * 简单对象和简单模的直和
% * 半单模的子模与商模
% * Schur 引理及除环；逆命题与有限维假设的反例
%

\subsection{简单对象和简单模的直和}
\label{b2:subsec:1301:01}

\begin{definition}\label{b2:def:simple-module}
设 \(R\) 为含幺环，所论左 \(R\) 模均为幺性模。非零左 \(R\) 模 \(S\) 只有 \(0,S\) 两个子模时，称 \(S\) 为简单模\termidx[jiandan-mo]{简单模}。左 \(R\) 模 \(M\) 可写为一族简单子模的内部直和
\[
M=\bigoplus_{i\in I}S_i,
\]
时，称 \(M\) 为\term[bandan-mo]{半单模}[semisimple module]。
\label{b2:def:semisimple-module}
允许 \(I=\varnothing\)，因此零模半单，但不是简单模。
\end{definition}

例如，域 \(k\) 上的简单模是一维向量空间，每个向量空间选基后都是一维空间的直和，所以半单。对素数 \(p\)，整数模 \(\mathbb Z/p\mathbb Z\) 简单，而短正合列
\[
0\longrightarrow\mathbb Z/p\mathbb Z\xrightarrow{\bar a\mapsto\overline{pa}}
\mathbb Z/p^2\mathbb Z\longrightarrow\mathbb Z/p\mathbb Z\longrightarrow0
\]
不分裂。右端被 \(p\) 零化，所以从右端到中项的任何同态，其像都落在中项被 \(p\) 零化的子模 \(p\mathbb Z/p^2\mathbb Z\) 中；商映射在这个子模上为零，故该同态不可能是截面。这个唯一的非零真子模没有补模；它也是中项唯一的简单子模，所以中项不能写成简单子模的直和，\(\mathbb Z/p^2\mathbb Z\) 不是半单模。它的两个简单合成因子是在商层中出现的，并非模内已经存在的两个简单直和分量。

\ProseParagraphBreak
模范畴中的简单对象就是上述简单模。给定 \(0\ne s\in S\)，子模 \(Rs\) 非零，故 \(Rs=S\)。因此每个简单模都循环，且形如 \(R/L\)，其中 \(L\) 是极大左理想：满射 \(R\to S\)、\(r\mapsto rs\) 的核为 \(L\)，中间子模对应迫使 \(L\) 极大。反过来，极大左理想的商模只有零子模和全模，因而简单。这里 \(L\) 只要求左理想，不必是双边理想，商对象通常只是模。

\ProseParagraphBreak
例如，令 \(k\) 为域、\(R=M_2(k)\)，并取第一列为零的矩阵组成的集合 \(L\)。左乘任意矩阵仍使第一列为零，所以 \(L\) 是左理想。映射 \(R\to k^2\) 取矩阵的第一列，是满射左 \(R\) 模同态，核恰为 \(L\)。列模 \(k^2\) 简单：给定非零列向量 \(v\) 与任意 \(w\)，选 \(v_j\ne0\)，令矩阵的第 \(j\) 列为 \(w/v_j\)、另一列为零，就有 \(Av=w\)。因此 \(R/L\cong k^2\) 简单，\(L\) 是极大左理想。但是，若 \(E_{ij}\) 表示第 \((i,j)\) 个位置为 \(1\)、其余位置为零的矩阵，则 \(E_{12}\in L\)，而 \(E_{12}E_{21}=E_{11}\notin L\)；故 \(L\) 不是双边理想，\(R/L\) 不能以代表元乘法成为商环。

\ProseParagraphBreak
“简单子模之和”暂时不要求和为直和。以下定理说明，只要简单子模的和已经覆盖整个模，总可以从中抽出一个直和分解；同时，它给出不依赖所选简单分解的半单性判据。

\begin{theorem}[半单模的等价刻画]\label{b2:thm:semisimple-module-characterizations}
设 \(R\) 为含幺环，\(M\) 为幺性左 \(R\) 模。下列条件等价：
\begin{equivalences}
\item \(M\) 是简单子模的直和。
\item \(M\) 是其简单子模之和。
\item 对每个子模 \(N\subseteq M\)，存在子模 \(K\) 使 \(M=N\oplus K\)。
\end{equivalences}
\end{theorem}
\begin{proof}
第一条立即蕴含第二条。设 \(M=\sum_{i\in I}S_i\)，其中每个 \(S_i\) 简单，固定 \(N\subseteq M\)。为了给 \(N\) 找补模，考察所有子集 \(J\subseteq I\)，使 \(\sum_{j\in J}S_j\) 是直和且与 \(N\) 交为零，按包含排序。空集满足条件。链的并仍满足这两个条件。确实，任何有限个指标都同属链中的某一项，所以这些分量之间若有和为零的关系，直和性使每个分量为零；这些分量之和若落入 \(N\)，这一项与 \(N\) 零交又使该和为零。由 Zorn 引理，取极大的 \(J\)，令 \(K=\bigoplus_{j\in J}S_j\)。

若 \(N+K\ne M\)，某个 \(S_i\) 不包含于 \(N+K\)。交 \(S_i\cap(N+K)\) 是简单模 \(S_i\) 的子模，只能为零或整个 \(S_i\)；后一情形已被排除，故此交为零。于是 \(K+S_i\) 仍为直和；若 \(k+s\in N\)，其中 \(k\in K,s\in S_i\)，则 \(s\in S_i\cap(N+K)=0\)，继而 \(k\in K\cap N=0\)。所以可将 \(i\) 加入 \(J\)，保持两个条件，与极大性矛盾。因此 \(M=N\oplus K\)，得到第三条。特别地，取 \(N=0\) 便已经得到第一条。

最后假设第三条。要把简单商层变成模内的简单子模，先注意补模条件传给每个子模 \(U\)：若 \(L\subseteq U\) 且 \(M=L\oplus C\)，将 \(u\in U\) 写成 \(u=l+c\) 后，由 \(l\in L\subseteq U\) 得 \(c\in U\cap C\)，所以 \(U=L\oplus(U\cap C)\)。取非零子模 \(U\) 中的 \(0\ne x\)。循环模 \(Rx\) 的真子模按包含排序；任一链的并仍不含生成元 \(x\)，故仍为真子模。Zorn 引理给出极大真子模 \(L\)，于是 \(Rx/L\) 简单。由传给 \(Rx\) 的补模条件，\(Rx=L\oplus S\)，商映射限制为同构 \(S\to Rx/L\)。这样，原先只在商模中出现的简单模被实现为 \(U\) 内的非零简单子模 \(S\)。

令 \(N\) 为 \(M\) 的全部简单子模之和。如果 \(N\ne M\)，取非零补模 \(K\)。刚才的结论给出 \(K\) 的简单子模 \(S\)，但 \(S\subseteq N\) 又使 \(S\subseteq N\cap K=0\)，矛盾。因此 \(N=M\)，得到第二条，完成等价性证明。
\end{proof}

\subsection{半单模的子模与商模}
\label{b2:subsec:1301:02}

\begin{proposition}\label{b2:prop:semisimple-subquotients}
设 \(R\) 为含幺环，所论左 \(R\) 模均为幺性模。半单左 \(R\) 模的每个子模与商模都半单，任意一族半单左 \(R\) 模的直和也半单。半单左 \(R\) 模有限生成，当且仅当它是有限多个简单子模的直和。对半单左 \(R\) 模 \(M\) 的子模 \(N\)，短正合列
\[
0\longrightarrow N\longrightarrow M\longrightarrow M/N\longrightarrow0
\]
分裂。
\end{proposition}
\begin{proof}
由\cref{b2:thm:semisimple-module-characterizations}，写 \(M=N\oplus K\)。取一个简单分解 \(M=\bigoplus_iS_i\)，投影 \(p:M\to N\) 将每个 \(S_i\) 映为零或简单模：限制映射的核为 \(0\) 或 \(S_i\)，非零像因此同构于 \(S_i\)。又 \(N=p(M)=\sum_i p(S_i)\)，由等价刻画得 \(N\) 半单。这里得到的是投影后的简单子模，不必是原来 \(S_i\) 中某些分量的和。同理投影到 \(K\) 说明 \(K\) 半单，而商映射限制为同构 \(K\to M/N\)，故商模半单，其逆后接 \(K\hookrightarrow M\) 就给出商映射的截面。

若各 \(M_a\) 都已分解为简单模直和，把全部指标并在一起，仍是直接和：任一元素只有有限多个非零外层分量，每个分量又只有有限多个非零内层分量。因此所有简单分量合成所需直和分解。

若 \(m_1,\ldots,m_t\) 生成半单模 \(M=\bigoplus_{i\in I}S_i\)，它们的支集之并是有限集 \(F\subseteq I\)。每个 \(S_i\) 是子模，故对生成元作标量乘法及求和不会产生 \(F\) 之外的直和坐标。因此 \(M=\bigoplus_{i\in F}S_i\)，只含有限多个简单分量。反过来，每个简单模由任一非零元素循环生成，有限个简单模的直和便由这些元素有限生成；零模对应空直和。
\end{proof}

两端半单不保证中项半单：前面的 \(\mathbb Z/p^2\mathbb Z\) 正是两个简单模之间不分裂的扩张。合成因子的名单没有记录这类扩张是否拆开。另一方面，命题中的有限生成判据将把半单环的简单分解限制为有限直和。

\subsection{Schur 引理及除环}
\label{b2:subsec:1301:03}

\begin{lemma}[Schur]\label{b2:lem:schur}
设 \(R\) 为含幺环，\(S,T\) 为幺性简单左 \(R\) 模。任何非零同态 \(f:S\to T\) 都是同构。因此若 \(S\not\cong T\)，则 \(\operatorname{Hom}_R(S,T)=0\)；而
\[
D=\operatorname{End}_R(S)
\]
是除环，其乘法为映射复合。
\end{lemma}
\begin{proof}
\(\ker f\) 是 \(S\) 的子模；因 \(f\ne0\)，它不能为 \(S\)，所以为零。\(\operatorname{im}f\) 是 \(T\) 的非零子模，所以为 \(T\)。故 \(f\) 双射；模同态的逆仍是模同态。取 \(S=T\)，每个非零自同态都可逆，恒等映射因 \(S\ne0\) 而非零，所以自同态环是除环。
\end{proof}

这一结果称为\term[Schur-yinli]{Schur 引理}[Schur's lemma]。结论是除环，不能在一般环上把它直接写成标量域。例如除环 \(D\) 的左正则模简单，而\cref{b2:prop:regular-hom-endomorphism}给出其自同态环为 \(D^{\mathrm{op}}\)。对含幺 \(k\) 代数 \(A\) 的非零幺性模 \(S\)，标量映射给出 \(k\) 到 \(\operatorname{End}_A(S)\) 中心的嵌入，但这个嵌入未必满射。例如 \(\mathbb C\) 作为实代数作用在自身上，也是简单模，其自同态环为 \(\mathbb C\)，大于原标量域 \(\mathbb R\)。

\begin{corollary}\label{b2:cor:schur-algebraically-closed}
设 \(k\) 为代数闭域、\(A\) 为结合含幺 \(k\) 代数，\(S\) 为非零、有限 \(k\) 维的幺性简单左 \(A\) 模，则
\(\operatorname{End}_A(S)=k\operatorname{id}_S\)。
\end{corollary}
\begin{proof}
取 \(f\in\operatorname{End}_A(S)\)。由于 \(S\) 有限维且 \(k\) 代数闭，\(f\) 的特征多项式有根 \(\lambda\in k\)，所以 \(f-\lambda\operatorname{id}_S\) 有非零核。该映射仍 \(A\) 线性，因为 \(k\) 的作用在中心。由\cref{b2:lem:schur}，它不可能是非零同态，故 \(f=\lambda\operatorname{id}_S\)。反过来，每个这样的标量映射当然与 \(A\) 作用相容。
\end{proof}

例如，若 \(S,T\) 非同构且简单，自同态 \(f:S\oplus T\to S\oplus T\) 的四个分块分别是两个对角自同态及 \(S\to T\)、\(T\to S\) 的映射。后两者由 Schur 引理为零，所以分块形式为
\[
\begin{pmatrix}a&0\\0&b\end{pmatrix},\qquad
a\in\operatorname{End}_R(S),\quad b\in\operatorname{End}_R(T),
\]
自同态环因此是 \(\operatorname{End}_R(S)\times\operatorname{End}_R(T)\)。如果改为 \(S\oplus S\)，四个分块都属于 \(D=\operatorname{End}_R(S)\)，映射可以写为
\[
f(s_1,s_2)=\bigl(a_{11}(s_1)+a_{12}(s_2),\,
a_{21}(s_1)+a_{22}(s_2)\bigr),\qquad a_{ij}\in D.
\]
复合时先施加后一映射，再施加前一映射，四个条目按矩阵乘法组合，故 \(\operatorname{End}_R(S\oplus S)\cong M_2(D)\)。这一区别正是 Artin–Wedderburn 分解中“不同块”与“同一块内的矩阵”各自出现的原因。

\begin{proposition}\label{b2:prop:schur-hypothesis-counterexamples}
幺性整数模 \(\mathbb Q\) 不是简单模，但 \(\operatorname{End}_{\mathbb Z}(\mathbb Q)\cong\mathbb Q\) 是除环。另设 \(k\) 为代数闭域、\(t\) 为 \(k\) 上的不定元，\(A=k(t)\)。左正则模 \(S={}_AA\) 是幺性简单左 \(A\) 模，具有无限 \(k\) 维数，且
\[
\operatorname{End}_A(S)\cong k(t)\supsetneq k.
\]
因此自同态环为除环不蕴含模简单，而代数闭基域上的简单模若没有有限维假设，其自同态也未必全是基域标量。
\end{proposition}
\begin{proof}
若 \(f:\mathbb Q\to\mathbb Q\) 为 \(\mathbb Z\) 线性，令 \(q=f(1)\)。对 \(m\in\mathbb Z\) 和正整数 \(n\)，有 \(nf(m/n)=f(m)=mq\)，故 \(f(m/n)=mq/n\)。因此 \(f\) 恰为乘以 \(q\) 的映射。反过来，乘以任意 \(q\in\mathbb Q\) 的映射都 \(\mathbb Z\) 线性，加法与复合分别对应有理数的加法与乘法，于是得到环同构 \(\operatorname{End}_{\mathbb Z}(\mathbb Q)\cong\mathbb Q\)。但 \(\mathbb Z\) 是 \(\mathbb Q\) 的非零真子模，所以 \(\mathbb Q\) 不简单。

因为 \(A=k(t)\) 是域，任意非零左理想含有一个可逆元，从而为全环，故其左正则模简单。向量 \(1,t,t^2,\ldots\) 在 \(k\) 上线性无关，所以 \(S\) 无限维。任意 \(A\) 线性自同态满足 \(f(a)=af(1)\)，而任意乘以 \(c\in A\) 的映射都 \(A\) 线性；由于 \(A\) 交换，这给出 \(\operatorname{End}_A(S)\cong A\)。其中乘以 \(t\) 的映射不是 \(k\) 标量映射：若它等于 \(\lambda\operatorname{id}_S\)，在 \(1\) 上求值就会得到 \(t=\lambda\in k\)，矛盾。事实上该映射没有 \(k\) 中的特征值，因为 \((t-\lambda)a=0\) 对每个 \(\lambda\in k\) 都迫使 \(a=0\)。这正说明前一推论中有限维假设保证特征值存在的作用。
\end{proof}

% !TeX root = ../../Book2.tex
% 纲要标题：### 13.2 半单环
% 纲要位置：计划纲要.md，第 1470 行。

\section{半单环}
\label{b2:sec:1302}

半单模可以定义在任意环上；半单环则要求环自身的正则模已经半单。因为每个模都由自由模取商得到，这个看似只涉及一个模的要求，实际上会迫使整个模范畴中的子模都能分裂。它还蕴含强烈的有限性，不能与仅仅“没有非零元素同时零化所有简单模”混淆。

% 纲要要点（供目录核对）：
%
% * 正则模半单的等价条件
% * 所有模的投射性与内射性
% * 半单性不等于半本原性
%

\subsection{正则模半单的等价条件}
\label{b2:subsec:1302:01}

\begin{definition}\label{b2:def:semisimple-ring}
设 \(R\) 为含幺环。若左正则模 \({}_RR\) 是半单模，就称 \(R\) 为\term[bandan-huan]{半单环}[semisimple ring]。零环的正则模为零，按此约定也是半单环。\secref{b2:sec:1303}将证明该定义与右正则模半单等价。
\end{definition}

\begin{theorem}\label{b2:thm:semisimple-ring-characterizations}
设 \(R\) 为含幺环，所论左 \(R\) 模均为幺性模。下列条件等价：
\begin{equivalences}
\item 左正则模 \({}_RR\) 半单。
\item 每个左理想都是 \({}_RR\) 的直和因子。
\item 每个左 \(R\) 模都是半单模。
\item 每条左 \(R\) 模的短正合列都分裂。
\end{equivalences}
若这些条件成立，\({}_RR\) 是有限多个简单左理想的直和，因而具有有限长度，特别是左 Noether 且左 Artin。
\end{theorem}
\begin{proof}
第一、二条的等价是\cref{b2:thm:semisimple-module-characterizations}应用于 \({}_RR\)，因为正则模的子模恰是左理想。假设第一条，要从这个正则模推广到任意左模 \(M\)，可先把它的生成元放进自由模：自由模 \(R^{(I)}\) 是若干份 \({}_RR\) 的直和，由\cref{b2:prop:semisimple-subquotients}仍半单。任意左模 \(M\) 都有满射 \(R^{(I)}\to M\)，例如按 \(M\) 的元素取自由生成元并将其送到对应元素。再用同一命题中商模保持半单的结论，得到第三条。第三条当然包含第一条。

若第三条成立，短正合列 \(0\to A\xrightarrow{i}B\xrightarrow{q}C\to0\) 中，\(i(A)\) 在 \(B\) 内有补模 \(K\)。\(q|_K\) 单射是因为 \(K\cap\ker q=0\)，满射是因为每个 \(b\) 可写成 \(i(a)+k\)。其逆后接 \(K\hookrightarrow B\) 便是截面，所以序列分裂。反过来，将第四条用于每条 \(0\to N\to M\to M/N\to0\)，由\cref{b2:thm:split-exact-sequence}得到每个子模都有补模，故每个 \(M\) 半单。

最后，\(R\) 由 \(1\) 生成，\cref{b2:prop:semisimple-subquotients}中的有限生成判据便使简单分量只有有限多个。具体地，若 \(R=\bigoplus_{i\in I}S_i\)，则 \(1=\sum_{i\in F}s_i\) 只涉及有限集 \(F\subseteq I\)。每个 \(S_i\) 是左理想，所以对任意 \(r\in R\)，有 \(r=r1=\sum_{i\in F}rs_i\in\bigoplus_{i\in F}S_i\)。因此这有限个分量已覆盖全部 \(R\)。逐个加上它们得到有限合成列，故 \({}_RR\) 有限长度。由\cref{b2:thm:finite-length-chain-conditions}，左理想的升链与降链都稳定，得到两项有限性结论。
\end{proof}

在此意义下，“半单 Artin 环”中的左 Artin 条件已由半单性推出。对域 \(k\)，无限直积环 \(R=\prod_{i\ge1}k\) 不会仅因每个因子是域就自动成为半单环：其正则模含有严格下降的左理想链
\[
R\supsetneq\{x:x_1=0\}\supsetneq
\{x:x_1=x_2=0\}\supsetneq\cdots.
\]
每一步都可由仅在相应坐标为 \(1\) 的元素见出严格包含。定理保证半单环必左 Artin，所以这个环不是半单环。它的各个坐标理想虽是简单模，但它们的直和只包含有限支集的元素，不能覆盖整个正则模。

\ProseParagraphBreak
这里讨论的是无限直积环在自身乘法下的正则模，不能据此断言同一个加法群在别的标量作用下也不半单。例如 \(\prod_{i\ge1}\mathbb F_2\) 作为整数模被 \(2\) 零化，因而是一个 \(\mathbb F_2\) 向量空间。选取向量空间基后，它就是一维 \(\mathbb F_2\) 空间的直和，每个分量作为整数模都同构于简单模 \(\mathbb Z/2\mathbb Z\)，故这个整数模半单。这里选取的基不是各个坐标向量组成的集合；后者只生成有限支集的向量。半单性取决于底环及其作用。

\subsection{所有模的投射性与内射性}
\label{b2:subsec:1302:02}

\begin{proposition}\label{b2:prop:semisimple-projective-injective}
设 \(R\) 为含幺环，所论左 \(R\) 模均为幺性模。若 \(R\) 半单，则每个左 \(R\) 模既投射又内射。反过来，若每个左 \(R\) 模都投射，或每个左 \(R\) 模都内射，则 \(R\) 半单。
\end{proposition}
\begin{proof}
设 \(R\) 半单。给定满射 \(q:B\to C\)，其核所确定的短正合列分裂，故有截面 \(s:C\to B\)。对任意 \(f:P\to C\)，映射 \(sf:P\to B\) 是提升，说明任意 \(P\) 投射。给定单射 \(i:A\to B\)，分裂性给出缩回 \(r:B\to A\)；任意 \(f:A\to E\) 都由 \(fr:B\to E\) 延拓，说明任意 \(E\) 内射。

反过来，若所有模都投射，对每条短正合列的商模使用提升性质，将其恒等映射提升为截面，所以所有短正合列都分裂。若所有模都内射，则将子模的恒等映射沿包含延拓，得到缩回，也使每条短正合列分裂。两种情形均由\cref{b2:thm:semisimple-ring-characterizations}推出半单性。
\end{proof}

这里要求“每个左模”都具有性质。给定模 \(M\) 的半单性只保证 \(M\) 自身的子模有补模；投射性却要求对任意满射提升从 \(M\) 出发的同态，内射性要求沿任意单射延拓以 \(M\) 为目标的同态，条件涉及其他模。因而一个给定模半单，不意味着它在原环上投射或内射。例如对素数 \(p\)，简单整数模 \(\mathbb Z/p\mathbb Z\) 不是投射模，因为 \(\mathbb Z\to\mathbb Z/p\mathbb Z\) 没有截面：它的任何截面都会把非零有限阶元素嵌入无挠模 \(\mathbb Z\)。它也不是内射模，因为从 \(p\mathbb Z\) 到它的同态 \(p\mapsto\bar1\) 不能延拓到 \(\mathbb Z\)：任何延拓都会把 \(p\) 送到 \(p\) 倍的某个元素，即零。

\ProseParagraphBreak
半单环上的直和分解通常不唯一到具体子模。例如 \(k^2\) 中 \(ke_1\) 的补空间可取任意 \(k(e_2+ae_1)\)。半单性保证补模存在，结构定理则描述简单类型及其重数；这两者都不要求为每个子模指定唯一补模。

\subsection{半单性与半本原性}
\label{b2:subsec:1302:03}

每个简单模 \(S\) 给出一个双边理想 \(\operatorname{Ann}_R(S)\)，由同时零化 \(S\) 全部元素的标量组成。把这些理想取交，可以检测一个环的所有简单模共同看不见哪些元素。

\ProseParagraphBreak
必须区别零化整个 \(S\) 与零化一个选定向量。对 \(0\ne s\in S\)，满射 \(R\to S\)、\(r\mapsto rs\) 的核 \(L_s=\{\,r\in R\mid rs=0\,\}\) 是极大左理想，而 \(\operatorname{Ann}_R(S)=\bigcap_{s\ne0}L_s\) 要求所有向量同时被零化。即使 \(S\cong R/L_s\)，也不能一般地写成 \(\operatorname{Ann}_R(S)=L_s\)。前面的 \(M_2(k)\) 列模就是例子：对 \(s=(1,0)^{\mathsf t}\)，\(L_s\) 是第一列为零的矩阵组成的非零左理想；但一个矩阵若零化全部列向量，其每一列都为零，故 \(\operatorname{Ann}_{M_2(k)}(k^2)=0\)。

\ProseParagraphBreak
所有简单模的同构类可以用一个集合表示：每个简单模都同构于某个 \(R/L\)，而极大左理想组成 \(R\) 的幂集中的一个集合。从这些商模中每个同构类选一个代表，便得到所需代表族；同构的模具有相同的零化子。

\begin{proposition}\label{b2:prop:simple-annihilator-intersection}
设 \(R\) 为非零含幺环，\(\Sigma\) 为幺性简单左 \(R\) 模的一组同构类代表。有
\[
\bigcap_{S\in\Sigma}\operatorname{Ann}_R(S)
=\bigcap_{L\text{ 为 }R\text{ 的极大左理想}}L.
\]
\end{proposition}
\begin{proof}
若 \(a\) 零化每个简单模，对每个极大左理想 \(L\)，商模 \(R/L\) 简单，故 \(a(1+L)=0\)，即 \(a\in L\)。这证明左边包含于右边。

反过来，设 \(a\) 属于全部极大左理想。固定任意简单模 \(S\)，但让 \(0\ne s\in S\) 遍历其所有非零向量。对每个这样的 \(s\)，满射 \(R\to S\)、\(r\mapsto rs\) 的核 \(L_s\) 都是极大左理想，因而包含 \(a\)，所以 \(as=0\)。这里可能需要随 \(s\) 改换 \(L_s\)，但 \(a\) 属于全部极大左理想的假设对每个 \(s\) 都适用。零向量也被零化，于是 \(a\in\operatorname{Ann}_R(S)\)，得到反向包含。
\end{proof}

若含幺环 \(R\) 的全部幺性简单左模的零化子之交为零，则称 \(R\) 具有\term[banbenyuan-xing]{半本原性}[semiprimitivity]。这就是说，简单模的作用族共同忠实：任何非零 \(a\in R\) 都在某个简单模的某个向量上作用非零；单个简单模的作用则未必忠实。第十四章将把这个交作为 Jacobson 根研究。

\ProseParagraphBreak
半单环一定满足这个条件。若 \(a\) 零化所有简单模，就零化左正则模的每个简单直和分量，因而零化整个 \(R\)；特别是 \(a=a1=0\)。但反过来不成立：\(\mathbb Z\) 的简单模为全部 \(\mathbb Z/p\mathbb Z\)，零化子为 \(p\mathbb Z\)，且
\[
\bigcap_{p\text{ 素数}}p\mathbb Z=0.
\]
若非零整数 \(a\) 属于此交，取 \(|a|+1\) 的任一素因子 \(p\)，它不整除 \(a\)，立即矛盾。但短正合列
\[
0\longrightarrow2\mathbb Z\longrightarrow\mathbb Z
\longrightarrow\mathbb Z/2\mathbb Z\longrightarrow0
\]
不分裂：若 \(2\mathbb Z\) 有补模 \(K\)，则 \(K\cong\mathbb Z/2\mathbb Z\)，而 \(\mathbb Z\) 中没有非零有限阶元素，不可能有这样的 \(K\)。因此 \(\mathbb Z\) 不半单。半本原性保证非零环元素能被某个简单作用检测，却没有保证子模与商模之间的扩张分裂。

\ProseParagraphBreak
零环没有非零简单模。若把空族理想的交按约定取为全环，上述交在零环中仍是零，半单性与半本原性在这个退化情形中均成立。后续结构定理将把零环对应于没有矩阵因子的空乘积，避免把它误列为除环。

% !TeX root = ../../Book2.tex
% 纲要标题：### 13.3 Artin–Wedderburn 定理
% 纲要位置：计划纲要.md，第 1476 行。

\section{Artin–Wedderburn 定理}
\label{b2:sec:1303}

半单环的正则模已分成有限多个简单部分。接下来考虑这些部分之间的全部同态：不同简单类型之间没有非零同态，同一类型的多个副本之间则组成矩阵。由于左正则模的自同态由右乘给出，最后必须通过反环把这个矩阵描述转回原环。

% 纲要要点（供目录核对）：
%
% * 半单环的矩阵块分解
% * 单 Artin 环与除环
% * 分解的唯一性、中心幂等元及左右约定
%

\subsection{半单环的矩阵块分解}
\label{b2:subsec:1303:01}

\begin{theorem}[Artin–Wedderburn]\label{b2:thm:artin-wedderburn}
设 \(R\) 为含幺环。\(R\) 半单，当且仅当存在有限多个除环 \(D_1,\ldots,D_t\) 与正整数 \(n_1,\ldots,n_t\)，使
\begin{equation}\label{b2:eq:artin-wedderburn}
R\cong\prod_{i=1}^tM_{n_i}(D_i).
\end{equation}
零环对应 \(t=0\) 的空乘积。每个这样的环也右半单，且左右均为 Artin 环和 Noether 环。
\end{theorem}
\begin{proof}
设 \(R\ne0\) 半单。由\cref{b2:thm:semisimple-ring-characterizations}，将正则模的有限简单分解按同构类型合并，写成
\[
{}_RR\cong\bigoplus_{i=1}^t S_i^{n_i},
\]
其中 \(S_i\) 两两非同构，\(n_i>0\)。令 \(E_i=\operatorname{End}_R(S_i)\)，由\cref{b2:lem:schur}它是除环。不同类型间的 Hom 为零：从 \(S_j^{n_j}\) 到 \(S_i^{n_i}\) 的映射在各副本间的坐标映射都为零。因此每个自同态保持每个 \(S_i^{n_i}\)，不同类型之间没有非对角块，而各块内的映射可以独立选取，得到
\[
\operatorname{End}_R({}_RR)
\cong\prod_i\operatorname{End}_R(S_i^{n_i})
\cong\prod_iM_{n_i}(E_i).
\]
第二个同构是在固定各副本与 \(S_i\) 的识别后，把自同态写成矩阵，条目表示“从第 \(j\) 个副本到第 \(i\) 个副本”的映射。若两个自同态的条目分别为 \(f_{ij},g_{ij}\)，先施加 \(g\)、再施加 \(f\)，从第 \(\ell\) 个副本到第 \(i\) 个副本的条目为 \(\sum_jf_{ij}\circ g_{j\ell}\)。这恰好是以复合作为 \(E_i\) 乘法的通常矩阵乘法，也解释了为何同一简单类型的副本形成一个矩阵块。

由\cref{b2:prop:regular-hom-endomorphism}，每个左线性自同态是某个右乘映射 \(r\mapsto ra\)；先右乘 \(b\)、再右乘 \(a\)，结果为右乘 \(ba\)。因此
\[
R^{\mathrm{op}}\cong\operatorname{End}_R({}_RR)
\cong\prod_iM_{n_i}(E_i).
\]
要回到原环 \(R\)，先对整个式子取反环，得到
\[
R\cong\prod_iM_{n_i}(E_i)^{\mathrm{op}}.
\]
这里得到的仍是矩阵环的反环。对任意除环 \(E\) 和 \(n\ge1\)，用转置把它写回通常的矩阵形式：
\[
\Phi:M_n(E)^{\mathrm{op}}\longrightarrow M_n(E^{\mathrm{op}}),
\qquad A^{\mathrm{op}}\longmapsto A^{\mathsf T}.
\]
来源反转整个矩阵的乘法，目标则反转每个条目的乘法；转置同时交换行列，才能把求和中的索引对应起来。确实，设 \(A=(a_{ij})\)、\(B=(b_{ij})\) 属于 \(M_n(E)\)，来源中 \(A^{\mathrm{op}}B^{\mathrm{op}}=(BA)^{\mathrm{op}}\)，目标中的乘积逐条目满足
\[
\bigl(A^{\mathsf T}\cdot_{E^{\mathrm{op}}}B^{\mathsf T}\bigr)_{i\ell}
=\sum_j a_{ji}\cdot_{\mathrm{op}}b_{\ell j}
=\sum_j b_{\ell j}a_{ji}
=(BA)_{\ell i}.
\]
这正是 \((BA)^{\mathsf T}\) 的第 \((i,\ell)\) 项。转置还保持加法与单位，并且再转置就恢复原矩阵，所以 \(\Phi\) 为含幺环同构。取 \(D_i=E_i^{\mathrm{op}}\)，便得到题述的原环分解。

反过来，设 \(D\) 为除环、\(n\ge1\)、\(A=M_n(D)\)。右乘矩阵单位 \(E_{jj}\) 保留第 \(j\) 列而清去其余列，因此按列得到左模直和
\[
{}_AA=AE_{11}\oplus\cdots\oplus AE_{nn}.
\]
每个左理想 \(AE_{jj}\) 由只有第 \(j\) 列可能非零的矩阵组成，提取该列便给出它与列模 \(D^n\) 的左 \(A\) 模同构。以标准列为基，\(D^n\) 是系数写在右侧的右 \(D\) 向量空间，而 \(A\) 通过矩阵乘法从左作用。这个左 \(A\) 模简单：若 \(v\ne0\)，选 \(v_j\ne0\)；要把 \(v\) 送到任意列向量 \(w\)，可读取第 \(j\) 坐标，用其逆消去 \(v_j\)，再把各 \(w_i\) 写入目标坐标。具体取只有第 \(j\) 列可能非零的矩阵 \(B\)，令 \(b_{ij}=w_i v_j^{-1}\)，则
\[
(Bv)_i=(w_i v_j^{-1})v_j=w_i.
\]
因此任何非零向量都生成全模。乘法次序不能改成 \(v_j^{-1}w_i\)，因为除环不必交换。故每个矩阵环左半单；有限乘积的正则模按各因子分解，仍为有限个简单左模的直和，其余因子在该分量上作用为零。

矩阵环的右正则模则按行分解为 \(E_{11}A\oplus\cdots\oplus E_{nn}A\)。对非零行向量 \(v\)，取 \(v_j\ne0\)；给定目标行 \(w\)，令 \(B\) 只有第 \(j\) 行可能非零，且 \(b_{j\ell}=v_j^{-1}w_\ell\)，便有 \((vB)_\ell=v_j(v_j^{-1}w_\ell)=w_\ell\)。所以各行理想是简单右模，有限乘积也右半单。两侧正则模都只有有限个简单分量，因而具有有限长度；分别把它们看作左 \(R\) 模与左 \(R^{\mathrm{op}}\) 模，\cref{b2:thm:finite-length-chain-conditions}给出两侧的升链与降链条件。零环的情形各项结论直接成立。
\end{proof}

这一定理称为\term[Artin-Wedderburn-dingli]{Artin–Wedderburn 定理}[Artin--Wedderburn theorem]。它把所有半单环归结为除环及其有限矩阵环。这些除环可以非交换；有限维的实代数就可能含有四元数除环这一类型。此前半单环只按左正则模定义，右半单性是在得到矩阵结构后由行分解证明的结论。

\subsection{单 Artin 环与除环}
\label{b2:subsec:1303:02}

\begin{definition}\label{b2:def:simple-ring}
设 \(R\) 为非零含幺环。若 \(R\) 只有 \(0,R\) 两个双边理想，就称为\term[jiandan-huan]{单环}[simple ring]。它与左正则模 \({}_RR\) 简单是不同条件：后者要求只有 \(0,R\) 两个左理想，因而更强。
\end{definition}

\begin{proposition}\label{b2:prop:simple-artinian}
非零含幺环 \(R\) 是单环且左 Artin，当且仅当 \(R\) 同构于某个 \(M_n(D)\)，其中 \(D\) 为除环、\(n\ge1\)。这些环也右 Artin。特别地，单半单环恰为单个除环上的全矩阵环。
\end{proposition}
\begin{proof}
先设 \(A=M_n(D)\)。上一定理已证明 Artin 性。若双边理想 \(I\) 含非零矩阵 \(B\)，取 \(b_{ij}\ne0\)。左右乘矩阵单位可以把这个非零条目搬到任意位置 \((k,\ell)\)：
\[
E_{ki}BE_{j\ell}=b_{ij}E_{k\ell}\in I.
\]
再左乘对角矩阵 \(b_{ij}^{-1}I_n\)，把该条目化成一，得到 \(E_{k\ell}\in I\)。对每个 \(a\in D\)，继续左乘 \(aI_n\)，还得到 \(aE_{k\ell}\in I\)。任意矩阵是这些矩阵的有限和，所以 \(I=A\)，证明单性。

反过来，设 \(R\) 单且左 Artin。非零左理想的降链条件保证存在极小非零左理想，它作为左模简单。要利用单性，我们还需要由这些简单左理想构造一个非零双边理想，所以令 \(U\) 为全部极小左理想之和。它首先是非零左理想。对简单左理想 \(L\) 和 \(r\in R\)，右乘映射
\[
\rho_r:L\longrightarrow R,\qquad x\longmapsto xr
\]
为左 \(R\) 线性，因为 \((ax)r=a(xr)\)。它的核是 \(L\) 的子模，简单性使核为零或全体 \(L\)；因此像 \(Lr\) 或者为零，或者由第一同构定理同构于 \(L\)，仍是简单左理想。故 \(Lr\subseteq U\)，对 \(U\) 中的有限和也有 \(Ur\subseteq U\)。于是 \(U\) 为非零双边理想，由单性 \(U=R\)。\cref{b2:thm:semisimple-module-characterizations}说明 \({}_RR\) 半单，再由\cref{b2:thm:artin-wedderburn}得到有限矩阵环乘积。若乘积有两个以上非零因子，其中一个因子便给出非零真双边理想，与单性矛盾。因此只剩一个因子。
\end{proof}

例如 \(M_n(D)\) 在 \(n>1\) 时是单环，却不是除环：\(E_{11}\) 非零而不可逆，且 \(AE_{11}\) 是非零真左理想。相反，若非零含幺环 \(R\) 的左正则模本身简单，对任意 \(0\ne a\in R\) 有 \(Ra=R\)，于是存在 \(b\) 使 \(ba=1\)。再对 \(b\ne0\) 取 \(c\) 使 \(cb=1\)，则 \(c=cba=a\)，所以 \(ab=1\)。因此每个非零元都有双侧逆，\(R\) 是除环。反过来，除环的每个非零左理想都含有单位，因而为全环，所以其左正则模简单。

\ProseParagraphBreak
上述证明还说明：在单环中，只要存在一个极小非零左理想，就已经可以推出矩阵环结构；左 Artin 条件提供了保证其存在的充分条件。随后用单性把全部极小左理想的和扩大为全环，环的双边理想结构与左模结构分别承担了不同步骤。

\subsection{分解的唯一性及左右约定}
\label{b2:subsec:1303:03}

一个半单分解可以改变简单分量的具体位置，却不应改变每种简单类型的副本数。固定简单模 \(S\)，从 \(S\) 到其他简单类型的同态为零，而到一个 \(S\) 副本的同态给出一个自同态除环。因此 \(\operatorname{Hom}_R(S,M)\) 可以检测这种类型的副本数。其标量来自来源 \(S\) 的自同态，作用方式是先改变输入，再施加同态，即预复合。

\begin{proposition}[半单分解中的重数]\label{b2:prop:semisimple-multiplicity}
设 \(R\) 为含幺环、\(M\) 为有限生成半单幺左 \(R\) 模、\(S\) 为简单幺左 \(R\) 模，\(E=\operatorname{End}_R(S)\)，其乘法为映射复合。则 \(\operatorname{Hom}_R(S,M)\) 通过以下运算成为右 \(E\) 向量空间（\(h\in\operatorname{Hom}_R(S,M),\ e\in E\)）：
\[
h\cdot e=h\circ e.
\]
任一简单直和分解中，\(S\) 的同构副本出现次数等于以下右维数：
\[
\dim_E\operatorname{Hom}_R(S,M).
\]
特别地，简单类型及各自重数不依赖所选分解。
\end{proposition}
\begin{proof}
由复合的结合律，
\[
(h\cdot e)\cdot e'=h\circ e\circ e'=h\cdot(ee'),
\]
而恒等映射给出单位作用，复合对加法分配，所以确为右作用。取简单分解 \(M=\bigoplus_{j\in J}T_j\)。由于 \(S\) 简单，任意 \(0\ne s\in S\) 都生成 \(S\)。给定 \(h:S\to M\)，\(h(s)\) 只含有限多个非零坐标，左标量作用不会增加这些坐标；由 \(h(rs)=rh(s)\)，整个 \(h(S)\) 都在同一个有限子直和中。因此各坐标映射组成有限支集的族，取坐标投影给出右 \(E\) 线性同构
\[
\operatorname{Hom}_R(S,M)\cong\bigoplus_{j\in J}\operatorname{Hom}_R(S,T_j).
\]
反向把有限多个坐标映射通过包含求和即可。这里即使 \(J\) 无限，得到的也是直和，而不是直积，原因在于来源 \(S\) 有限生成。又由于 \(M\) 有限生成，由\secref{b2:sec:1301}的有限支集论证知 \(J\) 有限。若 \(T_j\not\cong S\)，该项由 Schur 引理为零；若 \(u_j:S\to T_j\) 是同构，则每个同态唯一写为 \(u_j\circ e\)，所以该项是一维右 \(E\) 空间。维数恰为副本数。反过来，\(S\) 出现当且仅当这个 Hom 非零，因此出现的简单类型也由 \(M\) 本身决定。
\end{proof}

这里除环上的有限维数仍然唯一。基交换论证只需要非零系数可逆：若右线性展开含 \(v_j a\)、\(a\ne0\)，就在右侧乘 \(a^{-1}\) 解出 \(v_j\)，然后逐个替换生成元。一组无关向量不能多于一组有限生成元，交换两组基便得大小相同；不需要交换两个标量的次序。

\ProseParagraphBreak
将上一命题用于半单环的正则模，\(\operatorname{Hom}_R(S,R)\) 的右维数便恢复类型 \(S\) 的重数，而该类型还决定系数除环
\[
D_S=\operatorname{End}_R(S)^{\mathrm{op}}.
\]
这两种数据取自 \(R\) 的模与同态，不依赖矩阵坐标。不同矩阵因子则可以由环内部的中心幂等元区分。

\begin{proposition}\label{b2:prop:matrix-centers-central-idempotents}
设 \(D\) 为除环、\(n\ge1\)，则
\[
Z\bigl(M_n(D)\bigr)=Z(D)I_n.
\]
更一般地，设 \(t\ge0\)、\(D_1,\ldots,D_t\) 为除环、\(n_1,\ldots,n_t\ge1\)，并令 \(R=\prod_{i=1}^tM_{n_i}(D_i)\)，则
\[
Z(R)=\prod_{i=1}^t Z(D_i)I_{n_i}.
\]
对 \(1\le i\le t\)，令 \(c_i\in R\) 的第 \(i\) 坐标为 \(I_{n_i}\)，其余坐标为零。\(R\) 的中心幂等元恰为 \(\sum_{i\in F}c_i\)，其中 \(F\subseteq\{1,\ldots,t\}\)，总数为 \(2^t\)。按 \(e\le f\iff ef=e\) 比较中心幂等元时，极小非零者恰为 \(c_1,\ldots,c_t\)。\(t=0\) 时 \(R\) 为零环，只有幂等元零，没有非零中心幂等元。
\end{proposition}
\begin{proof}
设 \(C=(c_{ij})\in Z(M_n(D))\)。与各 \(E_{jj}\) 交换，迫使第 \(j\) 行、第 \(j\) 列的非对角元都为零，因此 \(C\) 为对角矩阵。再与每个 \(E_{ij}\) 交换，得到各对角元相同，故 \(C=dI_n\)。对任意 \(a\in D\)，与 \(aE_{11}\) 交换又给出 \(da=ad\)，所以 \(d\in Z(D)\)。反过来，\(d\in Z(D)\) 时，\(dI_n\) 与每个矩阵逐条目交换，证明第一个等式。这些论证在 \(n=1\) 时同样成立。

乘积环中的元素与全部元素交换，当且仅当每个坐标在相应因子的中心，得到第二个等式。\(Z(D_i)\) 是域：非零中心元的逆仍与每个元素交换。故其中的幂等元只能为零或一，中心幂等元便是在各坐标上分别选择零或单位矩阵，恰对应子集 \(F\)，共有 \(2^t\) 个。两个这样的幂等元相乘对应子集的交，所以 \(e\le f\) 对应子集包含；极小非空子集是单点集，给出各 \(c_i\)。空乘积时仅有空子集，对应零环中的唯一元素零，仍有 \(2^0=1\) 个幂等元。
\end{proof}

\begin{corollary}\label{b2:cor:artin-wedderburn-uniqueness}
设 \(R\) 为含幺环，且有两个含幺环分解
\[
R\cong\prod_{i=1}^tM_{n_i}(D_i)
\cong\prod_{j=1}^sM_{m_j}(F_j),
\]
其中 \(t,s\ge0\)，\(D_i,F_j\) 为除环，\(n_i,m_j\ge1\)。则 \(t=s\)，并可重排下标，使每个 \(i\) 都满足 \(n_i=m_i\) 及 \(D_i\cong F_i\)（含幺环同构）。零环的两个分解均为空乘积。
\end{corollary}
\begin{proof}
若 \(R=0\)，任何除环上的正阶矩阵环都非零，所以只能有 \(t=s=0\)。以下设 \(R\ne0\)。把两种乘积的各坐标单位通过同构送到 \(R\)，分别得到 \(c_1,\ldots,c_t\) 与 \(d_1,\ldots,d_s\)。由\cref{b2:prop:matrix-centers-central-idempotents}，它们分别列出了 \(R\) 的全部极小非零中心幂等元。因此 \(t=s\)，且重排后 \(c_i=d_i\)；各因子正是这些幂等元切出的环 \(c_iR\)，其单位为 \(c_i\)。

对任意幺左 \(R\) 模 \(T\)，各 \(c_iT\) 是子模，因为 \(c_i\) 位于中心。由 \(\sum_i c_i=1\) 及 \(c_ic_j=0\ (i\ne j)\)，有
\[
T=\bigoplus_{i=1}^t c_iT.
\]
若 \(T\) 简单，就只能有一个非零分量，该分量等于 \(T\)，其余因子在 \(T\) 上作用为零。因此不同矩阵因子的简单类型不会混合。对 \(A=M_n(D)\)，前面的按列分解给出 \({}_AA\cong(D^n)^n\)。每个简单左 \(A\) 模 \(T\) 都由任意非零元生成，因而是 \(A\) 的商；相应满射 \(A\to T\) 在某个列理想上的限制非零，Schur 引理使 \(T\cong D^n\)。故一个矩阵因子恰提供一种简单类型。

还要从这个类型恢复 \(D\)。令 \(e_i\) 为 \(D^n\) 的标准列，\(f\in\operatorname{End}_A(D^n)\)。由 \(E_{11}f(e_1)=f(E_{11}e_1)=f(e_1)\)，有 \(f(e_1)=e_1a\)，其中 \(a\in D\)。对任意列 \(v\)，取第一列为 \(v\)、其余列为零的矩阵 \(B\)，则 \(Be_1=v\)，左 \(A\) 线性给出
\[
f(v)=f(Be_1)=Bf(e_1)=B(e_1a)=va.
\]
反过来，每个右乘映射 \(r_a(v)=va\) 都左 \(A\) 线性，因为 \(B(va)=(Bv)a\)。又 \(r_a\circ r_b=r_{ba}\)，故
\[
\operatorname{End}_A(D^n)\cong D^{\mathrm{op}}.
\]
于是该简单类型的自同态除环再取反环，恰好恢复原来的 \(D\)。

对因子 \(c_iR\)，两种分解给出同一个简单左 \(R\) 模类型 \(S_i\)；因其余因子作用为零，左 \(R\) 线性与左 \(c_iR\) 线性在 \(S_i\) 上是同一条件。因此其自同态环在第一种描述下同构于 \(D_i^{\mathrm{op}}\)，在第二种描述下同构于 \(F_i^{\mathrm{op}}\)，取反环得到 \(D_i\cong F_i\)。按列分解又说明 \(S_i\) 在正则模中的重数分别为 \(n_i\) 与 \(m_i\)。\({}_RR\) 有限生成且半单，\cref{b2:prop:semisimple-multiplicity}使这两个重数都等于 \(\operatorname{Hom}_R(S_i,R)\) 的右 \(\operatorname{End}_R(S_i)\) 维数，因此 \(n_i=m_i\)。
\end{proof}

这里所谓唯一性是参数及因子的同构类型唯一，不是矩阵单位或具体同构映射唯一。改变同一简单类型的副本识别，会改变矩阵坐标。采用左模约定时，自同态除环必须再取反环才得到原环的矩阵系数；若采用右模约定，相应次序随之改变，不能在公式中无说明地省去 \(\mathrm{op}\)。

% !TeX root = ../../Book2.tex
% 纲要标题：### 13.4 有限维代数、群代数与有限环
% 纲要位置：计划纲要.md，第 1482 行。

\section{有限维代数、群代数与有限环}
\label{b2:sec:1304}

结构定理不仅给出抽象的环同构，也使维数、简单模和标量域之间的关系可以计算。本节先讨论有限维代数，再用最小的群代数例子比较不同特征，最后证明有限除环必交换。最后一个结论需要环本身只有有限多个元素，不能仅凭它在某个域上有限维就套用。

% 纲要要点（供目录核对）：
%
% * 有限维半单代数
% * 群代数的半单性例子与 Maschke 定理的陈述
% * 有限除环的交换性与 Wedderburn 小定理；循环模的无重因子判据
%
% ---
%

\subsection{有限维半单代数}
\label{b2:subsec:1304:01}

\begin{proposition}\label{b2:prop:finite-dimensional-semisimple-algebra}
设 \(k\) 为域、\(A\) 为有限维结合含幺半单 \(k\) 代数。存在非负整数 \(t\)、有限维除 \(k\) 代数 \(D_1,\ldots,D_t\) 及正整数 \(n_1,\ldots,n_t\)，使
\[
A\cong\prod_{i=1}^tM_{n_i}(D_i)
\]
成为 \(k\) 代数同构。相应的简单左模为列模 \(S_i=D_i^{n_i}\)，由第 \(i\) 个因子作用，其余因子作用为零，且
\[
\dim_kA=\sum_i n_i^2\dim_kD_i,
\qquad \dim_kS_i=n_i\dim_kD_i.
\]
若 \(k\) 代数闭，则所有 \(D_i=k\)，于是 \(A\cong\prod_iM_{n_i}(k)\)，且 \(\dim_kA=\sum_i(\dim_kS_i)^2\)。零代数对应 \(t=0\) 的空乘积，此时简单模列表为空，维数公式中的空和为零。
\end{proposition}
\begin{proof}
零代数的结论由空乘积约定成立。设 \(A\ne0\)。对任意简单左 \(A\) 模 \(S_i\)，取 \(0\ne s_i\in S_i\)，则 \(As_i=S_i\)，因为这是非零子模。因此 \(a\mapsto as_i\) 给出 \(k\) 线性满射 \(A\to S_i\)，从而 \(\dim_kS_i\le\dim_kA<\infty\)。\(\operatorname{End}_A(S_i)\) 是 \(\operatorname{End}_k(S_i)\) 的 \(k\) 子空间，故有限维；Schur 引理使它为除环，其反环 \(D_i\) 也是有限维除 \(k\) 代数。\cref{b2:thm:artin-wedderburn}的构造保持 \(k\) 标量：标量在正则模上作为乘 \(\lambda\) 作用，在各简单分量上仍是同一个标量自同态，取反环不会改变中心标量。因此得到 \(k\) 代数同构。矩阵有 \(n_i^2\) 个任意 \(D_i\) 条目，列模有 \(n_i\) 个条目，直接得到两条维数公式。若 \(k\) 代数闭，刚证明的简单模有限维性允许应用\cref{b2:cor:schur-algebraically-closed}，使各自同态除环等于 \(k\)，反环仍为 \(k\)，从而得到平方和公式。
\end{proof}

例如，代数闭域上的一个四维半单代数只能同构于 \(k^4\) 或 \(M_2(k)\)：其维数是正整数平方之和，只有 \(4=1+1+1+1\) 与 \(4=2^2\) 两种可能。前者有四种一维简单模，后者只有一种二维简单模。两个代数的总维数相同，模结构却不同。

\ProseParagraphBreak
一般基域上不能删去除代数的维数。比如 \(\mathbb C\) 作为实代数是半单的，其唯一简单模为 \(\mathbb C\)，实维数为二；公式是 \(\dim_{\mathbb R}\mathbb C=1^2\cdot2\)，而不是这个简单模实维数的平方。类似地，实四元数除代数维数为四，它的左正则模简单；有限维除代数仍可能非交换。

\subsection{群代数的半单性}
\label{b2:subsec:1304:02}

设 \(k\) 为域、\(G\) 为有限群。群代数 \(k[G]\) 以各 \(g\in G\) 为 \(k\) 基，乘法按群乘法线性延拓。第五卷第1.3节将通过对群作用平均投影，证明 Maschke 定理\footnote{马施克（Heinrich Maschke，1853--1908），德国数学家，研究有限群的线性表示和微分几何。}：\(k[G]\) 半单，当且仅当 \(|G|\cdot1_k\ne0\)，也就是 \(|G|\cdot1_k\) 在 \(k\) 中可逆。

\ProseParagraphBreak
平均要修正的是普通线性投影与群作用不相容的问题。若 \(\rho:G\to\operatorname{GL}_k(V)\) 是表示，\(U\subseteq V\) 在群作用下稳定，先选一个线性投影 \(p:V\to U\)，并通过包含把它看作 \(V\) 的自同态。在 \(|G|\cdot1_k\ne0\) 时，平均得到
\[
P=\frac1{|G|}\sum_{g\in G}\rho(g)p\rho(g)^{-1}.
\]
群作用会排列这些共轭项，使平均后的算子与群作用相容。每个共轭项在 \(U\) 上都是恒等，所以未经归一化的和在 \(U\) 上是 \(|G|\) 倍的恒等；只有能除以 \(|G|\)，才恢复所需的投影。这解释了特征零或 \(\operatorname{char}k=p>0\) 且 \(p\mathrel{\vcenter{\hbox{\scalebox{1.25}{\(\nmid\)}}}}|G|\) 的条件。

\ProseParagraphBreak
最小例子可以在当前知识范围内直接算出。设 \(G=C_2=\langle g\mid g^2=1\rangle\)，且 \(\operatorname{char}k\ne2\)。令
\[
e_+=\frac{1+g}{2},\qquad e_-=\frac{1-g}{2}.
\]
由 \(g^2=1\) 得 \(e_+^2=e_+\)、\(e_-^2=e_-\)、\(e_+e_-=0\)、\(e_++e_-=1\)，并且 \(ge_+=e_+\)、\(ge_-=-e_-\)。因此 \(k[C_2]=ke_+\oplus ke_-\)，左乘 \(e_+\) 与 \(e_-\) 分别是正则作用中 \(g\) 的 \(1\) 与 \(-1\) 特征子空间上的投影。这就是\cref{b2:prop:coprime-operator-decomposition}中用互素多项式构造投影的情形：\(t-1\) 与 \(t+1\) 互素，\((1+t)/2\) 与 \((1-t)/2\) 分别选出两个分量。具体地，
\[
k[C_2]\longrightarrow k\times k,
\quad a+bg\longmapsto(a+b,a-b)
\]
为代数同构，逆映射为 \((u,v)\mapsto u e_++v e_-\)。两种简单模分别使 \(g\) 作用为 \(1\) 与 \(-1\)，正则模就是这两个一维模的直和。

\ProseParagraphBreak
若 \(\operatorname{char}k=2\)，令 \(\varepsilon=g-1\)，则 \(\varepsilon^2=g^2-2g+1=0\)，从而
\[
k[C_2]\cong k[\varepsilon]/(\varepsilon^2).
\]
这里 \(\varepsilon\ne0\)，因为 \(1,g\) 仍是群代数的两个基向量。两个特征值 \(1,-1\) 在特征二时合并，留下非零的幂零元 \(\varepsilon\)。子模 \(k\varepsilon\) 与相应的商模都为一维，\(g\) 均作用为恒等，但这个子模没有直和补模；\cref{b2:prop:semisimple-cyclic-polynomial}将通过重复因子的一般判据证明这种不分裂，也给出该环不半单的结论。相同的群在不同特征下给出了不同的模分裂行为。

\ProseParagraphBreak
对 \(G=C_m\)、\(m\ge1\)，\cref{b2:prop:semisimple-cyclic-polynomial}的证明将给出同构 \(k[C_m]\cong k[t]/(t^m-1)\)，其中 \(t\) 的陪集对应循环群生成元 \(g\)。因而一个 \(C_m\) 表示就是一个满足 \(T^m=\operatorname{id}\) 的算子，群代数的模也正是由这个关系限制的 \(k[t]\) 模。这把问题接回第四章的算子模与主分量：不同不可约因子将分量分开，重复因子则可能留下不能分裂的幂零层。循环群的半单性也将在该命题中判定；一般有限群仍由第五卷的平均方法处理。

\subsection{Wedderburn 小定理}
\label{b2:subsec:1304:03}

\begin{theorem}[Wedderburn 小定理]\label{b2:thm:wedderburn-little}
每个只有有限多个元素的除环都是交换域。
\end{theorem}
\begin{proof}
设 \(D\) 为有限除环。中心 \(Z\) 是域：非零中心元 \(z\) 的逆仍与每个元素交换，因为由 \(zx=xz\) 左右乘以 \(z^{-1}\) 就得到 \(xz^{-1}=z^{-1}x\)。记 \(|Z|=q\ge2\)，把 \(D\) 视为有限维 \(Z\) 向量空间，令 \(n=\dim_ZD\)，则 \(|D|=q^n\)。若 \(n=1\)，每个元素都是中心中 \(1\) 的倍数，结论成立。下面假设 \(n>1\)，推出矛盾。

乘法群 \(D^\times\) 的中心为 \(Z^\times\)，含 \(q-1\) 个元素。为了计算非中心共轭类的大小，先考察代表元的中心化子。对 \(x\in D\setminus Z\)，其环中心化子
\[
C_D(x)=\{\,a\in D\mid ax=xa\,\}
\]
是包含 \(Z\) 的除子环：加、乘封闭，若 \(a\ne0\) 与 \(x\) 交换，则其逆也交换。设 \(m_x=\dim_Z C_D(x)\)，则 \(|C_D(x)|=q^{m_x}\)。虽然 \(C_D(x)\) 未必交换，\(D\) 仍是它的左向量空间，而 \(Z\) 位于 \(C_D(x)\) 的中心。若 \(c_1,\ldots,c_{m_x}\) 是 \(C_D(x)\) 的 \(Z\) 基，\(v_1,\ldots,v_{d_x}\) 是 \(D\) 的左 \(C_D(x)\) 基，则 \(c_i v_j\) 是 \(D\) 的 \(Z\) 基：先按 \(v_j\) 展开，再把左侧的系数按 \(c_i\) 展开，且两次展开都唯一。因此维数塔公式给出
\[
\dim_ZD=\dim_{C_D(x)}D\cdot\dim_ZC_D(x),
\qquad n=d_xm_x.
\]
这里的乘法次序始终是左系数乘基向量，证明不需要 \(C_D(x)\) 交换。故 \(m_x\mid n\)；又因 \(x\) 非中心，\(C_D(x)\) 是 \(D\) 的真 \(Z\) 子空间，所以 \(m_x<n\)。

群中心化子只取其中的可逆元：
\[
C_{D^\times}(x)=\{a\in D^\times\mid ax=xa\}=C_D(x)^\times,
\]
所以它的阶为 \(q^{m_x}-1\)。有限群的类方程把中心的单元素共轭类与其余共轭类分开，写成 \(|G|=|Z(G)|+\sum_j[G:C_G(x_j)]\)，其中 \(x_j\) 取遍非中心共轭类的代表。把第一卷命题4.2.2及式(4.1)的这一公式应用于 \(D^\times\)，取代表 \(x_1,\ldots,x_s\)，便得
\begin{equation}\label{b2:eq:wedderburn-class-equation}
q^n-1=(q-1)+\sum_{j=1}^s\frac{q^n-1}{q^{m_{x_j}}-1}.
\end{equation}
每个分数都是该代表的共轭类大小，而各分母的指数 \(m_{x_j}\) 都是 \(n\) 的真因子。

要从这个整数等式中得到约束，可以寻找一个同时整除左端与所有非中心共轭类大小的整数。分母的指数都是 \(n\) 的真因子，使第一卷定义21.1.1中的分圆因子 \(\Phi_n\) 正好具有这一性质。该卷式(21.1)及其后的整数系数归纳证明给出
\[
t^n-1=\prod_{d\mid n}\Phi_d(t),\qquad \Phi_d(t)\in\mathbb Z[t].
\]
对每个真因子 \(m\mid n\)，有整数系数多项式恒等式
\[
\frac{t^n-1}{t^m-1}=\prod_{\substack{d\mid n\\d\nmid m}}\Phi_d(t).
\]
因为 \(m<n\)，\(n\) 不整除 \(m\)，右端包含 \(\Phi_n(t)\)。代入 \(t=q\)，便知整数 \(\Phi_n(q)\) 整除 \(q^n-1\)，也整除式\eqref{b2:eq:wedderburn-class-equation}中的每个分数。从左端减去这些均被它整除的共轭类大小，剩下的正是 \(q-1\)，所以
\[
\Phi_n(q)\mid q-1.
\]

但 \(\Phi_n(q)>q-1\)。此时 \(\Phi_n(q)\) 已是一个整数；把同一个多项式放到复数域中分解，只是为了估计这个整数的大小。按其定义写
\[
\Phi_n(q)=\prod_{\zeta\text{ 为本原 }n\text{ 次单位根}}(q-\zeta).
\]
因 \(n>1\)，每个这样的 \(\zeta\ne1\)。单位圆上除 \(1\) 外的点都比 \(1\) 离正实数 \(q>1\) 更远，具体地，
\[
|q-\zeta|^2=(q-1)^2+2q(1-\operatorname{Re}\zeta)>(q-1)^2.
\]
非实根成共轭对，每对贡献 \((q-\zeta)(q-\overline\zeta)=|q-\zeta|^2>0\)；若有实本原根，则只能是 \(n=2\) 时的 \(-1\)，因子 \(q+1\) 也为正。因此 \(\Phi_n(q)>0\)，且它等于上述所有因子绝对值的乘积，严格大于 \((q-1)^{\varphi(n)}\ge q-1\)。这里 \(\varphi(n)\ge1\)，而 \(q-1\ge1\)；即使 \(q=2\)，每个因子绝对值也严格大于 \(1\)，严格不等式仍成立。这与正整数 \(\Phi_n(q)\) 整除 \(q-1\) 矛盾。故 \(n=1\)，\(D=Z\) 交换。
\end{proof}

上述结论称为\term[Wedderburn-xiaodingli]{Wedderburn 小定理}[Wedderburn's little theorem]。分圆多项式的整数性及因式分解给出关键的整除关系；有限性使中心和各中心化子的元素个数分别写成 \(q\) 的整数次幂，也使乘法群的类方程成为有限整数等式。本定理要求除环只有有限多个元素；例如四元数除代数 \(\mathbb H\) 虽在 \(\mathbb R\) 上四维，却有无限多个元素，并不交换。

\begin{corollary}\label{b2:cor:finite-semisimple-rings}
设 \(R\) 为含幺环。\(R\) 有限且半单，当且仅当存在非负整数 \(t\)、正整数 \(n_1,\ldots,n_t\) 及素数幂 \(q_1,\ldots,q_t\)，使
\[
R\cong\prod_{i=1}^t M_{n_i}(\mathbb F_{q_i}).
\]
\(t=0\) 时把空乘积解释为零环。
\end{corollary}
\begin{proof}
由\cref{b2:thm:artin-wedderburn}分解后，每个矩阵因子仍有限，其系数除环作为单个矩阵条目的取值集也有限。\cref{b2:thm:wedderburn-little}使每个除环都是有限域。反过来，有限个有限域上的有限阶矩阵环显然只有有限多个元素，且由结构定理半单。
\end{proof}

回到循环群代数，交换性使半单分解中的矩阵因子只能是一阶的域。对多项式商环，这种分解可以直接从不可约因子读出；若因子重复，问题则是相应的子模能否真正分裂，而不只是有几个简单的合成因子。

\begin{proposition}\label{b2:prop:semisimple-cyclic-polynomial}
设 \(k\) 为域，\(f\in k[t]\) 为非零非常数多项式。以下条件等价：\(k[t]/(f)\) 作为幺左 \(k[t]\) 模半单；它作为含幺环半单；\(f\) 的每个不可约因子在 \(k[t]\) 中只出现一次。

若 \(m\ge1\) 为整数、\(C_m\) 为 \(m\) 阶循环群，则群代数 \(k[C_m]\) 半单，当且仅当 \(\operatorname{char}k=0\)，或 \(\operatorname{char}k=p>0\) 且 \(p\) 不整除 \(m\)。
\end{proposition}
\begin{proof}
令 \(A=k[t]/(f)\)。因为 \(k[t]\to A\) 满射，\(A\) 的左 \(k[t]\) 子模与左 \(A\) 子模相同，简单性和直和分解也相同。因此 \(A\) 作为 \(k[t]\) 模半单，等价于左正则模 \({}_AA\) 半单，也就是环 \(A\) 半单。

若 \(f=c\,p_1\cdots p_r\)，其中 \(c\in k^\times\)，\(p_i\) 是两两不同的首一不可约多项式，则这些因子两两互素。第一卷定理10.4.3的中国剩余定理给出
\[
A\cong\prod_{i=1}^r k[t]/(p_i).
\]
每个因子都是域，也是简单 \(k[t]\) 模；有限乘积作为模就是有限直和，故 \(A\) 半单。

若某个不可约多项式 \(p\) 满足 \(p^2\mid f\)，则 \(A\) 有商模 \(Q=k[t]/(p^2)\)。设 \(A\) 半单，\cref{b2:prop:semisimple-subquotients}便使 \(Q\) 半单，故非零子模 \(pQ\) 应有直和补模，相应地存在 \(k[t]\) 线性投影 \(u:Q\to pQ\)，在 \(pQ\) 上为恒等。记 \(1,p\) 在 \(Q\) 中的陪集为 \(\overline1,\overline p\)。因 \(p^2Q=0\)，\(p\) 消去 \(pQ\)，而 \(u(\overline1)\in pQ\)，于是
\[
u(\overline p)=u(p\overline1)=p\,u(\overline1)=0.
\]
但 \(\overline p\ne0\)，且 \(\overline p\in pQ\)，所以投影又要求 \(u(\overline p)=\overline p\)，矛盾。故 \(A\) 不半单。

取 \(C_m=\langle g\rangle\)，映射 \(t\mapsto g\) 给出代数同构 \(k[t]/(t^m-1)\cong k[C_m]\)，因为它把前者的基 \(\overline1,\overline t,\ldots,\overline{t^{m-1}}\) 映到后者的基 \(1,g,\ldots,g^{m-1}\)。若 \(\operatorname{char}k=0\)，或特征 \(p\) 不整除 \(m\)，则 \(m\cdot1_k\ne0\)，而
\[
(t^m-1)'=mt^{m-1},
\qquad \gcd(t^m-1,mt^{m-1})=1.
\]
任何重复的不可约因子都会同时整除多项式及其导数，因此 \(t^m-1\) 没有重复不可约因子，群代数半单。若 \(\operatorname{char}k=p>0\) 且 \(p\mid m\)，则
\[
t^m-1=(t^{m/p}-1)^p.
\]
括号内仍为非常数多项式，至少有一个不可约因子，而它在右端的重数至少为 \(p\ge2\)，所以群代数不半单。
\end{proof}


\begin{chaptersummary}
半单模的关键性质是每个子模都有补模。简单子模之和可以抽成直和，半单模的子模与商模仍半单；但两端半单的扩张未必半单。Schur 引理把不同简单类型间的同态消去，并把同一简单类型的自同态组成除环。

\ProseParagraphBreak
半单环的正则模由有限多个简单左理想组成，这使整个左模范畴中的短正合列都分裂。把正则模的自同态按简单类型分块，再通过反环恢复原乘法，得到有限个矩阵环的乘积。简单类型决定系数除环，正则模中的重数决定矩阵大小；具体坐标可以改变，这些参数的同构类型不变。

\ProseParagraphBreak
半本原性只要求所有简单模共同零化的元素为零，它本身不保证正则模半单。有限维代数的结构还受基域影响，有限除环的交换性则依靠更强的“元素总数有限”。Wedderburn 小定理把有限半单环进一步归结为有限域上的矩阵环；下一章将用 Jacobson 根研究一般环偏离半单结构的部分。
\end{chaptersummary}

\BookExercises

\begin{exercise}\label{b2:ex:13-simple-polynomial-modules}
设 \(k\) 为域。分类 \(k[t]\) 的全部幺性简单左模，并说明哪些具有 \(k\) 维数一。在 \(k=\mathbb R\) 时给出一个二维简单模。
\end{exercise}
\begin{solution}
每个简单模是极大理想的商。\(k[t]\) 是主理想整环，极大理想恰为不可约多项式 \(f\) 生成的理想，所以简单模为 \(k[t]/(f)\)。取首一 \(f\) 后，零化子为 \((f)\)，不同的 \(f\) 给出不同同构类型。其 \(k\) 维数为 \(\deg f\)，因为 \(1,t,\ldots,t^{\deg f-1}\) 的陪集组成基。故一维情形正是 \(f=t-a\)。在实数域上，\(\mathbb R[t]/(t^2+1)\cong\mathbb C\) 为二维简单模，其中 \(t\) 作用为乘以 \(i\)。
\end{solution}

\begin{exercise}\label{b2:ex:13-squarefree-cyclic}
设 \(k\) 为域，\(V\ne0\) 为有限维 \(k\) 向量空间，\(T\in\operatorname{End}_k(V)\)，令 \(V_T\) 为由 \(t\cdot v=T(v)\) 定义的 \(k[t]\) 模。证明 \(V_T\) 半单当且仅当极小多项式 \(\mu_T(t)\) 不含重复不可约因子；\(T\) 在 \(k\) 上可对角化当且仅当 \(\mu_T\) 是互不相同的一次因子的乘积。分别分析实矩阵
\[
\begin{pmatrix}0&-1\\1&0\end{pmatrix},\qquad
\begin{pmatrix}0&1\\0&0\end{pmatrix}
\]
对应的模，说明半单性与在基域上可对角化的区别。
\end{exercise}
\begin{solution}
由\cref{b2:thm:rational-canonical-form,b2:prop:operator-characteristic-minimal}，存在首一非常数多项式 \(f_1\mid\cdots\mid f_s\)，使
\[
V_T\cong\bigoplus_{i=1}^s k[t]/(f_i),\qquad \mu_T=f_s.
\]
若 \(\mu_T\) 无重因子，每个 \(f_i\) 也无重因子，\cref{b2:prop:semisimple-cyclic-polynomial}使每个循环分量半单，故其直和半单。反过来，若 \(V_T\) 半单，则它的商 \(k[t]/(f_s)\) 半单，同一判据说明 \(f_s=\mu_T\) 无重因子。

若 \(T\) 可对角化，取由特征向量组成的基，每个基向量所在直线为简单 \(k[t]\) 模 \(k[t]/(t-\lambda)\)，零化全模的首一多项式中次数最小者正是全部不同特征值对应的 \(t-\lambda\) 之积。反过来，若 \(\mu_T\) 是不同一次因子的乘积，各 \(f_i\) 也如此；\secref{b2:sec:0403}的中国剩余分解把上述循环分量拆为一维模，合并这些一维模的基便给出特征向量基。

第一个实矩阵的极小多项式为 \(t^2+1\)：其平方为 \(-I\)，且它不是实标量矩阵。该多项式不可约，所以对应模简单、因而半单，但在 \(\mathbb R\) 上没有特征值，不能对角化。第二个矩阵非零而平方为零，极小多项式为 \(t^2\)，含重因子，因此对应模不半单。
\end{solution}

\begin{exercise}\label{b2:ex:13-graph-submodules}
设 \(R\) 为含幺环、\(S\) 为简单左 \(R\) 模，令 \(E=\operatorname{End}_R(S)\)。证明 \(S\oplus S\) 中与 \(S\oplus0\) 互补的子模恰为
\(K_f=\Bigl\{\,\bigl(f(s),s\bigr):s\in S\,\Bigr\}\)、\(f\in E\)。再分类 \(S\oplus S\) 的全部简单子模。
\end{exercise}
\begin{solution}
每个 \(K_f\) 是第二坐标投影的一个截面的像，且 \((x,y)=\bigl(x-f(y),0\bigr)+\bigl(f(y),y\bigr)\)，因此确为补模。反过来，对任一补模 \(K\)，第二投影 \(K\to S\) 单射，因为核包含于 \(K\cap(S\oplus0)\)；它满射来自直和覆盖。其逆形如 \(s\mapsto\bigl(f(s),s\bigr)\)，其中 \(f\) 线性，所以得到全部补模。

若简单子模 \(T\) 的第二投影为零，则 \(T\subseteq S\oplus0\)，由简单性 \(T=S\oplus0\)。若第二投影非零，Schur 引理使其成为同构，所以 \(T=K_f\)。这也说明一个固定直和分解不会列出所有简单子模。

更一般地，对任意正整数 \(n\)，\(\operatorname{Hom}_R(S,S^n)\cong E^n\) 是右 \(E\) 向量空间，非零同态均为单射；若非零同态 \(f,g:S\to S^n\) 的像相同，把它们视为到共同像的同构，则 \(e=g^{-1}f\in E^\times\) 且 \(f=g\circ e\)；反过来，这个等式也保证像相同。任一简单子模向某个坐标的投影非零，故由 Schur 引理同构于 \(S\)，所以全部简单子模都由这些同态得到。因此，\(S^n\) 的简单子模对应于这个右 \(E\) 向量空间的一维子空间，也就是右 \(E\) 空间的射影空间。
\end{solution}

\begin{exercise}\label{b2:ex:13-products-of-simple-modules}
令 \(p\) 遍历全部素数。证明整数模 \(\bigoplus_p\mathbb Z/p\mathbb Z\) 半单，但 \(\prod_p\mathbb Z/p\mathbb Z\) 不半单。若把全部因子都换成同一个 \(\mathbb Z/2\mathbb Z\)，直积是否半单？
\end{exercise}
\begin{solution}
直和按定义半单。每个简单整数模都为某个 \(\mathbb Z/p\mathbb Z\)，因此任意半单整数模中的元素只有有限支集，必有有限阶。不同素数的直积却含元素 \(x=(\bar1)_p\)，它不被任何非零整数 \(a\) 零化：取不整除 \(a\) 的素数 \(p\)，该坐标仍非零。因此直积不半单。若全部因子都是 \(\mathbb Z/2\mathbb Z\)，整个直积是 \(\mathbb F_2\) 向量空间，选取向量空间基后成为 \(\mathbb Z/2\mathbb Z\) 的直和，所以作为整数模半单。这个基不同于标准坐标向量族，后者只张成有限支集部分。同一个乘积作为其自身环的左正则模不半单，而作为整数模却半单；必须明确标量环，不能仅由“无限直积”判断半单性。
\end{solution}

\begin{exercise}\label{b2:ex:13-schur-converse}
令 \(M=\mathbb Q\oplus\mathbb Z\)，视为整数模。计算它的自同态环，并解释为什么从 \(\mathbb Q\) 到 \(\mathbb Z\) 的块必为零，而从 \(\mathbb Z\) 到 \(\mathbb Q\) 的块可以非零。求这个环中的全部幂等元。
\end{exercise}
\begin{solution}
由\cref{b2:prop:schur-hypothesis-counterexamples}，\(\operatorname{End}_{\mathbb Z}(\mathbb Q)=\mathbb Q\)；而整数模同态 \(\mathbb Z\to\mathbb Z\) 或 \(\mathbb Z\to\mathbb Q\) 都由 \(1\) 的像决定，分别得到 \(\mathbb Z\) 或 \(\mathbb Q\)。若 \(f:\mathbb Q\to\mathbb Z\)，则对每个正整数 \(n\)，有 \(f(1)=n f(1/n)\)，所以 \(f(1)\) 被所有正整数整除，只能为零。再由 \(b f(a/b)=a f(1)=0\) 和 \(\mathbb Z\) 无挠，得 \(f=0\)。因此
\[
\operatorname{End}_{\mathbb Z}(M)\cong
\left\{\begin{pmatrix}a&b\\0&d\end{pmatrix}:a,b\in\mathbb Q,\ d\in\mathbb Z\right\},
\]
其作用为 \((x,z)\mapsto(ax+bz,dz)\)，复合给出通常矩阵乘法。这里 \(\mathbb Q\) 与 \(\mathbb Z\) 并非一对简单模，不能借 Schur 引理把两个不同类型间的 Hom 都消去。

幂等条件为 \(a^2=a\)、\(d^2=d\)、\((a+d)b=b\)。前两式使 \(a,d\in\{0,1\}\)；若 \(a=d\)，则 \(b=0\)，得到零元和单位元。若 \((a,d)=(1,0)\) 或 \((0,1)\)，任意 \(b\in\mathbb Q\) 都可取，分别得到
\[
\begin{pmatrix}1&b\\0&0\end{pmatrix},\qquad
\begin{pmatrix}0&b\\0&1\end{pmatrix}.
\]
它们给出不同投影，表明即使两个固定分量间的一侧 Hom 为零，另一侧的非零 Hom 仍可改变补模。
\end{solution}

\begin{exercise}\label{b2:ex:13-schur-finite-dimension}
设 \(k\) 为域，将 \(M=k(t)\) 视为 \(k[t]\) 模，其中 \(t\) 按乘法作用。证明 \(\operatorname{End}_{k[t]}(M)\cong k(t)\)，但 \(M\) 不简单。比较把标量环扩大为 \(k(t)\) 后的简单性，并证明乘以 \(t\) 的自同态没有属于 \(k\) 的特征值。
\end{exercise}
\begin{solution}
对 \(f\in\operatorname{End}_{k[t]}(M)\)，令 \(c=f(1)\)。任取 \(a,b\in k[t]\)、\(b\ne0\)，线性性给出
\[
b f(a/b)=f(a)=ac.
\]
在 \(k(t)\) 中乘以 \(b\) 可逆，故 \(f(a/b)=(a/b)c\)。反过来，任意 \(c\in k(t)\) 的乘法映射都与 \(k[t]\) 作用相容，其加法和复合对应 \(k(t)\) 的加法与乘法，因此得到环同构。\(k[t]\) 是 \(M\) 的非零真 \(k[t]\) 子模，故 \(M\) 不简单；扩大为 \(k(t)\) 作用后，它成为域的一维向量空间，故简单。两种标量环给出相同的自同态环，却有不同的子模。

若乘以 \(t\) 有特征值 \(\lambda\in k\)，则存在 \(0\ne v\in k(t)\) 使 \((t-\lambda)v=0\)，与 \(k(t)\) 为域且 \(t-\lambda\ne0\) 矛盾。即使 \(k\) 代数闭，这个无限维算子也没有特征值；有限维 Schur 推论中的特征值步骤不能用于它。
\end{solution}

\begin{exercise}\label{b2:ex:13-endomorphism-blocks}
计算整数模 \(M=(\mathbb Z/2\mathbb Z)^3\oplus(\mathbb Z/3\mathbb Z)^2\) 的自同态环，并求其元素总数。
\end{exercise}
\begin{solution}
从二阶循环群到三阶循环群的同态像同时被二和三零化，所以只能为零，反方向也是如此。每个同类型分量中的同态由域中的一个标量决定，于是
\[
\operatorname{End}_{\mathbb Z}(M)\cong M_3(\mathbb F_2)\times M_2(\mathbb F_3).
\]
每个三阶矩阵有九个二值条目，每个二阶矩阵有四个三值条目，故元素总数为 \(2^9\cdot3^4\)。这里自同态环包含不可逆矩阵，因此它不是除环，符合 \(M\) 并非简单模的事实。
\end{solution}

\begin{exercise}\label{b2:ex:13-idempotent-summands}
设 \(k\) 为域，在 \(R=M_2(k)\) 中，令 \(I=RE_{11}\) 为第二列全零的左理想。求全部以 \(I\) 为像且在 \(I\) 上为恒等的左模投影 \(p:R\to I\)，写出各自的幂等元及补模。一个直和因子是否唯一决定它在正则模中的投影？
\end{exercise}
\begin{solution}
由\cref{b2:prop:regular-hom-endomorphism}，投影必为右乘 \(e=p(1)\in I\)，写 \(e=\begin{psmallmatrix}a&0\\b&0\end{psmallmatrix}\)。要求 \(p(E_{11})=E_{11}\) 给出 \(a=1\)；对任意 \(b\in k\)，所得
\[
e_b=\begin{pmatrix}1&0\\b&0\end{pmatrix}
\]
都幂等，并且右乘它在 \(I\) 上为恒等。若以 \((v,w)\) 表示一个矩阵的两列，则 \(p_b(v,w)=(v+bw,0)\)，故核为所有 \((-bw,w)\)、\(w\in k^2\) 组成的左理想。\cref{b2:prop:idempotent-module-decomposition}给出 \(R=I\oplus\ker p_b\)。不同 \(b\) 给出不同投影及补模，所以只指定直和因子 \(I\) 还不能唯一指定投影。
\end{solution}

\begin{exercise}\label{b2:ex:13-annihilator-vector-module}
设 \(D\) 为除环，\(n\ge1\)，令 \(R=M_n(D)\)、\(S=D^n\) 为标准幺性左列模。求 \(\operatorname{Ann}_R(S)\) 和标准列向量 \(e_1\) 的零化左理想，指出二者何时相同、何时不同。
\end{exercise}
\begin{solution}
若矩阵零化所有标准列，它的每列都为零，因此 \(\operatorname{Ann}_R(S)=0\)。而 \(Ae_1=0\) 当且仅当 \(A\) 的第一列为零，所以单个向量的零化子是这些矩阵组成的左理想。\(n=1\) 时这个理想也是零，两个零化子相同；对 \(n>1\)，它非零且不是双边理想，例如 \(E_{12}\) 的第一列为零，但 \(E_{12}E_{21}=E_{11}\) 的第一列非零。零化整个模得到双边理想，零化一个元素一般只得到左理想。这里 \(S\) 又是简单列模，且 \(\operatorname{Ann}_R(S)=0\)，所以 \(S\) 是忠实简单模；\chapref{b2:ch:14}将以忠实简单模定义本原环。
\end{solution}

\begin{exercise}\label{b2:ex:13-center-and-idempotents}
设 \(k\) 为含 \(q\) 个元素的有限域，\(R=M_2(k)\times k\)。分类 \(R\) 中幂等元的共轭类型，分别计算全部幂等元和中心幂等元的个数。这里共轭指由 \(R\) 中单位元实施的 \(e\mapsto ueu^{-1}\)。
\end{exercise}
\begin{solution}
若 \(P\in M_2(k)\) 幂等，线性投影给出 \(k^2=\operatorname{im}P\oplus\ker P\)。选择两部分的基，便使 \(P\) 共轭于 \(\operatorname{diag}(I_r,0)\)，其中 \(r=0,1,2\)，而秩在共轭下不变，所以恰有三种共轭类型。\(k\) 的幂等元只有零和一，且共轭不改变它们。因此 \(R\) 中的幂等元共有六种共轭类型。

秩一的投影由像直线 \(L\) 与核直线 \(K\ne L\) 唯一确定。\(k^2\) 有 \((q^2-1)/(q-1)=q+1\) 条直线，选定 \(L\) 后有 \(q\) 种 \(K\)，故秩一幂等矩阵共 \(q(q+1)\) 个。加上零矩阵和单位矩阵，再为 \(k\) 因子独立选择零或一，得到全部幂等元的个数
\[
2\bigl(2+q(q+1)\bigr).
\]
由\cref{b2:prop:matrix-centers-central-idempotents}，中心幂等元却只有 \((0,0),(I_2,0),(0,1),(I_2,1)\) 四个。秩一投影不在中心，不能把它切出的左模分量误当成乘积环的双边矩阵块。
\end{solution}

\begin{exercise}\label{b2:ex:13-concrete-finite-product}
设 \(R=M_2(\mathbb F_3)\times M_3(\mathbb F_2)\)。求全部简单左模的同构类型、各自元素个数、它们在左正则模中的重数，以及 \(R\) 的元素总数。
\end{exercise}
\begin{solution}
两种简单模分别为 \(S=\mathbb F_3^2\) 和 \(T=\mathbb F_2^3\)，由对应因子作用，另一因子作用为零。它们分别有九个、八个元素；通过中心幂等元的不同作用可知不可能同构。按列分解，\({}_RR\cong S^2\oplus T^3\)，所以重数为二和三，正则模长度为五。环本身的元素总数为 \(3^4\cdot2^9\)。两种模来自不同特征，这里不为它们虚构共同的基域。
\end{solution}

\begin{exercise}\label{b2:ex:13-dimension-nine}
设 \(k\) 为代数闭域。列出所有九维结合含幺半单 \(k\) 代数的同构类型，并为每个类型列出简单模的维数。
\end{exercise}
\begin{solution}
由\cref{b2:prop:finite-dimensional-semisimple-algebra}，维数是矩阵大小的平方之和。可能的正平方为 \(1,4,9\)，故全部类型为
\[
k^9,\qquad M_2(k)\times k^5,\qquad
M_2(k)\times M_2(k)\times k,\qquad M_3(k).
\]
相应简单模维数依次为九个一；一个二和五个一；两个二和一个一；一个三。Artin–Wedderburn 参数唯一性说明这些类型彼此不同，也说明没有遗漏。
\end{solution}

\begin{exercise}\label{b2:ex:13-commutative-semisimple}
证明交换含幺半单环恰为有限个域的乘积，包括零环的空乘积情形。指出为什么这里不能把“有限个”改为“任意多个”。
\end{exercise}
\begin{solution}
在 Artin–Wedderburn 分解中，若某个 \(n_i>1\)，则 \(E_{11}E_{12}=E_{12}\) 而 \(E_{12}E_{11}=0\)，不交换，所以所有 \(n_i=1\)。各 \(D_i\) 作为该环的商也交换，因此是域。反过来，有限个域的乘积由结构定理半单且显然交换。无限乘积 \(\prod_{j\ge1}k\) 中“前 \(r\) 个坐标为零”的理想给出无限严格降链，而半单环必须 Artin，所以不满足结论。
\end{solution}

\begin{exercise}\label{b2:ex:13-cyclic-three}
分别在 \(\mathbb Q\)、\(\mathbb C\) 及特征三的域 \(k\) 上分析循环群 \(C_3\) 的群代数。前两种情形写出半单分解；后一种情形证明不半单。不要调用一般 Maschke 定理。
\end{exercise}
\begin{solution}
群代数为 \(k[t]/(t^3-1)\)。在 \(\mathbb Q\) 上，\(t^3-1=(t-1)(t^2+t+1)\)，第二个因子没有有理根所以不可约，且两因子互素。\secref{b2:sec:0403}的中国剩余分解给出 \(\mathbb Q[C_3]\cong\mathbb Q\times\mathbb Q(\zeta_3)\)，其中 \(\zeta_3\) 为本原三次单位根，满足 \(\zeta_3^2+\zeta_3+1=0\)，故第二因子就是 \(\mathbb Q[t]/(t^2+t+1)\)。在 \(\mathbb C\) 上，三个单位根互异，故 \(\mathbb C[C_3]\cong\mathbb C^3\)。

特征三时 \(t^3-1=(t-1)^3\)，写 \(\varepsilon=t-1\) 得 \(R=k[\varepsilon]/(\varepsilon^3)\)。若 \((\varepsilon^2)\) 有模投影 \(p:R\to(\varepsilon^2)\)，则 \(p(\varepsilon^2)=\varepsilon^2p(1)=0\)，但投影在该非零理想上应为恒等，矛盾。故此环不半单。
\end{solution}

\begin{exercise}\label{b2:ex:13-finite-order-p-four}
设 \(p\) 为素数。分类所有恰有 \(p^4\) 个元素的含幺半单环，并指出哪些非交换。
\end{exercise}
\begin{solution}
由\cref{b2:cor:finite-semisimple-rings}，所有因子形如 \(M_n(\mathbb F_{p^f})\)，它对元素总数的指数贡献为 \(fn^2\)，各贡献之和为四。若 \(n>1\)，只能 \(n=2,f=1\)，且此因子已占全部指数，得到 \(M_2(\mathbb F_p)\)。其余情况全为域，按四的正整数分拆得到
\[
\mathbb F_{p^4},\quad \mathbb F_{p^3}\times\mathbb F_p,\quad
\mathbb F_{p^2}\times\mathbb F_{p^2},\quad
\mathbb F_{p^2}\times\mathbb F_p\times\mathbb F_p,\quad
\mathbb F_p^4.
\]
只有矩阵环这一种非交换。第一卷定理17.1.1与定理17.1.2保证各素数幂大小的有限域存在，且由元素数决定同构类型；再加上结构定理的唯一性，得到列表无重叠且穷尽。
\end{solution}

\begin{exercise}\label{b2:ex:13-smallest-noncommutative-simple}
分类所有有 \(p^8\) 个元素的含幺单环，其中 \(p\) 为素数。再求非交换有限单环可能具有的最小元素数。
\end{exercise}
\begin{solution}
有限环当然左 Artin，因此单环由\cref{b2:prop:simple-artinian}写为 \(M_n(D)\)。\(D\) 有限，由 Wedderburn 小定理为 \(\mathbb F_{p^f}\)，且 \(fn^2=8\)。于是只有 \(n=1,f=8\) 或 \(n=2,f=2\)，即 \(\mathbb F_{p^8}\) 与 \(M_2(\mathbb F_{p^2})\)。一般非交换有限单环必须 \(n\ge2\)，所以元素数至少为 \(2^{2^2}=16\)，且 \(M_2(\mathbb F_2)\) 达到这个值。
\end{solution}

\begin{exercise}\label{b2:ex:13-recover-multiplicities}
设 \(k\) 为域，\(R=M_2(k)\times k\)，\(M\) 为有限维幺性左 \(R\) 模，且 \(\dim_kM=8\)、\(\dim_k\operatorname{End}_R(M)=13\)。求 \(M\) 的简单分解。
\end{exercise}
\begin{solution}
两种简单模为 \(S=k^2\)、\(T=k\)，自同态环均为 \(k\)，不同类型间 Hom 为零。由\cref{b2:thm:artin-wedderburn}，\(R\) 是半单环，\cref{b2:thm:semisimple-ring-characterizations}于是保证每个左 \(R\) 模半单，包括题设的 \(M\)。因此 \(M\cong S^a\oplus T^b\)，其中 \(a,b\ge0\)。同类分量间的同态给出矩阵块，不同类型间的同态为零，故 \(\operatorname{End}_R(M)\cong M_a(k)\times M_b(k)\)，其维数为 \(a^2+b^2\)。题设两项维数因此给出
\(2a+b=8\)、\(a^2+b^2=13\)。第二式迫使 \((a,b)=(2,3)\) 或 \((3,2)\)，第一式选出 \((3,2)\)。因此 \(M\cong S^3\oplus T^2\)。这两项数值共同恢复重数，单用总维数则不能做到。
\end{solution}

\begin{exercise}\label{b2:ex:13-dual-numbers-semisimple-modules}
设 \(k\) 为域，\(R=k[\varepsilon]/(\varepsilon^2)\)。证明一个幺性左 \(R\) 模半单，当且仅当 \(\varepsilon\) 在其上作用为零。把结论翻译成满足 \(T^2=0\) 的线性算子的陈述。
\end{exercise}
\begin{solution}
若 \(S\) 简单，\(\varepsilon S\) 是子模，因为环交换。若它非零，则等于 \(S\)，于是 \(\varepsilon^2S=\varepsilon S=S\)，与 \(\varepsilon^2=0\) 矛盾。所以 \(\varepsilon S=0\)，每个半单模的全部简单分量都被 \(\varepsilon\) 零化，整个模也如此。反过来，若 \(\varepsilon M=0\)，作用经过 \(R/(\varepsilon)=k\)，选取 \(k\) 基后每个一维子空间都是简单 \(R\) 子模，故半单。若用 \(T\) 表示 \(\varepsilon\) 的作用，结论就是对应模半单当且仅当 \(T=0\)，而不仅仅是 \(T^2=0\)。
\end{solution}

\begin{exercise}\label{b2:ex:13-infinite-product-semiprimitivity}
令 \(R=\prod_{i\ge1}k_i\)，其中各 \(k_i\) 为域。对每个 \(i\)，令 \(S_i=k_i\) 通过第 \(i\) 个坐标接受 \(R\) 作用。证明 \(R\) 满足半本原性的零化子条件，却不是半单环。
\end{exercise}
\begin{solution}
各 \(S_i\) 简单，因为其非零向量在对应域作用下生成整个 \(k_i\)。\(\operatorname{Ann}_R(S_i)\) 正是第 \(i\) 个坐标为零的理想，这些零化子的交为零。因此所有简单模零化子的交作为其子集也为零，\(R\) 半本原。另一方面，依次要求前一个、前两个、前三个坐标为零，得到无限严格下降的左理想链。由\cref{b2:thm:semisimple-ring-characterizations}，半单环必左 Artin，故此 \(R\) 不半单。
\end{solution}

\begin{exercise}\label{b2:ex:13-projective-factors}
设 \(R\) 为含幺环，幺性左 \(R\) 模 \(M\) 有有限合成列，且该合成列中每个简单商因子都是投射模。证明 \(M\) 半单。反过来，半单模的简单因子是否必为原环上的投射模？
\end{exercise}
\begin{solution}
对合成列长度归纳。长度零时为零模。设末端为 \(0\to N\to M\xrightarrow{q}S\to0\)，其中 \(S\) 为投射简单模。将 \(\operatorname{id}_S\) 沿满射 \(q\) 提升，得到截面，故 \(M\cong N\oplus S\)。\(N\) 的较短合成列仍有投射简单商因子，归纳得它半单，所以 \(M\) 半单。逆命题不成立：取素数 \(p\)，\(\mathbb Z/p\mathbb Z\) 本身简单且半单，但整数模满射 \(\mathbb Z\to\mathbb Z/p\mathbb Z\) 不分裂，因此它不是投射模。对有限长度模，沿一条合成列逐层分裂，就能把全部简单商提升为实际直和分量；反过来，半单模的子模都有补模，合成列中的各次扩张都分裂。简单商因子投射是保证这些扩张分裂的充分条件，却不是半单性所必需的条件。
\end{solution}

% !TeX root = ../../Book2.tex
% 纲要标题：## 第14章　Jacobson 根、稠密性与本原环

\chapter{Jacobson 根、稠密性与本原环}
\label{b2:ch:14}
\BookChapterNumber{b2:ch:14}

\begin{chapterabstract}
Jacobson 根收集在每个简单模上都作用为零的环元素，其单位判据把这种作用性质转成可逆性问题。本章证明根在商环与矩阵环中的公式，再用非交换 Nakayama 引理提升有限生成模的生成元。随后区分本原环、半本原环与半单环，以简单模的自同态除环为正确标量建立 Jacobson 稠密性定理。代数闭域上的有限维情形进一步给出 Burnside 矩阵代数定理；具体反例说明底域、维数和量词各自的作用。
\end{chapterabstract}

\begin{learninggoals}
\begin{itemize}
\item 在极大左、右理想的交、简单模零化子与单位判据之间转换，计算商环和矩阵环的 Jacobson 根。
\item 证明并运用非交换 Nakayama 引理，分别在局部环与一般半单商上计算最少生成元个数。
\item 区分一个忠实简单模与一族简单模共同检测环元素的条件，理解半本原环的次直积表示。
\item 核对自同态除环的反环与右标量约定，使用稠密性实现有限组插值，辨明何时能得到全部矩阵代数。
\end{itemize}
\end{learninggoals}

在双数环 \(A=k[\varepsilon]/(\varepsilon^2)\) 中，\(\varepsilon\ne0\)，却在简单模 \(A/(\varepsilon)\) 上作用为零。另一方面，对任意 \(a\in A\)，\(1-a\varepsilon\) 的逆可直接写成 \(1+a\varepsilon\)。一个元素对所有简单模的作用与它参与单位判据的表现，竟然指向同一类环元素；Jacobson 根把这种现象从算例中抽出来。

\ProseParagraphBreak

对全部简单模考察作用，得到它们共同的零化子 \(J(R)\)；若改为固定一个简单模，问题就变成这个模能否忠实地反映环元素，以及环元素能在它上面实现哪些算子。前一个问题引出本原环，后一个问题引出稠密性定理。两者都从简单模出发，但“一族简单模共同忠实”与“一个简单模忠实”是不同的条件。

\ProseParagraphBreak

根也控制生成元的提升。对有限生成模 \(M\)，\(M/J(R)M\) 的任意一组生成元，无论怎样选择代表元，都会生成 \(M\)。Nakayama 引理说明为何根层不会妨碍这个提升；稠密性定理则说明为何简单模上的算子只需保持自同态除环给出的关系，就能实现任意有限组相容的指定像。

\ProseParagraphBreak

本章需要\chapref{b2:ch:12}的有限长度和局部环判据，以及\chapref{b2:ch:13}的简单模、Schur 引理和半单环结构。一般环部分不预设任何链条件；只有生成元提升、有限长度的半单性判据及有限维表示结论分别加入相应条件。单位判据中的可逆性要求“对每个环元素”成立；稠密性要求向量在自同态右除环上线性无关，不能改成较小底域上的线性无关。

\ProseParagraphBreak
本章的结论在后文有两种不同用途。\secref{b2:sec:1401}的根与单位判据及\secref{b2:sec:1402}的 Nakayama 引理支撑\chapref{b2:ch:15}的根滤过与投射覆盖，进而用于\chapref{b2:ch:16}的 Morita 理论。\secref{b2:sec:1403}的简单模和\secref{b2:sec:1404}的稠密性，尤其是自同态右除环上的有限维情形，则进入\chapref{b2:ch:17}的中心单代数和\chapref{b2:ch:18}的 Brauer 群。\secref{b2:sec:1405}的 Burnside 定理给出代数闭域上的进一步情形。

\ProseParagraphBreak
有限维代数的根可参阅 Etingof 等人的《Introduction to Representation Theory》\cite[Section 3.5, Definition 3.5.1, Proposition 3.5.2, p. 46]{EtingofEtAl2011RepresentationTheory}；代数闭域上有限维不可约表示的稠密性见\cite[Section 3.2, Corollary 3.2.1, Theorem 3.2.2(i), p. 43]{EtingofEtAl2011RepresentationTheory}。一般含幺环的 Jacobson 根与本原环可参阅 Lam 的《A First Course in Noncommutative Rings》\cite[Chapters 2 and 4]{Lam2001NoncommutativeRings}。

% !TeX root = ../../Book2.tex
% 纲要标题：### 14.1 Jacobson 根
% 纲要位置：计划纲要.md，第 1492 行。

\section{Jacobson 根}
\label{b2:sec:1401}

简单模给出环最小的非零作用对象。某个环元素可能在每个简单模上都作用为零；这些元素组成 Jacobson 根。本节把这个描述与极大左理想、可逆性及取商联系起来。环 \(R\) 均含单位元，模均幺性，不假设 \(R\) 交换或满足链条件。

% 纲要要点（供目录核对）：
%
% * 极大左理想的交及简单模作用的共同核；极大左理想商首先为模
% * 根中的可逆性判据及 Jacobson 逆公式；区分单个幂零元与理想中元素各自幂零
% * 商环、矩阵环与有限长度判据；矩阵环简单模对应、单位提升及上三角根计算
%

\subsection{极大左理想的交及等价刻画}
\label{b2:subsec:1401:01}

每个简单左模都同构于某个 \(R/L\)，其中 \(L\) 为极大左理想。\cref{b2:prop:simple-annihilator-intersection}已经证明：一个元素零化全部简单模，当且仅当它属于全部极大左理想。因而这些左理想的交正好把简单作用共同检测不到的部分收集起来。这里 \(R/L\) 首先是简单左模；\(L\) 未必为双边理想，不能一般地把它当作商环。

\begin{definition}\label{b2:def:jacobson-radical}
设 \(R\) 为含幺环。将 \(R\) 的全部极大左理想取交，所得集合称为 \(R\) 的\term[Jacobson-gen]{Jacobson 根}[Jacobson radical]，记为
\[
J(R)=\bigcap_{L\text{ 极大左理想}}L.
\]
对零环约定 \(J(0)=0\)。非零环的每个真左理想都包含于某个极大左理想。
\end{definition}
\symidx{\(J(R)\)：环 \(R\) 的 Jacobson 根}

最后的存在性可由第一卷附录A定理A.1.4的 Zorn 引理得到：包含给定真左理想的真左理想族非空，链的并仍是左理想，而且不含一，因而仍真；空链用原理想作上界。由此取到极大元。这里“极大”在所有真左理想之间比较，不要求它本身是双边理想。

\begin{proposition}\label{b2:prop:radical-simple-annihilators}
设 \(R\) 为含幺环，\(J(R)\) 为其 Jacobson 根。\(J(R)\) 等于全部幺性简单左 \(R\) 模的零化子的交，因而是双边理想。对每个这样的简单模 \(S\)，令其作用同态为
\[
\rho_S:R\longrightarrow\operatorname{End}_{\mathbb Z}(S),
\qquad \rho_S(r)(s)=rs,
\]
其中自同态取自 \(S\) 的加法群，则 \(J(R)=\bigcap_S\ker\rho_S\)。两个交都可只取简单模的同构类代表；取全部循环商 \(R/L\)、\(L\) 为极大左理想，也足以计算它们。
\end{proposition}
\begin{proof}
交的等式已在\cref{b2:prop:simple-annihilator-intersection}证明：一个元素属于所有极大左理想，恰好使它在每个简单模的每个向量上作用为零。每个简单模本来就同构于某个循环商 \(R/L\)，同构模又有相同的零化子，所以只需所述代表族。模零化子是双边理想，因为 \(xS=0\) 蕴含 \((rx)s=r(xs)=0\)、\((xr)s=x(rs)=0\)。因此这些零化子的交 \(J(R)\) 也为双边理想。作用的结合律使 \(\rho_S\) 保持乘法，且 \(\ker\rho_S=\operatorname{Ann}_R(S)\)，便得到共同核的表述。零环没有非零幺性简单模，空族的交按全环取值，等式也成立。
\end{proof}

例如，对域 \(k\) 和 \(n\ge1\)，\(R=M_n(k)\) 的列向量模 \(k^n\) 简单且忠实，所以 \(J(R)=0\)。但是 \(x=E_{12}\) 在 \(n\ge2\) 时是非零幂零元。取 \(r=E_{21}\)，就有 \(rx=E_{22}\)，它是非零幂等元，已经不再幂零。单个元素幂零，并不保证任意左倍仍幂零；这与整个理想中的元素都幂零是不同的条件。

\subsection{根中的可逆性判据}
\label{b2:subsec:1401:02}

\begin{theorem}[可逆性判据]\label{b2:thm:jacobson-unit-criterion}
设 \(R\) 为含幺环，\(J(R)\) 为其 Jacobson 根。对 \(x\in R\)，以下条件等价：
\begin{equivalences}
\item\label{b2:item:radical-member} \(x\in J(R)\)。
\item\label{b2:item:radical-unit-left} 对每个 \(r\in R\)，\(1-rx\) 为单位。
\item\label{b2:item:radical-unit-right} 对每个 \(r\in R\)，\(1-xr\) 为单位。
\end{equivalences}
此外，对任意 \(a,b\in R\)，\(1-ab\) 可逆当且仅当 \(1-ba\) 可逆；若 \(u=(1-ab)^{-1}\)，则
\[
(1-ba)^{-1}=1+bua.
\]
\end{theorem}
\begin{proof}
先证明附带的恒等事实。若 \(u=(1-ab)^{-1}\)，则 \((1-ab)u=u(1-ab)=1\) 分别给出 \(u-1-abu=0\) 与 \(u-1-uab=0\)。于是
\[
(1-ba)(1+bua)=1+b(u-1-abu)a=1,
\qquad
(1+bua)(1-ba)=1+b(u-1-uab)a=1.
\]
所以 \(1+bua\) 是 \(1-ba\) 的双侧逆；交换 \(a,b\) 便得另一方向。

对根中的元素，要先由极大左理想的定义得到左逆，再利用该左逆仍具有 \(1+J(R)\) 的形式补出右逆。设 \(j\in J(R)\)。若左理想 \(R(1-j)\) 为真，它包含在某个极大左理想 \(L\) 中，而 \(j\in L\)，导致 \(1\in L\)，矛盾。故存在 \(b\) 满足 \(b(1-j)=1\)。此时 \(b=1+bj\in1+J(R)\)。对 \(-bj\in J(R)\) 重复同一论证，有 \(c\) 使 \(cb=1\)。于是
\[
c=cb(1-j)=1-j,
\]
所以 \((1-j)b=1\)，\(1-j\) 确实可逆。若 \(x\in J(R)\)，双边理想性给出 \(rx\in J(R)\)，因而得到第二项。

若第二项成立而 \(x\notin J(R)\)，取极大左理想 \(L\) 不含 \(x\)。此时 \(L+Rx=R\)，存在 \(\ell\in L,r\in R\) 使 \(1=\ell+rx\)。所以单位 \(1-rx=\ell\) 落在真左理想中，左乘其逆即得 \(1\in L\)，矛盾。

最后将已证恒等事实用于 \(a=r,b=x\)，便得第二、三项等价。
\end{proof}

其中 \(1-ab\) 与 \(1-ba\) 的可逆性等价及逆公式通常称为 Jacobson 引理。单位判据要求 \(1-rx\) 对任意 \(r\) 都可逆，描述的是所有左倍 \(rx\) 都不破坏 \(1\) 的可逆性；它为模掉根后的单位提升提供了依据。

\begin{corollary}[根的左右对称性]\label{b2:cor:jacobson-left-right}
对任意含幺环 \(R\)，有
\[
J(R)=\bigcap_{L\text{ 极大左理想}}L
=\bigcap_{K\text{ 极大右理想}}K.
\]
\end{corollary}
\begin{proof}
极大右理想正是反环 \(R^{\mathrm{op}}\) 的极大左理想。把\cref{b2:thm:jacobson-unit-criterion}用于 \(R^{\mathrm{op}}\)，可知 \(x\) 属于全部极大右理想，当且仅当对每个 \(r\in R\)，元素 \(1-xr\) 在 \(R\) 中可逆；反环与原环的单位相同。再把同一判据用于 \(R\)，这恰是 \(x\in J(R)\) 的条件。零环按约定也满足等式。
\end{proof}

单位判据中的“对每个 \(r\)”不能省去。前面的矩阵 \(x=E_{12}\) 满足 \((1-x)^{-1}=1+x\)，但取 \(r=E_{21}\) 后，\(1-rx=1-E_{22}\) 的第二列为零，不可逆。这正是单个幂零元不必属于根的障碍。

\begin{corollary}\label{b2:cor:nil-ideal-radical}
设 \(R\) 为含幺环、\(I\subseteq R\) 为双边理想。若 \(I\) 的每个元素都幂零，则 \(I\subseteq J(R)\)。特别地，若 \(I^n=0\) 对某个正整数 \(n\) 成立，结论也成立。
\end{corollary}
\begin{proof}
对 \(x\in I,r\in R\)，理想的左乘封闭性保证 \(rx\in I\)，所以 \(rx\) 也幂零；这是单个幂零元的例子中缺少的保证。取正整数 \(m\) 使 \((rx)^m=0\)，有限几何级数给出
\[
(1-rx)^{-1}=1+rx+\cdots+(rx)^{m-1}.
\]
由可逆性判据，\(x\in J(R)\)。理想幂零时，每个元素的同一次幂都为零，故为特例。
\end{proof}

英文文献把理想中每个元素各自幂零的条件称为 nil ideal，把存在共同指数使 \(I^n=0\) 的条件称为 nilpotent ideal。本书将“幂零理想”用于后者；前者不要求各元素具有同一个幂零指数。

\ProseParagraphBreak
对域 \(k\) 和 \(n\ge1\)，在 \(k[t]/(t^n)\) 中，理想 \((\bar t)\) 幂零，故包含于根；商为域，\((\bar t)\) 又是极大理想，所以根恰为它。一般环的根未必幂零，例如 \(k[\![t]\!]\) 的单位恰为常数项非零的级数：逐次解系数可构造其逆。因此非单位元组成的理想 \((t)\) 是唯一极大理想，\(J\bigl(k[\![t]\!]\bigr)=(t)\)，但 \(t^n\ne0\) 对每个 \(n\) 成立。

\subsection{商环、矩阵环与有限长度判据}
\label{b2:subsec:1401:03}

\begin{proposition}\label{b2:prop:radical-quotient}
设 \(R\) 为含幺环、\(I\) 为其双边理想，\(\pi:R\to R/I\) 为商映射，则 \(\pi\bigl(J(R)\bigr)\subseteq J(R/I)\)。若 \(I\subseteq J(R)\)，则
\[
J(R/I)=J(R)/I.
\]
特别地，\(J\bigl(R/J(R)\bigr)=0\)。
\end{proposition}
\begin{proof}
\(R/I\) 的幺性简单左模通过 \(\pi\) 成为简单 \(R\) 模，并被 \(I\) 零化；反过来，满足 \(IS=0\) 的幺性简单 \(R\) 模 \(S\)，其作用也唯一经过 \(R/I\)。两种作用的子模相同，所以简单性不变。因此取商只保留原来简单模中被 \(I\) 零化的那些测试对象。\(J(R)\) 零化全部简单 \(R\) 模，当然也零化这些对象，由\cref{b2:prop:radical-simple-annihilators}得到包含。若 \(I\subseteq J(R)\)，则每个极大左理想都含 \(I\)；商与原像使它们与 \(R/I\) 的极大左理想一一对应。逐个取交便得等式，再取 \(I=J(R)\) 得最后断言。
\end{proof}

测试简单模的族缩小后，共同零化它们的元素可能增加，所以上述包含可以严格。例如 \(J(\mathbb Z)=0\)：若整数 \(x\ne0\)，则 \(1-3x\) 不是整数单位，单位判据排除 \(x\)。但 \(\mathbb Z/4\mathbb Z\) 的根为非零理想 \((\bar2)\)。

\begin{proposition}\label{b2:prop:matrix-radical}
设 \(R\) 为含幺环、\(n\ge1\) 为整数。每个幺性简单左 \(R\) 模 \(S\) 都给出按矩阵乘列作用的幺性简单左 \(M_n(R)\) 模 \(S^n\)。反过来，对任意幺性简单左 \(M_n(R)\) 模 \(N\)，\(E_{11}N\) 通过角环 \(E_{11}M_n(R)E_{11}\cong R\) 成为幺性简单左 \(R\) 模，且
\[
N\cong(E_{11}N)^n
\]
为左 \(M_n(R)\) 模同构。这两个构造给出两类简单模同构类的一一对应。这里 \(E_{ij}\) 为标准矩阵单位。相应地，有
\[
J\bigl(M_n(R)\bigr)=M_n\bigl(J(R)\bigr).
\]
\end{proposition}
\begin{proof}
先把两种环的简单模对应起来。取幺性简单左 \(R\) 模 \(S\)。任意非零列向量有非零坐标 \(s\)，矩阵单位可把它单独移到任意坐标，而 \(Rs=S\)，所以该向量生成全部 \(S^n\)，证明 \(S^n\) 简单。

反过来，取任意幺性简单左 \(M_n(R)\) 模 \(N\)，令 \(S=E_{11}N\)。角环 \(E_{11}M_n(R)E_{11}\) 由仅第 \((1,1)\) 个条目可能非零的矩阵组成，\(r\mapsto rE_{11}\) 是从 \(R\) 到该角环的含幺环同构，角环的单位为 \(E_{11}\)。因此按 \(r\cdot s=(rE_{11})s\) 定义的左 \(R\) 作用是幺性的。这里 \(rE_{ij}\) 表示第 \((i,j)\) 个条目为 \(r\)、其余条目为零的矩阵。

因为 \(1=\sum_iE_{ii}\)，非零 \(v\in N\) 总有某个 \(E_{ii}v\ne0\)。此时 \(E_{1i}v\ne0\)，否则再乘 \(E_{i1}\) 会得到 \(E_{ii}v=0\)，所以 \(S\ne0\)。若 \(0\ne U\subseteq S\) 是 \(R\) 子模，对任意矩阵 \(B=(b_{ji})\) 有
\[
B E_{i1}u=\sum_j E_{j1}\bigl((b_{ji}E_{11})u\bigr)
\qquad(u\in U).
\]
右端仍属于 \(\sum_jE_{j1}U\)，故这个非零子模在所有矩阵作用下稳定，由 \(N\) 简单得 \(\sum_iE_{i1}U=N\)。再施加 \(E_{11}\)，便有 \(U=S\)，所以 \(S\) 简单。

矩阵单位还给出坐标同构 \(S^n\to N\)、\((s_i)\mapsto\sum_iE_{i1}s_i\)。它的逆取 \(v\mapsto(E_{1i}v)_i\)：\(E_{1i}\) 读取第 \(i\) 个坐标，\(E_{i1}\) 将其放回原位置，由 \(E_{1j}E_{i1}=\delta_{ji}E_{11}\) 与 \(\sum_iE_{i1}E_{1i}=1\) 得到两边互逆。若 \(v=\sum_jE_{j1}s_j\)，则
\[
E_{1i}(Bv)=\sum_j(b_{ij}E_{11})s_j=\sum_j b_{ij}\cdot s_j,
\]
正是矩阵乘列的第 \(i\) 个坐标，所以这是左 \(M_n(R)\) 模同构。对原先由 \(S\) 构造的 \(S^n\)，\(E_{11}S^n\) 就是第一坐标上的 \(S\) 副本；模同构也保持这些构造，所以它们在同构类上互逆。

现在，一个矩阵属于 \(J(M_n(R))\)，当且仅当它零化全部简单模 \(S^n\)。把任意 \(s\in S\) 单独放在第 \(j\) 个坐标，就知这等价于每个条目 \(b_{ij}\) 零化全部简单 \(R\) 模。由\cref{b2:prop:radical-simple-annihilators}，又等价于所有 \(b_{ij}\in J(R)\)，得到根公式。零环的两边都为零，直接成立。
\end{proof}

环的根由极大左理想的交给出，模内全部极大子模的交也记录了所有简单商共同检测不到的部分。当这个交为零时，有限长度使这些简单商足以把整个模嵌入有限个简单模的直和，从而给出半单性；这一步中的有限性条件不能删除。

\begin{proposition}\label{b2:prop:finite-length-radical-zero}
设 \(R\) 为含幺环，\(M\) 为有限长度幺性左 \(R\) 模。若 \(M\) 的全部极大子模之交为零，则 \(M\) 半单。因此若正则左模 \({}_RR\) 有限长度，\(R/J(R)\) 为半单环。
\end{proposition}
\begin{proof}
零模直接成立。非零有限长度模有极大子模，由\cref{b2:thm:finite-length-chain-conditions,b2:prop:chain-maximal-minimal}的升链条件及极大元判据可得。为了把全部极大子模的交压缩为有限交，考察它们的所有非空有限交。降链条件使这一族中有一个极小者 \(K=\bigcap_{i=1}^r L_i\)。若 \(K\ne0\)，取 \(0\ne x\in K\)。全部极大子模交为零，故某个极大子模 \(L\) 不含 \(x\)，于是 \(K\cap L\subsetneq K\)，仍是一个有限交，矛盾。因此已有有限个极大子模满足 \(\bigcap_{i=1}^rL_i=0\)。

映射 \(M\to\bigoplus_{i=1}^rM/L_i\)、\(m\mapsto(m+L_i)_i\) 的核为 \(K=0\)，所以它单射。每个商模简单，目标半单，\cref{b2:prop:semisimple-subquotients}保证其子模也半单，故 \(M\) 半单。

令 \(B=R/J(R)\)。它作为左 \(R\) 模具有有限长度，因为有限长度对取商保持。其 \(R\) 作用已经经过 \(B\)，所以 \(R\) 子模与左 \(B\) 子模相同；这些子模在正则 \(B\) 模中又正是左理想。\cref{b2:prop:radical-quotient}给出 \(J(B)=0\)，故上述判据使它作为 \(R\) 模半单。同样因为两种标量作用的子模相同，简单性及直和分解也相同，\({}_BB\) 半单，亦即 \(B\) 是半单环。
\end{proof}

\begin{corollary}[单位的提升]\label{b2:cor:units-lift-radical}
设 \(R\) 为含幺环，\(I\) 为满足 \(I\subseteq J(R)\) 的双边理想。对任意 \(a\in R\)，\(a\) 是 \(R\) 的单位，当且仅当 \(a+I\) 是 \(R/I\) 的单位。更一般地，对任意整数 \(n\ge1\) 和矩阵 \(A\in M_n(R)\)，\(A\) 可逆，当且仅当逐条目取商所得矩阵 \(\overline A\in M_n(R/I)\) 可逆。
\end{corollary}
\begin{proof}
保持单位元的商映射将单位送到单位，所以只需证明反向。若 \(a+I\) 可逆，取其逆的一个提升 \(b\in R\)，有 \(ab=1+i\)、\(ba=1+j\)，其中 \(i,j\in I\)。由\cref{b2:thm:jacobson-unit-criterion}，\(1+i\) 与 \(1+j\) 都可逆。于是
\[
a\bigl(b(1+i)^{-1}\bigr)=1,
\qquad
\bigl((1+j)^{-1}b\bigr)a=1.
\]
一个左逆与一个右逆必相等：若 \(da=1\)、\(ac=1\)，则 \(d=d(ac)=(da)c=c\)。因此 \(a\) 有双侧逆。

逐条目取商的环同态 \(M_n(R)\to M_n(R/I)\) 满射，核为 \(M_n(I)\)，所以目标同构于 \(M_n(R)/M_n(I)\)。又由\cref{b2:prop:matrix-radical}，有 \(M_n(I)\subseteq M_n(J(R))=J(M_n(R))\)。将刚证的单位提升结论用于这个矩阵环及其理想，便得到矩阵版本。
\end{proof}

\begin{proposition}\label{b2:prop:triangular-radical}
设 \(k\) 为域、\(n\ge1\) 为整数，\(T_n(k)\) 为 \(n\times n\) 上三角矩阵组成的含幺 \(k\) 代数，\(N\) 为其中严格上三角矩阵组成的双边理想。则
\[
J(T_n(k))=N,\qquad N^n=0,\qquad T_n(k)/N\cong k^n.
\]
对 \(n\ge2\)，有 \(N^{n-1}\ne0\)，故 \(N\) 的最小幂零指数恰为 \(n\)；\(n=1\) 时 \(N=0\)。
\end{proposition}
\begin{proof}
取对角条目的映射 \(\delta:T_n(k)\to k^n\) 保持加法、乘法和单位元：上三角矩阵乘积的第 \((i,i)\) 个条目恰为两个第 \((i,i)\) 个条目的乘积。它满射，核正是 \(N\)，因此 \(N\) 是双边理想，且 \(T_n(k)/N\cong k^n\)。

严格上三角矩阵中非零条目只能沿严格递增的行列指标出现。\(n\) 个这样的矩阵相乘时，任一可能非零项都需要 \(n+1\) 个严格递增的指标，而指标总共只有 \(n\) 个，所以乘积为零，\(N^n=0\)。由\cref{b2:cor:nil-ideal-radical}，\(N\subseteq J(T_n(k))\)。\(k^n\) 的各坐标模 \(k\) 都简单，而且共同忠实，所以由\cref{b2:prop:radical-simple-annihilators}，\(J(k^n)=0\)。再用\cref{b2:prop:radical-quotient}，根在对角商中的像为零，故 \(J(T_n(k))\subseteq N\)，得到等式。

若 \(n\ge2\)，标准矩阵单位的乘积 \(E_{12}E_{23}\cdots E_{n-1,n}=E_{1n}\ne0\)，属于 \(N^{n-1}\)，所以没有小于 \(n\) 的幂零指数。\(n=1\) 时严格上三角部分为空，\(N=0\)，所述结论直接成立。
\end{proof}

% !TeX root = ../../Book2.tex
% 纲要标题：### 14.2 非交换 Nakayama 引理与生成元提升
% 纲要位置：计划纲要.md，第 1498 行。

\section{非交换 Nakayama 引理与生成元提升}
\label{b2:sec:1402}

模去 Jacobson 根后，可以先在剩余模中考察生成问题。Nakayama 引理把商模的生成性提升回来；它不要求环交换，但对一般的根理想需要模有限生成。

% 纲要要点（供目录核对）：
%
% * 对有限生成左 \(R\)-模 \(M\) 和双边理想 \(I\subseteq J(R)\)，证明 \(IM=M\Rightarrow M=0\)；整理系数核对右乘封闭性，删生成元用 \(1-a\) 的可逆性；理想幂零时无须有限生成
% * 证明 \(M=N+IM\Rightarrow M=N\) 的任意生成元提升形式；不要求子模有限生成，并区分满射判据与单射
% * 局部环中最少生成组与不可删减生成组的关系；附加半单商假设下按矩阵块容量计算最少生成元数
% * 第三卷第3.2节保留交换环行列式技巧及局部版本的完整处理；第二卷第15章只调用本节已证的非交换形式
%

\subsection{非交换 Nakayama 引理}
\label{b2:subsec:1402:01}

\begin{theorem}[非交换 Nakayama 引理]\label{b2:thm:noncommutative-nakayama}
设 \(R\) 为含幺环、\(M\) 为幺左 \(R\) 模，\(I\) 为双边理想。若 \(M\) 有限生成、\(I\subseteq J(R)\) 且 \(IM=M\)，则 \(M=0\)。若 \(I^r=0\) 对某个正整数 \(r\) 成立，则 \(IM=M\) 同样蕴含 \(M=0\)，不需要 \(M\) 有限生成。
\end{theorem}

这称为\term[Nakayama-yinli]{Nakayama 引理}[Nakayama's lemma]。\(IM\) 表示有限和 \(\sum_j a_jm_j\)、\(a_j\in I\) 构成的集合。它是左子模，因为对 \(c\in R\)，有 \(c\sum_j a_jm_j=\sum_j(ca_j)m_j\)，而 \(I\) 对左乘封闭。有限生成情形的证明要从一组最少生成元中删去一个元素；\(1-a\) 的可逆性正用来解出这个元素。

\begin{proof}
先设 \(M\) 有限生成且 \(I\subseteq J(R)\)。若 \(M\ne0\)，其有限生成组的元素个数构成非空自然数集，故可由自然数的良序性选出元素个数最少的一组 \(x_1,\ldots,x_n\)，此时 \(n\ge1\)，无须另加 Noether 条件。由 \(x_n\in IM\)，写成有限和 \(x_n=\sum_jb_jm_j\)，其中 \(b_j\in I\)，再将各 \(m_j\) 展开为 \(m_j=\sum_i r_{ji}x_i\)。于是
\[
x_n=\sum_{i=1}^n a_i x_i,\qquad
a_i=\sum_j b_jr_{ji}\in I.
\]
这里系数 \(b_j\) 位于 \(r_{ji}\) 左侧，故 \(a_i\in I\) 用的是 \(I\) 对右乘封闭。所以 \((1-a_n)x_n=\sum_{i<n}a_ix_i\)。\cref{b2:thm:jacobson-unit-criterion}保证 \(1-a_n\) 为单位，左乘其逆便把 \(x_n\) 写成其余生成元的线性组合，与最少性矛盾。\(n=1\) 时右侧为空和，这同样表明唯一生成元为零，矛盾。

若 \(I^r=0\)，则由有限和的定义及作用的结合律有 \(I(IM)=I^2M\)。迭代 \(M=IM\) 得
\[
M=IM=I^2M=\cdots=I^rM=0.
\]
这个论证不涉及生成组，因而适用于任意幺左模。
\end{proof}

对于一般的 \(I\subseteq J(R)\)，有限生成条件不能省去。设 \(k\) 为域、\(t\) 为不定元，取 \(R=k[t]_{(t)}\)，即分母在 \(t=0\) 处不为零的有理函数组成的局部环。其非单位恰为理想 \((t)\)，故 \(J(R)=(t)\)。作为 \(R\) 模，非零的 \(M=k(t)\) 却满足 \(tM=M\)。

\ProseParagraphBreak
它不有限生成。把非零有理函数写成 \(f=t^a g(t)/h(t)\)，其中 \(a\in\mathbb Z\)、\(g,h\in k[t]\) 且 \(g(0)h(0)\ne0\)，定义它在 \(t=0\) 处的阶为 \(v_t(f)=a\)，并取 \(v_t(0)=+\infty\)。乘积的阶相加，和的阶不低于各项阶的最小值：提出共同的最低次 \(t\) 幂后，加法抵消只可能使阶升高。\(R\) 的元素均有非负阶，所以若有限组 \(f_1,\ldots,f_n\) 的阶均不低于 \(-N\)，其中 \(N\ge0\)，则用 \(R\) 中元素作系数相乘并相加仍不能使阶低于 \(-N\)。它们因而不能生成 \(t^{-N-1}\)。

\subsection{生成元提升与有限生成商模}
\label{b2:subsec:1402:02}

\begin{corollary}[生成元提升]\label{b2:cor:nakayama-generators}
设 \(R\) 为含幺环、\(M\) 为幺左 \(R\) 模、\(N\subseteq M\) 为子模，\(I\) 为双边理想。假设 \(M\) 有限生成且 \(I\subseteq J(R)\)，或者 \(I^r=0\) 对某个正整数 \(r\) 成立。若 \(M=N+IM\)，则 \(M=N\)。因此，在上述任一假设下，\(M/IM\) 的一组生成元在 \(M\) 中的任意提升都生成 \(M\)；幂零理想的情形不要求 \(M\) 或该生成组有限。
\end{corollary}
\begin{proof}
在 \(M\) 有限生成且 \(I\subseteq J(R)\) 的情形，商模 \(M/N\) 由原有限生成组的像生成，所以仍有限生成，不需要 \(N\) 有限生成。等式 \(M=N+IM\) 给出 \(I(M/N)=M/N\)，由\cref{b2:thm:noncommutative-nakayama}，商模为零。若 \(I^r=0\)，则同一定理的幂零版本直接用于 \(M/N\)，也得商模为零，无须有限生成。

对于后一断言，任意选取所给商模生成元的提升 \(x_\lambda\)，并令 \(N=\sum_\lambda Rx_\lambda\)。每个 \(m+IM\) 都是这些商类的有限线性组合，故 \(m\) 减去相应提升的线性组合落在 \(IM\) 中。这就给出 \(M=N+IM\)，应用前一断言即可；整个论证没有要求提升具有额外性质。
\end{proof}

\begin{proposition}\label{b2:prop:nakayama-surjectivity-test}
设 \(R\) 为含幺环，\(M\) 为有限生成幺左 \(R\) 模、\(P\) 为任意幺左 \(R\) 模，\(f:P\to M\) 为 \(R\) 模同态，\(I\subseteq J(R)\) 为双边理想，则 \(f\) 满射当且仅当诱导的 \(P/IP\to M/IM\) 满射。
\end{proposition}
\begin{proof}
正向由代表元直接提升。反向，商映射满射意味着 \(M=f(P)+IM\)。对有限生成目标 \(M\) 应用\cref{b2:cor:nakayama-generators}，得 \(f(P)=M\)。
\end{proof}

这一判据不能用来检测单射。设 \(k\) 为域，取 \(R=k[t]/(t^2)\)、\(I=J(R)=(\bar t)\)、\(P=R\)、\(M=R/I\)，令 \(f:P\to M\) 为自然商映射。诱导映射 \(P/IP\to M/IM\) 是 \(k\to k\) 的恒等映射，但 \(\ker f=I\ne0\)。核恰好位于来源模的 \(IP\) 中，模去这一层后便完全消失，所以剩余映射的单射性不能恢复原映射的单射性。

\subsection{剩余模与极小生成集}
\label{b2:subsec:1402:03}

沿用\cref{b2:def:local-ring}的约定，局部环是非零环，且其非单位组成一个双边理想 \(\mathfrak m\)。此时 \(R/\mathfrak m\) 是除环：每个非零商类都有单位代表元。任意真左理想不含单位，所以包含于 \(\mathfrak m\)；后者已极大，因而是唯一极大左理想，\(J(R)=\mathfrak m\)。

\ProseParagraphBreak
对任意含幺环 \(R\) 及有限生成幺左 \(R\) 模 \(M\)，记 \(\mu_R(M)\) 为生成 \(M\) 所需的最少元素个数，约定 \(\mu_R(0)=0\)。元素数等于 \(\mu_R(M)\) 的生成组称为最少生成组；删去任一元素就不再生成的组称为不可删减生成组。局部环中，剩余模是除环上的向量空间，可以先求它的一组基，再由 Nakayama 引理提升回原模。
\symidx{\(\mu_R(M)\)：有限生成模所需的最少生成元个数}

\begin{proposition}\label{b2:prop:local-minimal-generators}
设 \(R\) 为非零含幺局部环、\(D=R/J(R)\) 为其剩余除环，\(M\) 为有限生成幺左 \(R\) 模。剩余模 \(M/J(R)M\) 是左 \(D\) 向量空间，以下 \(\dim_D\) 表示左 \(D\) 维数。生成 \(M\) 所需的最少元素个数为
\[
\mu_R(M)=\dim_D M/J(R)M.
\]
在上述局部环条件下，生成组不可删减，当且仅当其像为剩余模的一组左 \(D\) 基。因此不可删减生成组与最少生成组相同。
\end{proposition}
\begin{proof}
任意生成组的像都生成剩余模，所以元素数不少于右侧的维数。反向，提升一组有限左 \(D\) 基，由\cref{b2:cor:nakayama-generators}得到 \(M\) 的生成组，故最小数恰为该维数。若生成组的像线性相关，取一个有限非平凡关系 \(\sum_i a_i\bar x_i=0\)，其中 \(a_i\in D\)，并选 \(a_j\ne0\)，则
\[
\bar x_j=-a_j^{-1}\sum_{i\ne j}a_i\bar x_i.
\]
由于这是左 \(D\) 向量空间，逆元须乘在左侧。删去 \(x_j\) 后，其余元素仍生成剩余模，再用提升推论便仍生成 \(M\)，所以原组可再删减。若像为基，删去任意一项连剩余模都不能生成，原组因而不可删减。
\end{proof}

一般环中不可删减生成组未必最少。设 \(k\) 为域，取 \(R=k\times k\) 的正则左模。\((1,0),(0,1)\) 生成 \(R\)，删去任一个都会失去一个坐标方向，所以不可删减；但 \((1,1)\) 是单位，单独就生成 \(R\)。因此 \(\mu_R(R)=1\)，这组不可删减生成组却有两个元素。

\ProseParagraphBreak
对于一般环，\cref{b2:prop:radical-quotient}只保证 \(J(R/J(R))=0\)，并不保证 \(R/J(R)\) 半单。例如 \(J(\mathbb Z)=0\)，而 \(\mathbb Z\) 的正则模不是半单模：它不含非零简单子模，因为每个非零子模 \(n\mathbb Z\) 都含有真非零子模 \(2n\mathbb Z\)。下面的矩阵块公式因而需要附加半单商的假设。有限维代数的正则模有有限长度，由\cref{b2:prop:finite-length-radical-zero,b2:thm:artin-wedderburn}，其半单商才自动具有这种有限矩阵块分解。

\ProseParagraphBreak
先看 \(R=M_2(k)\)、\(S=k^2\) 及 \(M=S^3\)，其中 \(k\) 为域，\(R\) 同时左乘三个列向量。任何单个元素 \(v\) 生成的子模都是 \(R\to M\)、\(A\mapsto Av\) 的像，其 \(k\) 维数至多为 \(\dim_kR=4\)，而 \(\dim_kM=6\)，所以一个元素不能生成 \(M\)。另一方面，取标准列向量 \(e_1,e_2\)，则
\[
v=(e_1,e_2,0),\qquad w=(0,0,e_1)
\]
生成 \(M\)：\(Av=(Ae_1,Ae_2,0)\) 的前两个分量是 \(A\) 的两列，可以独立指定；\(Bw=(0,0,Be_1)\) 可以任意指定第三个分量。因此 \(\mu_R(S^3)=2\)，三个简单直和分量并不意味着需要三个模生成元。

\ProseParagraphBreak
这个现象来自正则模所含的简单模重数。若 \(B=\prod_{i=1}^sM_{n_i}(D_i)\) 为有限个矩阵块的直积，则一份正则左模 \(B\) 包含 \(n_i\) 份第 \(i\) 类简单列模 \(S_i\)。一个生成元给出一份 \(B\) 到目标模的同态，因而至多提供 \(n_i\) 份这种简单模。若目标含 \(m_i\) 份 \(S_i\)，生成元数至少为 \(\lceil m_i/n_i\rceil\)。同一组生成元要同时满足所有类型的重数要求，所以应取这些下界的最大值；下面的证明还会说明这个最大值确实足够。

\begin{proposition}\label{b2:prop:semisimple-top-generators}
设 \(R\) 为非零含幺环，\(M\) 为有限生成幺左 \(R\) 模。设正整数 \(s,n_1,\ldots,n_s\)、除环 \(D_1,\ldots,D_s\) 及非负整数 \(m_1,\ldots,m_s\) 满足
\[
R/J(R)\cong\prod_{i=1}^s M_{n_i}(D_i),\qquad
M/J(R)M\cong\bigoplus_{i=1}^s S_i^{m_i},
\]
其中 \(S_i=D_i^{n_i}\) 为相应矩阵块的简单列模，则
\[
\mu_R(M)=\max_i\left\lceil\frac{m_i}{n_i}\right\rceil.
\]
零模的最少生成元个数取零。
\end{proposition}
\begin{proof}
记半单商为 \(B\)。由\cref{b2:thm:artin-wedderburn}的正则模分解，\({}_BB\cong\bigoplus_i S_i^{n_i}\)，故 \(B^r\cong\bigoplus_i S_i^{rn_i}\)。\cref{b2:prop:semisimple-multiplicity}给出的简单模重数唯一性说明，\(\bigoplus_iS_i^{m_i}\) 是 \(B^r\) 的商，当且仅当每个 \(m_i\le rn_i\)：充分性取对应分量投影，必要性由满射的分裂及直和重数相加得到。这正是剩余模可由 \(r\) 个元素生成的条件。原模生成元映到剩余模生成元，而由\cref{b2:cor:nakayama-generators}，后者的提升又生成原模，所以两者的最少数相同，得到公式。
\end{proof}

交换局部环的常用 Nakayama 引理已经包含在以上结论中。交换性还允许另一种证明方法：用有限生成元的关系矩阵及其伴随矩阵，把若干线性关系转成一个同时零化全部生成元的标量等式。\cref{b2:prop:adjugate-identity}提供相应的矩阵恒等式；第三卷第3.2节将把这种行列式技巧及其整性应用完整展开。

\ProseParagraphBreak
对于有限维非交换代数，下一章先使根的幂稳定，再把稳定项看成有限生成左模，用\cref{b2:thm:noncommutative-nakayama}证明该项为零。

% !TeX root = ../../Book2.tex
% 纲要标题：### 14.3 本原环与半本原环
% 纲要位置：计划纲要.md，第 1505 行。

\section{本原环与半本原环}
\label{b2:sec:1403}

模去 Jacobson 根以后，每个非零环元素至少能被某个简单模检测到。更强的问题是，能否用同一个简单模同时检测全部元素？本原环要求后者，半本原环只要求前者。它们与半单环的区别，将通过具体作用和无限乘积表现出来。

% 纲要要点（供目录核对）：
%
% * 忠实简单模与本原环；区分单个忠实作用、一族共同忠实作用及半单性
% * 半本原环的次直积结构；取极大左理想中最大的双边理想形成真正的本原商环
% * 半本原但非半单的实例；有限坐标满射与无限直积满射的区别
%

\subsection{忠实简单模与本原环}
\label{b2:subsec:1403:01}

\begin{definition}\label{b2:def:primitive-semiprimitive}
设 \(R\) 为非零含幺环。若 \(R\) 有忠实的幺性简单左模，则称它为\term[zuo-benyuan-huan]{左本原环}[left primitive ring]。若 \(J(R)=0\)，则称它为\term[banbenyuan-huan]{半本原环}[semiprimitive ring]。本节简称“本原”时均指左本原；右本原是另一条件，一般不能自动互换\footnote{右本原而非左本原的例子见\cite[pp. 473--475]{Bergman1964PrimitiveOneSided}；对该环取反环得到相反方向的例子。}。
\end{definition}

由\cref{b2:prop:radical-simple-annihilators}，半本原性等价于：对每个 \(0\ne r\in R\)，都存在一个幺性简单左模 \(S\) 和其中的向量 \(s\)，使 \(rs\ne0\)。这里 \(S\) 可以随 \(r\) 改变。本原性则要求先找到同一个简单模 \(S\)，使对每个 \(0\ne r\in R\)，都能在这个 \(S\) 中找到 \(s\) 满足 \(rs\ne0\)。因此每个本原环都半本原。对任意简单左模 \(S\)，商环 \(R/\operatorname{Ann}_R(S)\) 在 \(S\) 上的作用忠实且简单，所以它本原。

\ProseParagraphBreak
这里的“本原环”与下一章\cref{b2:def:primitive-idempotent}中的“本原幂等元”有不同的定义。前者要求一个忠实简单模，后者要求一个非零幂等元不能继续分成两个非零正交幂等元之和，不能仅凭同一个名称互换条件。

\begin{proposition}\label{b2:prop:commutative-primitive-field}
非零交换含幺环 \(R\) 左本原，当且仅当 \(R\) 是域。
\end{proposition}
\begin{proof}
若 \(S\) 忠实且简单，取 \(0\ne s\in S\)，满射 \(R\to S\)、\(r\mapsto rs\) 的核为极大理想 \(\mathfrak m\)，所以 \(S\cong R/\mathfrak m\)。交换性使这个商模的零化子恰为 \(\mathfrak m\)，忠实性便给出 \(\mathfrak m=0\)，故 \(R\) 是域。反向，一个域在自身上的正则模简单且忠实。
\end{proof}

在非交换情形，对任意除环 \(D\) 和整数 \(n\ge1\)，矩阵环 \(M_n(D)\) 都本原：列模 \(D^n\) 简单且忠实，简单性由矩阵单位移动非零坐标并乘除环元素得到。它在 \(n>1\) 时不是除环，说明上一个命题的交换性不能去掉。

\ProseParagraphBreak
非零半单环都半本原：左正则模是简单模的直和，而 \(J(R)\) 零化每个简单分量，所以 \(J(R)R=0\)，取右侧元素为 \(1\) 得 \(J(R)=0\)。但半单环未必本原：对任意域 \(k\)，环 \(k\times k\) 是两个域的乘积，因而半单；它是交换环却不是域，由上一个命题可知它不本原。

\subsection{半本原环的次直积结构}
\label{b2:subsec:1403:02}

\begin{definition}\label{b2:def:subdirect-product}
设 \(I\) 为指标集，\(R\) 与 \(R_i\)、\(i\in I\) 为含幺环，\(\phi:R\to\prod_{i\in I}R_i\) 为保幺环同态。若 \(\phi\) 单射，且它与每个因子投影的复合都满射，则称 \(\phi\) 把 \(R\) 表示为这些环的一个\term[cizhiji]{次直积}[subdirect product]。这里不要求 \(\phi\) 到整个直积满射。
\end{definition}

单射使像成为直积的一个子环，坐标满射使每个单独的因子都能取遍；不同坐标仍可能受到兼容条件限制。例如对任意域 \(k\)，对角嵌入 \(k\to k\times k\)、\(a\mapsto(a,a)\) 的两个坐标映射均满射，但它的像只包含两个坐标相等的元素。

\begin{theorem}\label{b2:thm:semiprimitive-subdirect}
非零含幺环 \(R\) 半本原，当且仅当它可以表示为一族左本原环的次直积。
\end{theorem}
\begin{proof}
对每个极大左理想 \(L\)，\(R/L\) 是简单左模，但 \(L\) 未必是双边理想，不能直接把它当作商环。为了得到一个在这个简单模上忠实作用的环商，应当模去作用的核，于是令
\[
I_L=\operatorname{Ann}_R(R/L)
=\{a\in R:aR\subseteq L\}.
\]
这个集合对加减封闭；若 \(aR\subseteq L\)，则对 \(r,x\in R\)，有 \((ra)x=r(ax)\in L\) 以及 \((ar)x=a(rx)\in L\)，故 \(I_L\) 是双边理想。取 \(x=1\) 可知 \(I_L\subseteq L\)。若双边理想 \(I\subseteq L\)，则每个 \(a\in I\) 都满足 \(aR\subseteq I\subseteq L\)，所以 \(I\subseteq I_L\)。这证明 \(I_L\) 是 \(L\) 中最大的双边理想。

商环 \(R/I_L\) 在 \(R/L\) 上的作用忠实，因为模去的恰是原作用的核；商前后子模相同，所以作用仍简单，故 \(R/I_L\) 本原。全部商映射合成
\[
\phi:R\longrightarrow\prod_L R/I_L.
\]
每个坐标映射满射，整体的核是 \(\bigcap_LI_L=J(R)\)，由\cref{b2:prop:radical-simple-annihilators}已知。故 \(J(R)=0\) 时它就是次直积表示。极大左理想组成集合，因此这里的乘积是集合编号的乘积。

反过来，设 \(\phi:R\to\prod_iR_i\) 为次直积表示，各 \(R_i\) 本原且半本原。每个坐标满射把 \(J(R)\) 映进 \(J(R_i)=0\)，这是\cref{b2:prop:radical-quotient}对其核所给的包含。于是 \(\phi\bigl(J(R)\bigr)=0\)，而 \(\phi\) 单射，故 \(J(R)=0\)。
\end{proof}

整数的模素数映射给出一个同时满足所有有限坐标要求、却不能任意指定无限坐标的例子：
\[
\mathbb Z\longrightarrow\prod_{p\text{ 素数}}\mathbb F_p.
\]
它单射，因为一个非零整数不可能被所有素数整除；各坐标显然满射。对任意有限多个不同素数，第一卷定理10.4.3的中国剩余定理允许同时任意指定这些坐标。但一个无限坐标族未必来自同一个整数：将奇素数坐标全取零、模二坐标取一就不可能提升。任何可能提升的整数须被所有奇素数整除，只能为零，又与模二条件矛盾。因此有限组同余条件的可解性，不能直接推广成对整个无限直积的满射性。

\subsection{非半单的半本原环}
\label{b2:subsec:1403:03}

\(\mathbb Z\) 的根为零，已在\secref{b2:sec:1401}由单位判据算出；它的正则模不半单，也已在\secref{b2:sec:1402}用不存在非零简单子模说明。另一方面，它不是域，由\cref{b2:prop:commutative-primitive-field}可知也不本原。所以这个半本原环既不本原，也不半单。

\ProseParagraphBreak
本原也不蕴含半单。令 \(k\) 为域，\(V\) 为有可数基 \(e_1,e_2,\ldots\) 的 \(k\) 向量空间，取 \(R=\operatorname{End}_k(V)\)。任意 \(0\ne v\in V\) 都可扩充为基；把 \(v\) 送到任意给定的 \(w\)，把这组基中的其余向量送到零，就唯一确定一个线性算子。无限维时扩充基使用第一卷附录A定理A.1.4的 Zorn 引理，和构造一般向量空间的基相同。故 \(Rv=V\)，自然模 \(V\) 简单；作用按定义忠实，所以 \(R\) 本原。

\ProseParagraphBreak
要检测这个环是否左 Artin，可以要求算子杀掉越来越多的指定基向量。这样的条件给出左理想
\[
L_n=\bigl\{\,T\in R:T(e_1)=\cdots=T(e_n)=0\,\bigr\}
\]
条件每增加一个，允许的算子就减少，形成无限严格降链 \(R\supsetneq L_1\supsetneq L_2\supsetneq\cdots\)。这些集合对加减和左乘稳定，因为复合 \(ST\) 仍将指定的基向量送到零；\(L_n\supsetneq L_{n+1}\) 的严格性可由只将 \(e_{n+1}\) 送到 \(e_1\)、将其余基向量送到零的算子核对，\(R\supsetneq L_1\) 则由恒等算子核对。因此 \(R\) 不是左 Artin 环，而\cref{b2:thm:semisimple-ring-characterizations}中的半单环均为左 Artin 环。这个例子说明本原性也不蕴含半单性，并且忠实简单模可能无限维。这里在有限组向量上指定算子的行为，正是下一节有限拓扑所记录的条件。

% !TeX root = ../../Book2.tex
% 纲要标题：### 14.4 Jacobson 稠密性定理
% 纲要位置：计划纲要.md，第 1229 行。

\section{Jacobson 稠密性定理}
\label{b2:sec:1404}

在一个简单模上，环元素可以把任意非零向量送到任何向量。稠密性定理把这个单向量事实推广为有限组向量的同时插值。能够自由指定的向量组须对自同态除环线性无关，而不只是对某个较小的底域无关。

% 纲要要点（供目录核对）：
%
% * 简单模上的作用与自同态除环
% * 稠密性定理的代数形式
% * 有限维情形的矩阵代数结论
%

\subsection{简单模的作用与自同态除环}
\label{b2:subsec:1404:01}

设 \(S\) 为简单左 \(R\) 模。由\cref{b2:lem:schur}，\(\Delta=\operatorname{End}_R(S)\) 是除环。为了把自同态看作右标量，规定
\[
D=\Delta^{\mathrm{op}},\qquad v\cdot d=f_d(v),
\]
其中 \(f_d\) 是 \(d\) 对应的 \(R\) 模自同态。因为 \(f_{d_1d_2}=f_{d_2}f_{d_1}\)，这给出右 \(D\) 向量空间结构。左 \(R\) 作用与它交换，所以作用同态的目标可缩为
\[
\rho:R\longrightarrow\operatorname{End}_D(S_D).
\]
这里 \(\operatorname{End}_D\) 表示右 \(D\) 线性自同态。
\symidx{\(D=\operatorname{End}_R(S)^{\mathrm{op}}\)：简单左模的右标量除环}

一个有限向量组 \(v_1,\ldots,v_n\) 右 \(D\) 线性无关，意指 \(\sum_i v_i d_i=0\) 只可能在全部 \(d_i=0\) 时成立。若它们原本有关系，任何 \(r\in R\) 都必保留该关系，故不能任意指定它们的像。

\ProseParagraphBreak
例如，令 \(R=\mathbb C\) 在 \(S=\mathbb C\) 上左乘，并同时把 \(S\) 看成二维实空间。向量 \(1,i\) 对实数线性无关，却满足右复关系 \(1\cdot i-i\cdot1=0\)。不可能有同一个左乘把 \(1\) 送到零而把 \(i\) 送到一。这就是必须使用正确标量除环的原因。

\subsection{稠密性定理的代数形式}
\label{b2:subsec:1404:02}

\begin{definition}\label{b2:def:finite-topology}
设 \(D\) 为除环、\(S\) 为右 \(D\) 向量空间。对 \(T\in\operatorname{End}_D(S_D)\) 和有限组 \(v_1,\ldots,v_m\in S\)，令
\[
U(T;v_1,\ldots,v_m)
=\bigl\{\,F\in\operatorname{End}_D(S_D):F(v_i)=T(v_i)\text{ 对所有 }i\,\bigr\}.
\]
以这些集合为各 \(T\) 的基本邻域，所得拓扑称为\term[youxian-tuopu]{有限拓扑}[finite topology]。有限个这样的条件可以合并成一组向量条件，因此确实给出邻域基。
\end{definition}

\begin{theorem}[Jacobson 稠密性]\label{b2:thm:jacobson-density}
设 \(R\) 为含幺环、\(S\) 为简单左 \(R\) 模，\(D=\operatorname{End}_R(S)^{\mathrm{op}}\) 按 \(s\cdot f^{\mathrm{op}}=f(s)\) 作用于 \(S\) 的右侧。对任意右 \(D\) 线性无关的有限组 \(v_1,\ldots,v_n\)，以及任意 \(w_1,\ldots,w_n\in S\)，存在 \(r\in R\) 使
\[
rv_i=w_i\qquad(1\le i\le n).
\]
\end{theorem}

这称为\term[Jacobson-choumixing-dingli]{Jacobson 稠密性定理}[Jacobson density theorem]。存在的 \(r\) 通常不唯一；定理只在给定的有限组向量上指定其作用。

\begin{proof}
对 \(n\) 归纳。\(n=1\) 时，\(v_1\ne0\)，简单性给出 \(Rv_1=S\)，结论成立。设已对 \(n-1\) 成立。考虑左理想
\[
L=\{\,a\in R:av_1=\cdots=av_{n-1}=0\,\}.
\]
\(Lv_n\) 是 \(S\) 的子模，因此或为零，或为 \(S\)。先排除第一种情况。

若 \(Lv_n=0\)，归纳假设说明映射 \(a\mapsto(av_1,\ldots,av_{n-1})\) 从 \(R\) 满射到 \(S^{n-1}\)，因而可以定义
\[
F:S^{n-1}\longrightarrow S,
\qquad F(av_1,\ldots,av_{n-1})=av_n.
\]
若两个 \(a\) 给出相同输入，它们之差属于 \(L\)，所以输出之差也为零。这证明良定义。使用代表元的加法及左乘，得到 \(F\) 左 \(R\) 线性。令 \(f_i:S\to S\) 是 \(F\) 在第 \(i\) 个直和分量上的限制，则 \(f_i\in\Delta\)，且
\[
v_n=F(v_1,\ldots,v_{n-1})=\sum_{i<n}f_i(v_i)=\sum_{i<n}v_i d_i
\]
对某些 \(d_i\in D\) 成立，与给定向量组右线性无关矛盾。因此 \(Lv_n=S\)。

现在先用归纳假设找 \(r_0\) 使 \(r_0v_i=w_i\)、\(i<n\)。再取 \(\ell\in L\) 使 \(\ell v_n=w_n-r_0v_n\)。令 \(r=r_0+\ell\)，则前 \(n-1\) 个值因 \(\ell\in L\) 保持不变，最后一个值也达到要求，归纳完成。
\end{proof}

更一般地，给定任意右 \(D\) 线性算子 \(T\) 及有限组向量，先在它们的右线性生成空间中选基，再对这组基应用定理，即可找到 \(r\) 在原来的全部向量上与 \(T\) 一致。因此每个基本邻域 \(U(T;v_1,\ldots,v_m)\) 都与 \(\rho(R)\) 相交，即
\[
\overline{\rho(R)}^{\,\mathrm{fin}}=\operatorname{End}_D(S_D).
\]
有限拓扑中的稠密性只要求有限组上的同时一致；当 \(S_D\) 无限维时，它不要求某个 \(r\) 在整个空间上等于 \(T\)。

\subsection{有限维情形的矩阵代数}
\label{b2:subsec:1404:03}

\begin{corollary}\label{b2:cor:density-finite-dimensional}
设 \(R\) 为含幺环、\(S\) 为简单左 \(R\) 模，并令 \(D=\operatorname{End}_R(S)^{\mathrm{op}}\) 按 \(s\cdot f^{\mathrm{op}}=f(s)\) 作用于 \(S\) 的右侧。若 \(\dim_D S=n<\infty\)，则
\[
R/\operatorname{Ann}_R(S)\cong\operatorname{End}_D(S_D)\cong M_n(D).
\]
特别地，本原环若有一个对其自同态右除环有限维的忠实简单模，则它是单 Artin 环。
\end{corollary}
\begin{proof}
取右 \(D\) 基 \(v_1,\ldots,v_n\)。对任意 \(T\in\operatorname{End}_D(S)\)，在\cref{b2:thm:jacobson-density}中指定 \(w_i=T(v_i)\)，便找到 \(r\) 在全部基向量上与 \(T\) 相同。两者右线性，所以处处相同，作用同态满射。它的核为 \(\operatorname{Ann}_R(S)\)，由环同构定理得到第一个同构。

把向量写成右坐标列 \(\sum_i v_i d_i\)，若 \(T(v_j)=\sum_i v_i a_{ij}\)，则新坐标为 \((\sum_j a_{ij}d_j)_i\)。矩阵从左作用于右坐标列，复合对应通常的矩阵乘法，因此得到 \(M_n(D)\)，不再取反环。矩阵环的单性和 Artin 性由\cref{b2:prop:simple-artinian}及其矩阵块说明得到。
\end{proof}

这并没有证明每个本原环都为矩阵环；有限维条件是将有限组插值变成整个算子的等式所必需的。在前一节的无限维例子中，没有一个有限向量组能充当整个空间的基。

% !TeX root = ../../Book2.tex
% 纲要标题：### 14.5 Burnside 矩阵代数定理
% 纲要位置：计划纲要.md，第 1235 行。

\section{Burnside 矩阵代数定理}
\label{b2:sec:1405}

代数闭域上的有限维不可约作用具有更强的结论：自同态除环只剩底域标量，因而稠密性直接给出全部矩阵。这个结论通常称为 Burnside 矩阵代数定理；它对群作用也有用，但陈述的对象是线性算子组成的代数。

% 纲要要点（供目录核对）：
%
% * 代数闭域上有限维不可约矩阵代数
% * 与稠密性定理的联系
% * 区分轨道计数引理及有限群 p^a q^b 可解性定理
%
% ---
%

\subsection{代数闭域上的不可约矩阵代数}
\label{b2:subsec:1405:01}

\begin{theorem}[Burnside 矩阵代数]\label{b2:thm:burnside-matrix-algebra}
设 \(k\) 为代数闭域，\(V\ne0\) 有限维，\(A\subseteq\operatorname{End}_k(V)\) 为含恒等算子的 \(k\) 子代数。若 \(V\) 没有非零真 \(A\) 稳定子空间，则 \(A=\operatorname{End}_k(V)\)。
\end{theorem}
\begin{proof}
假设使 \(V\) 成为简单左 \(A\) 模。由\cref{b2:cor:schur-algebraically-closed}，\(\operatorname{End}_A(V)=k\operatorname{id}_V\)：有限维与代数闭性使每个自同态有一个特征值，Schur 引理再迫使该自同态等于相应标量。因此本节的右标量除环 \(D=\operatorname{End}_A(V)^{\mathrm{op}}\) 就是 \(k\)。

\cref{b2:cor:density-finite-dimensional}于是说明作用映射 \(A\to\operatorname{End}_k(V)\) 满射。而这个映射原本就是子代数的包含，故 \(A=\operatorname{End}_k(V)\)。
\end{proof}

这一\term[Burnside-juzhen-daishu-dingli]{Burnside 矩阵代数定理}[Burnside's matrix algebra theorem]没有假设 \(k\) 的特征为零。有限维保证特征多项式有定义且次数为正，代数闭性保证它有根，再由 Schur 引理迫使上述自同态为标量。

\ProseParagraphBreak
例如，对域 \(k\) 上的一族矩阵，若其生成的含幺代数作用不可约，且 \(k\) 代数闭，则这些矩阵的有限乘积与线性组合生成全部矩阵空间。结论是代数生成，不是说原先有限个矩阵本身已经线性生成了全部矩阵。

\subsection{与稠密性定理的联系}
\label{b2:subsec:1405:02}

一般稠密性使用右除环 \(D=\operatorname{End}_A(V)^{\mathrm{op}}\)。上一定理先把它识别为 \(k\)，然后在有限基上同时指定像。代数闭假设若删去，结果可以改变：复数左乘组成 \(M_2(\mathbb R)\) 中的子代数
\[
\left\{
\,\begin{pmatrix}a&-b\\b&a\end{pmatrix}
\ \middle|\ a,b\in\mathbb R\,
\right\}.
\]
它在实平面上不可约，因为任何非零真稳定子空间只能是一条实直线，却不可能对乘 \(i\) 的九十度旋转稳定。但这个代数只有二维，严格小于四维的 \(M_2(\mathbb R)\)。稠密性对此给出的目标是复线性自同态，恰好就是原代数。

\ProseParagraphBreak
有限维也不能直接删除。设 \(V\) 有可数基 \(e_1,e_2,\ldots\)，\(F\) 为全部有限秩算子组成的双边理想，取 \(A=k\operatorname{id}+F\)。它含有将任意非零向量送到任意目标向量的秩一算子，所以作用不可约。可是右移算子 \(T(e_i)=e_{i+1}\) 不属于 \(A\)：对每个 \(\lambda\in k\)，\(T-\lambda\operatorname{id}\) 都单射，因为有限支集向量的最后一个非零坐标在右移后产生一个无法被抵消的末项。因此其像无限维，不可能是有限秩算子。

\ProseParagraphBreak
这个例子仍能满足有限组插值：把有限组向量所生成的子空间扩充为直和分量，在该有限维部分按要求定义算子，在一个补空间上取零，就得到所需有限秩算子。因此“在任意指定的有限组上一致”与“作为整个空间上的算子相等”确实有不同的含义。

\ProseParagraphBreak

文献中 Burnside 的名字还用于有限群作用的轨道计数引理，以及“阶为 \(p^a q^b\) 的有限群可解”的定理。本节的矩阵代数定理与这两个陈述分别讨论不同对象，不能只凭同名互相引用。特别地，矩阵代数定理没有关于群阶只含两个素因子的条件，也不直接给出群的可解性。

\ProseParagraphBreak
本定理的有限维代数闭域版本，可对照 Etingof 等人的\cite[Section 3.2, Corollary 3.2.1, Theorem 3.2.2(i), p. 43]{EtingofEtAl2011RepresentationTheory}；该书将它列为稠密性定理的一部分。后续研究有限群表示时，群代数的像正是一类矩阵代数，因此这个结果会进入表示论；本章给出的证明只依赖 Schur 引理、自同态的特征值和\cref{b2:thm:jacobson-density}。


\begin{chaptersummary}
Jacobson 根既是全部极大左理想的交，也是全部极大右理想及简单左模零化子的交，所以是双边理想。\(x\in J(R)\) 的单位判据要求所有 \(1-rx\) 可逆；只检验 \(1-x\)，甚至只检验底域标量倍，都可能遗漏非交换环中的障碍。根在一般商下只保证像包含于商环的根，只有商掉的理想包含于原根时才得到等式。

\ProseParagraphBreak
Nakayama 引理把根的可逆性用于有限生成组的消减，继而给出商模生成元的提升。局部环上，最少生成元数由一个除环上的剩余维数计算；半单商有多个矩阵块时，则需分别比较简单因子重数与矩阵块大小。一般根理想下，有限生成性不能删除；若理想本身幂零，则可以反复代入直到理想的幂为零，不再需要有限生成。

\ProseParagraphBreak
半本原环可由一族本原商环的次直积描述，本原环则要求同一个简单模忠实。对简单左模 \(S\)，置 \(D=\operatorname{End}_R(S)^{\mathrm{op}}\)。稠密性定理及其关系判据说明，有限组指定像能够由同一个环元素实现，当且仅当它们保留原向量间的全部右 \(D\) 线性关系。因此作用像在有限拓扑中的闭包为 \(\operatorname{End}_D(S_D)\)，这个闭包未必等于作用像本身。只有当该右除环上的维数有限时，有限组条件才覆盖整个算子；Burnside 定理再利用代数闭性把右除环缩为底域，从而得到全部矩阵。下一章将把 Jacobson 根与有限维性结合，研究根滤过和投射覆盖。
\end{chaptersummary}

\BookExercises

\begin{exercise}\label{b2:ex:14-scalar-unit-test}
设 \(k\) 为任意域。给出 \(x\in M_2(k)\)，使 \(1-\lambda x\) 对每个 \(\lambda\in k\) 都可逆，但 \(x\notin J\bigl(M_2(k)\bigr)\)。明确指出全环单位判据怎样排除这个 \(x\)。
\end{exercise}
\begin{solution}
取 \(x=E_{12}\)，则 \(x^2=0\)，所以 \((1-\lambda x)^{-1}=1+\lambda x\) 对任意标量成立。但取环元素 \(r=E_{21}\)，有 \(rx=E_{22}\)，故 \(1-rx=\operatorname{diag}(1,0)\) 不可逆。由\cref{b2:thm:jacobson-unit-criterion}，\(x\notin J\bigl(M_2(k)\bigr)\)。标量倍只检验了全环中很小的一部分元素，不能替代对全部 \(r\) 的条件。
\end{solution}

\begin{exercise}\label{b2:ex:14-integer-quotient-radical}
设 \(n=\prod_{i=1}^s p_i^{a_i}\ge2\)，其中 \(p_i\) 是两两不同的素数，\(a_i\) 为正整数。求 \(J(\mathbb Z/n\mathbb Z)\)，并求使其幂为零的最小正整数。
\end{exercise}
\begin{solution}
整数商环的极大理想是 \((\bar p_i)\)：其商为 \(\mathbb F_{p_i}\)，反向每个极大理想的原像是包含 \(n\mathbb Z\) 的整数极大理想，故只能是这些素数理想。因此根由
\[
\overline{p_1\cdots p_s}
\]
生成。其 \(r\) 次幂为零当且仅当 \(n\) 整除 \((p_1\cdots p_s)^r\)，即每个 \(r\ge a_i\)。最小正整数为 \(\max_i a_i\)。尤其 \(n\) 无平方因子时，根为零，此时最小正整数为一。
\end{solution}

\begin{exercise}\label{b2:ex:14-product-radical}
设 \((R_i)_{i\in I}\) 为由集合 \(I\) 编号的含幺环族。证明其直积满足
\[
J\left(\prod_{i\in I}R_i\right)=\prod_{i\in I}J(R_i).
\]
论证应同时适用于无限指标集。
\end{exercise}
\begin{solution}
直积中的元素可逆，当且仅当每个坐标可逆：逆元只能逐坐标取逆，而这些逆确实组成直积元素。若 \(x_i\in J(R_i)\)，则对任意 \(r=(r_i)\)，各 \(1-r_ix_i\) 都可逆，故 \(1-rx\) 可逆，得到正向包含。反过来，若 \(x\) 属于直积环的根，固定 \(i\) 及 \(a\in R_i\)，令 \(r\) 仅第 \(i\) 坐标为 \(a\)，其余为零。\(1-rx\) 可逆迫使 \(1-ax_i\) 可逆。对任意 \(a\) 使用单位判据，得 \(x_i\in J(R_i)\)，证明反向包含。没有一步要求元素只有有限支集。
\end{solution}

\begin{exercise}\label{b2:ex:14-matrix-root-powers}
设 \(R\) 为含幺环，\(I\) 为其双边理想，\(n,r\) 为正整数。证明 \(M_n(I)^r=M_n(I^r)\)。据此计算 \(J\bigl(M_3(\mathbb Z/12\mathbb Z)\bigr)\) 及其幂零指数。
\end{exercise}
\begin{solution}
矩阵乘积的条目是 \(r\) 个 \(I\) 元素乘积的有限和，故有包含 \(M_n(I)^r\subseteq M_n(I^r)\)。反向，\(I^r\) 中的每个元素为 \(a_1\cdots a_r\) 的有限和，且
\[
(a_1E_{i1})\left(\prod_{\nu=2}^{r-1}(a_\nu E_{11})\right)(a_rE_{1j})
=a_1\cdots a_r E_{ij}
\]
在 \(r\ge2\) 时成立，其中 \(r=2\) 时中间的空积取单位矩阵；\(r=1\) 直接成立。把任意矩阵拆成这些矩阵单位位置的和，得到反向包含。

\(J(\mathbb Z/12\mathbb Z)=(\bar6)\)，非零且平方为零。由\cref{b2:prop:matrix-radical}，所求根为全部条目在 \((\bar6)\) 中的三阶矩阵。上式说明其平方为零，而 \(\bar6E_{11}\ne0\)，所以幂零指数为二。
\end{solution}

\begin{exercise}\label{b2:ex:14-radical-general-map}
根在满同态下的像包含于目标的根。给出一个非满的保单位元环同态，使这个包含失败。可以使用局部环 \(k[t]_{(t)}\) 到其分式域的包含。
\end{exercise}
\begin{solution}
令 \(R=k[t]_{(t)}\subseteq k(t)\)。在 \(R\) 中，分母在零点非零；一个分式可逆当且仅当分子也在零点非零。因此非单位为理想 \((t)\)，它是唯一极大理想，\(J(R)=(t)\)。目标 \(k(t)\) 为域，根为零；包含映射却把 \(t\in J(R)\) 送到非零元素 \(t\)。所以包含失败。满同态的证明依赖目标简单模限制后仍简单；一般包含映射没有这一保证。
\end{solution}

\begin{exercise}\label{b2:ex:14-units-lift}
设 \(k=\mathbb F_q\)，\(R=k[\varepsilon]/(\varepsilon^3)\)，以下仍用 \(\varepsilon\) 表示其商类。求 \(|R^\times|\) 与 \(|\operatorname{GL}_2(R)|\)，并写出矩阵
\[
A=\begin{pmatrix}1&\varepsilon\\\varepsilon&1+\varepsilon\end{pmatrix}
\]
的逆。说明从 \(\operatorname{GL}_2(R)\) 到 \(\operatorname{GL}_2(k)\) 的约化映射的核。
\end{exercise}
\begin{solution}
理想 \(I=(\varepsilon)\) 满足 \(I^3=0\)，故包含于 \(J(R)\)。由\cref{b2:cor:units-lift-radical}，元素为单位当且仅当其常数项非零，所以 \(|R^\times|=(q-1)q^2\)。同一推论使矩阵约化映射满射：任意可逆的 \(k\) 矩阵任选逐项提升都可逆。其核恰为
\[
1+M_2(I)=\{\,I_2+B\mid B\in M_2(I)\,\}.
\]
每个这样的矩阵确实可逆，因为 \(B^3=0\)。\(I\) 有 \(q^2\) 个元素，四个条目可独立选取，故核有 \(q^8\) 个元素。因此
\[
|\operatorname{GL}_2(R)|=q^8(q^2-1)(q^2-q),
\]
其中两因子分别数出非零第一列及不在其线性生成空间中的第二列。

写 \(A=I_2+\varepsilon C\)，其中 \(C=\begin{pmatrix}0&1\\1&1\end{pmatrix}\)。由 \(\varepsilon^3=0\)，
\[
A^{-1}=I_2-\varepsilon C+\varepsilon^2C^2
=\begin{pmatrix}
1+\varepsilon^2&-\varepsilon+\varepsilon^2\\
-\varepsilon+\varepsilon^2&1-\varepsilon+2\varepsilon^2
\end{pmatrix}.
\]
这个有限几何级数公式在特征二时也成立，整数系数按 \(k\) 的特征解释。
\end{solution}

\begin{exercise}\label{b2:ex:14-triangular-radical}
设 \(k\) 为域，考虑分块上三角矩阵环
\[
R=\left\{\,
\begin{pmatrix}A&B\\0&C\end{pmatrix}
\ \middle|\ A\in M_2(k),\ B\in M_{2\times3}(k),\ C\in M_3(k)
\,\right\}.
\]
求 \(J(R)\)、其幂零指数及 \(R/J(R)\)。比较它与通常的 \(T_5(k)\)：为什么矩阵总大小都是五，根的幂零指数却不同？
\end{exercise}
\begin{solution}
取 \(N\) 为两个对角块都为零的矩阵集合。它是双边理想，且 \(N^2=0\)。由\cref{b2:cor:nil-ideal-radical}，\(N\subseteq J(R)\)。对角块映射给出
\[
R/N\cong M_2(k)\times M_3(k).
\]
两块的自然列模分别简单且忠实。把它们看成直积环的两个简单模，零化子之交为零，故 \(J(R/N)=0\)。由\cref{b2:prop:radical-quotient}，\(J(R)/N=0\)，所以 \(J(R)=N\)。\(N\ne0\) 且 \(N^2=0\)，其幂零指数为二。

\cref{b2:prop:triangular-radical}给出 \(T_5(k)\) 的根幂零指数为五。它的严格上三角根中有从第一位置经第二、三、四到第五位置的四次非零乘积；在这里的两块环中，根只连接第二块到第一块，两个根元素的乘积已经为零。指数取决于根中允许的乘积，不由矩阵的总大小单独决定。
\end{solution}

\begin{exercise}\label{b2:ex:14-nakayama-radical-needed}
取 \(R=k\times k\)、\(I=k\times0\)、\(M=I\)，均采用坐标乘法。证明 \(M\) 为非零有限生成模且 \(IM=M\)。说明这为何不与 Nakayama 引理矛盾。
\end{exercise}
\begin{solution}
\(M\) 由 \((1,0)\) 生成，故非零且循环；\(I^2=I\)，所以 \(IM=M\)。但 \(R\) 的两个简单坐标模的零化子交为零，故 \(J(R)=0\)，非零理想 \(I\) 不包含于根。也可以直接取 \(x=(1,0)\in I\)，此时 \(1-x=(0,1)\) 不可逆，单位判据排除 \(x\in J(R)\)。Nakayama 引理的根理想假设未满足。
\end{solution}

\begin{exercise}\label{b2:ex:14-nakayama-finite-needed}
设 \(R=k[t]_{(t)}\)，\(M=k(t)/R\)。证明 \(M\ne0\)、\(tM=M\)，且每个元素都被某个 \(t\) 的幂零化。证明 \(M\) 不有限生成，并说明 \(M/J(R)M\) 的值。
\end{exercise}
\begin{solution}
\(t^{-1}\notin R\)，所以 \(M\ne0\)。任意 \(f+R\) 都是 \(t(f/t+R)\)，故 \(tM=M\)。把有理函数的分母写成 \(t^n\cdot q(t)\)、\(q(0)\ne0\)，便有 \(t^nf\in R\)，说明每个陪集被某个 \(t^n\) 零化。

若有限个陪集生成 \(M\)，取共同指数 \(N\) 零化这些生成元，交换性使 \(t^N\) 零化它们的全部 \(R\) 线性组合。但 \(t^{-N-1}+R\) 乘 \(t^N\) 后为非零的 \(t^{-1}+R\)，矛盾。所以不有限生成。\(J(R)=(t)\)，因此 \(J(R)M=M\)，剩余模为零；在缺少有限生成时，零剩余模并不能检测原模为零。
\end{solution}

\begin{exercise}\label{b2:ex:14-nilpotent-ideal-generators}
设 \(R\) 为含幺环、\(I\) 为满足 \(I^r=0\) 的双边理想，\(r\ge1\)。对任意幺性左 \(R\) 模 \(M\)，证明 \(M\) 有限生成当且仅当 \(M/IM\) 在 \(R/I\) 上有限生成；在此情形证明
\[
\mu_R(M)=\mu_{R/I}(M/IM).
\]
若只假设 \(I\subseteq J(R)\)，上述有限生成性反向蕴含未必成立。用 \(R=k[t]_{(t)}\)、\(I=(t)\)、\(M=R\oplus k(t)\) 验证这一点。
\end{exercise}
\begin{solution}
有限生成模的商仍有限生成。反过来，提升 \(M/IM\) 的有限生成组为 \(x_1,\ldots,x_n\)，记 \(N=\sum_iRx_i\)，则 \(M=N+IM\)。由\cref{b2:cor:nakayama-generators}的幂零理想版本，不需预设 \(M\) 有限生成便有 \(M=N\)。任意生成组取像给出商模的生成组，任意商模生成组提升又给出原模生成组，两方向比较元素数便得最少数相等；零模也包括在内。

在题给局部环中，\(J(R)=(t)\)，且 \(tk(t)=k(t)\)。所以 \(M/IM\cong R/(t)\cong k\)，只需一个生成元；但 \(M\) 不有限生成，因为它的商模 \(k(t)\) 不有限生成。后者在零点的极点阶无统一上界，而有限个有理函数的 \(R\) 线性组合的极点阶有共同上界。这里 \(I\) 的任何幂均不为零，幂零理想下的反向提升论证因而不能使用。
\end{solution}

\begin{exercise}\label{b2:ex:14-local-generator-computation}
设 \(k\) 为域，令 \(R=k[x,y]/(x,y)^3\)，仍用 \(x,y\) 表示其商类，记 \(\mathfrak m=(x,y)\)。确定 \(J(R)\)，计算模 \(\mathfrak m\)、\(\mathfrak m^2\)、\(\mathfrak m^2\oplus R/(x)\) 的最少生成元个数。说明 \(x+y^2,y+x^2\) 是否生成 \(\mathfrak m\)。
\end{exercise}
\begin{solution}
\(\mathfrak m^3=0\)，且 \(R/\mathfrak m=k\)，故由幂零理想包含于根及商环根公式，\(J(R)=\mathfrak m\)。\(\mathfrak m/\mathfrak m^2\) 有基 \(x,y\)，所以 \(\mu_R(\mathfrak m)=2\)。\(\mathfrak m\mathfrak m^2=0\)，而 \(\mathfrak m^2\) 有基 \(x^2,xy,y^2\)，所以 \(\mu_R(\mathfrak m^2)=3\)。最后一个模的剩余空间由前三个基元素再加 \(R/(x)\) 的常数类组成，维数为四，故最少需要四个生成元。

\(x+y^2,y+x^2\) 在 \(\mathfrak m/\mathfrak m^2\) 中的像仍为 \(x,y\)，是一组基。目标模 \(\mathfrak m\) 有限生成，由\cref{b2:cor:nakayama-generators}，这两个提升生成 \(\mathfrak m\)，并且是不可删减生成组。
\end{solution}

\begin{exercise}\label{b2:ex:14-semisimple-generator-packing}
设 \(k\) 为域，令 \(R=M_2(k)\times M_3(k)\)，以 \(S_1=k^2,S_2=k^3\) 表示相应的简单列模。求 \(M=S_1^5\oplus S_2^7\) 的最少生成元个数，并构造达到这个数的生成组。
\end{exercise}
\begin{solution}
由\cref{b2:prop:semisimple-top-generators}，
\[
\mu_R(M)=\max\{\lceil5/2\rceil,\lceil7/3\rceil\}=3.
\]
构造三个向量 \(g_1,g_2,g_3\)。在五个 \(S_1\) 坐标中，\(g_j\) 的第 \(2j-1,2j\) 个位置分别放标准向量 \(e_1,e_2\)，超过五的位置删去，其余位置为零。在七个 \(S_2\) 坐标中，\(g_j\) 的第 \(3j-2,3j-1,3j\) 个位置分别放 \(e_1,e_2,e_3\)，超过七的位置删去，其余为零。

对每个 \(g_j\)，二阶矩阵的两列可以独立指定其负责的两个 \(S_1\) 坐标，三阶矩阵的三列可以独立指定其负责的三个 \(S_2\) 坐标；直积环中这两个矩阵又可独立选取。各 \(j\) 负责的坐标不交，合起来覆盖全模，所以三者生成 \(M\)。
\end{solution}

\begin{exercise}\label{b2:ex:14-irredundant-not-minimum}
设 \(R=k^4\)，采用坐标乘法，\(g_1,\ldots,g_m\in R\)。用各 \(g_j\) 的非零坐标位置给出它们生成左正则模 \({}_RR\) 的判据，以及这组生成元不可删减的判据。据此构造不可删减的三元生成组，证明任一不可删减有限生成组至多含四个元素，并求 \(\mu_R(R)\)。
\end{exercise}
\begin{solution}
记 \(A_j=\{\,i\in\{1,2,3,4\}\mid(g_j)_i\ne0\,\}\)。每个系数的四个坐标可独立选取，因此这组元素生成 \(R\) 当且仅当 \(\bigcup_jA_j=\{1,2,3,4\}\)：若漏掉某个位置，所有线性组合在那里为零；若每个位置均覆盖，利用坐标幂等元及域中逆元便可得到各个标准坐标向量。

在生成条件成立时，删去 \(g_j\) 后不再生成，当且仅当 \(A_j\) 含有一个不属于任何其他 \(A_\ell\) 的位置。因此不可删减等价于每个生成元都有这样的独占坐标。可取
\[
g_1=(1,1,0,0),\qquad g_2=(0,1,1,0),\qquad g_3=(0,0,0,1).
\]
它们分别有独占坐标一、三、四，故不可删减。给每个不可删减生成元选择一个独占坐标，这些坐标两两不同，故 \(m\le4\)。但单位元 \((1,1,1,1)\) 单独生成 \(R\)，而 \(R\ne0\)，所以 \(\mu_R(R)=1\)。
\end{solution}

\begin{exercise}\label{b2:ex:14-polynomial-semiprimitive}
对任意域 \(k\)，证明 \(J\bigl(k[t]\bigr)=0\)，但 \(k[t]\) 既不本原，也不半单。不使用“极大理想的交显然为零”作为证明。
\end{exercise}
\begin{solution}
对任意非零多项式 \(f\)，取 \(r=t\)，则 \(1-tf\) 具有正次数，不是 \(k[t]\) 的单位。由\cref{b2:thm:jacobson-unit-criterion}，\(f\notin J\bigl(k[t]\bigr)\)，所以根为零。交换本原环必须为域，而 \(t\) 不可逆，所以不本原。降链 \((t)\supsetneq(t^2)\supsetneq\cdots\) 不稳定，故不是 Artin 环；半单环由\cref{b2:thm:artin-wedderburn}具有 Artin 性，所以它不半单。
\end{solution}

\begin{exercise}\label{b2:ex:14-infinite-product}
设 \(k\) 为域，\(R=\prod_{n\ge1}k\)。证明每个非零元素都能被某个简单左模检测到，因而 \(R\) 半本原；再证明它没有忠实简单模，也不是半单环。
\end{exercise}
\begin{solution}
第 \(n\) 个坐标投影使 \(k\) 成为简单 \(R\) 模。如果 \(r\ne0\)，某个坐标 \(r_n\ne0\)，它在该简单模上作用非零。于是全部简单模零化子的交为零，由\cref{b2:prop:radical-simple-annihilators}，\(J(R)=0\)。\(R\) 交换且有非平凡幂等元，不是域，故由\cref{b2:prop:commutative-primitive-field}不本原，没有忠实简单模。令 \(I_n\) 为前 \(n\) 个坐标都为零的理想，则 \(R\supsetneq I_1\supsetneq I_2\supsetneq\cdots\)，故不 Artin，也不半单。
\end{solution}

\begin{exercise}\label{b2:ex:14-evaluation-subdirect}
设 \(k\) 为无限域，考虑取值映射 \(k[t]\rightarrow\prod_{a\in k}k\)、\(f\mapsto\bigl(f(a)\bigr)_a\)。证明它给出一个严格小于全直积的次直积。若 \(k\) 为有限域，这个映射的核与像又是什么？
\end{exercise}
\begin{solution}
无限域上，一个非零多项式只有有限多个根，所以取值映射单射。任一坐标映射满射，因为常数多项式可以取任意值。它不到全直积满射：只在 \(a=0\) 处取一、在其他点取零的坐标族不可能来自多项式，否则这个多项式有无限多个根，须为零，又与在零点的值矛盾。

有限域上，令 \(P(t)=\prod_{a\in k}(t-a)\)。一个多项式在每点取零，当且仅当逐个一次因子都整除它；这些因子两两互素，故等价于 \(P\mid f\)，核为 \((P)\)。对每个 \(a\)，Lagrange 多项式
\[
\ell_a(t)=\prod_{\substack{b\in k\\b\ne a}}\frac{t-b}{a-b}
\]
在 \(a\) 处取一，其余点取零，所以任意坐标族由有限和 \(\sum_a c_a\ell_a\) 给出，映射满射。于是 \(k[t]/(P)\cong\prod_{a\in k}k\)，有限与无限乘积的结论不同。
\end{solution}

\begin{exercise}\label{b2:ex:14-right-division-scalars}
设 \(D_0\) 为任意除环，\(R=D_0\)、\(S={}_{D_0}D_0\)。计算 \(\Delta=\operatorname{End}_R(S)\) 及本章采用的右标量除环 \(D=\Delta^{\mathrm{op}}\)。对 \(0\ne v\in S,w\in S\)，写出实现 \(rv=w\) 的 \(r\)，并说明为何 \(S_D\) 的无关组最多只有一个向量。
\end{exercise}
\begin{solution}
每个左线性自同态由 \(f(1)=a\) 决定，并为右乘 \(f(x)=xa\)。两个右乘的复合反转参数乘法，所以 \(\Delta\cong D_0^{\mathrm{op}}\)，再取反环得到 \(D\cong D_0\)。按此识别，右标量作用就是通常的右乘，\(S_D\) 以一为基，右维数为一。

取 \(r=wv^{-1}\)，有 \(rv=w\)。任意两个非零向量满足 \(v_2=v_1(v_1^{-1}v_2)\)，所以右线性相关。即使 \(D_0\) 在某个中心子域上有更高维数，也不能对这种右相关组任意指定像；正确的插值条件使用 \(D\)。
\end{solution}

\begin{exercise}\label{b2:ex:14-density-relations}
设 \(D_0\) 为除环、\(R=M_2(D_0)\)、\(S=D_0^2\) 为简单列模，右标量作用采用通常右乘。固定 \(a,b\in D_0\)，令
\[
v_1=e_1,\qquad v_2=e_2,\qquad v_3=e_1a+e_2b.
\]
给定 \(w_1,w_2,w_3\in S\)，求同时满足 \(rv_i=w_i\) 的全部 \(r\in R\)。若只给定 \(rv_1=w_1\) 与 \(rv_3=w_3\)，分别讨论 \(b=0\) 和 \(b\ne0\)，并注意非交换系数的位置。
\end{exercise}
\begin{solution}
与全部矩阵单位及 \(D_0\) 条目作用交换的自同态，恰好在两个坐标上分别右乘同一个 \(D_0\) 元素，所以 \(\operatorname{End}_R(S)\cong D_0^{\mathrm{op}}\)，相应右标量除环确为 \(D_0\)。同时指定前两列已唯一确定矩阵 \(r=(w_1\ w_2)\)。右线性要求
\[
w_3=w_1a+w_2b,
\]
这也充分：满足时上述矩阵实现全部指定像。原向量的任意关系均为 \((d_1,d_2,d_3)=(-ad,-bd,d)\)，\(d\in D_0\)，所以这一条件也正是\cref{b2:cor:density-compatible-relations}的全部关系保持条件。

若只指定第一、第三个像，当 \(b=0\) 时必须且只需 \(w_3=w_1a\)，第二列可任取。当 \(b\ne0\) 时，第一、第三个向量右线性无关；第二列必须为
\[
(w_3-w_1a)b^{-1},
\]
故任意 \(w_1,w_3\) 都对应唯一矩阵。因为这里是右标量关系，\(b^{-1}\) 乘在向量的右侧；把它移到左侧在一般除环上既没有相同含义，也没有相同结果。
\end{solution}

\begin{exercise}\label{b2:ex:14-finite-rank-commutant}
设 \(V=\bigoplus_{i\ge1}ke_i\)，\(F\) 为全部有限秩算子，\(F_0\) 为相对于此基只有有限多个非零矩阵条目的算子。证明 \(F_0\) 是不含恒等算子的 \(k\) 子代数，且 \(F_0\subsetneq F\)。证明 \(F_0\) 在 \(\operatorname{End}_k(V)\) 的有限拓扑中已经稠密，并比较 \(k\operatorname{id}+F_0\) 与 \(k\operatorname{id}+F\)。
\end{exercise}
\begin{solution}
有限多个非零条目的矩阵相加、数乘和相乘仍只有有限多个非零条目，所以 \(F_0\) 为子代数；每个这样的算子像均有限维，故属于 \(F\)。恒等矩阵有无限多个非零对角项，不在 \(F_0\)。算子 \(U(e_i)=e_1\) 的像为 \(ke_1\)，但有无限多个非零条目，说明 \(F_0\subsetneq F\)。

给定目标算子 \(T\) 及有限个测试向量 \(v_1,\ldots,v_n\)，它们总共只使用有限个基位置，记其集合为 \(J\)。令 \(B(e_j)=T(e_j)\) 对 \(j\in J\) 成立，在其余基向量上取零。有限多个 \(T(e_j)\) 又各有有限支集，所以 \(B\) 的非零矩阵条目只有有限多个，\(B\in F_0\)；并且 \(B(v_i)=T(v_i)\)。于是每个基本邻域都与 \(F_0\) 相交，其闭包为全部自同态环。

有 \(k\operatorname{id}+F_0\subsetneq k\operatorname{id}+F\)：上面的 \(U\) 不能写成 \(c\operatorname{id}+B\)、\(B\in F_0\)。若 \(c=0\)，\(U\) 有无限多个非零条目；若 \(c\ne0\)，\(U-c\operatorname{id}\) 在所有 \(i\ge2\) 的对角位置都非零，同样不在 \(F_0\)。由\cref{b2:prop:finite-rank-density-counterexample}，后者还严格小于 \(\operatorname{End}_k(V)\)。这两个含幺真子代数却都有同样的有限拓扑闭包，说明稠密性不能区别它们在无限多基位置上的限制。
\end{solution}

\begin{exercise}\label{b2:ex:14-burnside-rational-counterexample}
令 \(J=\begin{pmatrix}0&2\\1&0\end{pmatrix}\)，\(A=\mathbb Q[J]\subseteq M_2(\mathbb Q)\)。证明 \(A\) 在 \(\mathbb Q^2\) 上作用不可约，但不等于全部矩阵环；求其中心化子 \(\operatorname{End}_A(\mathbb Q^2)\)，并说明稠密性给出的正确目标。
\end{exercise}
\begin{solution}
\(J^2=2I\)，而 \(t^2-2\) 在 \(\mathbb Q\) 上不可约，所以 \(A\cong\mathbb Q(\sqrt2)\)。把向量 \((a,b)\) 识别为 \(a+b\sqrt2\)，\(J\) 对应乘 \(\sqrt2\)。任意非零 \(A\) 稳定子空间都是这个域的非零左理想，故为全体，作用不可约。\(A\) 的有理维数为二，严格小于矩阵环的四。

令 \(C=\begin{pmatrix}u&v\\w&z\end{pmatrix}\)，等式 \(CJ=JC\) 给出 \(u=z,v=2w\)，所以 \(C=uI+wJ\in A\)。因此中心化子恰为 \(A\)，它是交换域；相应右标量除环也是 \(A\)，\(\mathbb Q^2\) 在它上面一维，稠密性给出的全部右线性自同态正是 \(A\)。缺少代数闭性时，不能把这个目标改成全部有理线性自同态。
\end{solution}

\begin{exercise}\label{b2:ex:14-algebra-generation}
设 \(k\) 为任意域，令 \(U=I+E_{12}\)、\(V=I+E_{21}\)。证明 \(U,V\) 生成的含幺 \(k\) 子代数是 \(M_2(k)\)，却不能仅用 \(I,U,V\) 的线性组合得到所有矩阵。说明这里的结论为什么不需要 \(k\) 代数闭。
\end{exercise}
\begin{solution}
由 \(U-I=E_{12}\)、\(V-I=E_{21}\)，生成代数包含这两个矩阵单位，继而包含乘积 \(E_{12}E_{21}=E_{11}\)、\(E_{21}E_{12}=E_{22}\)。四个矩阵单位是 \(M_2(k)\) 的基，所以生成代数为全体。\(I,U,V\) 的线性生成空间只有三维：它们与 \(I,E_{12},E_{21}\) 具有相同线性生成空间，后者显然无关。乘积提供了原线性组合中缺少的方向。这里是显式矩阵运算，对任意域成立；Burnside 定理的代数闭条件是其一般不可约性判据的条件，并非所有具体生成等式成立的必要条件。
\end{solution}

\begin{exercise}\label{b2:ex:14-group-degree-bound}
设有限群 \(G\) 在代数闭域 \(k\) 上有非零有限维不可约表示 \(\rho:G\rightarrow\operatorname{GL}(V)\)，记 \(d=\dim_kV\)。证明 \(\rho(G)\) 的线性生成空间为 \(\operatorname{End}_k(V)\)，并推出 \(d^2\le |G|\)。不要加入关于特征不整除群阶的假设。
\end{exercise}
\begin{solution}
记 \(A=\operatorname{span}_k\rho(G)\)。群乘法使 \(\rho(g)\rho(h)=\rho(gh)\)，所以 \(A\) 对乘法封闭，且包含恒等算子。子空间对全部群元素稳定，当且仅当对其线性生成空间稳定，故 \(A\) 的作用不可约。\cref{b2:thm:burnside-matrix-algebra}于是给出 \(A=\operatorname{End}_k(V)\)。左侧由至多 \(|G|\) 个矩阵生成，右侧维数为 \(d^2\)，所以 \(d^2\le|G|\)。这个不等式对任意特征成立。
\end{solution}

% !TeX root = ../../Book2.tex
% 纲要标题：## 第15章　有限维代数、根滤过与投射覆盖
% 纲要位置：计划纲要.md，第 1525 行。

\chapter{有限维代数、根滤过与投射覆盖}
\label{b2:ch:15}
\BookChapterNumber{b2:ch:15}

\begin{chapterabstract}
有限维代数的 Jacobson 根是幂零理想，而模掉根的商是半单代数。两项结论使一个模可以按根的作用分层，但这些半单层之间仍可能有不分裂的扩张。本章用根与基座刻画 Loewy 列，计算三角矩阵代数的具体层次，再通过幂等元的相容提升构造投射覆盖。覆盖的存在性、唯一性和不可分解投射模的分类均在此证明。最后选取基本角代数，以半单商给出箭图的顶点，以 \(J/J^2\) 给出箭头。
\end{chapterabstract}

\begin{learninggoals}
\begin{itemize}
\item 区分降链稳定与根幂零，准确说明 Nakayama 引理在证明中的作用。
\item 计算模的根、基座和 Loewy 层，并辨认这些层之间是否分裂。
\item 完成幂零理想模的单个及正交幂等元组提升，利用多余核构造投射覆盖。
\item 说明不可分解投射模与简单模的对应，保留一般基域上自同态除环的作用。
\item 构造基本满角代数，按左模约定理解根第一层的箭头方向。
\end{itemize}
\end{learninggoals}

令 \(A\) 为域 \(k\) 上的二阶上三角矩阵代数。严格上三角矩阵组成理想 \(J\)，且 \(J^2=0\)；商 \(A/J\) 只看见两个对角位置，因而同构于 \(k\times k\)。商映射保留两个对角位置，却把非零的矩阵单位 \(E_{12}\) 送到零，也丢去了对角元与它相乘的作用。若研究一个 \(A\) 模，先模掉 \(J\) 能读出顶层，却无法独自恢复各层如何接合。

\ProseParagraphBreak

沿 \(J\) 的幂逐层观察，最先消失的乘法信息就能成为层之间关系的线索。若要从半单商中的简单对象回到原代数，还需找到足够小、又能映满目标的投射模。幂等元的提升和投射覆盖正是为这种回返服务；在三角矩阵例子中，两个对角幂等元与非对角位置已经把顶点和连接它们的箭头显露出来。

\ProseParagraphBreak

本章需要\cref{b2:thm:projective-characterizations}中的投射模刻画、\cref{b2:thm:artin-wedderburn}中的半单代数结构，以及\cref{b2:thm:noncommutative-nakayama}。角环及其与自同态环的关系使用\cref{b2:prop:corner-endomorphism}。除幂等元提升处明确讨论一般环外，\(A\) 始终是域 \(k\) 上的有限维结合含幺代数，模默认是幺性左模；对模的有限性要求在陈述中写明。

\ProseParagraphBreak
有限维代数及投射覆盖可参阅 Etingof 等人的《Introduction to Representation Theory》\cite[\S 3.5; \S\S 9.1--9.2]{EtingofEtAl2011RepresentationTheory}。该书第9章明确假设基域代数闭。本章的主要构造保留一般基域，简单模的自同态可能是大于 \(k\) 的除环；在计数 Hom 空间维数时尤其要区别这两种假设。根幂零性的主要证明采用第十四章的 Nakayama 引理，与该书 Proposition 3.5.3 的合成列证明可以相互参照。

% !TeX root = ../../Book2.tex
% 纲要标题：### 15.1 根的幂零性
% 纲要位置：计划纲要.md，第 1527 行。

\section{根的幂零性}
\label{b2:sec:1501}

有限维性使左理想的升降链都停止，但“链停止”还不等于“某个幂为零”。根滤过的各层可以通过半单商来理解，而要让滤过最终到达零，还需将有限维性与非交换 Nakayama 引理合起来。始终令 \(A\) 为域 \(k\) 上的有限维结合含幺代数，记 \(J=J(A)\)。

% 纲要要点（供目录核对）：
%
% * 有限维代数的 Jacobson 根
% * 半单商、根滤过与伴随分次代数
% * 幂零根和一般根的区别；区分逐元素幂零与统一混合乘积为零，说明伴随分次乘法保留的层次
% * 对有限维代数 \(A\) 的降链 \(J(A)^n\) 先取稳定项，再用第14.2节 Nakayama 引理推出其为零；不把 Artin 性本身当作根幂零性的全部证明
%

\subsection{有限维代数的 Jacobson 根}
\label{b2:subsec:1501:01}

回顾\cref{b2:def:jacobson-radical,b2:prop:radical-simple-annihilators}，\(J(A)\) 既是全部极大左理想的交，也是全部简单幺左模的共同零化子，因而是双边理想。特别地，对每个简单左模 \(S\)，都有 \(JS=0\)。后面各层被 \(J\) 零化的计算，使用的正是这个性质。

\ProseParagraphBreak
在 \(A\) 有限维的当前条件下，每个简单幺左 \(A\) 模都是 \(A\) 的商，因而有限维。每个有限生成幺左 \(A\) 模也有限维：有限个生成元给出满射 \(A^n\to M\)。反过来，若 \(m_1,\ldots,m_d\) 是 \(M\) 的 \(k\) 基，则 \(\sum_i c_im_i=\sum_i(c_i1_A)m_i\)，其中 \(c_i1_A\in A\)，所以它们也是 \(A\) 模生成元。这说明此时有限生成与有限 \(k\) 维等价，并不说明 \(A\) 模的最少生成元数等于 \(k\) 维数。

\begin{definition}\label{b2:def:nilpotent-ideal}
设 \(R\) 为含幺环、\(I\) 为其双边理想。对正整数 \(N\)，\(I^N\) 是全部 \(N\) 个 \(I\) 中元素的乘积所生成的加法子群。若某个正整数 \(N\) 使 \(I^N=0\)，则称 \(I\) 为\term[mi-ling-lixiang]{幂零理想}[nilpotent ideal]，满足该式的最小正整数称为其\term[miling-zhishu]{幂零指数}[nilpotency index]。
\end{definition}

第十四章出现的 nil ideal 只要求 \(\forall x\in I\,\exists n_x\ge1\)，使 \(x^{n_x}=0\)。幂零理想则要求 \(\exists N\ge1\,\forall x_1,\ldots,x_N\in I\)，都有 \(x_1\cdots x_N=0\)。后者除了使用一个共同指数，还要控制不同元素的混合乘积；仅考察同一个元素反复相乘，不能直接得到 \(I^N=0\)。

\begin{proposition}\label{b2:prop:nilpotent-ideal-in-radical}
任意含幺环的幂零双边理想都包含在其 Jacobson 根中。
\end{proposition}
\begin{proof}
若 \(I^N=0\)，则每个 \(x\in I\) 都满足 \(x^N=0\)，故 \(I\) 是 nil 双边理想。\cref{b2:cor:nil-ideal-radical}已证明这种理想包含于 Jacobson 根；它允许各元素使用不同的幂次，因而也适用于这里的共同幂次 \(N\)。
\end{proof}

例如 \(A=k[t]/(t^m)\)、\(m\ge1\)，其理想 \((\bar t)\) 幂零，且商为域 \(k\)。它因此恰是 \(J(A)\)：包含关系由上命题得到，而商的根为零，再用\cref{b2:prop:radical-quotient}即得反向包含。这里的根不是全部零因子的任意集合，而是由理想及简单模刻画的特定对象。

\subsection{半单商与根滤过}
\label{b2:subsec:1501:02}

\begin{proposition}\label{b2:prop:finite-algebra-semisimple-quotient}
设 \(k\) 为域、\(A\) 为有限维结合含幺 \(k\) 代数，\(J=J(A)\) 为其 Jacobson 根。商代数 \(A/J\) 是半单代数。每个被 \(J\) 零化的有限维幺左 \(A\) 模都是半单模。
\end{proposition}
\begin{proof}
由\cref{b2:prop:radical-quotient}，\(J(A/J)=0\)。商代数的左正则模有限维，故具有有限长度；其极大子模正是极大左理想，交为零。应用\cref{b2:prop:finite-length-radical-zero}，可知这个正则模半单，所以 \(A/J\) 是半单代数。

若 \(JM=0\)，定义 \((a+J)m=am\)。由于 \(J\) 在 \(M\) 上为零，此定义与代表元无关，使 \(M\) 成为幺左 \(A/J\) 模。它是某个有限自由 \((A/J)^n\) 的商。有限直和与商保持半单性，见\cref{b2:prop:semisimple-subquotients}，故 \(M\) 半单；其 \(A/J\) 子模与 \(A\) 子模是同一批子空间。
\end{proof}

由\cref{b2:thm:artin-wedderburn}，可以写
\begin{equation}\label{b2:eq:finite-algebra-semisimple-quotient}
A/J\cong\prod_{i=1}^s M_{n_i}(D_i),
\end{equation}
其中各 \(D_i\) 是有限维 \(k\) 除代数。暂不假设 \(k\) 代数闭，所以不能把 \(D_i\) 一律写成 \(k\)。

\ProseParagraphBreak
只要 \(J\) 在一个幺左模上作用为零，\(A\) 的作用就能且只能通过 \(A/J\) 给出；反过来，每个幺左 \(A/J\) 模沿商映射限制标量，都被 \(J\) 零化。两种作用给出的子模和模同态完全相同，所以简单性及半单性也相同。这里 \(A/J\) 半单，是根为零与有限长度共同带来的结论；对任意环，仅有 Jacobson 根为零不足以推出半单性。

\ProseParagraphBreak
\(J\) 的各个幂给出\term[gen-lvguo]{根滤过}[radical filtration]
\[
A\supseteq J\supseteq J^2\supseteq\cdots.
\]
每个 \(J^r/J^{r+1}\) 都被 \(J\) 零化，因为 \(J J^r=J^{r+1}\)，所以由\cref{b2:prop:finite-algebra-semisimple-quotient}，它作为左模半单。类似地，对任何有限维左模 \(M\)，商 \(J^rM/J^{r+1}M\) 也是半单模。各层半单并不说明原模是这些层的直和；层与层之间的扩张可以不分裂。

\ProseParagraphBreak
若只把这些商空间逐个列出，\(A\) 的乘法就不在数据中了。但 \(J^rJ^s\subseteq J^{r+s}\)，所以第 \(r\) 层和第 \(s\) 层的代表元相乘，仍能在第 \(r+s\) 层读出一个商类。伴随分次代数把这种乘法也记录下来。

\begin{definition}\label{b2:def:associated-radical-algebra}
设 \(k\) 为域、\(A\) 为有限维结合含幺 \(k\) 代数，\(J=J(A)\)，并约定 \(J^0=A\)。在分次向量空间
\[
\operatorname{gr}_J(A)=\bigoplus_{r\ge0}J^r/J^{r+1}
\]
上，对非负整数 \(r,s\) 及 \(a\in J^r\)、\(b\in J^s\) 定义
\[
(a+J^{r+1})(b+J^{s+1})=ab+J^{r+s+1}.
\]
所得分次 \(k\) 代数称为\term[gen-lvguo-bansui-fenci-daishu]{根滤过的伴随分次代数}[associated graded algebra of the radical filtration]。
\end{definition}
\symidx{\(\operatorname{gr}_J(A)\)：根滤过的伴随分次代数}

乘积不依赖代表元：若分别把 \(a,b\) 改动一个 \(J^{r+1},J^{s+1}\) 中的元素，新增各项都在 \(J^{r+s+1}\) 中。结合律由 \(A\) 的乘法继承，单位为 \(1+J\)，而次数按和相加。这个乘法只记录 \(ab\) 模 \(J^{r+s+1}\) 的值；若 \(ab\) 实际落在更深的 \(J^{r+s+1}\) 中，即使 \(ab\ne0\)，它在第 \(r+s\) 次的乘积也为零。因此较深层的修正项会被丢掉，不能仅凭伴随分次代数恢复 \(A\) 的全部乘法。

\ProseParagraphBreak
在 \(A=k[t]/(t^m)\) 中，\(J^r\) 的基为 \(\bar t^r,\ldots,\bar t^{m-1}\)，\(0\le r<m\)。这个例子恰有与滤过相容的分次结构：令 \(u\) 的次数为一，映射 \(k[u]/(u^m)\to\operatorname{gr}_J(A)\) 将 \(u\) 送到 \(\bar t+J^2\)，将 \(u^r\) 送到第 \(r\) 层的 \(\bar t^r+J^{r+1}\)。它保持乘法，并把各次的一维基对应起来，因而是分次 \(k\) 代数同构。这里能恢复截断多项式代数，依赖于这一具体乘法，不能作为一般情形的结论。每个非零层又是一维简单模 \(k\)，但当 \(m>1\) 时正则模不是这些层的直和：\(\bar t\) 在任何这样的简单模直和上均作用为零，而在正则模上并非如此。

\subsection{幂零根与一般根}
\label{b2:subsec:1501:03}

对有限维代数，本节将证明 \(J\) 是幂零理想。这个结论比“\(J\) 中每个元素各自幂零”更强，因为同一个整数要使任意若干根元素的乘积都为零。它也不意味着所有幂零元素都在根中。

\ProseParagraphBreak
例如 \(M_2(k)\) 是半单代数，故根为零，但矩阵单位 \(E_{12}\ne0\) 满足 \(E_{12}^2=0\)。幂零的单个元素生成的双边理想未必幂零；这里 \(E_{21}E_{12}=E_{22}\)，所以它生成的双边理想已经含有非零幂等元。由于 \(E_{22}^n=E_{22}\ne0\) 对每个正整数 \(n\) 都成立，这个双边理想的各个幂都非零。\cref{b2:prop:nilpotent-ideal-in-radical}的假设是整个双边理想幂零。

\ProseParagraphBreak
若去掉有限维性，Jacobson 根可以完全不幂零。例如形式幂级数环 \(k[\![t]\!]\) 中，常数项非零的级数可以逐次解出逆级数的系数，因而可逆；常数项为零的级数不可能可逆。故非单位元恰为理想 \((t)\)，它是唯一极大理想，等于 Jacobson 根。但 \(t^n\notin(t^{n+1})\)，所以 \((t)^n\supsetneq(t)^{n+1}\) 对所有 \(n\ge0\) 都成立，降链根本没有稳定项。Nakayama 引理仍适用；只是没有 \(J^{n+1}=J^n\) 这个等式，就不能把它用于稳定项来推出零。

\subsection{根幂零性的 Nakayama 证明}
\label{b2:subsec:1501:04}

\begin{theorem}\label{b2:thm:finite-algebra-radical-nilpotent}
设 \(k\) 为域、\(A\) 为有限维结合含幺 \(k\) 代数，\(J=J(A)\) 为其 Jacobson 根，则 \(J\) 幂零。若 \(A\ne0\)，可以取正整数 \(N\le\dim_k A\) 使 \(J^N=0\)。因此 \(J\) 是 \(A\) 中最大的幂零双边理想。
\end{theorem}
\begin{proof}
降链 \(A\supseteq J\supseteq J^2\supseteq\cdots\) 是有限维向量空间的降链，故存在 \(n\ge0\) 使 \(J^n=J^{n+1}\)。令 \(M=J^n\)，将它看成左 \(A\) 模。它有限维，因而有限生成；等式正是
\[
JM=J J^n=J^{n+1}=J^n=M.
\]
对理想 \(J\subseteq J(A)\) 和有限生成左模 \(M\) 使用\cref{b2:thm:noncommutative-nakayama}，得到 \(M=0\)。所以某个 \(J^n\) 为零，证明幂零性。

当 \(A\ne0\) 时，\(J\ne A\)，因为单位不属于任何真极大左理想。若在到达零以前某两个相邻幂相等，上面的同一论证会迫使它们为零，矛盾。因此从 \(A\) 到第一个零幂，每一步维数都严格下降，步数不超过 \(\dim_k A\)，得到所述界。最后，任意幂零双边理想由\cref{b2:prop:nilpotent-ideal-in-radical}包含于 \(J\)，而 \(J\) 本身已证幂零。
\end{proof}

证明用到了两个不同的事实：有限维性使某一项稳定，Nakayama 引理使这个稳定项为零。仅凭降链停止，只能得到 \(J^n=J^{n+1}\)，不能省去后一步。界 \(N\le\dim_k A\) 只把每个非零层的维数估为至少一，是一个统一的粗界，并不说最小幂零指数等于代数的维数。这个幂零性也将保证稍后提升幂等元的过程在有限步后结束。

% !TeX root = ../../Book2.tex
% 纲要标题：### 15.2 模的根、基座与 Loewy 列
% 纲要位置：计划纲要.md，第 1252 行。

\section{模的根、基座与 Loewy 列}
\label{b2:sec:1502}

现在把代数的根用于单个模。一个方向是不断取最大半单商，另一个方向是不断取出全部简单子模；它们给出从上到下和从下到上的两列。本节的 \(M\) 均为有限维左 \(A\) 模，\(J=J(A)\)。

% 纲要要点（供目录核对）：
%
% * 模的根与基座（socle；所有简单子模之和）
% * Loewy 列及半单层
% * 具体三角矩阵代数的计算
%

\subsection{模的根与基座}
\label{b2:subsec:1502:01}

\begin{definition}\label{b2:def:module-radical-socle}
设 \(M\) 为左 \(R\) 模。将 \(M\) 的全部极大真子模取交，所得子模称为 \(M\) 的\term[mo-de-gen]{根}[radical of a module]，记为 \(\operatorname{rad}M\)；若不存在极大真子模，该交按约定取 \(M\)。将全部简单子模求和，所得子模称为 \(M\) 的\term[jizuo]{基座}[socle]，记为 \(\operatorname{Soc}M\)；若没有简单子模，该和为零。对零模约定 \(\operatorname{rad}0=\operatorname{Soc}0=0\)。
\end{definition}
\symidx{\(\operatorname{rad}M\)、\(\operatorname{Soc}M\)：模的根与基座}

非零有限维模总有极大真子模及简单子模：分别取维数最大的真子模和维数最小的非零子模即可。基座是半单模，因为它是简单子模之和，适用\cref{b2:thm:semisimple-module-characterizations}。

\begin{theorem}\label{b2:thm:module-radical-jm}
设 \(k\) 为域、\(A\) 为有限维结合含幺 \(k\) 代数，\(J=J(A)\)，\(M\) 为有限维左 \(A\) 模，则
\[
\operatorname{rad}M=JM.
\]
商 \(M/JM\) 半单，而且对每个半单左 \(A\) 模 \(T\)，每个 \(A\) 模同态 \(M\to T\) 都唯一通过这个商分解。
\end{theorem}
\begin{proof}
对极大子模 \(L\)，商 \(M/L\) 简单，故被 \(J\) 零化。这说明 \(JM\subseteq L\)，对全部 \(L\) 取交得 \(JM\subseteq\operatorname{rad}M\)。

另一方面，\(M/JM\) 是有限维 \(A/J\) 模，由\cref{b2:prop:finite-algebra-semisimple-quotient}半单。若 \(m\notin JM\)，把它的非零商类投影到某个非零的简单直和分量，得到一个不将 \(m\) 送到零的满同态 \(M\to S\)。其核为不含 \(m\) 的极大子模，所以 \(m\notin\operatorname{rad}M\)。这给出反向包含。

最后，半单模是简单模的直和，所以 \(J\) 在其上作用为零。任意同态 \(f:M\to T\) 因而满足 \(f(JM)=J\cdot f(M)=0\)。由商模泛性质，它唯一通过 \(M/JM\) 分解。
\end{proof}

\begin{proposition}\label{b2:prop:module-socle-annihilator}
设 \(k\) 为域、\(A\) 为有限维结合含幺 \(k\) 代数，\(J=J(A)\)，\(M\) 为有限维左 \(A\) 模，则其基座为
\[
\operatorname{Soc}M=\{\,m\in M\mid Jm=0\,\}.
\]
\end{proposition}
\begin{proof}
每个简单子模均被 \(J\) 零化，所以基座包含于右侧。右侧是子模：若 \(Jm=0\)，则对 \(a\in A,j\in J\)，有 \(j(am)=(ja)m=0\)，因为 \(ja\in J\)。这个子模被 \(J\) 零化，由\cref{b2:prop:finite-algebra-semisimple-quotient}半单，故是若干简单子模之和，包含于基座。
\end{proof}

在 \(M=A=k[t]/(t^m)\)、\(m>1\) 中，\(\operatorname{rad}M=(\bar t)\)，而 \(\operatorname{Soc}M=k\bar t^{m-1}\)。前者收集不能通过简单商检测到的元素，后者收集已经构成简单子模的部分；它们通常不是同一个子空间。

\subsection{Loewy 列与半单层}
\label{b2:subsec:1502:02}

\begin{definition}\label{b2:def:loewy-filtrations}
设 \(M\) 为左 \(R\) 模。令 \(\operatorname{rad}^0M=M\)、\(\operatorname{rad}^{r+1}M=\operatorname{rad}(\operatorname{rad}^rM)\)（\(r\ge0\)），所得降列称为\term[gen-Loewy-lie]{根 Loewy 列}[radical Loewy series]。另令 \(\operatorname{Soc}^0M=0\)，并令 \(\operatorname{Soc}^{r+1}M\) 是 \(M/\operatorname{Soc}^rM\) 的基座在 \(M\) 中的原像，得到\term[jizuo-Loewy-lie]{基座 Loewy 列}[socle Loewy series]。
\end{definition}

以上两列对任意环上的模都有定义；下面用根的幂及其零化子写出的公式，适用于有限维代数上的有限维模。

\begin{proposition}\label{b2:prop:loewy-series}
设 \(k\) 为域、\(A\) 为有限维结合含幺 \(k\) 代数，\(J=J(A)\)，\(M\) 为有限维左 \(A\) 模。对所有 \(r\ge0\)，有
\[
\operatorname{rad}^rM=J^rM,\qquad
\operatorname{Soc}^rM=\{\,m\in M\mid J^rm=0\,\}.
\]
两列的相邻层均半单。若 \(J^N=0\)，根列至多在第 \(N\) 步到达零，基座列至多在第 \(N\) 步到达 \(M\)。
\end{proposition}
\begin{proof}
第一式由\cref{b2:thm:module-radical-jm}逐次用于有限维子模得到。第二式对 \(r\) 归纳；\(r=0\) 时 \(J^0=A\)，被 \(A\) 零化的元素只有零，符合定义。假设第 \(r\) 步成立，由\cref{b2:prop:module-socle-annihilator}，\(m\) 属于下一步当且仅当 \(Jm\subseteq\operatorname{Soc}^rM\)，即 \(J^rJm=0\)，也就是 \(J^{r+1}m=0\)。

根列的层 \(J^rM/J^{r+1}M\) 被 \(J\) 零化，故半单；基座列的层按定义是一个商模的基座，也半单。最后由\cref{b2:thm:finite-algebra-radical-nilpotent}取 \(J^N=0\)，代入两个公式便得到终止性。
\end{proof}

对有限维 \(k\) 代数 \(A\) 的有限维左模 \(M\)，记 \(J=J(A)\)。使 \(J^\ell M=0\) 的最小非负整数 \(\ell\) 称为 \(M\) 的\term[Loewy-changdu]{Loewy 长度}[Loewy length]。由命题，它也等于使 \(\operatorname{Soc}^\ell M=M\) 的最小整数。零模长度为零；非零半单模长度为一。

\ProseParagraphBreak
Loewy 层记录半单的商，不保证能在原模中为每一层同时选出子模代表。对双数代数 \(k[\varepsilon]/(\varepsilon^2)\) 的正则模，两层均为简单模 \(k\)，但层间扩张不分裂，因为若正则模为 \(k\oplus k\)，则 \(\varepsilon\) 应作用为零。这类层间关系正是稍后投射覆盖与箭图表示需要保存的信息。

\subsection{三角矩阵代数的计算}
\label{b2:subsec:1502:03}

令 \(A=T_3(k)\) 为上三角矩阵代数，矩阵单位记为 \(E_{ij}\)、\(1\le i\le j\le3\)。严格上三角矩阵组成理想 \(N\)，有 \(N^3=0\)，而 \(A/N\cong k^3\) 半单。由\cref{b2:prop:nilpotent-ideal-in-radical}与\cref{b2:prop:radical-quotient}，得到
\[
J=N=kE_{12}+kE_{13}+kE_{23},\qquad J^2=kE_{13},\qquad J^3=0.
\]
简单模 \(S_i\) 是一维空间，矩阵 \(a\) 通过对角元 \(a_{ii}\) 作用。它们给出全部简单模，因为简单模被 \(J\) 零化，因而来自 \(k^3\) 的一个坐标因子。

\ProseParagraphBreak
取 \(e_j=E_{jj}\)、\(P_j=Ae_j\)。它是第 \(j\) 列可能非零的矩阵组成的左理想，基为 \(E_{1j},\ldots,E_{jj}\)。矩阵单位乘法
\[
E_{ab}E_{ij}=\delta_{bi}E_{aj}
\]
给出
\begin{equation}\label{b2:eq:triangular-projective-layers}
J^rP_j=\operatorname{span}_k\{\,E_{ij}\mid1\le i\le j-r\,\},
\qquad 0\le r\le j,
\end{equation}
空指标集表示零。每个非零层由 \(E_{j-r,j}\) 的陪集生成；左乘矩阵 \(a\) 后，较低行号的项属于下一层，只剩系数 \(a_{j-r,j-r}\)，所以这个层同构于 \(S_{j-r}\)。

\ProseParagraphBreak
因此，从顶层向下读，\(P_1,P_2,P_3\) 的简单层依次为
\[
P_1:\ S_1;\qquad P_2:\ S_2,\ S_1;\qquad
P_3:\ S_3,\ S_2,\ S_1.
\]
它们的 Loewy 长度分别为一、二、三。另一方面，\(J\) 零化的向量恰为第一行号的部分，所以
\[
\operatorname{Soc}P_j=kE_{1j}\cong S_1,
\qquad
\operatorname{Soc}^rP_j=\operatorname{span}_k\{\,E_{ij}\mid i\le\min(r,j)\,\}.
\]
这些公式可直接从\cref{b2:prop:loewy-series}及矩阵乘法核验。

\ProseParagraphBreak
正则模分解为 \(A=P_1\oplus P_2\oplus P_3\)。于是
\[
A/JA\cong S_1\oplus S_2\oplus S_3,
\qquad \operatorname{Soc}A\cong S_1^{\oplus3}.
\]
最大半单商与最大半单子模甚至可以含有完全不同的重数。自然列向量模 \(k^3\) 则与 \(P_3\) 同构，映射为 \(E_{i3}\mapsto e_i\)；其根列是熟悉的坐标旗标，只是排列方向与从零开始的升旗标相反。

% !TeX root = ../../Book2.tex
% 纲要标题：### 15.3 投射覆盖
% 纲要位置：计划纲要.md，第 1541 行。

\section{投射覆盖}
\label{b2:sec:1503}

自由满射总存在，但其来源可能带有可以删去的投射分量。投射覆盖要求满射在模论意义下已经不能缩小来源。为了构造它，先要把半单商中的列模提升回来；这一步依赖幂等元提升，不能用商中的分解代替原代数中的分裂。本节的 \(A\) 仍有限维，\(J=J(A)\)；幂等元提升的两个定理则只需所写的环与幂零理想假设。

% 纲要要点（供目录核对）：
%
% * 极小生成与本质满射
% * 在构造投射覆盖前，就地证明幂零理想模的幂等元提升，并处理有限组正交幂等元的相容提升
% * 明确本质满射的模论含义为核是多余子模；从提升后的本原幂等元构造不可分解投射模及其到简单模的覆盖
% * 有限维代数中的投射覆盖；目标有限生成时来源的有限性，投射覆盖与最少自由生成的区别
% * 不可分解投射模与简单模
% * 由有限维性、半单商及根的幂零性逐项验证覆盖的存在和唯一性，避免把半单商中的分解直接当作原模中的分裂
% * 正文给出分别提升不保持正交的反例；唯一性包含目标同构时每个提升可逆，极小性说明投射目标的覆盖为同构
%

\subsection{极小生成与本质满射}
\label{b2:subsec:1503:01}

\begin{definition}\label{b2:def:superfluous-submodule}
设 \(R\) 为含幺环、\(P\) 为幺左 \(R\) 模、\(N\le P\) 为子模。若对每个子模 \(L\le P\)，等式 \(L+N=P\) 都蕴含 \(L=P\)，则称 \(N\) 为 \(P\) 的\term[duoyu-zimo]{多余子模}[superfluous submodule]，记为 \(N\ll P\)。若幺左 \(R\) 模的满同态 \(p:P\to M\) 满足 \(\ker p\ll P\)，则称 \(p\) 为\term[benzhi-manshe]{本质满射}[essential epimorphism]。
\end{definition}
\symidx{\(N\ll P\)：\(N\) 是 \(P\) 的多余子模}

这一条件也可写成：对每个子模 \(L\le P\)，若 \(p(L)=M\)，则 \(L=P\)。事实上，\(p(L)=M\) 等价于每个 \(x\in P\) 可写成 \(\ell+n\)，其中 \(\ell\in L,n\in\ker p\)，也就是 \(L+\ker p=P\)。因此任何仍能满射到目标的子模都已经是整个来源。多余性总是相对于指定的母模而言，判定时量化的是这个母模的全部子模。

\ProseParagraphBreak

等价地，对任意左 \(R\) 模同态 \(g:Q\to P\)，若复合 \(pg:Q\to M\) 满射，则 \(g\) 满射：此时 \(p(g(Q))=M\)，上一条件给出 \(g(Q)=P\)。反过来，把 \(g\) 取为任意子模 \(L\le P\) 的包含映射，就得到上一条件。

\begin{definition}\label{b2:def:projective-cover}
设 \(R\) 为含幺环、\(M\) 为幺左 \(R\) 模。若幺左 \(R\) 模的满同态 \(p:P\to M\) 的来源 \(P\) 投射，且 \(\ker p\ll P\)，则称 \(p\) 为 \(M\) 的一个\term[toushe-fugai]{投射覆盖}[projective cover]。定义中保留这个满同态，而不仅记录 \(P\) 的同构类型。
\end{definition}

若 \(M\) 有限生成，它的任意投射覆盖的来源也有限生成：取 \(M\) 的有限生成元，逐个提升到 \(P\)，令 \(L\) 为这些提升所生成的子模，则 \(p(L)=M\)。本质满射的条件给出 \(L=P\)。因此，在有限维代数 \(A\) 上，有限生成目标模的覆盖来源有限生成，进而有限维；这适用于本章讨论的有限维目标模。

\ProseParagraphBreak
\cref{b2:cor:nakayama-generators}说明，\(M/JM\) 的一组 \(A/J\) 模生成元可以提升为 \(M\) 的生成元。局部代数中，剩余模是在除环 \(A/J\) 上的向量空间，选基可得到最少的生成元。但一般情形下，由最少生成元给出的自由满射仍未必是投射覆盖。

\ProseParagraphBreak
例如，设 \(k\) 为域、\(A=M_2(k)\)、\(S=k^2\)。\(S\) 由一个非零向量生成，取第一列的映射 \(A\to S\) 因而用了最少的一个模生成元。不过真左理想 \(AE_{11}\) 已经满射到 \(S\)，故这个自由满射不是本质满射。限制 \(AE_{11}\to S\) 则是同构；\(AE_{11}\) 是正则模的直和因子，因而投射，这才给出一个投射覆盖。\(AE_{11}\) 的 \(k\) 维数为二，不能是 \(A\) 上的自由模。因此，最少生成元使自由来源的秩最小，来源仍可能缩小为一个非自由的投射模。

\subsection{幂等元与正交幂等元组的提升}
\label{b2:subsec:1503:02}

\begin{theorem}\label{b2:thm:idempotent-lifting}
设 \(R\) 为含幺环，\(I\) 为幂零双边理想。每个幂等元 \(\bar e\in R/I\) 都能提升为幂等元 \(e\in R\)。此外，两个这样的提升 \(e,f\) 共轭：存在 \(u\in1+I\)，使 \(f=ueu^{-1}\)。幂零假设不能一般省去：在 \(R=\mathbb Z\)、\(I=6\mathbb Z\) 中，幂等元 \(3+6\mathbb Z\) 不能提升为整数幂等元。
\end{theorem}
\begin{proof}
先处理 \(I^2=0\)。任取 \(a\) 提升 \(\bar e\)，令 \(d=a^2-a\in I\)，这是 \(a\) 未必幂等的误差。希望取 \(e=a+c\)，其中 \(c\in I\) 且与 \(a\) 交换。因为 \(I^2=0\) 使 \(c^2=0\)，所需条件化为
\[
e^2-e=d+ac+ca-c=d+(2a-1)c=0.
\]
这个误差修正方程对 \(c\) 是线性的。又有 \((2a-1)^2=1+4d\)；对任意 \(x\in I\)，\(dx\in I^2=0\)，所以 \((2a-1)^2x=x\)。因此在 \(I\) 上左乘 \(2a-1\) 的算子平方为恒等，可以取
\[
c=-(2a-1)d=(1-2a)d.
\]
所取 \(d,c\) 都是 \(a\) 的多项式，故确实与 \(a\) 交换；它们也都在 \(I\) 中。代入得到
\[
e^2-e=d+(2a-1)c
=d-(2a-1)^2d.
\]
由于 \((2a-1)^2=1+4d\)，右侧等于 \(-4d^2=0\)。因此 \(e\) 幂等，且与 \(a\) 模 \(I\) 相同。这里没有除以二，对任意特征都成立。

一般地，取 \(I^N=0\)。已在 \(R/I^r\) 中得到的幂等元，可以提升到 \(R/I^{r+1}\)：自然满射的核为 \(I^r/I^{r+1}\)，它的平方为零，因为 \(2r\ge r+1\)、\(r\ge1\)。将平方零情形用于这个商环，依次作 \(r=1,\ldots,N-1\) 的提升，最后得到 \(R/I^N=R\) 中的幂等元。

为证共轭性，先把 \(e,f\) 看成正则右模 \(R_R\) 上的左乘投影；左乘是右 \(R\) 线性的，因为 \((ex)r=e(xr)\)。它们分别给出分解
\[
R_R=eR\oplus(1-e)R=fR\oplus(1-f)R.
\]
我们寻找的左乘映射应把第一组分解的两个部分分别送到第二组对应的部分，从而满足 \(ue=fu\)。对 \(eR\) 中的元素左乘 \(f\)，对 \((1-e)R\) 中的元素左乘 \(1-f\)，便得到候选元素
\[
u=fe+(1-f)(1-e).
\]
事实上，对 \(x\in eR\) 有 \(ux=fx\in fR\)，对 \(x\in(1-e)R\) 有 \(ux=(1-f)x\in(1-f)R\)；直接相乘也得到 \(ue=fe=fu\)。这个公式先完成分解的匹配，还须保证匹配映射可逆。因 \(\bar f=\bar e\)，模 \(I\) 后有 \(\bar u=\bar e+(1-\bar e)=1\)，故 \(u\in1+I\)。幂零性使 \(u\) 可逆，可由有限几何和写出其逆。于是 \(ue=fu\) 给出 \(f=ueu^{-1}\)。

最后，\(3^2-3=6\)，所以 \(3+6\mathbb Z\) 幂等。整数幂等元满足 \(e(e-1)=0\)，而 \(\mathbb Z\) 是整环，故只有零或一；它们模六都不等于三，说明这个幂等元不能提升。这里 \((6\mathbb Z)^r=6^r\mathbb Z\ne0\) 对每个正整数 \(r\) 成立，恰好缺少幂零假设。
\end{proof}

要提升一组分解，所选提升还须保持彼此正交和总和为一；分别任意提升每个幂等元并不足够。\cref{b2:prop:idempotent-lift-compatibility-counterexample}在本节末给出各自幂等而不相容的完整计算。下一定理把这种相容性一并保证。

\begin{theorem}\label{b2:thm:orthogonal-idempotent-lifting}
设 \(R\) 为含幺环、\(I\subseteq R\) 为幂零双边理想，\(m\) 为正整数。若 \(\bar e_1,\ldots,\bar e_m\in R/I\) 满足
\[
\bar e_i\bar e_j=\delta_{ij}\bar e_i,
\qquad \sum_{i=1}^m\bar e_i=1,
\]
则存在提升 \(e_1,\ldots,e_m\in R\)，同样满足 \(e_ie_j=\delta_{ij}e_i\) 与 \(\sum_i e_i=1\)。
\end{theorem}
\begin{proof}
对 \(m\) 归纳。\(m=1\) 时取 \(e_1=1\)。设 \(m>1\)，先用\cref{b2:thm:idempotent-lifting}提升 \(\bar e_1\)，并令 \(c=1-e_1\)。余下部分的单位应为 \(c\)，所以转到角环 \(cRc\) 中提升其余幂等元。映射
\[
cRc\longrightarrow\bar c(R/I)\bar c
\]
满射：右侧元素可由任意代表元两边乘 \(c\) 得到。其核为 \(cIc\)，因为同时属于 \(I\) 与 \(cRc\) 的元素 \(x\) 满足 \(x=cxc\)。这个核幂零，\((cIc)^N\subseteq cI^Nc=0\)。

余下的 \(\bar e_2,\ldots,\bar e_m\) 都在该商角环中，仍相互正交，总和为 \(\bar c\)。在 \(cRc\) 中应用归纳假设，得到正交提升 \(e_2,\ldots,e_m\)，总和为 \(c\)。它们均被 \(e_1\) 在左右两侧乘为零，所以连同 \(e_1\) 得到所求完整正交组。
\end{proof}

平方零修正与角环归纳可参阅 Etingof 等人的《Introduction to Representation Theory》\cite[Proposition 9.1.1; Corollary 9.1.3, pp. 209--210]{EtingofEtAl2011RepresentationTheory}。本书已把单个提升和整组相容提升分别展开；后面的构造将同时使用这两个结论。

\subsection{多余子模与简单模的投射覆盖}
\label{b2:subsec:1503:03}

\begin{proposition}\label{b2:prop:finite-superfluous-radical}
设 \(k\) 为域、\(A\) 为有限维结合含幺 \(k\) 代数，\(J=J(A)\)。有限维幺左 \(A\) 模 \(P\) 的子模 \(N\) 多余，当且仅当 \(N\subseteq JP=\operatorname{rad}P\)。因此，对有限生成投射幺左 \(A\) 模的满射 \(p:P\to M\)，以下条件等价：它是投射覆盖；\(\ker p\subseteq JP\)；诱导映射 \(P/JP\to M/JM\) 是同构。若这个满射是投射覆盖且目标 \(M\) 本身投射，则 \(p\) 为同构。
\end{proposition}
\begin{proof}
若 \(N\ll P\)，而某个极大子模 \(L\) 不包含 \(N\)，则 \(L\subsetneq L+N\)，极大性迫使 \(L+N=P\)，与 \(L\ne P\) 矛盾。因此 \(N\) 包含于所有极大子模的交，由\cref{b2:thm:module-radical-jm}得到 \(N\subseteq JP=\operatorname{rad}P\)。

反过来，若 \(N\subseteq JP\) 且 \(L+N=P\)，则 \(L+JP=P\)。商 \(P/L\) 有限生成，并满足 \(J(P/L)=P/L\)，由\cref{b2:thm:noncommutative-nakayama}，\(P/L=0\)，即 \(L=P\)。

对满射 \(p\)，有 \(p(JP)=JM\)。诱导的商映射满射，核为 \((\ker p+JP)/JP\)：若 \(p(x)\in JM\)，取 \(y\in JP\) 使 \(p(y)=p(x)\)，便有 \(x-y\in\ker p\)。所以该核为零恰等于 \(\ker p\subseteq JP\)，得到其余等价性。

若 \(p\) 是覆盖且 \(M\) 投射，取截面 \(s:M\to P\)，便有 \(P=\ker p\oplus s(M)\)。核多余使 \(s(M)=P\)，从而 \(\ker p=0\)，\(p\) 是同构。
\end{proof}

\begin{definition}\label{b2:def:primitive-idempotent}
设 \(k\) 为域、\(A\) 为结合含幺 \(k\) 代数、\(e\in A\) 为非零幂等元。若 \(e\) 不能写成两个非零正交幂等元之和，则称它为\term[benyuan-midengyuan]{本原幂等元}[primitive idempotent]。这里的正交要求两个方向的乘积都为零。
\end{definition}

这里“本原”修饰幂等元，表示它不能再作正交分解；\cref{b2:def:primitive-semiprimitive}中的本原环则要求环有忠实简单模，两者的定义不同。

\begin{proposition}\label{b2:prop:primitive-projective-indecomposable}
设 \(k\) 为域、\(A\) 为结合含幺 \(k\) 代数、\(e\in A\) 为非零幂等元。幺左 \(A\) 模 \(Ae\) 有限生成且投射；它不可分解，当且仅当 \(e\) 本原。
\end{proposition}
\begin{proof}
\(Ae\) 由单个元素 \(e\) 生成。正则左模具有分解 \(A=Ae\oplus A(1-e)\)，投影为右乘 \(e\)，所以 \(Ae\) 是有限自由模的直和因子，因而投射。若 \(e=f+g\) 是非零正交幂等元之和，则 \(Ae=Af\oplus Ag\)，给出非平凡分解。

反过来，若 \(Ae\) 有非平凡分解，对应投影是 \(\operatorname{End}_A(Ae)\) 中一个非零且非恒等的幂等元。由\cref{b2:prop:corner-endomorphism}，该自同态环同构于 \((eAe)^{\mathrm{op}}\)，所以角环中也有 \(f^2=f\)、\(f\ne0,e\)。于是 \(e=f+(e-f)\)，且 \(f(e-f)=(e-f)f=0\)，这与本原性矛盾。
\end{proof}

\begin{proposition}\label{b2:prop:primitive-projective-cover}
设 \(k\) 为域、\(A\) 为有限维结合含幺 \(k\) 代数，\(J=J(A)\)。取 \(A/J\) 的本原幂等元 \(\bar e\)，并取其在 \(A\) 中的幂等元提升 \(e\)，则 \(e\) 本原，\(Ae/Je\) 是简单幺左 \(A\) 模，且自然满射
\[
Ae\longrightarrow Ae/Je
\]
是投射覆盖。每个简单幺左 \(A\) 模都能由这样的满射得到覆盖。
\end{proposition}
\begin{proof}
先看单个矩阵环 \(M_n(D)\)，其中 \(D\) 为除环。幂等矩阵作为右 \(D\) 向量空间上的算子，其像与核给出直和；选择相应基后，它共轭于 \(\operatorname{diag}(I_r,0)\)。后者是 \(r\) 个秩一坐标投影之和，所以 \(r>1\) 时不本原。若 \(r=1\)，两个非零正交幂等元的像会给出两个非零直和分量，不能同时放在这个一维像中。因此非零幂等元本原当且仅当 \(r=1\)。再看\eqref{b2:eq:finite-algebra-semisimple-quotient}中的有限矩阵环直积，本原幂等元只能在一个矩阵因子中具有秩一，其余因子为零。相应左理想同构于一个列模，且简单：若列 \(v\) 的第 \(j\) 个坐标非零，为把它映到任意给定列 \(w\)，取矩阵 \(C\) 的第 \(j\) 列为 \((w_iv_j^{-1})_i\)，其余列为零，便有 \(Cv=w\)。

\(J\) 零化每个简单左 \(A\) 模，所以它们都可看作 \(A/J\) 的简单模。半单商的正则模是上述简单列模的有限直和；任意简单模都是这个正则模的商。相应满射在至少一个列模上的限制非零，由\cref{b2:lem:schur}，该限制是同构。这也说明上述构造遍及全部简单模。

若 \(e=f+g\) 是正交幂等元分解，模 \(J\) 后，\(\bar e\) 的本原性使 \(\bar f,\bar g\) 至少一个为零。该幂等元便属于幂零理想 \(J\)，只能为零，因为幂等元的每个正幂都等于自身。因此 \(e\) 本原。最后，\(Ae\to(A/J)\bar e\) 的核为 \(Je\)：核元素既在 \(J\) 中，又满足 \(x=xe\)，故在 \(Je\) 中。于是 \(Ae/Je\) 是刚才的简单列模，且核为 \(J(Ae)\)。由\cref{b2:prop:finite-superfluous-radical}，满射为投射覆盖。
\end{proof}

\subsection{投射覆盖的存在性}
\label{b2:subsec:1503:04}

从 \(A/J\) 的每个矩阵因子选出全部对角矩阵单位，利用\cref{b2:thm:orthogonal-idempotent-lifting}得到 \(A\) 中的完整正交提升。每个提升给出上述简单列模的覆盖；同一矩阵因子的列模同构，但暂不预设其提升模已经同构。

\ProseParagraphBreak
从每一种简单模的同构类型中取一个代表 \(S_1,\ldots,S_s\)，并固定由\cref{b2:prop:primitive-projective-cover}构造的覆盖 \(p_i:P_i\to S_i\)。于是 \(P_i/JP_i\cong S_i\)，且 \(P_i\) 不可分解。

\begin{theorem}\label{b2:thm:projective-cover-existence}
设 \(k\) 为域、\(A\) 为有限维结合含幺 \(k\) 代数，\(J=J(A)\)。从全部简单幺左 \(A\) 模的同构类中各选一个代表 \(S_1,\ldots,S_s\)，并给每个 \(S_i\) 选定投射覆盖 \(p_i:P_i\to S_i\)，则每个有限维幺左 \(A\) 模 \(M\) 都有投射覆盖。若
\[
M/JM\cong\bigoplus_{i=1}^s S_i^{\oplus m_i},
\]
可以取覆盖的来源为 \(P=\bigoplus_iP_i^{\oplus m_i}\)。
\end{theorem}
\begin{proof}
由\cref{b2:thm:module-radical-jm}，所写的有限半单分解存在。合并 \(p_i\) 的相应份数并接上这个同构，得到满同态 \(h:P\to M/JM\)，其在商上诱导同构 \(P/JP\cong M/JM\)。\(P\) 是有限个投射模的直和，因而投射，可以沿商映射 \(\pi:M\to M/JM\) 提升 \(h\)，得到 \(p:P\to M\) 且 \(\pi p=h\)。

因为 \(h\) 满射，有 \(p(P)+JM=M\)。对有限生成的 \(M\) 及其子模 \(p(P)\) 应用\cref{b2:cor:nakayama-generators}的生成元提升形式，得到 \(p(P)=M\)。又因为 \(p\) 在模 \(J\) 后是同构，\(\ker p\subseteq JP\)。由\cref{b2:prop:finite-superfluous-radical}，这个满射正是投射覆盖。
\end{proof}

这里并未把 \(M/JM\) 的简单直和分解直接提升成 \(M\) 的简单直和分解。被提升的是它们的投射覆盖，以及一个从投射模出发的映射；满射性由 Nakayama 引理重新保证。

\ProseParagraphBreak
例如在 \(A=k[t]/(t^3)\) 中，\(M=A/(\bar t^2)\) 的最大半单商为 \(k\)，对应投射覆盖是 \(A\to M\)。核 \((\bar t^2)\) 包含在 \(JA\) 中，所以多余；这个覆盖保留了 \(M\) 的两层结构，而来源 \(A\) 本身有三层。

\subsection{不可分解投射模与简单模}
\label{b2:subsec:1503:05}

\begin{theorem}\label{b2:thm:indecomposable-projective-simple}
设 \(k\) 为域、\(A\) 为有限维结合含幺 \(k\) 代数，\(J=J(A)\)。从全部简单幺左 \(A\) 模的同构类中各选一个代表 \(S_1,\ldots,S_s\)，并按\cref{b2:prop:primitive-projective-cover}给每个 \(S_i\) 选定形如 \(p_i:P_i=Ae_i\to S_i\) 的投射覆盖。于是每个有限生成投射幺左 \(A\) 模都是 \(P_1,\ldots,P_s\) 的有限直和。有限生成不可分解投射模恰为这些 \(P_i\) 的同构类型，而
\[
P\longmapsto P/JP
\]
给出有限生成不可分解投射模与简单模的同构类型之间的一一对应。
\end{theorem}
\begin{proof}
对有限生成投射模 \(Q\)，按\cref{b2:thm:projective-cover-existence}构造覆盖 \(p:P\to Q\)，其中 \(P\) 是若干 \(P_i\) 的直和。\(Q\) 投射，而 \(p\) 是投射覆盖，由\cref{b2:prop:finite-superfluous-radical}，\(p\) 本身就是同构，故 \(Q\cong P\)。

各 \(P_i\) 已由\cref{b2:prop:primitive-projective-cover}证明不可分解；若 \(Q\) 不可分解，上述直和只能有一个非零分量。不同的 \(P_i\) 不同构，因为其商 \(P_i/JP_i\cong S_i\) 不同构。每个简单模都有其中一个覆盖，故得到双射。
\end{proof}

因此，这里对应的是三种对象的同构类型：
\[
S_i\quad\longleftrightarrow\quad P_i\cong Ae_i
\quad\longleftrightarrow\quad
\{\,e\text{ 为本原幂等元}\mid Ae\cong Ae_i\,\}.
\]
右端是一类本原幂等元，而不是某一个唯一的幂等元。

\ProseParagraphBreak
直和中的重数也可以直接从最大半单商读出。对有限生成投射左 \(A\) 模 \(Q\)，令 \(\Delta_i=\operatorname{End}_A(S_i)\)，它由\cref{b2:lem:schur}为除环。\(\operatorname{Hom}_A(Q,S_i)\) 通过像端的复合成为左 \(\Delta_i\) 向量空间，具体为 \(\alpha\cdot f=\alpha\circ f\)。因为 \(J\) 零化 \(S_i\)，每个同态都通过 \(Q/JQ\) 分解，因而
\begin{equation}\label{b2:eq:projective-simple-multiplicity}
Q\cong\bigoplus_iP_i^{\oplus m_i}
\quad\Longrightarrow\quad
m_i=\dim_{\Delta_i}\operatorname{Hom}_A(Q,S_i).
\end{equation}
理由是不同简单模之间没有非零同态，\(\operatorname{Hom}_A(S_i,S_i)=\Delta_i\) 在自身上的左维数为一。这也证明这些重数唯一，且没有假设 \(\Delta_i=k\)。\cref{b2:prop:semisimple-multiplicity}从来源简单模的自同态作预复合，得到右除环向量空间；这里简单模位于像端，作后复合，因而得到左除环向量空间。

\subsection{投射覆盖的唯一性与极小性}
\label{b2:subsec:1503:06}

存在性已经由\cref{b2:thm:projective-cover-existence}得到。以下唯一性须比较连同映射的覆盖；只比较来源的维数，不足以说明它们在 \(M\) 上相容。

\begin{theorem}\label{b2:thm:projective-cover-uniqueness}
设 \(k\) 为域、\(A\) 为有限维结合含幺 \(k\) 代数，\(M,N\) 为有限维幺左 \(A\) 模，\(h:M\to N\) 为 \(A\) 模同构，\(p:P\to M\) 与 \(q:Q\to N\) 为投射覆盖。则存在满足 \(qa=hp\) 的 \(A\) 模同态 \(a:P\to Q\)，且每个满足这一等式的 \(a\) 都是同构。取 \(N=M\)、\(h=1_M\)，便得到两个覆盖在目标 \(M\) 上同构；这样的同构本身未必唯一。
\end{theorem}
\begin{proof}
由 \(P\) 的投射性，可以沿满射 \(q\) 提升 \(hp\)，所以满足 \(qa=hp\) 的 \(a\) 存在。现在固定任意一个这样的 \(a\)。由 \(Q\) 的投射性，沿 \(p\) 提升 \(h^{-1}q\)，得到 \(b:Q\to P\)，满足 \(pb=h^{-1}q\)。于是 \(pba=p\)、\(qab=q\)。对任意 \(x\in P\)，有 \(x-ba(x)\in\ker p\)，所以
\[
ba(P)+\ker p=P.
\]
核多余，故 \(ba(P)=P\)。同样由 \(qab=q\) 及 \(\ker q\ll Q\) 得 \(ab(Q)=Q\)。这两步只用了覆盖的多余核条件。

目标 \(M,N\) 有限生成，本质满射性使来源 \(P,Q\) 也有限生成；又因 \(A\) 有限维，\(P,Q\) 都有限维，因而具有有限长度。\cref{b2:prop:finite-length-endomorphism}使满自同态 \(ba\) 与 \(ab\) 可逆。于是 \((ba)^{-1}b\) 是 \(a\) 的左逆，而 \(b(ab)^{-1}\) 是其右逆；一个映射的左逆与右逆相同，故给定的 \(a\) 是同构。

例如 \(A=k[t]/(t^2)\) 的覆盖 \(A\to k\) 上，恒等映射和乘以 \(1+\bar t\) 的映射都是同构，且模 \(\bar t\) 后都诱导恒等，二者不同。因此唯一的是覆盖在目标上同构的类型。
\end{proof}

\begin{proposition}\label{b2:prop:projective-cover-minimality}
设 \(k\) 为域、\(A\) 为有限维结合含幺 \(k\) 代数，\(M\) 为幺左 \(A\) 模。对有限生成投射幺左 \(A\) 模的满射 \(p:P\to M\)，它是投射覆盖，当且仅当每个满足 \(pu=p\) 的 \(A\) 模自同态 \(u:P\to P\) 都可逆。若 \(p\) 是覆盖，且 \(q:Q\to M\) 是任意投射幺左 \(A\) 模的满射，则 \(P\) 是 \(Q\) 的直和因子，不要求 \(Q\) 有限生成。
\end{proposition}
\begin{proof}
覆盖情形中，\(u(P)+\ker p=P\)，所以 \(u\) 满射。由命题假设，\(P\) 在有限维 \(k\) 代数 \(A\) 上有限生成，故它本身有限维，\(u\) 因而可逆。反过来，若 \(L+\ker p=P\)，则 \(p|_L:L\to M\) 满射。利用 \(P\) 的投射性，把 \(p\) 提升到 \(L\)，再复合包含映射得到 \(u:P\to P\)，其像在 \(L\) 中且 \(pu=p\)。若它必可逆，则 \(L=P\)，所以核多余。

对第二个断言，投射性给出 \(a:P\to Q\)、\(b:Q\to P\)，满足 \(qa=p\)、\(pb=q\)。于是 \(pba=p\)，第一部分使 \(ba\) 可逆。映射 \(a(ba)^{-1}:P\to Q\) 是 \(b\) 的截面，因此 \(Q\cong P\oplus\ker b\)。
\end{proof}

这使“覆盖极小”有了可检验的含义：任何其他投射满射的来源都含有它作为直和因子；若一个来源还能通过自同态缩小却仍覆盖目标，就不满足这种极小性。半单代数上的有限维模都投射，所以它们的覆盖不会增添额外的层，来源经覆盖映射与目标同构。关于同构类的选择虽然唯一，关于具体提升映射的选择仍须在后续构造中保留。

\begin{proposition}\label{b2:prop:idempotent-lift-compatibility-counterexample}
设 \(k\) 为域、\(R=M_2\bigl(k[\varepsilon]/(\varepsilon^2)\bigr)\)，\(\varepsilon\) 的商类仍记为 \(\varepsilon\)，\(E_{ij}\) 为矩阵单位，\(I=\varepsilon R\)。虽然 \(I^2=0\)，幂等元
\[
e=E_{11}+\varepsilon E_{12},\qquad f=E_{22}
\]
分别提升了 \(R/I\) 中正交且和为一的 \(\bar E_{11},\bar E_{22}\)，所选 \(e,f\) 却不正交。若改取 \(f'=1-e=E_{22}-\varepsilon E_{12}\)，则 \(e,f'\) 为同一商幂等元组的正交提升，且和为一。
\end{proposition}
\begin{proof}
\(I\) 为双边理想，且 \(\varepsilon^2=0\) 给出 \(I^2=0\)。由矩阵单位乘法，\(E_{11}E_{12}=E_{12}\)、\(E_{12}E_{11}=0\)，所以
\[
e^2=E_{11}+\varepsilon E_{12}=e,\qquad f^2=f.
\]
模 \(I\) 后，它们分别为 \(E_{11},E_{22}\)，其乘积在两个方向都为零，和为单位。但是在 \(R\) 中，
\[
ef=\varepsilon E_{12}\ne0,\qquad fe=0.
\]
因此分别取出的幂等提升不相互正交，它们的和也为 \(1+\varepsilon E_{12}\)，而不是一。取 \(f'=1-e\)，由 \(e^2=e\) 直接得到 \((f')^2=f'\)、\(ef'=f'e=0\) 及 \(e+f'=1\)；同时 \(f'-f\in I\)，所以没有改变第二个商幂等元。所有计算均适用于任意特征。
\end{proof}

% !TeX root = ../../Book2.tex
% 纲要标题：### 15.4 基本代数的动机
% 纲要位置：计划纲要.md，第 1267 行。

\section{基本代数的动机}
\label{b2:sec:1504}

正则模中同一种不可分解投射模可能出现多次；矩阵大小的一部分信息正来自这种重复。本节从每一种同构类型中只保留一个，构造一个较小的角代数，并说明它为什么适合记录简单模之间的关系。本节只证明这个角代数的基本性质，模范畴的等价留在下一章建立。

% 纲要要点（供目录核对）：
%
% * 本原幂等元
% * 简单模的自同态除环
% * 基本与分裂基本代数；第五卷箭图表示的代数来源
%
% ---
%

\subsection{本原幂等元}
\label{b2:subsec:1504:01}

\cref{b2:def:primitive-idempotent}与\cref{b2:prop:primitive-projective-indecomposable}已经把本原幂等元与不可分解投射左理想联系起来。本原性不要求幂等元位于中心：例如上三角矩阵代数的 \(E_{ii}\) 本原，却通常不是中心元。正交分解一的用途，是分解左正则模；只有各项还位于中心时，才能把它解释成代数的直积。

\begin{proposition}\label{b2:prop:corner-radical}
设 \(k\) 为域、\(A\) 为有限维结合含幺 \(k\) 代数，\(J=J(A)\)。对幂等元 \(e\in A\)，记 \(\bar e=e+J\in A/J\)。角代数 \(B=eAe\) 满足
\[
J(B)=eJe,
\qquad
B/J(B)\cong\bar e(A/J)\bar e.
\]
若 \(e\ne0\)，则 \(e\) 本原当且仅当 \(\bar e\) 本原；此时右侧是除环，\(eAe\) 是局部环。
\end{proposition}
\begin{proof}
向商代数映射并取角给出满射 \(eAe\to\bar e(A/J)\bar e\)，核为 \(eJe\)，这与\cref{b2:thm:orthogonal-idempotent-lifting}证明中的角环商计算相同。\(eJe\) 幂零，因为其 \(N\) 次乘积包含在 \(eJ^Ne\) 中，而 \(J^N=0\)。因此 \(eJe\subseteq J(B)\)。

在\eqref{b2:eq:finite-algebra-semisimple-quotient}的每个矩阵块中，对幂等算子 \(\bar e\) 选择像与核的基，使它成为 \(\operatorname{diag}(I_{r_i},0)\)。相应角环为 \(M_{r_i}(D_i)\)，零秩的块省略。所以商角环是半单代数，根为零。由\cref{b2:prop:radical-quotient}，\(J(B)/eJe=0\)，故 \(J(B)=eJe\)。

若 \(\bar e\) 本原，\cref{b2:prop:primitive-projective-cover}已证 \(e\) 本原。反过来，若 \(e\) 本原，则 \(Ae\) 不可分解，由\cref{b2:thm:indecomposable-projective-simple}，\(Ae/Je\cong(A/J)\bar e\) 简单。上述矩阵块描述迫使恰有一个 \(r_i=1\)，其余为零，故 \(\bar e\) 本原，商角环是除环。

最后，在含幺环 \(B\) 中，模 \(J(B)\) 后可逆的元素本身可逆。确实，若 \(xy\) 与 \(yx\) 都模 \(J(B)\) 等于单位，则它们由\cref{b2:thm:jacobson-unit-criterion}可逆，从而 \(x\) 同时有右逆 \(y(xy)^{-1}\) 与左逆 \((yx)^{-1}y\)，两者相同。于是当 \(B/J(B)\) 是除环时，\(B\) 的非单位元恰为 \(J(B)\)，它是唯一极大左理想，故 \(B\) 局部。此处角环的单位是 \(e\)。
\end{proof}

本原性在这里能传过商，依赖有限维情形下 \(J\) 的幂零性；对任意理想取商，幂等元的提升与上述结论都不能直接沿用。

\ProseParagraphBreak

例如 \(A=k[t]/(t^m)\) 的单位 \(1\) 是本原幂等元：该代数局部，模根后只有一个除环 \(k\)。当 \(m>1\) 时，角代数就是 \(A\) 自身，它虽局部却不是除环。局部性允许存在非零根，这与简单商是除环并不矛盾。

\subsection{简单模的自同态除环}
\label{b2:subsec:1504:02}

固定简单模 \(S_i\) 及其覆盖 \(P_i\)，仍记 \(\Delta_i=\operatorname{End}_A(S_i)\)。Schur 引理使 \(\Delta_i\) 成为除环，而不是预先成为基域。它控制\eqref{b2:eq:projective-simple-multiplicity}中的重数计数。

\begin{proposition}\label{b2:prop:projective-endomorphism-residue}
设 \(k\) 为域、\(A\) 为有限维结合含幺 \(k\) 代数，\(J=J(A)\)。令 \(S_i\) 为简单左 \(A\) 模，\(p_i:P_i\to S_i\) 为其投射覆盖，并记 \(\Delta_i=\operatorname{End}_A(S_i)\)。沿 \(P_i/JP_i\cong S_i\) 取商，从 \(P_i\) 的自同态得到满同态
\[
\operatorname{End}_A(P_i)\longrightarrow\operatorname{End}_A(S_i)=\Delta_i.
\]
其核是把 \(P_i\) 映入 \(JP_i\) 的全部自同态，恰为 \(\operatorname{End}_A(P_i)\) 的 Jacobson 根。
\end{proposition}
\begin{proof}
每个自同态保持 \(JP_i\)，所以诱导商上的自同态。若给定 \(\alpha:S_i\to S_i\)，利用 \(P_i\) 的投射性，把 \(\alpha p_i:P_i\to S_i\) 沿 \(p_i:P_i\to S_i\) 提升，就得到所需原像，证明满射。

核的描述来自商映射的定义。记此核为 \(K\)。若 \(u\in K\)，则对每个 \(r\) 有 \(u(J^rP_i)\subseteq J^r\cdot u(P_i)\subseteq J^{r+1}P_i\)。因而 \(N\) 个核中自同态的复合把 \(P_i\) 映入 \(J^NP_i=0\)，所以 \(K\) 是幂零理想。它包含于该自同态环的根，而模掉它的商是除环 \(\Delta_i\)，根为零；再用\cref{b2:prop:radical-quotient}，得到核恰为根。
\end{proof}

若 \(P_i=Ae\)，由\cref{b2:prop:corner-endomorphism}，这个自同态环是 \((eAe)^{\mathrm{op}}\)。因此半单商矩阵块里的除环与简单模自同态除环相差一个反环；按本书的左模约定，可在\eqref{b2:eq:finite-algebra-semisimple-quotient}中选取
\[
\Delta_i=\operatorname{End}_A(S_i),\qquad D_i=\Delta_i^{\mathrm{op}}.
\]
非交换除环上不能省略这个反环。

\begin{proposition}\label{b2:prop:split-simple-endomorphisms}
设 \(k\) 为代数闭域、\(A\) 为有限维结合含幺 \(k\) 代数。每个有限维简单左 \(A\) 模 \(S\) 都满足 \(\operatorname{End}_A(S)=k\operatorname{id}_S\)。
\end{proposition}
\begin{proof}
这是\cref{b2:cor:schur-algebraically-closed}的应用：本章的简单模有限维，基域代数闭，故每个简单自同态都是标量。该推论不要求 \(A\) 本身半单，因此可以直接用于这里的有限维代数。
\end{proof}

在一般基域上，此结论不成立。例如把 \(\mathbb C\) 看成实代数，简单左模 \(S=\mathbb C\) 的自同态除环是 \(\mathbb C\)，并不等于 \(\mathbb R\)。相应的 Hom 空间在基域上的维数为二，在这个自同态除环上的维数才为一。

\subsection{箭图表示的代数来源}
\label{b2:subsec:1504:03}

\begin{definition}\label{b2:def:basic-algebra}
设 \(k\) 为域、\(B\) 为有限维结合含幺 \(k\) 代数。若 \(B/J(B)\) 是有限个除 \(k\) 代数的直积，即其 Artin--Wedderburn 分解的每个矩阵块大小都为一，则称 \(B\) 为\term[jiben-daishu]{基本代数}[basic algebra]。
若进一步有 \(B/J(B)\cong k^s\)，则称 \(B\) 为\term[fenlie-jiben-daishu]{分裂基本代数}[split basic algebra]。
\end{definition}

\begin{proposition}\label{b2:prop:basic-corner}
设 \(k\) 为域、\(A\) 为有限维结合含幺 \(k\) 代数。可以选出幂等元 \(e\in A\)，使 \(eAe\) 为基本 \(k\) 代数，并满足 \(AeA=A\)。具体地，可在半单商 \(A/J(A)\) 的每个矩阵因子中只保留一个对角本原幂等元，再把这些幂等元相容提升并求和。
\end{proposition}
\begin{proof}
把\eqref{b2:eq:finite-algebra-semisimple-quotient}中全部对角矩阵单位组成完整正交组，由\cref{b2:thm:orthogonal-idempotent-lifting}提升为 \(e_{ij}\)，其中 \(i\) 标记矩阵因子，\(1\le j\le n_i\)。令 \(e=\sum_i e_{i1}\)，正交性保证它幂等。由\cref{b2:prop:corner-radical}，
\[
eAe/J(eAe)\cong\bar e(A/J)\bar e\cong\prod_iD_i,
\]
所以这个角代数基本。

在每个矩阵块中，一个对角矩阵单位生成整个双边理想：\(E_{p1}E_{11}E_{1q}=E_{pq}\)。因此 \(\bar e\) 在 \(A/J\) 中生成全部商代数，得到 \(AeA+J=A\)。写 \(1=x+j\)，其中 \(x\in AeA,j\in J\)。由根的可逆性判据，\(x=1-j\) 可逆；包含单位元的双边理想只能是整个 \(A\)，故 \(AeA=A\)。
\end{proof}

构造中同时得到 \(AeA=A\)：只保留每种列模的一份，仍能由角上的元素生成整个代数。这正是下一章比较模范畴所需的条件。设 \(A\) 为含幺代数、\(e\in A\) 为幂等元。若 \(AeA=A\)，则称 \(e\) 为\term[man-midengyuan]{满幂等元}[full idempotent]。例如 \(M_n(k)\) 中取 \(e=E_{11}\)，角代数就是 \(k\)；上三角矩阵代数和 \(k[t]/(t^m)\) 则已经是基本代数。基本性并不要求根为零。

\begin{proposition}\label{b2:prop:full-corner-radical-powers}
设 \(k\) 为域、\(A\) 为有限维结合含幺 \(k\) 代数，\(J=J(A)\)，\(e\in A\) 为满足 \(AeA=A\) 的幂等元，\(B=eAe\)。对每个整数 \(r\ge0\)，有
\[
J(B)^r=eJ^re.
\]
\end{proposition}
\begin{proof}
由\cref{b2:prop:corner-radical}，\(J(B)=eJe\)，而 \(r=0\) 时两边都等于 \(eAe\)。设 \(r\ge1\) 且已知 \((eJe)^r=eJ^re\)。显然 \((eJ^re)(eJe)\subseteq eJ^{r+1}e\)。反过来，由 \(AeA=A\) 可写 \(1=\sum_t a_teb_t\)。对 \(x\in J^r\)、\(y\in J\)，有
\[
exye=\sum_t(exa_te)(eb_tye)\in(eJ^re)(eJe).
\]
所以 \(eJ^{r+1}e\subseteq(eJ^re)(eJe)\)，归纳即得结论。
\end{proof}

对分裂基本代数，半单商 \(k^s\) 的坐标幂等元给出顶点，\(J/J^2\) 则记录顶点间的一次连接。一般基本代数的半单商仍可能含有非平凡除环，此时仅以普通 \(k\)-箭头记录维数会丢失顶点除环和箭头的双模结构。

\begin{definition}\label{b2:def:basic-algebra-quiver}
设 \(k\) 为域、\(B\) 为有限维分裂基本 \(k\) 代数，即 \(B/J(B)\cong k^s\)。取完整正交本原幂等元组 \(e_1,\ldots,e_s\)，并记 \(J=J(B)\)。以 \(1,\ldots,s\) 为顶点，从 \(i\) 到 \(j\) 画 \(\dim_k e_j(J/J^2)e_i\) 条箭头，得到的允许重边与圈的有向图称为 \(B\) 的\term[jiantu]{箭图}[quiver]。箭头方向配合左模约定。
\end{definition}

这张图的同构类型与所选幂等元组无关。若另选完整正交本原幂等元组 \(f_1,\ldots,f_s\)，由\cref{b2:prop:corner-radical}可知两组元素模 \(J\) 后都是 \(k^s\) 的坐标幂等元；重新编号后有 \(f_i-e_i\in J\)。对 \(x\in J\)，于是 \(f_jxf_i-e_jxe_i\in J^2\)，两组选择在 \(J/J^2\) 中给出相同的分量，箭头数相同。

\begin{proposition}\label{b2:prop:basic-algebra-arrow-generators}
设 \(k\) 为域、\(B\) 为有限维基本 \(k\) 代数，\(J=J(B)\)，并假设 \(B/J\cong k^s\)。取完整正交本原幂等元组 \(e_1,\ldots,e_s\)。对每个 \(e_j(J/J^2)e_i\) 选一组 \(k\) 基，并为每个基向量选取提升 \(x\in e_jJe_i\)，则这些提升连同 \(e_1,\ldots,e_s\) 生成代数 \(B\)。对每个左 \(B\) 模 \(M\)，有向量空间分解 \(M=\bigoplus_i e_iM\)，且提升 \(x\in e_jJe_i\) 的作用给出 \(e_iM\to e_jM\) 的线性映射。
\end{proposition}
\begin{proof}
令 \(L\) 为全部箭头提升的线性生成空间，\(C\) 为它们与全部 \(e_i\) 生成的子代数。由于提升的商类构成 \(J/J^2\) 的基，有 \(J=L+J^2\)。用 \(L^r\) 表示 \(r\) 个 \(L\) 中元素乘积的线性生成空间；展开 \(r\) 个这样的和的乘积，得
\[
J^r\subseteq L^r+J^{r+1}\subseteq C+J^{r+1}.
\]
又因 \(B/J=k^s\) 由 \(e_i\) 的商类生成，有 \(B=C+J\)。逐次把余项从 \(J\) 推到 \(J^2,J^3,\ldots\)，最终用 \(J^N=0\) 得 \(B=C\)。

对模块，\(\sum_i e_i=1\) 给出 \(m=\sum_i e_im\)，正交性使这项表达唯一，所以是向量空间直和。若 \(m=e_im\)，则 \(xm=e_jxe_im\in e_jM\)，得到所述线性映射。一般的 \(e_iM\) 未必是 \(B\) 子模；保存所有箭头的作用才能恢复各坐标之间的联系。
\end{proof}

例如 \(T_3(k)\) 的 \(J/J^2\) 有基 \(\overline E_{12},\overline E_{23}\)，给出箭头 \(2\to1\)、\(3\to2\)。元素 \(E_{13}=E_{12}E_{23}\) 来自两条箭头的复合，不再需要第三条独立箭头。一般代数中，不同路径的乘积还可能满足关系，所以箭图本身并不总能确定代数。第五卷第9.1节将引入路径代数，第9.5节进一步讨论基本代数与关系；本卷\chapref{b2:ch:16}先说明满幂等元的角代数如何与原代数的模范畴联系。


\begin{chaptersummary}
有限维性使根的幂构成的降链稳定，Nakayama 引理再使稳定项为零；半单商则把各个根滤过层变成半单模。对有限维模，\(\operatorname{rad}M=JM\)，而基座由被 \(J\) 零化的元素组成。两列的终止步数相同，但各层的具体位置和重数可能不同。

\ProseParagraphBreak
投射覆盖需要投射来源及多余核。幂等元提升把半单商中的列模转化为不可分解投射模；从最大半单商提升映射，再用 Nakayama 引理恢复满射性，便得到任意有限维模的覆盖。两个覆盖在目标上同构，不过具体同构未必唯一。简单模决定不可分解投射模的类型，简单自同态除环决定相应的重数计数。

\ProseParagraphBreak
选取每种不可分解投射类型的一份，给出基本满角代数。它的半单商去除了重复的矩阵大小；因所选幂等元满足 \(AeA=A\)，角代数的根幂还满足 \(J(eAe)^r=eJ(A)^re\)。分裂基本代数的 \(J/J^2\) 可以用箭头记录，较高次乘积则对应路径及关系；这些构造将分别在下一章的 Morita 理论和第五卷的箭图表示中继续展开。
\end{chaptersummary}

\BookExercises

\begin{exercise}\label{b2:ex:15-block-triangular-radical}
令 \(A\) 为所有分块矩阵 \(\begin{psmallmatrix}B&X\\0&C\end{psmallmatrix}\) 组成的代数，其中 \(B\in M_r(k)\)、\(C\in M_s(k)\)、\(X\in M_{r\times s}(k)\)，且 \(r,s\ge1\)。求 \(J(A)\)、半单商和根的幂零指数，并说明根以外仍可有非零幂零元素。
\end{exercise}
\begin{solution}
仅右上块可能非零的矩阵组成理想 \(I\)，任意两个这样的矩阵相乘为零，所以 \(I^2=0\)。商为 \(M_r(k)\times M_s(k)\)，是半单代数；由\cref{b2:prop:nilpotent-ideal-in-radical}及\cref{b2:prop:radical-quotient}，\(J(A)=I\)。因为 \(r,s\ge1\)，\(I\ne0\)，幂零指数为二。若 \(r\ge2\)，在左上块放 \(E_{12}\)、其余块为零，便得到不属于 \(J(A)\) 的非零平方零元素；\(s\ge2\) 时同样可用右下块。若 \(r=s=1\)，所有幂零元素恰在严格上三角理想内。
\end{solution}

\begin{exercise}\label{b2:ex:15-loewy-length-bounds}
设 \(M\ne0\) 为有限维左 \(A\) 模。证明其 Loewy 长度不超过 \(\dim_kM\)。再证明有限直和的 Loewy 长度等于各分量 Loewy 长度的最大值（全零时为零），而子模与商模的 Loewy 长度都不超过原模。
\end{exercise}
\begin{solution}
若 \(J^rM=J^{r+1}M\)，将 Nakayama 引理用于有限生成模 \(J^rM\)，得到它为零。因此到达零以前，每个相邻根层都使维数至少下降一，步数至多 \(\dim_kM\)。

对直和，\(J^r(\bigoplus_iM_i)=\bigoplus_iJ^rM_i\)，所以第 \(r\) 步为零恰好等于每个分量在该步为零，得到最大值公式。若 \(N\le M\)，则 \(J^rN\subseteq J^rM\)；商中有 \(J^r(M/N)=(J^rM+N)/N\)。原模的某次幂为零时，这两处也为零，证明所需界。
\end{solution}

\begin{exercise}\label{b2:ex:15-associated-radical-algebra}
令 \(A=T_3(k)\)、\(J=J(A)\)。计算根滤过的伴随分次代数 \(\operatorname{gr}_J(A)\)，并说明它同构于给 \(E_{ij}\) 赋次数 \(j-i\) 的原上三角矩阵代数。
\end{exercise}
\begin{solution}
对上三角矩阵，\(J^r\) 由满足 \(j-i\ge r\) 的 \(E_{ij}\) 生成；所以第 \(r\) 个分次片以 \(j-i=r\) 的矩阵单位的商类为基。把 \(E_{ij}\) 送到相应商类，是分次向量空间同构。非零乘积 \(E_{ij}E_{j\ell}=E_{i\ell}\) 的次数正好相加，其他矩阵单位乘积为零，故这也是代数同构。
\end{solution}

\begin{exercise}\label{b2:ex:15-radical-socle-subquotient}
设 \(k\) 为域。令 \(A=T_3(k)\)、\(P_i=AE_{ii}\)、\(M=P_2\oplus P_3\)，并取子模
\[
N=kE_{12}\oplus kE_{13}=\operatorname{Soc}M.
\]
计算 \(\operatorname{rad}N\)、\(\operatorname{rad}(M/N)\) 与 \(\operatorname{Soc}(M/N)\)。比较 \(N\cap\operatorname{rad}M\) 与 \(\operatorname{rad}N\)，以及商映射 \(M\to M/N\) 所诱导的 \(\operatorname{Soc}M\to\operatorname{Soc}(M/N)\) 是否满射。
\end{exercise}
\begin{solution}
记 \(J=J(A)\)。由\eqref{b2:eq:triangular-projective-layers}取 \(r=1\) 的计算，
\[
JM=kE_{12}\oplus(kE_{13}+kE_{23}),\qquad JN=0.
\]
所以 \(\operatorname{rad}N=0\)，而 \(N\cap\operatorname{rad}M=N\ne0\)。由\cref{b2:prop:radical-socle-subquotients}，商模的根为
\[
\operatorname{rad}(M/N)=(JM+N)/N=k(E_{23}+N).
\]
这里第一分量 \(P_2/kE_{12}\) 是简单模 \(S_2\)；第二分量 \(P_3/kE_{13}\) 的根为 \(k(E_{23}+N)\)，顶层为 \(S_3\)。具体地，商模中 \(E_{22}+N\) 和 \(E_{23}+N\) 都被 \(J\) 零化，而 \(E_{23}(E_{33}+N)=E_{23}+N\ne0\)。因此
\[
\operatorname{Soc}(M/N)=k(E_{22}+N)\oplus k(E_{23}+N)\cong S_2^{\oplus2}.
\]
原基座就是 \(N\)，在商中像为零，故基座上的诱导映射不满射。取商会产生原模中并非简单子模的新的基座部分；对子模有交公式，并不意味着对商也能直接取原基座的像。
\end{solution}

\begin{exercise}\label{b2:ex:15-dual-numbers-loewy}
令 \(A=k[\varepsilon]/(\varepsilon^2)\)，\(M=A^{\oplus r}\oplus k^{\oplus s}\)，其中 \(r,s\ge0\)，\(\varepsilon\) 在 \(k\) 上作用为零。计算模的根、基座、最大半单商及 Loewy 长度，并说明在 \(r>0\) 时根列不可能作为简单子模的直和分解实现。
\end{exercise}
\begin{solution}
\(J=(\varepsilon)\)，故 \(JM=(k\varepsilon)^{\oplus r}\oplus0\)，维数为 \(r\)。被 \(J\) 零化的部分是 \((k\varepsilon)^{\oplus r}\oplus k^{\oplus s}\)，故基座维数为 \(r+s\)；商 \(M/JM\cong k^{\oplus(r+s)}\)。若 \(r>0\)，Loewy 长度为二；若 \(r=0,s>0\)，为一；全零时为零。\(r>0\) 时，\(\varepsilon\) 在 \(M\) 上非零，而在任何简单模直和上均为零，所以不能把各层拼成 \(M\) 的简单直和分解。
\end{solution}

\begin{exercise}\label{b2:ex:15-change-of-scalars-truncated}
设 \(k\) 为域、\(m\ge2\)、\(A=k[t]/(t^m)\)。分别把 \(A\) 看作 \(k[t]\) 模和左 \(A\) 模。比较两种模的有限展示性、投射性与平坦性；求它们的子模、合成因子及模的根。哪些结论依赖标量环，哪些只由这里共同的子模格决定？
\end{exercise}
\begin{solution}
作为 \(k[t]\) 模，\(A\) 有展示
\[
0\longrightarrow k[t]\xrightarrow{\,t^m\,}k[t]\longrightarrow A\longrightarrow0,
\]
所以有限展示。它不投射：整环 \(k[t]\) 上的投射模是自由模的直和因子，因而无挠，但 \(\bar1\ne0\) 被 \(t^m\) 零化。它也不平坦：把单射 \(k[t]\xrightarrow{t^m}k[t]\) 与 \(A\) 张量后，得到 \(A\xrightarrow{0}A\)，不再单射。作为左 \(A\) 模，它是秩一自由模，因此有限展示、投射且平坦。

\(k[t]\) 对 \(A\) 的作用经过商映射 \(k[t]\to A\)，所以两种模的子模完全相同。它们恰为 \((\bar t^r)\)、\(0\le r\le m\)，形成一条严格下降链；每个相邻商都同构于 \(k\)，在 \(k[t]\) 模中即 \(k[t]/(t)\)，在 \(A\) 模中即 \(A/(\bar t)\)。两种模都只有一个极大子模 \((\bar t)\)，故根相同。标量环改变了投射性和平坦性，却没有改变此例的子模格、合成因子层数与根。
\end{solution}

\begin{exercise}\label{b2:ex:15-triangular-column-submodules}
求 \(T_3(k)\) 的自然列向量模 \(V=k^3\) 的全部子模，并判断它是否不可分解。
\end{exercise}
\begin{solution}
若子模 \(W\) 含向量 \(v\)，对角矩阵单位 \(E_{ii}\) 使其每个坐标向量 \(v_ie_i\) 也属于 \(W\)，因此 \(W\) 是若干坐标轴的直和。若 \(e_j\in W\)，上三角矩阵单位 \(E_{ij}\)、\(i<j\) 又使 \(e_i\in W\)。所以只有
\[
0,\quad ke_1,\quad ke_1+ke_2,\quad V
\]
这四个子模。任意两个非零子模都包含 \(ke_1\)，不可能交为零，因此 \(V\) 不可分解。它的三层依次为 \(S_3,S_2,S_1\)，与\eqref{b2:eq:triangular-projective-layers}相符。
\end{solution}

\begin{exercise}\label{b2:ex:15-triangular-projective-hom}
仍令 \(A=T_3(k)\)、\(P_i=AE_{ii}\)。计算全部 \(\operatorname{Hom}_A(P_i,P_j)\)，给出非零同态的公式，并解释为什么 \(\operatorname{End}_A(P_3)=k\) 不表示 \(P_3\) 简单。
\end{exercise}
\begin{solution}
同态由 \(E_{ii}\) 的像 \(x\) 决定，且必须满足 \(x=E_{ii}x\in E_{ii}AE_{jj}\)。反过来，这样的 \(x\) 给出 \(aE_{ii}\mapsto ax\)；因为 \(x=E_{ii}x\)，该公式与表达方式无关并且左线性。\(E_{ii}AE_{jj}\) 在 \(i\le j\) 时由 \(E_{ij}\) 生成，否则为零。因此 Hom 空间的维数表为
\[
\begin{pmatrix}1&1&1\\0&1&1\\0&0&1\end{pmatrix},
\]
第 \((i,j)\) 项的非零同态都是右乘 \(\lambda E_{ij}\) 的映射。特别地各 \(\operatorname{End}_A(P_i)=k\)。然而 \(P_3\) 有真子模 \(kE_{13}\)，所以不简单；Schur 引理的逆命题在这里不成立。
\end{solution}

\begin{exercise}\label{b2:ex:15-idempotent-nonnilpotent-counterexample}
设整数 \(n\ge2\) 有分解 \(n=\prod_{i=1}^r p_i^{a_i}\)，其中 \(p_i\) 为不同素数、\(a_i\ge1\)。分类 \(\mathbb Z/n\mathbb Z\) 的全部幂等元，求其个数，并判断哪些能提升为 \(\mathbb Z\) 的幂等元。以 \(n=12\) 具体核验，再说明“全部商幂等元可提升”是否迫使理想 \(n\mathbb Z\) 幂零。
\end{exercise}
\begin{solution}
由第一卷定理10.4.3的中国剩余定理，
\[
\mathbb Z/n\mathbb Z\cong\prod_{i=1}^r\mathbb Z/p_i^{a_i}\mathbb Z.
\]
在模 \(p^a\) 的商中，幂等条件是 \(p^a\mid e(e-1)\)。相邻整数互素，故 \(p\) 不会同时整除 \(e\) 与 \(e-1\)，整个素数幂必须整除其中一个。因此每个坐标的幂等元只有零或一。各坐标可独立选取，全部幂等元由 \(\{0,1\}^r\) 参数化，共有 \(2^r\) 个。

整数环只有零、一两个幂等元，所以可提升的商幂等元恰为各坐标全零或全一者。于是全部可提升当且仅当 \(r=1\)，也就是 \(n\) 为素数幂。模十二时，四个幂等元为 \(0,1,4,9\)：后两者分别满足 \(4^2-4=12\)、\(9^2-9=72\)，并分别在模四、模三的坐标上取 \((0,1)\)、\((1,0)\)，不能提升。

另一方面，无论 \(n\ge2\) 取何值，都有 \((n\mathbb Z)^m=n^m\mathbb Z\ne0\) 对每个正整数 \(m\) 成立。当 \(n\) 为素数幂时全部幂等元可提升，但这个理想仍不幂零。因此幂零性是正文提升定理的充分假设，不能由某个具体商中的全部幂等元可提升反推出它。
\end{solution}

\begin{exercise}\label{b2:ex:15-square-zero-matrix-lift}
令 \(R=M_2\bigl(k[\varepsilon]/(\varepsilon^2)\bigr)\)、\(I=\varepsilon M_2(k)\)。对
\[
a=E_{11}+\varepsilon\begin{pmatrix}u&v\\w&z\end{pmatrix},
\]
用平方零修正公式求出幂等提升 \(e\)。再将它写成 \((1+\varepsilon X)E_{11}(1+\varepsilon X)^{-1}\)。
\end{exercise}
\begin{solution}
相乘得 \(d=a^2-a=\varepsilon\operatorname{diag}(u,-z)\)。使用\cref{b2:thm:idempotent-lifting}中平方零情形的修正 \(e=a+(1-2a)d\)。因 \(\varepsilon^2=0\)，有 \((a-E_{11})d=0\)，故 \((1-2a)d=(1-2E_{11})d=\varepsilon\operatorname{diag}(-u,-z)\)，得到
\[
e=E_{11}+\varepsilon\begin{pmatrix}0&v\\w&0\end{pmatrix}.
\]
取 \(X=\begin{psmallmatrix}0&-v\\w&0\end{psmallmatrix}\)，则 \(XE_{11}-E_{11}X=\begin{psmallmatrix}0&v\\w&0\end{psmallmatrix}\)。又 \((1+\varepsilon X)^{-1}=1-\varepsilon X\)，展开共轭即得 \(e\)。所有计算对任意特征成立。
\end{solution}

\begin{exercise}\label{b2:ex:15-compatible-idempotent-lifts}
设 \(k\) 为域，\(R=M_2(k[\varepsilon]/(\varepsilon^2))\)、\(I=\varepsilon M_2(k)\)。求 \(E_{11},E_{22}\in R/I\) 的全部幂等元提升，并确定哪些有序提升对 \((e,f)\) 正交且满足 \(e+f=1\)。若 \(k=\mathbb F_q\)，求这样的完整正交提升组有多少个。
\end{exercise}
\begin{solution}
写 \(e=E_{11}+\varepsilon\begin{psmallmatrix}u&v\\w&z\end{psmallmatrix}\)。直接计算得
\[
e^2-e=\varepsilon\begin{pmatrix}u&0\\0&-z\end{pmatrix}.
\]
所以所有幂等提升恰为 \(e=E_{11}+\varepsilon(vE_{12}+wE_{21})\)。同样，\(E_{22}\) 的全部幂等提升为 \(f=E_{22}+\varepsilon(cE_{12}+dE_{21})\)。两方向的乘积分别为
\[
ef=\varepsilon(v+c)E_{12},\qquad fe=\varepsilon(w+d)E_{21}.
\]
于是正交恰好要求 \(c=-v,d=-w\)，此时 \(f=1-e\)，总和也为一。反过来，总和为一也给出这两个条件。因此完整正交组由任意 \((v,w)\in k^2\) 唯一确定；有限域情形共有 \(q^2\) 组。这里不仅检查单个幂等方程，还须同时满足两方向的零乘积条件。
\end{solution}

\begin{exercise}\label{b2:ex:15-superfluous-ambient}
证明多余子模的任意子模仍在同一母模中多余。若 \(N\ll P\) 且 \(N\subseteq Q\subseteq P\)，是否必有 \(N\ll Q\)？给出有限维代数中的反例。
\end{exercise}
\begin{solution}
若 \(N'\subseteq N\) 且 \(L+N'=P\)，则 \(L+N=P\)，由 \(N\ll P\) 得 \(L=P\)。所以 \(N'\ll P\)。

缩小母模却未必保持多余性。取 \(A=k[\varepsilon]/(\varepsilon^2)\)、\(P=A\)、\(N=Q=k\varepsilon\)。\(N=JP\)，由\cref{b2:prop:finite-superfluous-radical}，\(N\ll P\)。但在 \(Q=N\ne0\) 中，\(0+N=Q\) 而 \(0\ne Q\)，故 \(N\not\ll Q\)。
\end{solution}

\begin{exercise}\label{b2:ex:15-truncated-projective-cover}
令 \(A=k[t]/(t^n)\)、\(1\le s\le n\)。构造 \(M=A/(\bar t^s)\) 的投射覆盖，并将其核识别为一个循环商模。分别求来源、目标和非零核的 Loewy 长度。
\end{exercise}
\begin{solution}
自然满射 \(A\to M\) 的来源自由，核 \((\bar t^s)\subseteq JA\)，所以是覆盖。映射 \(A\to(\bar t^s)\)、\(a\mapsto a\bar t^s\) 的核为 \((\bar t^{n-s})\)，对截断多项式逐项比较系数即可看出；因此核模同构于 \(A/(\bar t^{n-s})\)。\(s=n\) 时此模为零。来源与目标的 Loewy 长度分别为 \(n,s\)，核非零时为 \(n-s\)。覆盖保留的是目标的最小投射来源，并不要求来源和目标的长度相同。
\end{solution}

\begin{exercise}\label{b2:ex:15-matrix-algebra-covers}
令 \(A=M_m(k)\)、\(S=k^m\)、\(M=S^{\oplus r}\)，其中 \(m\ge1,r\ge0\)。确定 \(M\) 的投射覆盖，并判断它何时是自由 \(A\) 模。
\end{exercise}
\begin{solution}
\(A\) 半单，\(S\cong AE_{11}\) 投射，所以 \(M\) 本身投射。恒等映射 \(M\to M\) 的核为零，是投射覆盖。正则模按列分解为 \(A\cong S^{\oplus m}\)，因而 \(m\mid r\) 时 \(M\cong A^{r/m}\)。反过来，若 \(M\) 自由，其有限 \(k\) 维数使自由基有限，设有 \(d\) 个，则 \(mr=m^2d\)，故 \(m\mid r\)。这包含 \(r=0\) 的情形，并给出许多不能选为自由模的投射覆盖。
\end{solution}

\begin{exercise}\label{b2:ex:15-cover-automorphisms}
令 \(A=k[t]/(t^n)\)、\(n\ge1\)，\(p:A\to k\) 为自然覆盖。求全部满足 \(pu=p\) 的左模自同构 \(u:A\to A\)。若 \(k=\mathbb F_q\)，求这样的自同构个数。
\end{exercise}
\begin{solution}
自同态由 \(u(1)=a\) 决定，公式为 \(u(x)=xa\)。条件 \(pu=p\) 等价于 \(a\) 的常数项为一，即 \(a\in1+(\bar t)\)。这些元素均可逆，所以它们都给出自同构，也没有别的可能。域有 \(q\) 个元素时，\(\bar t,\ldots,\bar t^{n-1}\) 的系数各有 \(q\) 种选择，个数为 \(q^{n-1}\)。因此覆盖的唯一同构类型不意味着覆盖上的自同构只能是恒等。
\end{solution}

\begin{exercise}\label{b2:ex:15-lifting-between-covers}
设 \(k\) 为域、\(A=k[t]/(t^4)\)，对 \(1\le i\le4\) 令 \(M_i=A/(\bar t^i)\)，自然投射覆盖记为 \(p_i:A\to M_i\)。固定 \(1\le r,s\le4\)，描述所有同态 \(h:M_s\to M_r\)，并求使 \(p_rH=hp_s\) 的全部 \(H:A\to A\)；判断这些提升何时可逆。分别考察 \(M_1\to M_2\)、\(1\mapsto\bar t\)，以及自然商映射 \(M_2\to M_1\)，说明目标间的映射不是同构时，提升可能具有怎样不同的行为。
\end{exercise}
\begin{solution}
同态 \(h\) 由 \(z=h(1)\in M_r\) 决定。关系 \(\bar t^s\cdot1=0\) 要求 \(\bar t^sz=0\)，即
\[
z\in\bar t^{\max(r-s,0)}M_r.
\]
比较截断多项式的系数即可得到这一条件。每个来源上的自同态形如 \(H(x)=xa\)，其中 \(a=H(1)\in A\)。提升条件恰为 \(a+(\bar t^r)=z\)；固定一个代表元 \(a_0\) 后，全部提升由 \(a\in a_0+(\bar t^r)\) 给出。这样的 \(H\) 可逆，当且仅当 \(a\) 的常数项非零：常数项为零时它在模根后为零，不能可逆；常数项非零时可用有限几何和求逆。因为 \(r\ge1\)，所有提升的常数项相同，所以对给定 \(h\)，全部提升同时可逆或同时不可逆。

对 \(M_1\to M_2\)、\(1\mapsto\bar t\)，此映射单射，所有提升都有 \(a\equiv\bar t\pmod{\bar t^2}\)，常数项为零，故都不可逆。对自然满射 \(M_2\to M_1\)，所有提升满足 \(a\equiv1\pmod{\bar t}\)，故都可逆。这两个目标映射均非同构；正文的同构提升定理给出充分条件，不能把它倒过来作为必要条件。
\end{solution}

\begin{exercise}\label{b2:ex:15-triangular-idempotents}
求 \(T_2(k)\) 的全部幂等元，判断哪些本原，并证明所有非平凡者分别共轭于 \(E_{11}\) 或 \(E_{22}\)。
\end{exercise}
\begin{solution}
写幂等元为 \(\begin{psmallmatrix}a&b\\0&c\end{psmallmatrix}\)，方程为 \(a^2=a,c^2=c,(a+c-1)b=0\)。域上的前两式使 \(a,c\in\{0,1\}\)。若两者同为零或同为一，必有 \(b=0\)；其余两种情形中 \(b\) 任意。因此全部幂等元为 \(0,1,E_{11}+bE_{12},E_{22}+bE_{12}\)。

后两族模根后是 \(k^2\) 中的坐标幂等元，由\cref{b2:prop:primitive-projective-cover}本原；\(1=E_{11}+E_{22}\) 不本原，零不计本原。又
\[
(1+cE_{12})E_{11}(1-cE_{12})=E_{11}-cE_{12},
\]
对 \(E_{22}\) 的对应公式为 \(E_{22}+cE_{12}\)。分别取 \(c=-b\) 或 \(c=b\)，就得到所需共轭。
\end{solution}

\begin{exercise}\label{b2:ex:15-corner-filtration}
设 \(k\) 为域、\(n\ge1\)。令 \(A=T_n(k)\)，选取 \(1\le q_1<\cdots<q_s\le n\)，并令 \(e=\sum_{a=1}^sE_{q_aq_a}\)。识别 \(B=eAe\)，对所有 \(r\ge0\) 分别计算 \(J(B)^r\) 与 \(eJ(A)^re\)，并判断 \(e\) 何时为满幂等元。再取 \(n=4\)、\((q_1,q_2,q_3)=(1,3,4)\)，比较二次和三次根幂。
\end{exercise}
\begin{solution}
角代数的基为 \(E_{q_aq_b}\)、\(a\le b\)，乘法保持矩阵单位规则，故按所选指标的次序同构于 \(T_s(k)\)。在 \(T_n(k)\) 中，\(J(A)^r\) 由距主对角线至少 \(r\) 的矩阵单位组成；一条长度为 \(r\) 的非零乘积，恰需选择 \(r\) 个严格增加的指标间隔。因此
\[
\begin{aligned}
J(B)^r&=\operatorname{span}_k\{E_{q_aq_b}\mid b-a\ge r\},\\
eJ(A)^re&=\operatorname{span}_k\{E_{q_aq_b}\mid q_b-q_a\ge r\}.
\end{aligned}
\]
对 \(r=0\)，两式均包含对角项，给出 \(B\)；空指标集按约定给零。前式在所选指标中数可连接的步数，后式允许在原代数中经过未选的中间指标。

若漏选某个 \(i\)，任何 \(aeb\) 的第 \(i\) 个对角元都为零，所以 \(AeA\) 不含 \(E_{ii}\)。故 \(e\) 满，当且仅当所有指标均被选中，即 \(e=1\)。在指定的四阶例子中，
\[
J(B)^2=kE_{14},\quad eJ(A)^2e=kE_{13}+kE_{14},\qquad
J(B)^3=0,\quad eJ(A)^3e=kE_{14}.
\]
虽然一次根取角的公式仍成立，较高根幂可以在不同次数都出现差异。
\end{solution}

\begin{exercise}\label{b2:ex:15-basic-corner-example}
令 \(A=M_2(k)\times k[t]/(t^2)\)。选择一个满幂等元 \(e\)，使 \(eAe\) 基本；求半单商，并给出正则左模中各类不可分解投射模的出现次数。
\end{exercise}
\begin{solution}
取 \(e=(E_{11},1)\)，则 \(eAe\cong k\times k[t]/(t^2)\)，根为 \(0\times(\bar t)\)，半单商为 \(k\times k\)，所以基本。又
\[
e+(E_{21},0)e(E_{12},0)=(I_2,1),
\]
说明 \(AeA\) 包含单位，故 \(e\) 满。两类不可分解投射模分别是 \(P_1=k^2\times0\) 与 \(P_2=0\times k[t]/(t^2)\)，作用分别经两个因子。正则模分解为 \(P_1^{\oplus2}\oplus P_2\)。第二个分量的根非零，说明选取基本角代数没有抹去这些非半单结构。
\end{solution}

\begin{exercise}\label{b2:ex:15-real-simple-division-ring}
把 \(A=M_2(\mathbb C)\) 看成实代数，并令 \(S=\mathbb C^2\) 为左模。求 \(\operatorname{End}_A(S)\) 及其在 \(\mathbb R\) 上的维数，给出 \(S\) 的投射覆盖和 \(A\) 的一个基本角代数。
\end{exercise}
\begin{solution}
任何 \(A\) 线性映射都与标量矩阵 \(iI_2\) 交换，故首先是复线性的。与 \(E_{11},E_{22}\) 交换使它为对角矩阵，再与 \(E_{12}\) 交换使两个对角元相同。因此 \(\operatorname{End}_A(S)\cong\mathbb C\)，实维数为二，在自身除环上的维数为一。

\(S\cong AE_{11}\)，且 \(A\) 半单，所以 \(S\) 投射，恒等映射是其覆盖。取 \(e=E_{11}\)，得到基本角代数 \(eAe\cong\mathbb C\)。它的半单商就是 \(\mathbb C\)，不能在实基域上改成 \(\mathbb R\)。
\end{solution}

\begin{exercise}\label{b2:ex:15-quiver-with-zero-composition}
设 \(k\) 为域、\(B=T_4(k)/(kE_{14})\)，矩阵单位的商类沿用原符号。求 \(J(B)\) 的各次幂和箭图，写出由 \(P_4=BE_{44}\) 给出的 \(S_4\) 的投射覆盖，并计算 \(P_4\) 的 Loewy 层。比较相邻箭头的二步复合与从第四顶点到第一顶点的三步复合，说明关系会在哪一层改变列模。
\end{exercise}
\begin{solution}
矩阵单位的乘法说明 \(kE_{14}\) 是双边理想。严格上三角理想的像为
\[
J(B)=kE_{12}+kE_{13}+kE_{23}+kE_{24}+kE_{34},\qquad
J(B)^2=kE_{13}+kE_{24},\quad J(B)^3=0.
\]
这是幂零理想，商为 \(k^4\)，所以它确为 Jacobson 根。\(J(B)/J(B)^2\) 的基是 \(E_{12},E_{23},E_{34}\) 的商类，箭图为 \(4\to3\to2\to1\)。两个相邻箭头的复合分别给出非零的 \(E_{24}\) 与 \(E_{13}\)，但
\[
E_{12}E_{23}E_{34}=E_{14}=0.
\]
因此该关系不消掉二步路径，而只使这条三步路径为零。

列模 \(P_4\) 的基为 \(E_{44},E_{34},E_{24}\)，根链为
\[
P_4\supset kE_{34}+kE_{24}\supset kE_{24}\supset0.
\]
自然满射 \(P_4\to S_4\) 的核是 \(J(B)P_4\)，故为投射覆盖；非零 Loewy 层依次是 \(S_4,S_3,S_2\)。原 \(T_4(k)\) 的第四列还有最底层 \(S_1\)，商掉 \(E_{14}\) 正是移除了这一层，箭头数和二步乘积均未改变。
\end{solution}

\begin{exercise}\label{b2:ex:15-projective-hom-composition-factors}
假设 \(k\) 代数闭。设 \(P_i\to S_i\) 为简单模的投射覆盖，\(M\) 为有限维模。对 \(M\) 的任一合成列，证明其中同构于 \(S_i\) 的因子个数等于 \(\dim_k\operatorname{Hom}_A(P_i,M)\)。
\end{exercise}
\begin{solution}
对简单模 \(S_j\)，所有 \(P_i\to S_j\) 的同态都通过 \(P_i/JP_i\cong S_i\) 分解。由 Schur 引理及\cref{b2:prop:split-simple-endomorphisms}，Hom 的 \(k\) 维数在 \(i=j\) 时为一，否则为零。

取合成列 \(0=M_0\subset M_1\subset\cdots\subset M_\ell=M\)。将 \(\operatorname{Hom}_A(P_i,-)\) 用于每个 \(0\to M_{r-1}\to M_r\to M_r/M_{r-1}\to0\)。\(P_i\) 投射，由\cref{b2:cor:projective-hom-exact}得到短正合列，所以这些 Hom 空间的维数逐项相加。每个因子贡献一或零，恰得到题中计数。这也表明该个数不依赖所选合成列。该计数公式可对照 Etingof 等人\cite[Proposition 9.2.3, p. 211]{EtingofEtAl2011RepresentationTheory}。
\end{solution}

